% Two editions from one source. The Kindle edition is selected by defining
% \KindleEdition before the class is read, which the Makefile does on the
% command line:  pdflatex "\def\KindleEdition{}% Two editions from one source. The Kindle edition is selected by defining
% \KindleEdition before the class is read, which the Makefile does on the
% command line:  pdflatex "\def\KindleEdition{}% Two editions from one source. The Kindle edition is selected by defining
% \KindleEdition before the class is read, which the Makefile does on the
% command line:  pdflatex "\def\KindleEdition{}% Two editions from one source. The Kindle edition is selected by defining
% \KindleEdition before the class is read, which the Makefile does on the
% command line:  pdflatex "\def\KindleEdition{}\input{main.tex}"
\ifdefined\KindleEdition
  \documentclass[10pt,oneside]{book}
\else
  \documentclass[11pt,a4paper,twoside,openright]{book}
\fi
\input{preamble}
\renewcommand*{\thisvolume}{functions}

\title{\Huge Functions and Their Graphs\\[0.4em]
  \large A Refresher with Exercises and Worked Solutions\\[0.8em]
  \normalsize Domain \(\cdot\) Range \(\cdot\) Graphs \(\cdot\)
  Transformations \(\cdot\) Regions \(\cdot\) Trigonometry}
\author{}
\date{}

\begin{document}
\frontmatter
\maketitle

\chapter*{About this sheet}
\addcontentsline{toc}{chapter}{About this sheet}

The algebra volume manipulates \emph{expressions}. This one is about
\emph{functions}: an expression together with a decision about which numbers
may be substituted into it, a record of which numbers come out, and above
all a picture.

Everything that follows in the later volumes leans on this material and
rarely stops to explain it. A limit is taken at a point of the domain or at
its boundary; a derivative exists only where the function does; an
asymptote is read off a graph; a constrained optimum is the point where two
families of curves touch. A mathematics sequence for economics or business
takes this for granted: one is expected to know what the graphs of the usual
functions look like, and to be able to sketch curves defined by equations in
two variables. That is the brief of this volume.

Five chapters, in the order they build on each other. \textbf{Domain and
range} fixes what a function is and where it lives. \textbf{The standard
functions and their graphs} is the catalogue --- powers, roots, exponentials,
logarithms, rational functions with their poles and asymptotes --- together
with the curves that are not graphs at all: circles, ellipses, hyperbolas,
isoquants. \textbf{Transforming functions} covers shifting, stretching,
reflecting, composing and inverting, and closes with the question an
economist meets most often, which is the reverse one: given the shape you
need, which formula has it? \textbf{Inequalities in two variables} turns
curves into regions, the skill behind every feasible set. And
\textbf{trigonometric functions} supplies the periodic shapes that later material
assumes and the calculus chapters use as test cases.

The structure is the same in every chapter: \textbf{Core Rules} states the
rules compactly; \textbf{Practice Problems} drills them; \textbf{Solutions}
works every problem through in full; each chapter closes with
\emph{additional problems} for a second pass. Seventy problems per chapter,
every one with a worked solution.

Work the problems before reading the solutions, and draw the pictures by
hand. A curve you have sketched yourself is one you will recognise in an
examination.

\tableofcontents

\mainmatter
\include{chapters/ch01}
\include{chapters/ch02}
\include{chapters/ch03}
\include{chapters/ch04}
\include{chapters/ch05}

\end{document}
"
\ifdefined\KindleEdition
  \documentclass[10pt,oneside]{book}
\else
  \documentclass[11pt,a4paper,twoside,openright]{book}
\fi
% Shared preamble for the Functions training sheet.
\usepackage[T1]{fontenc}
\usepackage[utf8]{inputenc}
\usepackage{lmodern}
\usepackage[a4paper,margin=2.4cm]{geometry}
\usepackage{amsmath,amssymb,amsthm}
\usepackage{microtype}
\usepackage{enumitem}
\usepackage{xcolor}
\usepackage[most]{tcolorbox}
\usepackage{fancyhdr}
\usepackage[hidelinks]{hyperref}

\definecolor{accent}{HTML}{1F4E79}
\definecolor{soft}{HTML}{F4F7FA}

% --- function graphs ------------------------------------------------------
\usepackage{pgfplots}
\pgfplotsset{compat=1.18}
\usetikzlibrary{arrows.meta}

% One consistent look for every function graph in these books.
\pgfplotsset{
  fnaxis/.style={
    width=5.9cm, height=4.6cm,
    axis lines=middle,
    axis line style={-{Stealth[length=4pt]}, gray!65!black, line width=0.3pt},
    tick style={gray!65!black},
    tick label style={font=\tiny, gray!50!black},
    label style={font=\scriptsize},
    xlabel={$x$}, ylabel={$y$},
    every axis x label/.style={at={(ticklabel* cs:1.0)}, anchor=west},
    every axis y label/.style={at={(ticklabel* cs:1.0)}, anchor=south east,
      xshift=-1pt},
    every axis plot/.append style={accent, very thick, smooth, samples=201},
    clip=true, enlargelimits=false,
    title style={font=\scriptsize, yshift=4pt},
  },
  fnaxis/.append style={/pgfplots/legend style={font=\tiny, draw=none,
    fill=none, at={(0.98,0.03)}, anchor=south east}},
}

% A row of small graphs, evenly spaced with \hfill between them, plus an
% optional caption line: \begin{gallery} <pic>\hfill<pic> \gcap{...} \end{gallery}
\newenvironment{gallery}
  {\par\medskip\noindent\ignorespaces}
  {\par\medskip\noindent}
\newcommand{\gcap}[1]{\par\smallskip{\centering\footnotesize #1\par}}

% --- plane regions, implicit curves, surfaces -----------------------------
\usepgfplotslibrary{fillbetween}
\definecolor{regionfill}{HTML}{9DB4D8}
\pgfplotsset{colormap={house}{color=(soft) color=(accent!45) color=(accent)}}

\pgfplotsset{
  % Squared, gridded axis for plane regions and implicitly given curves.
  geoaxis/.style={
    fnaxis,
    width=5.9cm, height=5.9cm,
    grid=both,
    grid style={gray!22, line width=0.2pt},
    minor tick num=1,
  },
  % Level-curve maps: many thin curves, no single curve emphasised.
  contouraxis/.style={
    geoaxis,
    every axis plot/.append style={accent!75, line width=0.5pt, smooth,
      samples=201},
  },
  % Surfaces: \addplot3[surf, domain=..., y domain=...] {expr};
  surfaxis/.style={
    width=6.0cm, height=5.0cm,
    view={130}{22},
    colormap name=house,
    ticks=none,
    xlabel={$x$}, ylabel={$y$}, zlabel={$f$},
    label style={font=\tiny, gray!50!black},
    title style={font=\scriptsize},
    every axis plot/.append style={samples=25, shader=faceted interp,
      draw=accent!35, line width=0.1pt},
  },
}

% Shorthand for a shaded region between two named paths.
% \region{<path A>}{<path B>}{<soft clip options>}
\newcommand{\region}[3]{%
  \addplot[regionfill, opacity=0.55, forget plot, sharp plot]
    fill between[of=#1 and #2, #3];}

% --- rule boxes -----------------------------------------------------------
\newtcolorbox{rulebox}[1]{%
  breakable, enhanced, colback=soft, colframe=accent,
  boxrule=0.4pt, left=7pt, right=7pt, top=6pt, bottom=6pt,
  fonttitle=\bfseries, coltitle=white, colbacktitle=accent,
  title={#1}}

% --- worked examples ------------------------------------------------------
% One fully narrated solution per chapter, carrying the chapter's central
% technique and the mistakes that technique invites:
%   \begin{worked}{Task} ... \wstep{Step 1.} ... \wtrap{...} ... \end{worked}
\newtcolorbox{worked}[1]{%
  breakable, enhanced, colback=white, colframe=accent!65,
  boxrule=0.4pt, left=7pt, right=7pt, top=6pt, bottom=6pt,
  fonttitle=\bfseries, coltitle=accent, colbacktitle=soft,
  title={#1}}

% Run-in step heading inside a worked example.
\newcommand{\wstep}[1]{\par\smallskip\noindent\textbf{#1}\enspace\ignorespaces}

% A pitfall, set off from the running solution so it can be found again.
\newcommand{\wtrap}[1]{%
  \par\smallskip
  {\leftskip=1.3em\noindent
   \llap{\textcolor{accent}{\rule[0.22ex]{0.9em}{0.6pt}}\hspace{0.4em}}%
   \textbf{Trap.}\enspace\itshape #1\par}%
  \smallskip}

% --- problem / solution lists --------------------------------------------
\newlist{problems}{enumerate}{1}
\setlist[problems]{label=\textbf{\arabic*.}, leftmargin=2.1em,
  itemsep=0.55ex, parsep=0pt, topsep=0.8ex}
\newlist{solutions}{enumerate}{1}
\setlist[solutions]{label=\textbf{\arabic*.}, leftmargin=2.1em,
  itemsep=1.0ex, parsep=0pt, topsep=0.8ex}

% --- running heads --------------------------------------------------------
\pagestyle{fancy}
\fancyhf{}
\fancyhead[LE,RO]{\thepage}
\fancyhead[RE]{\small\itshape\nouppercase{\leftmark}}
\fancyhead[LO]{\small\itshape\nouppercase{\rightmark}}
\setlength{\headheight}{14pt}
\renewcommand{\headrulewidth}{0.3pt}
\renewcommand{\chaptermark}[1]{\markboth{\chaptername\ \thechapter.\ #1}{}}
\renewcommand{\sectionmark}[1]{\markright{\thesection\ #1}}

% --- shorthands used throughout -------------------------------------------
\newcommand{\dd}[2]{\frac{\mathrm{d}#1}{\mathrm{d}#2}}
\newcommand{\pd}[2]{\frac{\partial #1}{\partial #2}}
\newcommand{\pdd}[3]{\frac{\partial^{2} #1}{\partial #2\,\partial #3}}
\newcommand{\R}{\mathbb{R}}
\newcommand{\grad}{\nabla}
\DeclareMathOperator{\sgn}{sgn}
\DeclareMathOperator{\Hess}{H}

% German gloss, used on first mention of a technical term
\newcommand{\de}[1]{\emph{#1}}

\renewcommand*{\thisvolume}{functions}

\title{\Huge Functions and Their Graphs\\[0.4em]
  \large A Refresher with Exercises and Worked Solutions\\[0.8em]
  \normalsize Domain \(\cdot\) Range \(\cdot\) Graphs \(\cdot\)
  Transformations \(\cdot\) Regions \(\cdot\) Trigonometry}
\author{}
\date{}

\begin{document}
\frontmatter
\maketitle

\chapter*{About this sheet}
\addcontentsline{toc}{chapter}{About this sheet}

The algebra volume manipulates \emph{expressions}. This one is about
\emph{functions}: an expression together with a decision about which numbers
may be substituted into it, a record of which numbers come out, and above
all a picture.

Everything that follows in the later volumes leans on this material and
rarely stops to explain it. A limit is taken at a point of the domain or at
its boundary; a derivative exists only where the function does; an
asymptote is read off a graph; a constrained optimum is the point where two
families of curves touch. A mathematics sequence for economics or business
takes this for granted: one is expected to know what the graphs of the usual
functions look like, and to be able to sketch curves defined by equations in
two variables. That is the brief of this volume.

Five chapters, in the order they build on each other. \textbf{Domain and
range} fixes what a function is and where it lives. \textbf{The standard
functions and their graphs} is the catalogue --- powers, roots, exponentials,
logarithms, rational functions with their poles and asymptotes --- together
with the curves that are not graphs at all: circles, ellipses, hyperbolas,
isoquants. \textbf{Transforming functions} covers shifting, stretching,
reflecting, composing and inverting, and closes with the question an
economist meets most often, which is the reverse one: given the shape you
need, which formula has it? \textbf{Inequalities in two variables} turns
curves into regions, the skill behind every feasible set. And
\textbf{trigonometric functions} supplies the periodic shapes that later material
assumes and the calculus chapters use as test cases.

The structure is the same in every chapter: \textbf{Core Rules} states the
rules compactly; \textbf{Practice Problems} drills them; \textbf{Solutions}
works every problem through in full; each chapter closes with
\emph{additional problems} for a second pass. Seventy problems per chapter,
every one with a worked solution.

Work the problems before reading the solutions, and draw the pictures by
hand. A curve you have sketched yourself is one you will recognise in an
examination.

\tableofcontents

\mainmatter
\chapter{Functions, Domain and Range}

Up to here every chapter has manipulated \emph{expressions}. A
\de{Funktion} is the next step: an expression together with a decision
about which numbers may be substituted into it, and a record of which
numbers come out. Those two sets, the \de{Definitionsbereich} (domain) and
the \de{Wertebereich} (range), are not decoration. A limit is taken at a
point of the domain or at its boundary; a derivative exists only where the
function does; an optimisation problem is solved over a domain; and a
demand function that returns a negative price has left the region where it
means anything.

Finding a maximal domain is pure algebra of the kind the previous chapters
drilled: a denominator must not vanish, an even root needs a non-negative
radicand, a logarithm needs a positive argument. Each condition is an
equation or an inequality, and the domain is what survives all of them at
once. Finding a range is slightly different work: either the function is
built from pieces whose ranges are known, or the equation $y=f(x)$ is
solved for $x$ and one asks which $y$ admit a solution.

\section{Core Rules}

\begin{rulebox}{Function, domain, range}
A function $f\colon D\to\R$ assigns to every $x\in D$ exactly one value
$f(x)\in\R$. Two sets belong to the description:
\begin{itemize}[leftmargin=1.6em,itemsep=0.3ex]
  \item the \textbf{domain} $D(f)$ \de{(Definitionsbereich)}: the $x$ that
        may be substituted;
  \item the \textbf{range} $W(f)=\{f(x): x\in D(f)\}$
        \de{(Wertebereich, Bildmenge)}: the values actually attained.
\end{itemize}
The \textbf{graph} of $f$ is the set $\{(x,f(x)) : x\in D(f)\}$ in the
plane. A curve is the graph of a function exactly when every vertical line
meets it at most once (\textbf{vertical line test}). Reading off a graph:
$D(f)$ is its shadow on the $x$-axis, $W(f)$ its shadow on the $y$-axis.
\end{rulebox}

\begin{gallery}
\begin{tikzpicture}
  \begin{axis}[fnaxis, xmin=-0.6, xmax=5.4, ymin=-0.6, ymax=3.4,
      xtick={1,3,5}, ytick={1,2,3}, title={$f(x)=\sqrt{x-1}$}]
    \addplot[domain=1:5]{sqrt(x-1)};
  \end{axis}
\end{tikzpicture}\hfill
\begin{tikzpicture}
  \begin{axis}[fnaxis, xmin=-3.4, xmax=3.4, ymin=-0.6, ymax=4.4,
      xtick={-2,2}, ytick={2,4}, title={$f(x)=x^{2}$}]
    \addplot[domain=-2.1:2.1]{x^2};
  \end{axis}
\end{tikzpicture}\hfill
\begin{tikzpicture}
  \begin{axis}[fnaxis, xmin=-3.4, xmax=3.4, ymin=-1.2, ymax=4.4,
      xtick={-2,2}, ytick={1,3}, title={$f(x)=\mathrm{e}^{x}$}]
    \addplot[domain=-3:1.45]{exp(x)};
  \end{axis}
\end{tikzpicture}
\gcap{$D=[1,\infty)$, $W=[0,\infty)$ \quad\textbullet\quad
$D=\R$, $W=[0,\infty)$ \quad\textbullet\quad $D=\R$, $W=(0,\infty)$.}
\end{gallery}

\begin{rulebox}{The maximal domain: three obstructions}
Unless a domain is prescribed, $D(f)$ is the \textbf{maximal domain}: every
real $x$ for which the formula produces a real number. Only three things
can go wrong, and each is an (in)equality to solve:
\begin{itemize}[leftmargin=1.6em,itemsep=0.3ex]
  \item \textbf{Denominator:} $\dfrac{\;\cdot\;}{g(x)}$ requires
        $g(x)\neq 0$. Solve $g(x)=0$ and delete those points.
  \item \textbf{Even root:} $\sqrt[2k]{g(x)}$ requires $g(x)\geq 0$.
        Odd roots ($\sqrt[3]{\;}$) are defined for all reals.
  \item \textbf{Logarithm:} $\ln g(x)$ requires $g(x)>0$ (strictly).
\end{itemize}
A power $g(x)^{r}$ with non-integer $r$ counts as a root: it requires
$g(x)\geq 0$ for $r>0$ and $g(x)>0$ for $r<0$. If several obstructions
occur, \textbf{intersect} the resulting sets. A root in a denominator
combines two conditions into one: $\dfrac{1}{\sqrt{g(x)}}$ needs
$g(x)>0$.
\end{rulebox}

\begin{rulebox}{Domains prescribed by context}
In economics the domain is usually narrower than the maximal domain and
comes from the model, not from the formula:
quantities and inputs satisfy $q\geq 0$, $K,L\geq 0$; a price must satisfy
$p\geq 0$; a share or probability lives in $[0,1]$; an average cost
$AC(q)=C(q)/q$ needs $q>0$. State the economic domain explicitly, because
the formula will happily return values outside it: the linear demand
$p=120-3q$ is defined for every real $q$, but only $q\in[0,40]$ keeps both
$q\geq0$ and $p\geq 0$.
\end{rulebox}

\begin{rulebox}{Finding the range}
Four routes, in increasing order of generality:
\begin{enumerate}[leftmargin=1.7em,itemsep=0.3ex]
  \item \textbf{Standard building blocks.}
        $x^{2}\colon[0,\infty)$, \;$\sqrt{x}\colon[0,\infty)$,
        \;$|x|\colon[0,\infty)$, \;$\mathrm{e}^{x}\colon(0,\infty)$,
        \;$\ln x\colon\R$, \;$1/x\colon\R\setminus\{0\}$,
        \;$x^{3}\colon\R$, \;$\sin x,\cos x\colon[-1,1]$.
        Then track what the outer operations do: adding $c$ shifts the
        range by $c$, multiplying by $a>0$ scales it, multiplying by $a<0$
        scales \emph{and} flips it.
  \item \textbf{Complete the square} for a quadratic:
        $a(x-h)^{2}+k$ has range $[k,\infty)$ if $a>0$ and
        $(-\infty,k]$ if $a<0$.
  \item \textbf{Monotonicity.} If $f$ is strictly increasing on $[a,b]$,
        then $W=[f(a),f(b)]$; open ends stay open, and an end at
        $\pm\infty$ is replaced by the limiting value.
  \item \textbf{Solve $y=f(x)$ for $x$} (the general method). The range is
        the set of $y$ for which the solution exists and lies in $D(f)$.
        For $f(x)=\dfrac{x}{x+1}$: \;$y(x+1)=x\Rightarrow x=\dfrac{y}{1-y}$,
        solvable for every $y\neq 1$, hence $W=\R\setminus\{1\}$.
\end{enumerate}
\end{rulebox}

\begin{rulebox}{Composite functions}
$(g\circ f)(x)=g\bigl(f(x)\bigr)$ is evaluated \textbf{inside out}. Its
domain carries two conditions at once:
\[
  D(g\circ f)=\{\,x\in D(f)\;:\;f(x)\in D(g)\,\}.
\]
So $x$ must be admissible for the inner function, and its \emph{output}
must be admissible for the outer one. For $g(u)=\ln u$ and $f(x)=x^{2}-1$:
$x^{2}-1>0$, hence $D=(-\infty,-1)\cup(1,\infty)$. Order matters:
$g\circ f\neq f\circ g$ in general, and the two domains usually differ.
\end{rulebox}

\begin{rulebox}{Symmetry, monotonicity, injectivity}
\begin{itemize}[leftmargin=1.6em,itemsep=0.3ex]
  \item $f$ is \textbf{even} \de{(gerade)} if $f(-x)=f(x)$ for all $x\in D$
        (graph symmetric about the $y$-axis), \textbf{odd}
        \de{(ungerade)} if $f(-x)=-f(x)$ (symmetric about the origin).
        Only even powers $\Rightarrow$ even; only odd powers
        $\Rightarrow$ odd; a mix is neither.
  \item $f$ is \textbf{strictly increasing} if $x_{1}<x_{2}$ implies
        $f(x_{1})<f(x_{2})$, strictly decreasing if the inequality
        reverses.
  \item $f$ is \textbf{injective} \de{(eineindeutig)} if $f(x_{1})=f(x_{2})$
        forces $x_{1}=x_{2}$: no horizontal line meets the graph twice.
        Strictly monotone $\Rightarrow$ injective $\Rightarrow$ invertible
        on its range (Chapter~3).
\end{itemize}
\end{rulebox}

\section{Worked Example}

\begin{worked}{Find the maximal domain of $f(x)=\dfrac{\ln(x+3)}{\sqrt{4-x^{2}}}$, and the range of $g(x)=\dfrac{x+3}{x-2}$}
\wstep{Step 1. List the obstructions, one per structure in the formula.}
Three of them appear here, which is why this example was chosen:
\begin{itemize}[leftmargin=1.6em,itemsep=0.2ex]
  \item the logarithm needs $x+3>0$;
  \item the even root needs $4-x^{2}\ge0$;
  \item the denominator needs $\sqrt{4-x^{2}}\neq0$.
\end{itemize}
\wtrap{The last two combine into the \emph{strict} inequality $4-x^{2}>0$.
A square root sitting in a denominator is not allowed to be zero, so the
usual ``$\ge$'' for an even root tightens to ``$>$''. Handling the root and
the denominator as separate conditions and then forgetting to intersect them
is the standard way to hand in $[-2,2]$ instead of the right answer.}

\wstep{Step 2. Solve each condition separately.}
\[
  x+3>0 \iff x>-3 ;
  \qquad
  4-x^{2}>0 \iff x^{2}<4 \iff -2<x<2 .
\]

\wstep{Step 3. Intersect --- never unite.}
The function is defined where \emph{all} conditions hold at once:
\[
  D_{f}=(-3,\infty)\cap(-2,2)=\boxed{(-2,2)} .
\]
\wtrap{``And'' means intersection. Each obstruction removes points; nothing
in the formula ever gives points back. If your answer is larger than one of
the individual conditions, you took a union by mistake.}

\wstep{Step 4. Check the boundary behaviour, which confirms the strictness.}
As $x\to2^{-}$ the denominator $\to0^{+}$ while $\ln5>0$, so $f\to+\infty$.
The endpoint really must be excluded. At $x=-2$ the same thing happens.

\wstep{Step 5. The range of $g$: solve $y=g(x)$ for $x$.}
The range is the set of $y$ for which the equation $y=g(x)$ has a solution.
So treat $y$ as given and solve:
\[
  y=\frac{x+3}{x-2}
  \;\Longrightarrow\;
  y(x-2)=x+3
  \;\Longrightarrow\;
  yx-x=2y+3
  \;\Longrightarrow\;
  x(y-1)=2y+3 .
\]

\wstep{Step 6. Read off which $y$ fail.}
\[
  x=\frac{2y+3}{y-1},
\]
which exists for every $y$ except $y=1$; and at $y=1$ the previous line reads
$0=5$, impossible. Hence
\[
  \boxed{\;\text{range}(g)=\R\setminus\{1\}\;}
\]

\wstep{Step 7. Confirm it against the graph.}
$g(x)=\frac{x+3}{x-2}=1+\frac{5}{x-2}$, so the graph is the hyperbola
$5/x$ shifted, with horizontal asymptote $y=1$. The value $1$ is approached
but never attained --- exactly what the algebra said. Two independent
routes agreeing is worth the extra line.
\wtrap{Do not assume the range of a rational function is all of $\R$, and do
not assume it is the same as the domain. The excluded input ($x=2$) and the
excluded output ($y=1$) are different numbers, read off different parts of
the formula.}
\end{worked}

\section{Practice Problems}

State the maximal domain $D(f)\subseteq\R$ in problems 1--18, using
interval notation.

\begin{problems}
  \item $f(x)=\dfrac{x+3}{x-2}$
  \item $f(x)=\sqrt{x-5}$
  \item $f(x)=\sqrt{9-x^{2}}$
  \item $f(x)=\ln(x-4)$
  \item $f(x)=\ln(x^{2}-4)$
  \item $f(x)=\dfrac{\sqrt{x}}{x-1}$
  \item $f(x)=\dfrac{1}{\sqrt{x-3}}$
  \item $f(x)=\sqrt{x^{2}-x-6}$
  \item $f(x)=\dfrac{\ln x}{x-2}$
  \item $f(x)=\mathrm{e}^{1/x}$
  \item $f(x)=\sqrt{x+2}+\sqrt{4-x}$
  \item $f(x)=\dfrac{x^{2}-1}{x^{2}-3x+2}$
  \item $f(x)=\ln(\ln x)$
  \item $f(x)=x^{3/4}$
  \item $f(x)=\dfrac{1}{1-\mathrm{e}^{x}}$
  \item $f(x)=\sqrt{\ln x}$
  \item $f(x)=\sqrt[3]{x-7}$
  \item $f(x)=\dfrac{\sqrt{4-x}}{x^{2}-x}$
\end{problems}

\begin{problems}[resume]
  \item A firm has cost $C(q)=50+2q$. State the economic domain of $C$
        and of the average cost $AC(q)=C(q)/q$.
  \item For the demand schedule $p(q)=120-3q$, state the largest interval
        of quantities on which both $q\geq0$ and $p\geq0$ hold.
  \item State the domain of the Cobb--Douglas product
        $Y(K,L)=K^{0.3}L^{0.7}$ as a subset of the $(K,L)$-plane, and say
        why the exponents force the restriction.
  \item Utility is $U(x)=\ln x$. State the economic domain and explain what
        happens to $U$ as $x\to 0^{+}$.
\end{problems}

Determine the range $W(f)$ in problems 23--36. The domain is the maximal
one unless stated otherwise.

\begin{problems}[resume]
  \item $f(x)=x^{2}+4$
  \item $f(x)=x^{2}-6x+11$
  \item $f(x)=-2x^{2}+8x-3$
  \item $f(x)=\mathrm{e}^{x}-3$
  \item $f(x)=5-\mathrm{e}^{x}$
  \item $f(x)=\ln(x+2)$
  \item $f(x)=\dfrac{1}{x-2}+3$
  \item $f(x)=\sqrt{x-1}$
  \item $f(x)=\sqrt{16-x^{2}}$
  \item $f(x)=|x-3|$
  \item $f(x)=\dfrac{2x+1}{x-3}$
  \item $f(x)=\dfrac{x^{2}}{x^{2}+1}$
  \item $f(x)=x^{2}$ on $D=[-1,3]$
  \item $f(x)=2x-5$ on $D=(1,4]$
\end{problems}

\begin{problems}[resume]
  \item With $f(u)=\sqrt{u}$ and $g(x)=x-4$: give $(f\circ g)(x)$ and
        $(g\circ f)(x)$ with both domains.
  \item With $f(u)=1/u$ and $g(x)=x+2$: give $(f\circ g)(x)$ and its
        domain.
  \item With $f(u)=\ln u$ and $g(x)=9-x^{2}$: give $(f\circ g)(x)$ and its
        domain.
  \item With $f(u)=\sqrt{u}$ and $g(x)=1-x^{2}$: give $(f\circ g)(x)$, its
        domain and its range.
  \item Classify as even, odd, or neither:
        (a) $x^{4}-3x^{2}$, \;(b) $x^{3}-x$, \;(c) $x^{2}+x$,
        \;(d) $\dfrac{1}{x}$.
  \item Which of $x^{3}$, $x^{2}$, $\mathrm{e}^{x}$, $|x|$ are injective on
        $\R$?
  \item Give domain and range of
        $f(x)=\begin{cases} x^{2}, & x<1,\\ 2x-1, & x\geq 1.\end{cases}$
  \item Does $x^{2}+y^{2}=25$ define $y$ as a function of $x$? Justify with
        the vertical line test, and give the two functions obtained by
        solving for $y$, with their domains.
  \item A graph consists of the segment from $(-2,1)$ to $(0,5)$ together
        with the horizontal ray from $(0,5)$ to the right. Read off
        $D$ and $W$.
\end{problems}

\section{Solutions}

\begin{solutions}
  \item Denominator $x-2\neq0$: \;$D=\R\setminus\{2\}
        =(-\infty,2)\cup(2,\infty)$.
  \item $x-5\geq0$: \;$D=[5,\infty)$.
  \item $9-x^{2}\geq0\iff x^{2}\leq9\iff|x|\leq3$: \;$D=[-3,3]$.
  \item $x-4>0$: \;$D=(4,\infty)$.
  \item $x^{2}-4>0\iff|x|>2$: \;$D=(-\infty,-2)\cup(2,\infty)$.
  \item Two conditions: $x\geq0$ for the root and $x\neq1$ for the
        denominator: \;$D=[0,1)\cup(1,\infty)$.
  \item The root sits in the denominator, so $x-3>0$ (not merely
        $\geq0$): \;$D=(3,\infty)$.
  \item $x^{2}-x-6=(x-3)(x+2)\geq0$; the parabola opens upwards, so the
        product is non-negative outside the roots:
        \;$D=(-\infty,-2]\cup[3,\infty)$.
  \item $\ln x$ needs $x>0$, the denominator needs $x\neq2$:
        \;$D=(0,2)\cup(2,\infty)$.
  \item Only $x\neq0$; the exponential itself is unrestricted:
        \;$D=\R\setminus\{0\}$.
  \item $x+2\geq0$ and $4-x\geq0$, i.e. $x\geq-2$ \emph{and} $x\leq4$:
        \;$D=[-2,4]$.
  \item $x^{2}-3x+2=(x-1)(x-2)$: \;$D=\R\setminus\{1,2\}$. (The numerator
        also vanishes at $x=1$, but a cancelling factor does not repair the
        domain: $x=1$ stays excluded.)
  \item Inner: $x>0$; outer: $\ln x>0\iff x>1$. The stronger condition
        wins: \;$D=(1,\infty)$.
  \item Exponent $3/4$ is not an integer, so the base must satisfy
        $x\geq0$: \;$D=[0,\infty)$.
  \item $1-\mathrm{e}^{x}=0\iff \mathrm{e}^{x}=1\iff x=0$:
        \;$D=\R\setminus\{0\}$.
  \item $\ln x$ needs $x>0$, and $\ln x\geq0$ needs $x\geq1$; the
        intersection is \;$D=[1,\infty)$.
  \item An odd root accepts every real radicand: \;$D=\R$.
  \item Root: $4-x\geq0\iff x\leq4$. Denominator: $x^{2}-x=x(x-1)\neq0$,
        so $x\neq0$ and $x\neq1$. Together:
        \;$D=(-\infty,0)\cup(0,1)\cup(1,4]$.
  \item $C$: production cannot be negative, so $D(C)=[0,\infty)$; at $q=0$
        the value $C(0)=50$ is the fixed cost. $AC$: dividing by $q$
        excludes the origin, so $D(AC)=(0,\infty)$.
  \item $p\geq0\iff 120-3q\geq0\iff q\leq40$; with $q\geq0$ this gives
        $q\in[0,40]$.
  \item $D=\{(K,L): K\geq0,\;L\geq0\}$, the non-negative quadrant. The
        exponents $0.3$ and $0.7$ are not integers, so $K^{0.3}$ and
        $L^{0.7}$ require non-negative bases; negative input quantities
        would be meaningless anyway.
  \item $D=(0,\infty)$: consumption must be strictly positive, since
        $\ln 0$ is undefined. As $x\to0^{+}$, $U(x)\to-\infty$: the
        logarithmic utility punishes zero consumption infinitely hard,
        which is exactly why interior solutions appear in Cobb--Douglas
        problems.
  \item $x^{2}\geq0$, so $x^{2}+4\geq4$: \;$W=[4,\infty)$.
  \item $x^{2}-6x+11=(x-3)^{2}+2$: \;$W=[2,\infty)$.
  \item $-2x^{2}+8x-3=-2(x^{2}-4x)-3=-2(x-2)^{2}+5$; the parabola opens
        downwards: \;$W=(-\infty,5]$.
  \item $\mathrm{e}^{x}\in(0,\infty)$, shifted down by $3$:
        \;$W=(-3,\infty)$.
  \item $\mathrm{e}^{x}\in(0,\infty)$, negated to $(-\infty,0)$, shifted up
        by $5$: \;$W=(-\infty,5)$.
  \item $\ln$ attains every real value, and $x+2$ runs through all of
        $(0,\infty)$ on $D=(-2,\infty)$: \;$W=\R$.
  \item $1/(x-2)$ attains every value except $0$; adding $3$ shifts this:
        \;$W=\R\setminus\{3\}$.
  \item A square root returns exactly the non-negative numbers, and
        $x-1$ runs through $[0,\infty)$: \;$W=[0,\infty)$.
  \item $D=[-4,4]$, on which $16-x^{2}$ runs through $[0,16]$; the root of
        that is \;$W=[0,4]$.
  \item An absolute value is non-negative and attains every non-negative
        value (at $x=3\pm w$): \;$W=[0,\infty)$.
  \item Solve $y=\dfrac{2x+1}{x-3}$: \;$y(x-3)=2x+1\Rightarrow
        x(y-2)=3y+1\Rightarrow x=\dfrac{3y+1}{y-2}$, which exists exactly
        for $y\neq2$: \;$W=\R\setminus\{2\}$. ($y=2$ is the horizontal
        asymptote.)
  \item Solve $y=\dfrac{x^{2}}{x^{2}+1}$: \;$y(x^{2}+1)=x^{2}\Rightarrow
        x^{2}=\dfrac{y}{1-y}$. A real $x$ exists iff
        $\dfrac{y}{1-y}\geq0$, i.e. $0\leq y<1$: \;$W=[0,1)$.
  \item On $[-1,3]$ the square is not monotone: it falls to $0$ at $x=0$
        and rises to $9$ at $x=3$. Minimum $0$, maximum $9$:
        \;$W=[0,9]$.
  \item $f$ is strictly increasing, so the endpoints map to the endpoints:
        $f(1)=-3$ (excluded, as $1\notin D$) and $f(4)=3$ (included):
        \;$W=(-3,3]$.
  \item $(f\circ g)(x)=\sqrt{x-4}$ with $D=[4,\infty)$;
        $(g\circ f)(x)=\sqrt{x}-4$ with $D=[0,\infty)$. Same ingredients,
        different functions and different domains.
  \item $(f\circ g)(x)=\dfrac{1}{x+2}$, requiring $x+2\neq0$:
        \;$D=\R\setminus\{-2\}$.
  \item $(f\circ g)(x)=\ln(9-x^{2})$, requiring $9-x^{2}>0\iff|x|<3$:
        \;$D=(-3,3)$.
  \item $(f\circ g)(x)=\sqrt{1-x^{2}}$, requiring $1-x^{2}\geq0$:
        \;$D=[-1,1]$. On it $1-x^{2}$ runs through $[0,1]$, so
        $W=[0,1]$ (the upper unit semicircle).
  \item (a) even (only even powers); (b) odd (only odd powers);
        (c) neither: $f(-x)=x^{2}-x$ equals neither $f(x)$ nor $-f(x)$;
        (d) odd: $1/(-x)=-1/x$.
  \item $x^{3}$ and $\mathrm{e}^{x}$ are injective (both strictly
        increasing). $x^{2}$ and $|x|$ are not: $f(-1)=f(1)$ in both cases.
  \item $D=\R$. First piece: $x$ runs through $(-\infty,1)$, so $x^{2}$
        attains every value in $[0,\infty)$. Second piece: $2x-1$ on
        $[1,\infty)$ gives $[1,\infty)$. Union: \;$W=[0,\infty)$.
  \item No. The vertical line $x=0$ meets the circle at $(0,5)$ and
        $(0,-5)$, so the assignment is not unique. Solving gives two
        functions, $y=\sqrt{25-x^{2}}$ (upper semicircle) and
        $y=-\sqrt{25-x^{2}}$ (lower), each with domain $[-5,5]$ and ranges
        $[0,5]$ and $[-5,0]$.
  \item $D=[-2,\infty)$: the segment covers $x\in[-2,0]$, the ray covers
        $x\in[0,\infty)$. On the segment $y$ climbs from $1$ to $5$, on the
        ray $y$ stays at $5$: \;$W=[1,5]$.
\end{solutions}

\section{Additional Problems}

\begin{problems}
  \item $f(x)=\dfrac{3x}{2x+7}$: maximal domain.
  \item $f(x)=\sqrt{2x+8}$: maximal domain and range.
  \item $f(x)=\ln(5-2x)$: maximal domain.
  \item $f(x)=\dfrac{1}{\ln x}$: maximal domain.
  \item $f(x)=\dfrac{x-1}{\sqrt{x^{2}-4}}$: maximal domain.
  \item $f(x)=\sqrt{\dfrac{x-1}{x+2}}$: maximal domain.
  \item $f(x)=\ln\!\bigl(x^{2}+1\bigr)$: maximal domain and range.
  \item $f(x)=\mathrm{e}^{-x^{2}}$: maximal domain and range.
  \item $f(x)=\dfrac{1}{x^{2}-5x+6}$: maximal domain.
  \item $f(x)=\sqrt{4-x}+\dfrac{1}{x+1}$: maximal domain.
  \item $f(x)=(x-3)^{-1/2}$: maximal domain.
  \item $f(x)=\dfrac{|x|}{x}$: maximal domain and range.
  \item $f(x)=3-\dfrac{2}{x+1}$: range.
  \item $f(x)=\dfrac{x-1}{x+1}$: range.
  \item $f(x)=4-\sqrt{x+9}$: range.
  \item $f(x)=x^{2}-4x+7$ on $D=[0,3]$: range.
  \item $f(x)=\mathrm{e}^{2x}+1$ on $D=[0,\ln 2]$: range.
  \item $f(x)=\ln x$ on $D=[1,\mathrm{e}^{3}]$: range.
  \item With $f(u)=u^{2}$ and $g(x)=\sqrt{x-3}$: give $(f\circ g)(x)$ and
        its domain, and explain why the simplified formula alone would give
        the wrong answer.
  \item With $f(u)=\dfrac{1}{u-1}$ and $g(x)=\mathrm{e}^{x}$: give
        $(f\circ g)(x)$ and its domain.
  \item Revenue is $R(q)=q\,(200-4q)$. State the economic domain on which
        both $q\ge 0$ and the price stay non-negative, and the range of $R$
        on it.
  \item A production function is $Q(L)=10\sqrt{L}$ with $L\in[0,100]$.
        Give the range, and the average product $Q(L)/L$ with its domain.
  \item Classify as even, odd, or neither: (a) $\mathrm{e}^{-x^{2}}$,
        \;(b) $x\,|x|$, \;(c) $\ln x$, \;(d) $\dfrac{x^{3}}{x^{2}+1}$.
  \item Show that $f(x)=\dfrac{2x-1}{x+4}$ is injective on its domain by
        solving $f(x_{1})=f(x_{2})$.
  \item Give domain and range of
        $f(x)=\begin{cases} -x, & x<0,\\ \sqrt{x}, & 0\le x\le 4,\\
        2, & x>4.\end{cases}$
\end{problems}

\section{Additional Solutions}

\begin{solutions}
  \item $2x+7\neq0\iff x\neq-\tfrac{7}{2}$:
        \;$D=\R\setminus\{-\tfrac{7}{2}\}$.
  \item $2x+8\geq0\iff x\geq-4$: \;$D=[-4,\infty)$, and since $2x+8$ then
        runs through $[0,\infty)$, \;$W=[0,\infty)$.
  \item $5-2x>0\iff x<\tfrac{5}{2}$: \;$D=(-\infty,\tfrac{5}{2})$.
  \item $x>0$ for the logarithm and $\ln x\neq0\iff x\neq1$:
        \;$D=(0,1)\cup(1,\infty)$.
  \item The root is in the denominator, so $x^{2}-4>0\iff|x|>2$:
        \;$D=(-\infty,-2)\cup(2,\infty)$.
  \item Need $\dfrac{x-1}{x+2}\geq0$ with $x\neq-2$. The quotient is
        non-negative outside the interval between the two critical points
        $-2$ and $1$: \;$D=(-\infty,-2)\cup[1,\infty)$. (At $x=1$ the
        quotient is $0$, which is allowed; $x=-2$ is not.)
  \item $x^{2}+1\geq1>0$ always, so $D=\R$; the argument runs through
        $[1,\infty)$, hence \;$W=[0,\infty)$.
  \item $D=\R$. The exponent $-x^{2}$ runs through $(-\infty,0]$, and
        $\mathrm{e}^{u}$ maps that to $(0,1]$: \;$W=(0,1]$ (the bell curve,
        maximum $1$ at $x=0$).
  \item $x^{2}-5x+6=(x-2)(x-3)$: \;$D=\R\setminus\{2,3\}$.
  \item $4-x\geq0\iff x\leq4$, and $x\neq-1$:
        \;$D=(-\infty,-1)\cup(-1,4]$.
  \item A negative non-integer exponent means a root in a denominator:
        $x-3>0$, \;$D=(3,\infty)$.
  \item $D=\R\setminus\{0\}$. The quotient equals $1$ for $x>0$ and $-1$
        for $x<0$: \;$W=\{-1,1\}$, a range consisting of two points.
  \item $2/(x+1)$ attains every value except $0$; negating and adding $3$
        keeps exactly one value out of reach: \;$W=\R\setminus\{3\}$.
  \item $y=\dfrac{x-1}{x+1}\Rightarrow y(x+1)=x-1\Rightarrow
        x(y-1)=-1-y\Rightarrow x=\dfrac{-(1+y)}{y-1}=\dfrac{1+y}{1-y}$,
        which exists for $y\neq1$: \;$W=\R\setminus\{1\}$.
  \item $D=[-9,\infty)$; $\sqrt{x+9}$ runs through $[0,\infty)$, and
        $4-\sqrt{x+9}$ therefore through \;$W=(-\infty,4]$.
  \item $x^{2}-4x+7=(x-2)^{2}+3$. On $[0,3]$ the vertex $x=2$ lies inside,
        giving the minimum $3$; the endpoints give $f(0)=7$ and $f(3)=4$,
        so the maximum is $7$: \;$W=[3,7]$.
  \item Strictly increasing: $f(0)=2$ and
        $f(\ln 2)=\mathrm{e}^{2\ln 2}+1=4+1=5$: \;$W=[2,5]$.
  \item Strictly increasing: $\ln 1=0$ and $\ln \mathrm{e}^{3}=3$:
        \;$W=[0,3]$.
  \item $(f\circ g)(x)=\bigl(\sqrt{x-3}\bigr)^{2}=x-3$, but only where the
        inner function is defined: \;$D=[3,\infty)$. The simplified formula
        $x-3$ suggests $D=\R$; the domain is inherited from the inner
        function and is not restored by cancellation.
  \item $(f\circ g)(x)=\dfrac{1}{\mathrm{e}^{x}-1}$, requiring
        $\mathrm{e}^{x}\neq1\iff x\neq0$: \;$D=\R\setminus\{0\}$.
  \item Price $200-4q\ge0\iff q\le50$, so the domain is $[0,50]$.
        $R(q)=200q-4q^{2}=-4(q-25)^{2}+2500$, a downward parabola with
        vertex at $q=25$; on $[0,50]$ it runs from $0$ up to $2500$ and
        back to $0$: \;$W=[0,2500]$.
  \item $\sqrt{L}$ runs through $[0,10]$ on $[0,100]$, so
        $W=[0,100]$. The average product is
        $Q(L)/L=10/\sqrt{L}$, which needs $L>0$: domain $(0,100]$.
  \item (a) even: $(-x)^{2}=x^{2}$; (b) odd:
        $(-x)|-x|=-x|x|$; (c) neither, and in fact the test cannot even be
        applied: the domain $(0,\infty)$ is not symmetric about $0$;
        (d) odd: numerator odd, denominator even.
  \item $\dfrac{2x_{1}-1}{x_{1}+4}=\dfrac{2x_{2}-1}{x_{2}+4}$
        $\Rightarrow(2x_{1}-1)(x_{2}+4)=(2x_{2}-1)(x_{1}+4)$
        $\Rightarrow 2x_{1}x_{2}+8x_{1}-x_{2}-4
        =2x_{1}x_{2}+8x_{2}-x_{1}-4$
        $\Rightarrow 9x_{1}=9x_{2}\Rightarrow x_{1}=x_{2}$. Injective on
        $\R\setminus\{-4\}$.
  \item $D=\R$. First piece: $-x$ on $(-\infty,0)$ gives $(0,\infty)$.
        Second: $\sqrt{x}$ on $[0,4]$ gives $[0,2]$. Third: the constant
        $2$. The union is \;$W=[0,\infty)$.
\end{solutions}

\chapter{The Standard Functions and Their Graphs}

A first course in mathematics for economists assumes that the graphs of the
usual functions are known by sight: which way a cubic runs off at its ends,
where a rational function has a pole, that $\ln x$ never turns around, that
an isoquant is a hyperbola. That knowledge is not decoration either. Every
optimisation problem later on begins with a picture, every limit is checked
against one, and a wrong sketch produces a wrong answer faster than any
arithmetic slip.

This chapter is the catalogue, and the drill that makes it stick. First the
standard shapes; then the features that can be read straight off a formula
without any calculus: intercepts, roots, symmetry, poles and asymptotes,
and the behaviour at the far ends of the axis. The last block leaves the
world of $y=f(x)$ and sketches curves given by an \emph{equation} in $x$
and $y$: circles, ellipses, hyperbolas, budget lines and indifference
curves.

\section{Core Rules}

\begin{rulebox}{The power functions}
Whole-number powers alternate: even powers are symmetric about the
$y$-axis and non-negative, odd powers are symmetric about the origin and
run from $-\infty$ to $+\infty$. Negative powers have a pole at $0$; the
root $\sqrt{x}$ is the upper half of a sideways parabola.
\end{rulebox}

\begin{gallery}
\begin{tikzpicture}
  \begin{axis}[fnaxis, xmin=-2.6, xmax=2.6, ymin=-2.6, ymax=2.6,
      xtick={-2,2}, ytick={-2,2}, title={$y=x$}]
    \addplot[domain=-2.4:2.4]{x};
  \end{axis}
\end{tikzpicture}\hfill
\begin{tikzpicture}
  \begin{axis}[fnaxis, xmin=-2.6, xmax=2.6, ymin=-1.0, ymax=4.4,
      xtick={-2,2}, ytick={2,4}, title={$y=x^{2}$}]
    \addplot[domain=-2.1:2.1]{x^2};
  \end{axis}
\end{tikzpicture}\hfill
\begin{tikzpicture}
  \begin{axis}[fnaxis, xmin=-2.6, xmax=2.6, ymin=-4.4, ymax=4.4,
      xtick={-2,2}, ytick={-4,4}, title={$y=x^{3}$}]
    \addplot[domain=-1.62:1.62]{x^3};
  \end{axis}
\end{tikzpicture}
\end{gallery}

\begin{gallery}
\begin{tikzpicture}
  \begin{axis}[fnaxis, xmin=-0.6, xmax=4.6, ymin=-0.8, ymax=2.6,
      xtick={2,4}, ytick={1,2}, title={$y=\sqrt{x}$}]
    \addplot[domain=0:4.4]{sqrt(x)};
  \end{axis}
\end{tikzpicture}\hfill
\begin{tikzpicture}
  \begin{axis}[fnaxis, xmin=-3.4, xmax=3.4, ymin=-3.4, ymax=3.4,
      xtick={-2,2}, ytick={-2,2}, title={$y=1/x$}]
    \addplot[domain=-3.2:-0.31]{1/x};
    \addplot[domain=0.31:3.2]{1/x};
  \end{axis}
\end{tikzpicture}\hfill
\begin{tikzpicture}
  \begin{axis}[fnaxis, xmin=-3.4, xmax=3.4, ymin=-0.6, ymax=4.4,
      xtick={-2,2}, ytick={2,4}, title={$y=1/x^{2}$}]
    \addplot[domain=-3.2:-0.48]{1/x^2};
    \addplot[domain=0.48:3.2]{1/x^2};
  \end{axis}
\end{tikzpicture}
\end{gallery}

\begin{rulebox}{Exponential, logarithm, absolute value}
$\mathrm{e}^{x}$ is positive everywhere, strictly increasing, and flattens
onto the asymptote $y=0$ as $x\to-\infty$; it passes through $(0,1)$.
$\ln x$ is its mirror image at the line $y=x$: defined only for $x>0$,
strictly increasing, passing through $(1,0)$, with the $y$-axis as a
vertical asymptote. A general exponential $a^{x}$ rises for $a>1$ and falls
for $0<a<1$, always through $(0,1)$; note $\bigl(\tfrac12\bigr)^{x}
=2^{-x}$, so the decreasing case is the increasing one reflected in the
$y$-axis. The absolute value $|x|$ is a V with its vertex at the origin.
\end{rulebox}

\begin{gallery}
\begin{tikzpicture}
  \begin{axis}[fnaxis, xmin=-3.4, xmax=2.4, ymin=-0.8, ymax=4.4,
      xtick={-2,2}, ytick={1,3}, title={$y=\mathrm{e}^{x}$}]
    \addplot[domain=-3.2:1.45]{exp(x)};
  \end{axis}
\end{tikzpicture}\hfill
\begin{tikzpicture}
  \begin{axis}[fnaxis, xmin=-0.8, xmax=5.4, ymin=-3.4, ymax=2.4,
      xtick={1,3,5}, ytick={-2,1}, title={$y=\ln x$}]
    \addplot[domain=0.045:5.2]{ln(x)};
  \end{axis}
\end{tikzpicture}\hfill
\begin{tikzpicture}
  \begin{axis}[fnaxis, xmin=-3.4, xmax=3.4, ymin=-0.8, ymax=3.4,
      xtick={-2,2}, ytick={1,3}, title={$y=|x|$}]
    \addplot[domain=-3:3]{abs(x)};
  \end{axis}
\end{tikzpicture}
\end{gallery}

\begin{rulebox}{What can be read off a formula without calculus}
\begin{itemize}[leftmargin=1.6em,itemsep=0.3ex]
  \item \textbf{$y$-intercept:} evaluate $f(0)$ (if $0\in D$).
  \item \textbf{Roots} \de{(Nullstellen)}: solve $f(x)=0$. For a product,
        set each factor to zero. A factor $(x-a)^{k}$ with $k$ odd makes
        the graph \emph{cross} the axis at $a$; with $k$ even it only
        \emph{touches} it.
  \item \textbf{Sign:} between consecutive roots and poles the sign cannot
        change. Pick one test point per interval; that fixes the whole
        picture.
  \item \textbf{Symmetry:} $f(-x)=f(x)$ gives mirror symmetry in the
        $y$-axis, $f(-x)=-f(x)$ point symmetry in the origin.
  \item \textbf{End behaviour of a polynomial:} only the leading term
        $a_{n}x^{n}$ matters as $x\to\pm\infty$. Even $n$: both ends go the
        same way (up if $a_{n}>0$). Odd $n$: the ends go opposite ways
        (down--left, up--right if $a_{n}>0$).
\end{itemize}
\end{rulebox}

\begin{rulebox}{Rational functions: poles and asymptotes}
Write $f(x)=\dfrac{p(x)}{q(x)}$ with $p,q$ polynomials and
\textbf{cancel common factors first}.
\begin{itemize}[leftmargin=1.6em,itemsep=0.3ex]
  \item A zero of $q$ that does \emph{not} cancel is a \textbf{pole}
        \de{(Polstelle)}: the vertical line $x=a$ is a
        \textbf{vertical asymptote}, and $|f|\to\infty$ there.
  \item A zero of $q$ that \emph{does} cancel is a \textbf{hole}: the graph
        is that of the reduced function, minus the single point $x=a$.
  \item \textbf{Horizontal asymptote}, by comparing degrees:
        $\deg p<\deg q\Rightarrow y=0$; \;$\deg p=\deg q\Rightarrow
        y=\dfrac{\text{leading coeff. of }p}{\text{leading coeff. of }q}$;
        \;$\deg p>\deg q\Rightarrow$ none.
  \item If $\deg p=\deg q+1$, polynomial division gives
        $f(x)=(\text{linear})+\dfrac{\text{remainder}}{q(x)}$, and the
        linear part is an \textbf{oblique asymptote}
        \de{(schiefe Asymptote)}, because the remainder term dies out.
\end{itemize}
\end{rulebox}

\begin{gallery}
\begin{tikzpicture}
  \begin{axis}[fnaxis, xmin=-2.4, xmax=6.4, ymin=-4.4, ymax=6.4,
      xtick={2,4,6}, ytick={-2,2,4}, title={$y=\dfrac{2x+1}{x-3}$}]
    \addplot[domain=-2.2:2.55]{(2*x+1)/(x-3)};
    \addplot[domain=3.42:6.2]{(2*x+1)/(x-3)};
    \addplot[gray, dashed, thin, forget plot, samples=2, domain=-2.2:6.2]{2};
    \addplot[gray, dashed, thin, forget plot, samples=2] coordinates {(3,-4.2)(3,6.2)};
  \end{axis}
\end{tikzpicture}\hfill
\begin{tikzpicture}
  \begin{axis}[fnaxis, xmin=-4.4, xmax=4.4, ymin=-6.4, ymax=6.4,
      xtick={-2,2}, ytick={-4,4}, title={$y=\dfrac{x^{2}+1}{x}$}]
    \addplot[domain=-4.2:-0.17]{(x^2+1)/x};
    \addplot[domain=0.17:4.2]{(x^2+1)/x};
    \addplot[gray, dashed, thin, forget plot, samples=2, domain=-4.2:4.2]{x};
  \end{axis}
\end{tikzpicture}\hfill
\begin{tikzpicture}
  \begin{axis}[fnaxis, xmin=-3, xmax=42, ymin=-3, ymax=32,
      xtick={10,20,30,40}, ytick={5,15,25},
      title={$AC(q)=\dfrac{100+5q}{q}$}]
    \addplot[domain=3.9:40]{(100+5*x)/x};
    \addplot[gray, dashed, thin, forget plot, samples=2, domain=-2.8:40]{5};
  \end{axis}
\end{tikzpicture}
\gcap{Horizontal asymptote $y=2$ and pole $x=3$ \quad\textbullet\quad
oblique asymptote $y=x$ \quad\textbullet\quad average cost falling towards
the asymptote $AC=5$.}
\end{gallery}

\begin{rulebox}{Curves given by an equation in two variables}
Not every curve is a graph $y=f(x)$. The standard ones:
\begin{itemize}[leftmargin=1.6em,itemsep=0.3ex]
  \item \textbf{Line:} $ax+by=c$. Sketch it from the two intercepts
        $x=c/a$ and $y=c/b$; solved for $y$ the slope is $-a/b$. This is
        the budget line $p_{x}x+p_{y}y=M$, with slope $-p_{x}/p_{y}$.
  \item \textbf{Circle:} $(x-x_{0})^{2}+(y-y_{0})^{2}=r^{2}$, centre
        $(x_{0},y_{0})$, radius $r$. A general
        $x^{2}+y^{2}+Dx+Ey+F=0$ is brought to this form by
        \textbf{completing the square} in $x$ and in $y$.
  \item \textbf{Ellipse:} $\dfrac{x^{2}}{a^{2}}+\dfrac{y^{2}}{b^{2}}=1$,
        crossing the axes at $(\pm a,0)$ and $(0,\pm b)$.
  \item \textbf{Hyperbola:} $xy=c$ with $c>0$ lies in the first and third
        quadrant with both axes as asymptotes. Solved for $y$ it is
        $y=c/x$. Every Cobb--Douglas isoquant $K^{\alpha}L^{\beta}=\bar Q$
        and every indifference curve $xy=\bar U$ has this shape.
\end{itemize}
To sketch, either solve for $y$, remembering that a square root splits the
curve into \textbf{two branches} $y=\pm\sqrt{\cdots}$, or read off centre,
radius and intercepts directly.
\end{rulebox}

\section{Worked Example}

\begin{worked}{Sketch $f(x)=\dfrac{2x+1}{x-3}$ and $g(x)=(x-1)^{2}(x+3)$ from the formulas alone, and say what $h(x)=\dfrac{x^{2}-1}{x-1}$ really is}
\wstep{Step 1. The rational function: find the poles.}
The denominator of $f$ vanishes at $x=3$ and the numerator does not
($2\cdot3+1=7\neq0$), so $x=3$ is a genuine \textbf{vertical asymptote}.

\wstep{Step 2. Find the horizontal asymptote by comparing degrees.}
Numerator and denominator both have degree $1$, so the limit at $\pm\infty$
is the ratio of leading coefficients:
\[
  \lim_{x\to\pm\infty}\frac{2x+1}{x-3}=\frac{2}{1}=2 ,
\]
a horizontal asymptote $y=2$. Equivalently, dividing out gives
$f(x)=2+\dfrac{7}{x-3}$, which shows the whole graph at once: the hyperbola
$7/x$ shifted right $3$ and up $2$.

\wstep{Step 3. Intercepts and signs.}
$f(x)=0$ when $2x+1=0$, i.e.\ $x=-\tfrac12$; and $f(0)=\tfrac{1}{-3}=-\tfrac13$.
The sign changes at $x=-\tfrac12$ (zero) and at $x=3$ (pole), giving
$+,-,+$ on the three intervals.

\wstep{Step 4. Approach the pole from both sides.}
As $x\to3^{+}$, $x-3\to0^{+}$ and the numerator is near $7>0$, so
$f\to+\infty$; as $x\to3^{-}$, $f\to-\infty$. Now the sketch is determined:
two branches, one below-left rising to the asymptote $y=2$ from below, one
above-right falling to it from above.
\wtrap{Test the sign on \emph{each side} of a pole rather than assuming
symmetry. A pole of odd order flips the sign, one of even order does not,
and drawing both branches on the same side of the horizontal asymptote is
the usual consequence of guessing.}

\wstep{Step 5. The polynomial: read the roots with their multiplicities.}
$g(x)=(x-1)^{2}(x+3)$ has a root of multiplicity $2$ at $x=1$ and a simple
root at $x=-3$. Even multiplicity means the graph \textbf{touches} the axis
and turns back; odd multiplicity means it \textbf{crosses}. So the curve
crosses at $-3$ and bounces off the axis at $1$.

\wstep{Step 6. End behaviour is decided by the leading term alone.}
Expanding just the leading term: $x^{2}\cdot x=x^{3}$, degree $3$ with a
positive coefficient, so $g\to-\infty$ as $x\to-\infty$ and $g\to+\infty$ as
$x\to+\infty$. With $g(0)=1\cdot3=3>0$, the shape is now forced: up through
$-3$, over a local maximum, down to touch the axis at $1$, then away
upward.
\wtrap{Degree and leading sign come from multiplying the leading terms of
every factor, counted with multiplicity. $(x-1)^{2}(x+3)$ is cubic, not
quadratic; miscounting the degree reverses one of the two ends.}

\wstep{Step 7. The impostor: $h(x)=\dfrac{x^{2}-1}{x-1}$.}
Here the numerator \emph{also} vanishes at $x=1$. Factor before concluding
anything:
\[
  h(x)=\frac{(x-1)(x+1)}{x-1}=x+1\qquad\text{for }x\neq1 .
\]
So $h$ has no asymptote at all. Its graph is the straight line $y=x+1$ with a
single point punched out at $(1,2)$ --- a \textbf{removable discontinuity},
a hole, not a pole.
\wtrap{A zero denominator does not by itself make an asymptote. Always
factor numerator and denominator first: a shared factor gives a hole, an
unshared one gives a pole. This is the same distinction that decides whether
a limit in Volume~4 exists or diverges.}
\end{worked}

\section{Practice Problems}

For each polynomial in problems 1--8, give the roots (with multiplicity
where it matters), the $y$-intercept, the symmetry, and the end behaviour
as $x\to\pm\infty$.

\begin{problems}
  \item $f(x)=x^{3}-4x$
  \item $f(x)=x^{4}-5x^{2}+4$
  \item $f(x)=-2x^{3}+6x$
  \item $f(x)=(x-1)^{2}(x+3)$
  \item $f(x)=x^{2}-6x+8$
  \item $f(x)=-x^{2}+4x-4$
  \item $f(x)=x^{5}-x$
  \item A polynomial satisfies $f(x)\to-\infty$ both as $x\to+\infty$ and
        as $x\to-\infty$. What does that say about its degree and its
        leading coefficient?
\end{problems}

In problems 9--20, give all vertical asymptotes, holes, and horizontal or
oblique asymptotes.

\begin{problems}[resume]
  \item $f(x)=\dfrac{1}{x-3}$
  \item $f(x)=\dfrac{2x+1}{x-3}$
  \item $f(x)=\dfrac{x^{2}-1}{x-1}$
  \item $f(x)=\dfrac{3x}{x^{2}-4}$
  \item $f(x)=\dfrac{x^{2}+1}{x}$
  \item $f(x)=\dfrac{x^{2}-4}{x^{2}-9}$
  \item $f(x)=\dfrac{x^{2}+3x+2}{x+1}$
  \item $f(x)=\dfrac{x^{2}-x-6}{x-3}$
  \item $f(x)=\dfrac{2x^{2}+3}{x-1}$
  \item $f(x)=\dfrac{5}{x^{2}+1}$
  \item $AC(q)=\dfrac{100+5q}{q}$ \;(average cost, $q>0$)
  \item $f(x)=\dfrac{x+3}{x^{2}+x-6}$
\end{problems}

\begin{problems}[resume]
  \item For $y=2^{x}$: give the $y$-intercept, the asymptote, the
        monotonicity and the range.
  \item Explain why the graph of $y=\bigl(\tfrac12\bigr)^{x}$ is the graph
        of $y=2^{x}$ reflected in the $y$-axis.
  \item For $y=\ln x$: give the $x$-intercept, the asymptote, the domain
        and the range.
  \item Sketch $y=\mathrm{e}^{-x}$ and state its asymptote and range.
  \item Both $\ln x$ and $\log_{10}x$ pass through $(1,0)$ and rise. What
        exactly distinguishes their graphs?
  \item The adjustment path $y=1-\mathrm{e}^{-x}$, $x\ge0$, is used for
        saturation effects. Give $y(0)$, the asymptote, and say in words
        what the curve does.
\end{problems}

\begin{problems}[resume]
  \item Match each graph (A)--(F) below to one of the formulas
        $\;\ln x$, \;$|x|$, \;$\sqrt{x}$, \;$x^{3}-x$,
        \;$\mathrm{e}^{-x}$, \;$1/x^{2}$.
\end{problems}

\begin{gallery}
\begin{tikzpicture}
  \begin{axis}[fnaxis, xmin=-2.6, xmax=2.6, ymin=-2.2, ymax=2.6,
      xtick={-2,2}, ytick={-2,2}, title={(A)}]
    \addplot[domain=-1.85:1.85]{x^3-x};
  \end{axis}
\end{tikzpicture}\hfill
\begin{tikzpicture}
  \begin{axis}[fnaxis, xmin=-0.7, xmax=4.6, ymin=-0.8, ymax=2.6,
      xtick={2,4}, ytick={1,2}, title={(B)}]
    \addplot[domain=0:4.4]{sqrt(x)};
  \end{axis}
\end{tikzpicture}\hfill
\begin{tikzpicture}
  \begin{axis}[fnaxis, xmin=-3.4, xmax=3.4, ymin=-0.6, ymax=4.4,
      xtick={-2,2}, ytick={2,4}, title={(C)}]
    \addplot[domain=-3.2:-0.48]{1/x^2};
    \addplot[domain=0.48:3.2]{1/x^2};
  \end{axis}
\end{tikzpicture}
\end{gallery}

\begin{gallery}
\begin{tikzpicture}
  \begin{axis}[fnaxis, xmin=-2.4, xmax=3.4, ymin=-0.8, ymax=4.4,
      xtick={-2,2}, ytick={1,3}, title={(D)}]
    \addplot[domain=-1.45:3.2]{exp(-x)};
  \end{axis}
\end{tikzpicture}\hfill
\begin{tikzpicture}
  \begin{axis}[fnaxis, xmin=-3.4, xmax=3.4, ymin=-0.8, ymax=3.4,
      xtick={-2,2}, ytick={1,3}, title={(E)}]
    \addplot[domain=-3:3]{abs(x)};
  \end{axis}
\end{tikzpicture}\hfill
\begin{tikzpicture}
  \begin{axis}[fnaxis, xmin=-0.8, xmax=5.4, ymin=-3.4, ymax=2.4,
      xtick={1,3,5}, ytick={-2,1}, title={(F)}]
    \addplot[domain=0.045:5.2]{ln(x)};
  \end{axis}
\end{tikzpicture}
\end{gallery}

\begin{problems}[resume]
  \item Sketch $f(x)=x^{2}-4x+3$: give the vertex, both intercepts and the
        direction of opening.
  \item Sketch $f(x)=\dfrac{x+1}{x-1}$: give both asymptotes and both
        intercepts.
  \item Sketch $f(x)=\sqrt{x+4}$: give the domain, the range and the
        starting point of the curve.
  \item Sketch $f(x)=|x|-2$: give the vertex and the roots.
  \item Sketch $f(x)=x^{3}-3x^{2}$: give the roots with multiplicity and
        the end behaviour.
  \item Classify as even, odd, or neither: (a) $x^{4}+2$, \;(b) $x^{3}+x$,
        \;(c) $1/x^{2}$, \;(d) $\mathrm{e}^{x}$.
  \item Which of these have a horizontal asymptote, and what is it?
        (a) $\dfrac{3x+1}{x-2}$, \;(b) $\dfrac{x^{2}}{x+1}$,
        \;(c) $\mathrm{e}^{-x}$ as $x\to+\infty$, \;(d) $\ln x$.
  \item On which intervals is $f(x)=\dfrac{(x-1)(x+2)}{x-3}$ positive?
  \item Describe the graph of $f(x)=\sqrt{x^{2}-1}$: domain, symmetry,
        range, and the two points where it meets the $x$-axis.
\end{problems}

In problems 37--45, identify the curve, and sketch it.

\begin{problems}[resume]
  \item $x^{2}+y^{2}=16$. Also solve for $y$ and give the two branches with
        their domains.
  \item $(x-2)^{2}+(y+1)^{2}=9$. Give centre and radius.
  \item $\dfrac{x^{2}}{9}+\dfrac{y^{2}}{4}=1$. Give the four axis
        intercepts.
  \item $xy=6$. Solve for $y$; in which quadrants does the curve lie?
  \item $y^{2}=x$. Is $y$ a function of $x$? Is $x$ a function of $y$?
  \item $2x+3y=12$. Give both intercepts and the slope.
  \item $x^{2}+y^{2}-4x+6y+9=0$. Complete the square; give centre and
        radius.
  \item The isoquant $K^{0.5}L^{0.5}=10$. Solve for $L$ and describe the
        shape for $K>0$.
  \item The budget line $4x+2y=40$ and the indifference curve $xy=50$.
        Give the slope of the budget line and $y$ as a function of $x$ on
        the indifference curve.
\end{problems}

\section{Solutions}

\begin{solutions}
  \item $x^{3}-4x=x(x^{2}-4)=x(x-2)(x+2)$: roots $-2,0,2$, all simple, all
        crossings; $f(0)=0$; odd (only odd powers); degree $3$ with
        positive leading coefficient, so $f\to-\infty$ as $x\to-\infty$ and
        $f\to+\infty$ as $x\to+\infty$.
  \item $x^{4}-5x^{2}+4=(x^{2}-1)(x^{2}-4)$: roots $\pm1,\pm2$;
        $f(0)=4$; even; degree $4$, positive leading coefficient, so both
        ends go to $+\infty$ (a W-shape).
  \item $-2x^{3}+6x=-2x(x^{2}-3)$: roots $0,\pm\sqrt3$; $f(0)=0$; odd;
        degree $3$ with \emph{negative} leading coefficient, so
        $f\to+\infty$ as $x\to-\infty$ and $f\to-\infty$ as $x\to+\infty$.
  \item Roots $x=1$ (double, the graph touches and turns back) and $x=-3$
        (simple, a crossing); $f(0)=(-1)^{2}\cdot3=3$; neither even nor
        odd; degree $3$, leading coefficient $+1$: down on the left, up on
        the right.
  \item $(x-2)(x-4)$: roots $2$ and $4$; $f(0)=8$; vertex at
        $x=3$, $f(3)=-1$; neither even nor odd; parabola opening upwards,
        both ends $+\infty$.
  \item $-(x-2)^{2}$: a double root at $x=2$ (the vertex touches the axis);
        $f(0)=-4$; opens downwards, both ends $-\infty$.
  \item $x^{5}-x=x(x^{4}-1)=x(x^{2}-1)(x^{2}+1)$: roots $-1,0,1$;
        $f(0)=0$; odd; degree $5$, positive leading coefficient: down on
        the left, up on the right.
  \item Both ends behave alike, so the degree is \textbf{even}; both go
        \emph{down}, so the leading coefficient is \textbf{negative}.
  \item Vertical asymptote $x=3$; $\deg p=0<1=\deg q$, so the horizontal
        asymptote is $y=0$. No holes.
  \item Vertical asymptote $x=3$; equal degrees, so
        $y=\tfrac{2}{1}=2$ is the horizontal asymptote.
  \item $\dfrac{(x-1)(x+1)}{x-1}=x+1$ for $x\neq1$: no pole but a
        \textbf{hole} at $(1,2)$. The graph is the line $y=x+1$ with that
        one point punched out; no horizontal asymptote.
  \item $x^{2}-4=(x-2)(x+2)$: poles at $x=\pm2$; $\deg p=1<2=\deg q$, so
        $y=0$.
  \item Pole at $x=0$. Degrees differ by one, so divide:
        $\dfrac{x^{2}+1}{x}=x+\dfrac{1}{x}$, giving the oblique asymptote
        $y=x$.
  \item Poles at $x=\pm3$; equal degrees with leading coefficients $1$ and
        $1$, so the horizontal asymptote is $y=1$.
  \item $\dfrac{(x+1)(x+2)}{x+1}=x+2$ for $x\neq-1$: a hole at $(-1,1)$,
        no pole, no horizontal asymptote.
  \item $\dfrac{(x-3)(x+2)}{x-3}=x+2$ for $x\neq3$: a hole at $(3,5)$.
  \item Pole at $x=1$. Division: $2x^{2}+3=(x-1)(2x+2)+5$, so
        $f(x)=2x+2+\dfrac{5}{x-1}$ and the oblique asymptote is $y=2x+2$.
  \item $x^{2}+1$ never vanishes: no poles at all. $\deg p=0<2$, so
        $y=0$; the graph is a bell-shaped hump with maximum $5$ at $x=0$,
        even.
  \item $AC(q)=\dfrac{100}{q}+5$: pole at $q=0$ and horizontal asymptote
        $AC=5$. Average cost falls towards the marginal cost $5$ as the
        fixed cost $100$ is spread over more units.
  \item $x^{2}+x-6=(x+3)(x-2)$, so
        $f(x)=\dfrac{x+3}{(x+3)(x-2)}=\dfrac{1}{x-2}$ for $x\neq-3$:
        a hole at $\bigl(-3,-\tfrac15\bigr)$, a pole at $x=2$, and the
        horizontal asymptote $y=0$.
  \item $y$-intercept $2^{0}=1$; asymptote $y=0$ (approached as
        $x\to-\infty$); strictly increasing; range $(0,\infty)$.
  \item Because $\bigl(\tfrac12\bigr)^{x}=(2^{-1})^{x}=2^{-x}$: the
        argument $x$ has been replaced by $-x$, and replacing $x$ by $-x$
        mirrors any graph in the $y$-axis.
  \item $x$-intercept $(1,0)$; the $y$-axis $x=0$ is a vertical asymptote
        (as $x\to0^{+}$, $\ln x\to-\infty$); $D=(0,\infty)$, $W=\R$.
  \item The decreasing exponential: it falls from $+\infty$ through
        $(0,1)$ towards the asymptote $y=0$ as $x\to+\infty$;
        $W=(0,\infty)$. It is $\mathrm{e}^{x}$ mirrored in the $y$-axis.
  \item Only the vertical scale: $\log_{10}x=\dfrac{\ln x}{\ln 10}$, so the
        common logarithm is the natural one compressed by the constant
        factor $\ln 10\approx2.303$. Same domain, same asymptote, same
        intercept, same shape.
  \item $y(0)=1-1=0$; as $x\to\infty$, $\mathrm{e}^{-x}\to0$, so the
        asymptote is $y=1$. The curve starts at the origin, rises steeply
        at first and then ever more slowly, approaching the saturation
        level $1$ without ever reaching it.
  \item (A) $x^{3}-x$ \;(odd, three roots); \;(B) $\sqrt{x}$;
        \;(C) $1/x^{2}$ \;(even, positive, pole at $0$);
        \;(D) $\mathrm{e}^{-x}$; \;(E) $|x|$; \;(F) $\ln x$.
  \item $x^{2}-4x+3=(x-1)(x-3)=(x-2)^{2}-1$: opens upwards, vertex
        $(2,-1)$, roots $1$ and $3$, $y$-intercept $3$.
  \item $\dfrac{x+1}{x-1}$: pole $x=1$, horizontal asymptote $y=1$ (equal
        degrees); root $x=-1$; $y$-intercept $f(0)=-1$. Two branches, one
        in each half-plane of the pole.
  \item $D=[-4,\infty)$, $W=[0,\infty)$; the curve starts at $(-4,0)$ and
        rises, ever more slowly: the $\sqrt{\;}$ shape moved four units to
        the left.
  \item A V with vertex $(0,-2)$; roots where $|x|=2$, i.e. $x=\pm2$;
        $y$-intercept $-2$.
  \item $x^{3}-3x^{2}=x^{2}(x-3)$: a double root at $x=0$ (touch) and a
        simple root at $x=3$ (crossing); degree $3$, positive leading
        coefficient: down on the left, up on the right.
  \item (a) even; (b) odd; (c) even, since $1/(-x)^{2}=1/x^{2}$;
        (d) neither.
  \item (a) yes, $y=3$ (equal degrees); (b) no: the numerator degree
        exceeds the denominator degree by one, so there is an oblique
        asymptote $y=x-1$ instead; (c) yes, $y=0$; (d) no: $\ln x$ grows
        without bound, however slowly.
  \item Critical points $-2,1,3$. Test one point per interval:
        $x=-3$ gives $\frac{(-4)(-1)}{-6}<0$; $x=0$ gives
        $\frac{(-1)(2)}{-3}>0$; $x=2$ gives $\frac{(1)(4)}{-1}<0$;
        $x=4$ gives $\frac{(3)(6)}{1}>0$. Hence $f>0$ on
        $(-2,1)\cup(3,\infty)$.
  \item $D=(-\infty,-1]\cup[1,\infty)$ (from $x^{2}-1\ge0$); even, so
        symmetric about the $y$-axis; $W=[0,\infty)$; it meets the axis at
        $(-1,0)$ and $(1,0)$ and rises from there, approaching the lines
        $y=\pm x$ from below for large $|x|$.
  \item Circle, centre $(0,0)$, radius $4$. Solving:
        $y=\pm\sqrt{16-x^{2}}$, two branches, each with domain $[-4,4]$;
        the plus sign is the upper semicircle, the minus sign the lower.
  \item Circle with centre $(2,-1)$ and radius $3$.
  \item Ellipse with $a=3$ and $b=2$: it crosses the axes at
        $(\pm3,0)$ and $(0,\pm2)$.
  \item $y=6/x$, a hyperbola with both axes as asymptotes. Since the
        product is positive, $x$ and $y$ share a sign: the curve lies in
        the first and third quadrants.
  \item A parabola opening to the right. $y$ is \emph{not} a function of
        $x$: the line $x=4$ meets the curve at $y=\pm2$. But $x=y^{2}$
        \emph{is} a function of $y$.
  \item Intercepts $(6,0)$ and $(0,4)$; solved for $y$,
        $y=4-\tfrac23 x$, so the slope is $-\tfrac23$.
  \item $(x^{2}-4x+4)+(y^{2}+6y+9)=-9+4+9$, i.e.
        $(x-2)^{2}+(y+3)^{2}=4$: centre $(2,-3)$, radius $2$.
  \item $K^{0.5}L^{0.5}=10\Rightarrow (KL)^{0.5}=10\Rightarrow KL=100
        \Rightarrow L=\dfrac{100}{K}$. A hyperbola in the first quadrant,
        falling and convex, with both axes as asymptotes: the familiar
        isoquant, along which more capital buys ever less labour saving.
  \item Budget line: $y=20-2x$, slope $-2$ (the price ratio
        $-p_{x}/p_{y}=-4/2$), intercepts $(10,0)$ and $(0,20)$.
        Indifference curve: $y=50/x$, a hyperbola in the first quadrant.
\end{solutions}

\section{Additional Problems}

\begin{problems}
  \item Asymptotes of $f(x)=\dfrac{4x-1}{2x+6}$.
  \item Asymptotes and holes of $f(x)=\dfrac{x^{2}-9}{x+3}$.
  \item Asymptotes and symmetry of $f(x)=\dfrac{x^{3}}{x^{2}-1}$.
  \item Asymptotes and holes of $f(x)=\dfrac{x^{2}-2x}{x^{2}-4}$.
  \item Asymptotes of $f(x)=\dfrac{6}{x^{2}-x-2}$.
  \item Asymptotes of $f(x)=\dfrac{x-4}{x^{2}+4}$.
  \item Roots, symmetry and end behaviour of $f(x)=x^{4}-4x^{2}$.
  \item Roots, sign and end behaviour of $f(x)=(x+2)(x-1)(x-3)$.
  \item Symmetry and asymptote of $f(x)=\dfrac{x}{x^{2}+1}$.
  \item Where do $y=3^{x}$ and $y=3^{-x}$ intersect, and what asymptote do
        they share?
  \item The graphs of $y=\mathrm{e}^{x}$ and $y=\ln x$ are mirror images.
        In which line, and which points correspond to each other?
  \item Domain, range, symmetry and asymptotes of $y=1/x^{2}$.
  \item Domain, symmetry and shape of $y=\sqrt[3]{x}$.
  \item Vertex and roots of $y=|x-2|$.
  \item For the cost function $C(q)=q^{3}-6q^{2}+15q$, tabulate
        $C(0),C(1),\dots,C(5)$ and describe the shape of the curve.
  \item Tabulate $f(x)=x^{2}-2x-3$ for $x=-2,\dots,4$ and sketch it from
        the table; give vertex and roots.
  \item Identify $9x^{2}+4y^{2}=36$ and give its axis intercepts.
  \item Identify $x^{2}-y^{2}=9$ and solve it for $y$.
  \item What curve is $(x+1)^{2}+(y-2)^{2}=0$?
  \item Identify $x^{2}+y^{2}+2x-8=0$: centre and radius.
  \item Compare $y=\sqrt{9-x^{2}}$ with $x^{2}+y^{2}=9$: which is a
        function, and what are its domain and range?
  \item The isocost line $5K+2L=100$: give both intercepts and $L$ as a
        function of $K$.
  \item For the indifference curve $xy=36$, tabulate $y$ at
        $x=1,2,3,4,6,9,12$ and describe the shape.
  \item The demand curve $p=100/q$ has the property that revenue $pq$ is
        constant. Give that constant, the asymptotes, and the economic
        reading of each.
  \item A curve passes through the origin, rises throughout, is concave,
        and flattens onto the horizontal line $y=4$. Which of
        $y=4-\mathrm{e}^{-x}$, \;$y=4(1-\mathrm{e}^{-x})$,
        \;$y=4\mathrm{e}^{-x}$ fits, and why do the other two fail?
\end{problems}

\section{Additional Solutions}

\begin{solutions}
  \item Pole where $2x+6=0$, i.e. $x=-3$; equal degrees, so the horizontal
        asymptote is $y=\tfrac42=2$.
  \item $\dfrac{(x-3)(x+3)}{x+3}=x-3$ for $x\neq-3$: a hole at
        $(-3,-6)$, no pole, no horizontal asymptote (the graph is a line).
  \item Poles at $x=\pm1$. Division: $x^{3}=x(x^{2}-1)+x$, so
        $f(x)=x+\dfrac{x}{x^{2}-1}$ and the oblique asymptote is $y=x$.
        Odd: numerator odd, denominator even.
  \item $\dfrac{x(x-2)}{(x-2)(x+2)}=\dfrac{x}{x+2}$ for $x\neq2$: a hole at
        $\bigl(2,\tfrac12\bigr)$, a pole at $x=-2$, horizontal asymptote
        $y=1$.
  \item $x^{2}-x-2=(x-2)(x+1)$: poles at $x=2$ and $x=-1$; horizontal
        asymptote $y=0$.
  \item $x^{2}+4>0$ always: no poles. Numerator degree smaller, so $y=0$.
  \item $x^{4}-4x^{2}=x^{2}(x^{2}-4)=x^{2}(x-2)(x+2)$: double root at $0$
        (touch), simple roots at $\pm2$; even; degree $4$ with positive
        leading coefficient, so both ends rise to $+\infty$.
  \item Roots $-2,1,3$, all simple. Degree $3$, positive leading
        coefficient: down on the left, up on the right. Signs alternate
        between the roots: negative on $(-\infty,-2)$, positive on
        $(-2,1)$, negative on $(1,3)$, positive on $(3,\infty)$.
  \item Odd (odd numerator over even denominator); no poles since
        $x^{2}+1>0$; horizontal asymptote $y=0$ at both ends. The graph is
        a hump above the axis on the right and its mirror image below on
        the left.
  \item $3^{x}=3^{-x}\iff x=-x\iff x=0$: they meet at $(0,1)$. Both have
        the horizontal asymptote $y=0$, the first as $x\to-\infty$, the
        second as $x\to+\infty$; each is the other reflected in the
        $y$-axis.
  \item In the line $y=x$. The point $(a,\mathrm{e}^{a})$ on the
        exponential corresponds to $(\mathrm{e}^{a},a)$ on the logarithm;
        in particular $(0,1)\leftrightarrow(1,0)$ and
        $(1,\mathrm{e})\leftrightarrow(\mathrm{e},1)$. The two functions
        are inverse to each other.
  \item $D=\R\setminus\{0\}$, $W=(0,\infty)$; even, hence symmetric about
        the $y$-axis; vertical asymptote $x=0$ and horizontal asymptote
        $y=0$. Both branches lie above the axis.
  \item $D=\R$: an odd root accepts negative radicands. It is odd, since
        $\sqrt[3]{-x}=-\sqrt[3]{x}$, strictly increasing, steep near the
        origin and flattening out: the graph of $x^{3}$ mirrored in
        $y=x$.
  \item A V with vertex $(2,0)$; a single (double) root at $x=2$;
        $y$-intercept $|0-2|=2$.
  \item $C(0)=0$, $C(1)=10$, $C(2)=14$, $C(3)=18$, $C(4)=28$, $C(5)=50$.
        Increasing throughout; the increments $10,4,4,10,22$ first fall and
        then rise, so the curve is concave up to about $q=2$ and convex
        afterwards: the S-shaped total cost curve of the textbooks, with
        its inflection at $q=2$.
  \item $f(-2)=5$, $f(-1)=0$, $f(0)=-3$, $f(1)=-4$, $f(2)=-3$, $f(3)=0$,
        $f(4)=5$. The table is symmetric about $x=1$: an upward parabola
        with vertex $(1,-4)$ and roots $-1$ and $3$.
  \item Divide by $36$: $\dfrac{x^{2}}{4}+\dfrac{y^{2}}{9}=1$, an ellipse
        with $a=2$, $b=3$, crossing the axes at $(\pm2,0)$ and
        $(0,\pm3)$.
  \item A hyperbola opening left and right. Solved for $y$:
        $y=\pm\sqrt{x^{2}-9}$, defined for $|x|\ge3$, two branches.
  \item A sum of two squares is zero only if both vanish: the ``curve'' is
        the single point $(-1,2)$.
  \item $(x+1)^{2}-1+y^{2}=8$, so $(x+1)^{2}+y^{2}=9$: a circle with centre
        $(-1,0)$ and radius $3$.
  \item $y=\sqrt{9-x^{2}}$ is a function, the upper semicircle, with
        $D=[-3,3]$ and $W=[0,3]$. The full equation $x^{2}+y^{2}=9$ is not
        a function of $x$: it also contains the lower half.
  \item Intercepts $(20,0)$ on the $K$-axis and $(0,50)$ on the $L$-axis;
        $L=50-2.5K$, a falling line with slope $-2.5$, the ratio of the
        input prices.
  \item $y=36/x$: $36,18,12,9,6,4,3$. A falling, convex hyperbola branch in
        the first quadrant, approaching both axes without touching them:
        halving $x$ doubles $y$.
  \item $pq=100$. The asymptotes are $p=0$ and $q=0$: as the quantity
        grows, the price falls towards zero without reaching it, and no
        finite price clears an arbitrarily small quantity. Revenue stays at
        $100$ everywhere, which is what unit elasticity means.
  \item $y=4(1-\mathrm{e}^{-x})$. It gives $y(0)=0$ and rises to the
        asymptote $4$. \;$y=4-\mathrm{e}^{-x}$ has the right asymptote but
        starts at $y(0)=3$, not at the origin; \;$y=4\mathrm{e}^{-x}$ falls
        from $4$ towards $0$ instead of rising.
\end{solutions}

\chapter{Transforming Functions}

Knowing the catalogue of standard shapes is worth much more once it is
combined with the second half of the skill: recognising that an unfamiliar
formula is a familiar function in disguise. The graph of
$y=-2\mathrm{e}^{-(x-3)}+5$ never has to be plotted point by point. It is
the exponential, mirrored twice, stretched, and moved, and all of this is
read off the formula without any arithmetic.

Three related skills belong here. \textbf{Transforming}: what the constants
in $a\,f\bigl(b(x-c)\bigr)+d$ do to a graph, with the standing trap that
everything \emph{inside} the function acts in the opposite direction and by
the reciprocal factor. \textbf{Composing and decomposing}: building a
function out of simpler ones, and taking one apart again, which is exactly
the reading that makes the chain rule mechanical later. And
\textbf{inverting}: swapping the roles of $x$ and $y$, the operation that
turns an exponential into a logarithm and a demand curve into an inverse
demand curve.

The chapter ends with the reverse question, the one an economist meets
most often: not ``what does this formula look like'', but ``which formula has
the shape I need?'' Given a vertex, an asymptote, a saturation level or two
data points, determine the parameters.

\section{Core Rules}

\begin{rulebox}{The four elementary moves}
Start from a known graph $y=f(x)$ and read the general form
\[
  y=a\,f\bigl(b(x-c)\bigr)+d .
\]
\begin{itemize}[leftmargin=1.6em,itemsep=0.3ex]
  \item \textbf{$d$: vertical shift.} $f(x)+d$ moves the graph \emph{up}
        by $d$ (down if $d<0$). Outside the function, and in the direction
        you expect.
  \item \textbf{$c$: horizontal shift.} $f(x-c)$ moves the graph
        \emph{right} by $c$. Inside the function, and \emph{against}
        the sign: $f(x+2)$ moves two units to the \emph{left}.
  \item \textbf{$a$: vertical scaling.} $a\,f(x)$ stretches the graph
        vertically by the factor $|a|$ ($0<|a|<1$ compresses), and
        \textbf{reflects it in the $x$-axis} if $a<0$. Roots stay put;
        every $y$-value is multiplied by $a$.
  \item \textbf{$b$: horizontal scaling.} $f(bx)$ \emph{compresses} the
        graph horizontally by the factor $|b|$ --- again the reciprocal of
        what the number suggests --- and \textbf{reflects it in the
        $y$-axis} if $b<0$. The $y$-intercept stays put.
\end{itemize}
Order matters: inside operations act on $x$ before $f$ is applied, outside
operations act on the result. Apply them as
\emph{horizontal scaling and shifting first, then $f$, then vertical
scaling and shifting}. And factor before reading: $f(2x-6)=f\bigl(2(x-3)\bigr)$
is a compression by $2$ together with a shift of $3$ to the right, not a
shift of $6$.
\end{rulebox}

\begin{gallery}
\begin{tikzpicture}
  \begin{axis}[fnaxis, xmin=-2.6, xmax=4.6, ymin=-1.2, ymax=4.6,
      xtick={-2,2,4}, ytick={2,4}, title={$y=(x-2)^{2}+1$}]
    \addplot[gray!60, dashed, thin, samples=101, domain=-2.1:2.1]{x^2};
    \addplot[domain=0.15:3.85]{(x-2)^2+1};
  \end{axis}
\end{tikzpicture}\hfill
\begin{tikzpicture}
  \begin{axis}[fnaxis, xmin=-2.6, xmax=2.6, ymin=-1.2, ymax=4.6,
      xtick={-2,2}, ytick={2,4}, title={$y=2x^{2}$}]
    \addplot[gray!60, dashed, thin, samples=101, domain=-2.1:2.1]{x^2};
    \addplot[domain=-1.5:1.5]{2*x^2};
  \end{axis}
\end{tikzpicture}\hfill
\begin{tikzpicture}
  \begin{axis}[fnaxis, xmin=-3.4, xmax=3.4, ymin=-1.2, ymax=4.6,
      xtick={-2,2}, ytick={2,4}, title={$y=\mathrm{e}^{-x}$}]
    \addplot[gray!60, dashed, thin, domain=-3.2:1.5]{exp(x)};
    \addplot[domain=-1.5:3.2]{exp(-x)};
  \end{axis}
\end{tikzpicture}
\gcap{Shift right and up \quad\textbullet\quad vertical stretch
\quad\textbullet\quad reflection in the $y$-axis. The original is dashed.}
\end{gallery}

\begin{rulebox}{Rewrite first, then read: other bases, negative inner factors, higher powers}
The four moves can be read off only once the formula has the shape
$a\,f\bigl(b(x-c)\bigr)+d$ for a base function you know. Four rewritings
get it there.
\begin{itemize}[leftmargin=1.6em,itemsep=0.3ex]
  \item \textbf{Exponentials with any base.} For $q>0$, $q\neq1$, the graph
        of $q^{x}$ passes through $(0,1)$ and has the asymptote $y=0$; it
        rises if $q>1$ and falls if $0<q<1$. Since $q^{-x}=(1/q)^{x}$, a
        negative exponent is a reflection in the $y$-axis, turning a rising
        exponential into a falling one. Collect everything into a single
        base before sketching:
        \[
          \frac{a^{x}}{b^{x}}=\Bigl(\frac{a}{b}\Bigr)^{x},\qquad
          a^{2x}=\bigl(a^{2}\bigr)^{x},\qquad
          \text{e.g. }\ \frac{3^{2x}}{2^{x}}=\Bigl(\frac{9}{2}\Bigr)^{x}.
        \]
        Then $y=c\,q^{x}+d$ has domain $\R$, the horizontal asymptote
        $y=d$, and range $(d,\infty)$ if $c>0$, $(-\infty,d)$ if $c<0$.
  \item \textbf{A negative factor inside.} In $\sqrt{k-mx}$ with $m>0$,
        factor out $-m$, not $m$:
        \[
          \sqrt{k-mx}=\sqrt{-m\Bigl(x-\frac{k}{m}\Bigr)} .
        \]
        The factor $-m$ is a reflection in the $y$-axis together with a
        horizontal compression by $m$; then comes a shift \emph{right} by
        $k/m$. The domain is $x\le k/m$: the curve starts at $x=k/m$ and
        runs to the \emph{left}. The same factoring handles $\ln(k-mx)$,
        with domain $x<k/m$ and the asymptote $x=k/m$.
  \item \textbf{Logarithms with an inner factor.}
        $c\ln(ax+b)+d=c\ln\bigl(a(x+\tfrac{b}{a})\bigr)+d$: domain
        $ax+b>0$, vertical asymptote $x=-b/a$, range $\R$. For $a>0$ there
        is a second reading, $c\ln a+c\ln\bigl(x+\tfrac{b}{a}\bigr)+d$: the
        horizontal compression by $a$ is the same as a vertical shift by
        $c\ln a$.
  \item \textbf{Higher powers in a denominator.} $y=\dfrac{c}{(x-a)^{n}}+b$
        has the asymptotes $x=a$ and $y=b$. For \emph{even} $n$ the
        denominator is positive on both sides of the pole, so both branches
        lie on the \emph{same} side of $y=b$ (above it if $c>0$, below if
        $c<0$), and the range is $(b,\infty)$ or $(-\infty,b)$. For
        \emph{odd} $n$ the denominator changes sign at $x=a$, the branches
        lie on \emph{opposite} sides, as for $1/x$, and the range is
        $\R\setminus\{b\}$.
\end{itemize}
Factor powers too: $(2x-6)^{2}=4(x-3)^{2}$, so
$\dfrac{8}{(2x-6)^{2}}=\dfrac{2}{(x-3)^{2}}$, with the pole at $x=3$.
\end{rulebox}

\begin{gallery}
\begin{tikzpicture}
  \begin{axis}[fnaxis, xmin=-3.4, xmax=3.4, ymin=-1.2, ymax=4.6,
      xtick={-2,2}, ytick={2,4}, title={$y=2^{-x}=\bigl(\tfrac12\bigr)^{x}$}]
    \addplot[gray!60, dashed, thin, domain=-3.2:2.1]{2^x};
    \addplot[domain=-2.1:3.2]{2^(-x)};
  \end{axis}
\end{tikzpicture}\hfill
\begin{tikzpicture}
  \begin{axis}[fnaxis, xmin=-3.4, xmax=4.4, ymin=-1.2, ymax=4.6,
      xtick={-2,2,4}, ytick={2,4}, title={$y=\sqrt{4-2x}$}]
    \addplot[gray!60, dashed, thin, samples=101, domain=0:4.2]{sqrt(x)};
    \addplot[samples=201, domain=-3.2:2]{sqrt(4-2*x)};
  \end{axis}
\end{tikzpicture}\hfill
\begin{tikzpicture}
  \begin{axis}[fnaxis, xmin=-2.4, xmax=3.4, ymin=-3.2, ymax=4.6,
      xtick={-2,2}, ytick={-2,2,4},
      title={$\tfrac{1}{(x-1)^{2}}$ and $\tfrac{1}{(x-1)^{3}}$}]
    \addplot[samples=101, domain=-2.3:0.53]{1/(x-1)^2};
    \addplot[samples=101, domain=1.47:3.3]{1/(x-1)^2};
    \addplot[gray!60, dashed, thin, samples=101, domain=-2.3:0.31]{1/(x-1)^3};
    \addplot[gray!60, dashed, thin, samples=101, domain=1.62:3.3]{1/(x-1)^3};
  \end{axis}
\end{tikzpicture}
\gcap{A negative exponent reflects in the $y$-axis \quad\textbullet\quad
$\sqrt{-2(x-2)}$ starts at $x=2$ and runs left \quad\textbullet\quad even
power (solid): both branches above; odd power (dashed): opposite sides.}
\end{gallery}

\begin{rulebox}{Absolute values: two different operations}
\begin{itemize}[leftmargin=1.6em,itemsep=0.3ex]
  \item $y=|f(x)|$ acts on the \textbf{output}: everything below the
        $x$-axis is folded \emph{up}, everything above stays. Roots become
        points where the graph touches the axis and turns.
  \item $y=f(|x|)$ acts on the \textbf{input}: the part of the graph with
        $x\geq0$ is kept and \emph{mirrored} to the left, and whatever the
        function did for $x<0$ is discarded. The result is always even.
\end{itemize}
The two coincide only by accident. Note also $|ax+b|=|a|\,
\bigl|x+\tfrac{b}{a}\bigr|$, which turns any such expression into a shifted
and stretched V.
\end{rulebox}

\begin{gallery}
\begin{tikzpicture}
  \begin{axis}[fnaxis, xmin=-2.6, xmax=2.6, ymin=-2.6, ymax=2.6,
      xtick={-2,2}, ytick={-2,2}, title={$y=f(x)=x^{3}-3x$}]
    \addplot[domain=-2.1:2.1]{x^3-3*x};
  \end{axis}
\end{tikzpicture}\hfill
\begin{tikzpicture}
  \begin{axis}[fnaxis, xmin=-2.6, xmax=2.6, ymin=-2.6, ymax=2.6,
      xtick={-2,2}, ytick={2}, title={$y=|f(x)|$}]
    \addplot[domain=-2.1:2.1]{abs(x^3-3*x)};
  \end{axis}
\end{tikzpicture}\hfill
\begin{tikzpicture}
  \begin{axis}[fnaxis, xmin=-2.6, xmax=2.6, ymin=-2.6, ymax=2.6,
      xtick={-2,2}, ytick={-2,2}, title={$y=f(|x|)$}]
    \addplot[domain=-2.1:2.1]{abs(x)^3-3*abs(x)};
  \end{axis}
\end{tikzpicture}
\end{gallery}

\begin{rulebox}{Composing and decomposing}
$(f\circ g)(x)=f\bigl(g(x)\bigr)$: the \textbf{inner} function $g$ is
applied first. To \emph{decompose} a formula, ask what you would compute
first with a calculator; that is the inner function:
\[
  \begin{aligned}
    \sqrt{3x+1}\;&:\quad g(x)=3x+1, \quad f(u)=\sqrt{u},\\
    \mathrm{e}^{x^{2}}\;&:\quad g(x)=x^{2}, \quad f(u)=\mathrm{e}^{u},\\
    (\ln x)^{3}\;&:\quad g(x)=\ln x, \quad f(u)=u^{3}.
  \end{aligned}
\]
Decompositions are not unique, but the ``innermost bracket first'' reading
is the one the chain rule wants. Note that all four elementary moves are
themselves compositions with linear functions: $f(b(x-c))$ is $f$ composed
with a linear inner function.
\end{rulebox}

\begin{rulebox}{Inverse functions}
$f^{-1}$ undoes $f$: \;$f^{-1}\bigl(f(x)\bigr)=x$ and
$f\bigl(f^{-1}(y)\bigr)=y$. It exists exactly when $f$ is
\textbf{injective}; if it is not, restrict the domain first (as with
$x^{2}$ on $[0,\infty)$). Three consequences and one recipe:
\begin{itemize}[leftmargin=1.6em,itemsep=0.3ex]
  \item $D(f^{-1})=W(f)$ and $W(f^{-1})=D(f)$: domain and range swap.
  \item The graph of $f^{-1}$ is the graph of $f$ \textbf{mirrored in the
        line $y=x$}; the point $(p,q)$ becomes $(q,p)$.
  \item $f^{-1}$ is \emph{not} $1/f$.
  \item \textbf{Recipe:} write $y=f(x)$, solve for $x$, then rename the
        variables. \;$y=\mathrm{e}^{2x}\Rightarrow \ln y=2x\Rightarrow
        x=\tfrac12\ln y$, so $f^{-1}(x)=\tfrac12\ln x$ with
        $D(f^{-1})=(0,\infty)$.
\end{itemize}
This is the operation behind every \emph{inverse demand} function: the
demand $q(p)$ solved for $p$ gives $p(q)$, the same curve read along the
other axis.
\end{rulebox}

\begin{gallery}
\begin{tikzpicture}
  \begin{axis}[fnaxis, xmin=-3.2, xmax=4.4, ymin=-3.2, ymax=4.4,
      xtick={-2,2,4}, ytick={-2,2,4},
      title={$\mathrm{e}^{x}$ and $\ln x$}]
    \addplot[domain=-3:1.45]{exp(x)};
    \addplot[domain=0.05:4.2]{ln(x)};
    \addplot[gray!60, dashed, thin, samples=2, domain=-3:4.2]{x};
  \end{axis}
\end{tikzpicture}\hfill
\begin{tikzpicture}
  \begin{axis}[fnaxis, xmin=-0.8, xmax=4.4, ymin=-0.8, ymax=4.4,
      xtick={2,4}, ytick={2,4},
      title={$x^{2}$ on $[0,\infty)$ and $\sqrt{x}$}]
    \addplot[domain=0:2.05]{x^2};
    \addplot[domain=0:4.2]{sqrt(x)};
    \addplot[gray!60, dashed, thin, samples=2, domain=0:4.2]{x};
  \end{axis}
\end{tikzpicture}
\gcap{Each pair is symmetric about the dashed line $y=x$.}
\end{gallery}

\begin{rulebox}{Reverse engineering: choosing the parameters}
Given the shape you want, pick the family and fit the constants.
\begin{itemize}[leftmargin=1.6em,itemsep=0.3ex]
  \item \textbf{Prescribed vertex:} use $y=a(x-h)^{2}+k$ with the vertex
        $(h,k)$, and determine $a$ from one further point.
  \item \textbf{Prescribed asymptotes:} $y=\dfrac{k}{x-c}+d$ has the
        vertical asymptote $x=c$ and the horizontal one $y=d$; a rational
        $\dfrac{ax+b}{x-c}$ has $y=a$ as horizontal asymptote.
  \item \textbf{Prescribed saturation:} $y=S\bigl(1-\mathrm{e}^{-kx}\bigr)$
        starts at $0$ and rises to the ceiling $S$; $y=S\mathrm{e}^{-kx}$
        starts at $S$ and decays to $0$.
  \item \textbf{Two data points, exponential family:} $y=A\,q^{x}$ with
        $A=y(0)$, and $q$ from the ratio of the two values. Halving every
        $T$ units means $q=2^{-1/T}$.
  \item \textbf{Economic shifts:} a change in fixed cost shifts $C(q)$
        \emph{vertically}; a per-unit tax $t$ adds $tq$ and therefore
        changes the \emph{slope}, not the intercept; measuring the input in
        thousands instead of units rescales $x$ horizontally.
\end{itemize}
\end{rulebox}

\section{Worked Example}

\begin{worked}{Describe the moves taking $y=f(x)$ to $y=-3f(2x-4)+1$, apply them to $f(x)=\sqrt{x}$, and invert $g(x)=2\ln(x-1)+3$}
\wstep{Step 1. Factor the inside before saying anything about the shift.}
The inside is $2x-4$, and it must be written as
\[
  2x-4=2(x-2)
\]
before it can be read. In this form the horizontal instructions are
unambiguous: \textbf{compress by a factor $2$}, and \textbf{shift right by
$2$} --- not by $4$.
\wtrap{This is the chapter's central trap. $f(2x-4)$ shifts by $2$, because
what matters is the value of $x$ that makes the argument zero. Reading the
shift straight off the $-4$ is the most common error with
transformations.}

\wstep{Step 2. Separate inside from outside, and get the directions right.}
\begin{center}\footnotesize
\begin{tabular}{lll}
  \hline
  where & operation & effect on the graph\\
  \hline
  inside, $2(\cdot)$ & multiply $x$ by $2$ & horizontal compression by $2$ \\
  inside, $(\cdot-2)$ & subtract $2$ from $x$ & shift \emph{right} $2$ \\
  outside, $-3(\cdot)$ & multiply by $-3$ & vertical stretch $3$, flip in the $x$-axis\\
  outside, $+1$ & add $1$ & shift \emph{up} $1$ \\
  \hline
\end{tabular}
\end{center}
Inside changes act \textbf{opposite} to how they read (subtracting moves
right, multiplying compresses); outside changes act \textbf{as} they read.
The reason is that inside operations act on the input before $f$ sees it,
so to keep the output the same you must compensate.

\wstep{Step 3. Apply the moves in the right order.}
Inside: shift and scale, in the order that the factored form dictates ---
compress first, then shift, or equivalently shift the compressed picture by
$2$. Outside: stretch and reflect first, then translate. Applying the
vertical shift before the stretch would multiply the shift by $-3$ as well,
landing at $-3f(\cdot)-3$ instead.

\wstep{Step 4. Apply it to $f(x)=\sqrt{x}$.}
\[
  y=-3\sqrt{2x-4}+1=-3\sqrt{2(x-2)}+1 .
\]
The domain moves with the graph: $\sqrt{\cdot}$ needs $2x-4\ge0$, so
$x\ge2$. The starting point $(0,0)$ of $\sqrt{x}$ goes to $(2,1)$: $x=2$
because the shift is right by~$2$, and $y=1$ because $-3\cdot0+1=1$. From
there the curve falls (the reflection), so its range is $(-\infty,1]$.

\wstep{Step 5. Track one more point, as a check.}
The point $(4,2)$ on $y=\sqrt{x}$: solve $2x-4=4$ for the new abscissa,
giving $x=4$; the new ordinate is $-3\cdot2+1=-5$. So $(4,2)\mapsto(4,-5)$,
and substituting $x=4$ into the final formula gives
$-3\sqrt{4}+1=-5$. \checkmark

\wstep{Step 6. Now the inverse of $g(x)=2\ln(x-1)+3$: undo the operations in reverse order.}
Forward, $g$ does: subtract $1$, take $\ln$, double, add $3$. Backwards:
subtract $3$, halve, exponentiate, add $1$. In symbols, set $y=g(x)$ and
solve for $x$:
\[
  y-3=2\ln(x-1)
  \;\Longrightarrow\;
  \frac{y-3}{2}=\ln(x-1)
  \;\Longrightarrow\;
  \mathrm{e}^{(y-3)/2}=x-1
  \;\Longrightarrow\;
  x=1+\mathrm{e}^{(y-3)/2} .
\]
Renaming the variable,
\[
  \boxed{\;g^{-1}(x)=1+\mathrm{e}^{(x-3)/2}\;}
\]

\wstep{Step 7. State the domain and range, which swap.}
$g$ has domain $(1,\infty)$ and range $\R$. Therefore $g^{-1}$ has domain
$\R$ and range $(1,\infty)$ --- visible in the formula, since
$\mathrm{e}^{(x-3)/2}>0$ always, so $g^{-1}>1$ always.
\wtrap{The domain of the inverse is the \emph{range} of the original. This is
not bookkeeping: it is what makes $g^{-1}$ a function at all. And check
$g^{-1}(g(x))=x$ on one value --- $g(1+\mathrm{e}^{0})=g(2)=2\ln1+3=3$ and
$g^{-1}(3)=1+\mathrm{e}^{0}=2$. \checkmark}
\end{worked}

\section{Practice Problems}

In problems 1--10, describe in words how the graph of $y=f(x)$ has to be
moved, scaled or reflected.

\begin{problems}
  \item $y=f(x-3)$
  \item $y=f(x)+2$
  \item $y=-f(x)$
  \item $y=f(-x)$
  \item $y=2f(x)$
  \item $y=f(2x)$
  \item $y=f\!\left(\dfrac{x}{2}\right)+1$
  \item $y=-3f(x-1)+4$
  \item $y=f(x+2)-5$
  \item $y=\tfrac12 f(-x)$
\end{problems}

In problems 11--18, write down the formula.

\begin{problems}[resume]
  \item $y=x^{2}$ shifted $3$ to the right and $1$ up.
  \item $y=\sqrt{x}$ reflected in the $x$-axis and shifted $4$ up.
  \item $y=\mathrm{e}^{x}$ compressed horizontally by the factor $2$.
  \item $y=\ln x$ shifted $5$ to the left.
  \item $y=1/x$ shifted $2$ to the right and $3$ down.
  \item $y=|x|$ shifted $1$ to the left and stretched vertically by $3$.
  \item $y=x^{3}$ reflected in the $y$-axis and shifted $2$ up.
  \item A downward parabola with vertex $(-1,5)$, twice as steep as
        $y=x^{2}$.
\end{problems}

In problems 19--24, name the base function and the transformations, and
give the features asked for.

\begin{problems}[resume]
  \item $y=2(x+1)^{2}-8$: vertex, roots, range.
  \item $y=-\sqrt{x-2}+3$: domain, starting point, range.
  \item $y=3\mathrm{e}^{x+1}$: asymptote, value at $x=-1$.
  \item $y=\ln(2x)$: give \emph{two} readings of this graph, one
        horizontal and one vertical.
  \item $y=\dfrac{1}{x-4}+2$: both asymptotes.
  \item $y=-|x+3|+5$: vertex, roots, direction of opening.
\end{problems}

\begin{problems}[resume]
  \item For $f(x)=x^{2}-2x-3$, describe the graph of $y=|f(x)|$: which
        part moves, and what happens at the roots?
  \item Sketch $y=|2x-6|$ by first writing it as a stretched, shifted V.
  \item For $f(x)=\sqrt{x}$, describe $y=f(|x|)$: give its domain and its
        symmetry.
  \item With $f(u)=u^{2}$ and $g(x)=x-5$: give $(f\circ g)(x)$ and
        $(g\circ f)(x)$, and say which transformation of $u^{2}$ each one
        is.
  \item With $f(u)=\sqrt{u}$ and $g(x)=x^{2}+9$: give $(f\circ g)(x)$ and
        its domain.
  \item Decompose $h(x)=\sqrt{3x+1}$ into $f\circ g$.
  \item Decompose $h(x)=\mathrm{e}^{x^{2}}$ into $f\circ g$.
  \item Decompose $h(x)=(\ln x)^{3}$ into $f\circ g$.
  \item Decompose $h(x)=\dfrac{1}{x^{2}+1}$ into $f\circ g$.
\end{problems}

In problems 34--39, find the inverse together with its domain and range.

\begin{problems}[resume]
  \item $f(x)=3x-7$
  \item $f(x)=\dfrac{x+2}{x-1}$
  \item $f(x)=x^{2}$ on $D=[0,\infty)$
  \item $f(x)=\ln(x-1)$
  \item $f(x)=\sqrt{x+4}$
  \item $f(x)=\dfrac{2x+1}{x-3}$
  \item Verify for $f(x)=3x-7$ that $f\bigl(f^{-1}(x)\bigr)=x$ and
        $f^{-1}\bigl(f(x)\bigr)=x$.
  \item The demand function is $q(p)=50-2p$. Give the inverse demand
        function $p(q)$, and state on which interval of quantities it is
        economically meaningful.
\end{problems}

\begin{problems}[resume]
  \item Determine $a$ so that the parabola with vertex $(2,-3)$ passes
        through $(0,1)$, and write the function.
  \item An exponential $y=A\,q^{x}$ passes through $(0,5)$ and $(2,45)$.
        Determine $A$ and $q$.
  \item Write a rational function with vertical asymptote $x=2$,
        horizontal asymptote $y=3$, and a root at the origin.
  \item Total cost is $C(q)=5q+100$. Write the new cost function and name
        the transformation when (a) the fixed cost rises by $40$;
        (b) a tax of $2$ per unit is levied.
\end{problems}

\section{Solutions}

\begin{solutions}
  \item Shifted $3$ units to the \emph{right} (inside, against the sign).
  \item Shifted $2$ units up.
  \item Reflected in the $x$-axis: every $y$-value changes sign.
  \item Reflected in the $y$-axis: left and right are swapped.
  \item Stretched vertically by the factor $2$; roots and $x$-intercepts
        stay where they are.
  \item Compressed horizontally by the factor $2$: what happened at $x=4$
        now happens at $x=2$. The $y$-intercept is unchanged.
  \item Stretched horizontally by the factor $2$ (the inner factor
        $\tfrac12$ is a compression by $\tfrac12$, i.e. a stretch), then
        moved $1$ up.
  \item Right $1$, stretched vertically by $3$, reflected in the $x$-axis,
        then up $4$ --- inside first, outside afterwards.
  \item Left $2$ and down $5$.
  \item Reflected in the $y$-axis and compressed vertically to half its
        height.
  \item $y=(x-3)^{2}+1$.
  \item $y=-\sqrt{x}+4$.
  \item $y=\mathrm{e}^{2x}$ (a factor $2$ \emph{inside} compresses).
  \item $y=\ln(x+5)$, with the vertical asymptote moved to $x=-5$.
  \item $y=\dfrac{1}{x-2}-3$, asymptotes $x=2$ and $y=-3$.
  \item $y=3|x+1|$.
  \item $y=(-x)^{3}+2=-x^{3}+2$.
  \item $y=-2(x+1)^{2}+5$.
  \item Base $x^{2}$: left $1$, stretched by $2$, down $8$. Vertex
        $(-1,-8)$; roots from $2(x+1)^{2}=8\iff(x+1)^{2}=4\iff x=1$ or
        $x=-3$; range $[-8,\infty)$.
  \item Base $\sqrt{x}$: right $2$, reflected in the $x$-axis, up $3$.
        $D=[2,\infty)$, the curve starts at $(2,3)$ and falls,
        $W=(-\infty,3]$.
  \item Base $\mathrm{e}^{x}$: left $1$, stretched by $3$. Asymptote
        $y=0$; $y(-1)=3\mathrm{e}^{0}=3$. (Equivalently
        $3\mathrm{e}\cdot\mathrm{e}^{x}$: a horizontal shift of an
        exponential is also a vertical stretch.)
  \item Horizontally: $\ln x$ compressed by the factor $2$. Vertically:
        since $\ln(2x)=\ln 2+\ln x$, the same graph is $\ln x$ shifted up
        by $\ln2\approx0.69$. Both readings are correct --- a peculiarity
        of the logarithm.
  \item Base $1/x$: right $4$, up $2$. Vertical asymptote $x=4$,
        horizontal asymptote $y=2$.
  \item Base $|x|$: left $3$, reflected (so it opens downwards), up $5$.
        Vertex $(-3,5)$; roots where $|x+3|=5$, i.e. $x=2$ and $x=-8$.
  \item $f(x)=(x-3)(x+1)$ is negative exactly on $(-1,3)$, and that arc is
        folded up above the axis; outside it the graph is unchanged. At the
        roots $x=-1$ and $x=3$ the curve now touches the axis and turns
        instead of crossing it, and the old minimum $f(1)=-4$ becomes a
        local maximum of height $4$.
  \item $|2x-6|=2|x-3|$: a V with vertex $(3,0)$, twice as steep as
        $|x|$, i.e. with slopes $-2$ and $+2$; $y$-intercept $6$.
  \item $y=\sqrt{|x|}$: $D=\R$ now, since $|x|\geq0$ always. The right half
        is unchanged and mirrored to the left, so the function is even,
        with range $[0,\infty)$ and a sharp point at the origin.
  \item $(f\circ g)(x)=(x-5)^{2}$: the parabola shifted $5$ to the right.
        $(g\circ f)(x)=x^{2}-5$: the parabola shifted $5$ down. Inner
        shifts move horizontally, outer shifts vertically.
  \item $(f\circ g)(x)=\sqrt{x^{2}+9}$, and since $x^{2}+9\geq9>0$ for
        every $x$, $D=\R$. (Its range is $[3,\infty)$.)
  \item $g(x)=3x+1$, $f(u)=\sqrt{u}$.
  \item $g(x)=x^{2}$, $f(u)=\mathrm{e}^{u}$.
  \item $g(x)=\ln x$, $f(u)=u^{3}$.
  \item $g(x)=x^{2}+1$, $f(u)=\dfrac{1}{u}$.
  \item $y=3x-7\Rightarrow x=\dfrac{y+7}{3}$, so
        $f^{-1}(x)=\dfrac{x+7}{3}$, with $D=W=\R$.
  \item $y=\dfrac{x+2}{x-1}\Rightarrow y(x-1)=x+2\Rightarrow x(y-1)=y+2
        \Rightarrow x=\dfrac{y+2}{y-1}$. Hence
        $f^{-1}(x)=\dfrac{x+2}{x-1}=f(x)$: the function is its own
        inverse, and its graph is symmetric about $y=x$. Domain and range
        are both $\R\setminus\{1\}$.
  \item $f^{-1}(x)=\sqrt{x}$, with $D(f^{-1})=[0,\infty)=W(f)$ and
        $W(f^{-1})=[0,\infty)=D(f)$. The restriction to $[0,\infty)$ was
        needed to make $f$ injective in the first place.
  \item $y=\ln(x-1)\Rightarrow \mathrm{e}^{y}=x-1\Rightarrow
        x=\mathrm{e}^{y}+1$: $f^{-1}(x)=\mathrm{e}^{x}+1$, with $D=\R$ and
        $W=(1,\infty)$ --- exactly the domain $(1,\infty)$ of $f$.
  \item $y=\sqrt{x+4}\Rightarrow y^{2}=x+4\Rightarrow x=y^{2}-4$, and only
        $y\geq0$ occurs: $f^{-1}(x)=x^{2}-4$ with $D=[0,\infty)$ and
        $W=[-4,\infty)$. Without the restriction $D=[0,\infty)$ the formula
        $x^{2}-4$ would not be the inverse.
  \item $y=\dfrac{2x+1}{x-3}\Rightarrow y(x-3)=2x+1\Rightarrow
        x(y-2)=3y+1\Rightarrow x=\dfrac{3y+1}{y-2}$. So
        $f^{-1}(x)=\dfrac{3x+1}{x-2}$, with $D=\R\setminus\{2\}$ (the
        range of $f$) and $W=\R\setminus\{3\}$ (the domain of $f$). The two
        asymptotes have swapped roles.
  \item $f\bigl(f^{-1}(x)\bigr)=3\cdot\dfrac{x+7}{3}-7=x+7-7=x$, and
        $f^{-1}\bigl(f(x)\bigr)=\dfrac{(3x-7)+7}{3}=\dfrac{3x}{3}=x$.
  \item $q=50-2p\Rightarrow 2p=50-q\Rightarrow p(q)=25-\dfrac{q}{2}$. It
        is meaningful for $q\in[0,50]$: at $q=50$ the price has fallen to
        $0$, and larger quantities would require a negative price.
  \item $y=a(x-2)^{2}-3$ with $1=a\cdot4-3$, so $a=1$ and
        $y=(x-2)^{2}-3$.
  \item $A=y(0)=5$. Then $45=5q^{2}\Rightarrow q^{2}=9\Rightarrow q=3$
        (positive base), so $y=5\cdot3^{x}$.
  \item $y=\dfrac{3x}{x-2}$: the denominator gives the pole $x=2$, equal
        degrees give the horizontal asymptote $y=3$, and the numerator
        vanishes at $x=0$, so the graph passes through the origin.
  \item (a) $C(q)=5q+140$: a pure vertical shift by $40$; the curve keeps
        its slope. (b) $C(q)=7q+100$: the tax adds $2q$, so the
        \emph{slope} changes from $5$ to $7$ while the intercept stays.
        A vertical shift and a change of slope are different operations,
        and only the first one is a translation.
\end{solutions}

\section{Additional Problems}

\begin{problems}
  \item $y=(x+4)^{2}-2$: base function, transformations, vertex, range.
  \item $y=-2\sqrt{x}+1$: transformations, domain, range.
  \item $y=\mathrm{e}^{-(x-1)}$: transformations, asymptote, value at
        $x=1$.
  \item $y=\ln(x+3)$: transformation, vertical asymptote, domain.
  \item $y=\dfrac{3}{x+2}-1$: transformations and both asymptotes.
  \item $y=\tfrac12|x-6|$: transformations and vertex.
  \item Write: $y=1/x$ reflected in the $x$-axis and shifted $3$ up.
  \item Write: $y=\mathrm{e}^{x}$ compressed horizontally by the factor
        $3$ and shifted $1$ down.
  \item Write: $y=\sqrt{x}$ stretched vertically by $5$ and shifted $1$ to
        the right.
  \item Write: a downward parabola with vertex $(3,-4)$, twice as steep as
        $y=x^{2}$.
  \item Describe the graph of $y=|x^{2}-1|$: which arc is folded, and what
        does the graph do at $x=\pm1$?
  \item Describe the graph of $y=|\ln x|$: domain, roots, and behaviour as
        $x\to0^{+}$.
  \item For $f(x)=x^{3}-3x$, give the formula for $f(|x|)$ and state its
        symmetry.
  \item With $f(u)=u^{2}$ and $g(x)=x+3$: give both composites and say why
        they differ.
  \item Decompose $h(x)=\sqrt{1-x^{2}}$ into $f\circ g$ and give its
        domain.
  \item Decompose $h(x)=\mathrm{e}^{\sqrt{x}}$ into $f\circ g$.
  \item Decompose $h(x)=(2x+5)^{10}$ into $f\circ g$.
  \item Inverse of $f(x)=5-2x$.
  \item Inverse of $f(x)=x^{3}+1$.
  \item Inverse of $f(x)=\dfrac{1}{x-2}$, with domain and range.
  \item Inverse of $f(x)=10^{x}$, with domain and range.
  \item Inverse of $f(x)=\ln(3x)$, with domain and range.
  \item Why does $f(x)=x^{2}$ have no inverse on $\R$, and which two
        restrictions repair it?
  \item A quantity is $200$ at $x=0$ and halves every $5$ units. Write it
        both as $A\,q^{x}$ and in the form $A\,\mathrm{e}^{-kx}$.
  \item Cost is $C(q)=3q+40$. Fixed cost rises to $60$ and the unit cost
        falls to $2.5$. Write the new function and name both changes as
        transformations of the graph.
\end{problems}

In problems 26--36, rewrite the formula first (single base, factored
inside, factored power), then give the chain of moves from the base
function, the maximal domain, the range and every asymptote.

\begin{problems}[resume]
  \item $y=3-2\cdot5^{-x}$
  \item $y=1+3\cdot\bigl(\tfrac25\bigr)^{x}$
  \item $y=\dfrac{12^{x}}{3^{x}}-2$
  \item $y=4-\dfrac{2^{2x}}{8^{x}}$
  \item $y=2\sqrt{6-2x}+1$
  \item $y=2-\sqrt{5-x}$
  \item $y=3\ln(2x+6)-1$
  \item $y=1-2\ln(4-x)$
  \item $y=\dfrac{2}{(x-1)^{2}}-3$
  \item $y=1-\dfrac{1}{(x+2)^{3}}$
  \item $y=\dfrac{4}{(2x-6)^{2}}+1$
\end{problems}

\section{Additional Solutions}

\begin{solutions}
  \item Base $x^{2}$, shifted $4$ to the left and $2$ down; vertex
        $(-4,-2)$; range $[-2,\infty)$.
  \item Base $\sqrt{x}$: stretched by $2$, reflected in the $x$-axis, then
        up $1$. $D=[0,\infty)$, and since $-2\sqrt{x}\in(-\infty,0]$,
        $W=(-\infty,1]$.
  \item $\mathrm{e}^{-(x-1)}$: reflect $\mathrm{e}^{x}$ in the $y$-axis,
        then shift $1$ to the right (the inner expression is
        $-(x-1)$). Asymptote $y=0$; $y(1)=\mathrm{e}^{0}=1$.
  \item $\ln x$ shifted $3$ to the left; the asymptote moves with it to
        $x=-3$; $D=(-3,\infty)$.
  \item Base $1/x$: stretched by $3$, left $2$, down $1$. Vertical
        asymptote $x=-2$, horizontal asymptote $y=-1$.
  \item Base $|x|$: right $6$ and compressed vertically to half the height;
        vertex $(6,0)$, slopes $\pm\tfrac12$.
  \item $y=-\dfrac{1}{x}+3$.
  \item $y=\mathrm{e}^{3x}-1$.
  \item $y=5\sqrt{x-1}$.
  \item $y=-2(x-3)^{2}-4$.
  \item The arc between the roots, where $x^{2}-1<0$, i.e. over
        $(-1,1)$, is folded up; it becomes a hump of height $1$ at $x=0$.
        At $x=\pm1$ the graph touches the axis and turns, forming two
        corners.
  \item $D=(0,\infty)$; the single root is $x=1$; on $(0,1)$ the logarithm
        is negative and is folded up, so the graph falls from $+\infty$ to
        $0$ there and rises again afterwards, with a corner at $(1,0)$. As
        $x\to0^{+}$ it goes to $+\infty$.
  \item $f(|x|)=|x|^{3}-3|x|$. It is even: the right half of the original
        cubic is kept and mirrored into the left half-plane, so the old
        left branch is discarded.
  \item $(f\circ g)(x)=(x+3)^{2}$, the parabola shifted $3$ to the left;
        $(g\circ f)(x)=x^{2}+3$, the parabola shifted $3$ up. Composition
        is not commutative: an inner shift moves the graph horizontally, an
        outer one vertically.
  \item $g(x)=1-x^{2}$, $f(u)=\sqrt{u}$; $D=[-1,1]$ from $1-x^{2}\geq0$.
  \item $g(x)=\sqrt{x}$, $f(u)=\mathrm{e}^{u}$; $D=[0,\infty)$.
  \item $g(x)=2x+5$, $f(u)=u^{10}$.
  \item $y=5-2x\Rightarrow x=\dfrac{5-y}{2}$, so
        $f^{-1}(x)=\dfrac{5-x}{2}$.
  \item $y=x^{3}+1\Rightarrow x=\sqrt[3]{y-1}$, so
        $f^{-1}(x)=\sqrt[3]{x-1}$, with $D=W=\R$.
  \item $y=\dfrac{1}{x-2}\Rightarrow x-2=\dfrac1y\Rightarrow
        x=2+\dfrac1y$: $f^{-1}(x)=2+\dfrac{1}{x}$, with
        $D=\R\setminus\{0\}$ and $W=\R\setminus\{2\}$.
  \item $f^{-1}(x)=\log_{10}x$, with $D=(0,\infty)$ and $W=\R$.
  \item $y=\ln(3x)\Rightarrow \mathrm{e}^{y}=3x\Rightarrow
        x=\dfrac{\mathrm{e}^{y}}{3}$: $f^{-1}(x)=\dfrac{\mathrm{e}^{x}}{3}$,
        with $D=\R$ and $W=(0,\infty)$.
  \item It is not injective: $f(-2)=f(2)=4$, so a horizontal line meets the
        graph twice and the inverse assignment would not be unique.
        Restricting to $[0,\infty)$ gives the inverse $\sqrt{x}$;
        restricting to $(-\infty,0]$ gives $-\sqrt{x}$.
  \item Halving every $5$ units means $q=2^{-1/5}$, so
        $y=200\cdot2^{-x/5}$. Since $2^{-x/5}
        =\mathrm{e}^{-(\ln2/5)x}$, the same function is
        $y=200\,\mathrm{e}^{-kx}$ with
        $k=\dfrac{\ln2}{5}\approx0.1386$.
  \item $C(q)=2.5q+60$. Raising the fixed cost from $40$ to $60$ shifts the
        line $20$ units up; lowering the unit cost from $3$ to $2.5$
        flattens it, a change of slope rather than a translation. Only the
        first is a shift of the graph.
  \item Rewrite $5^{-x}=\bigl(\tfrac15\bigr)^{x}$. From $y=5^{x}$:
        reflect in the $y$-axis (the graph now falls), stretch vertically by
        $2$, reflect in the $x$-axis (it now rises again), shift up $3$.
        $D=\R$. Since $2\cdot5^{-x}$ takes every value in $(0,\infty)$,
        $-2\cdot5^{-x}$ takes every value in $(-\infty,0)$ and
        $W=(-\infty,3)$. Horizontal asymptote $y=3$, approached from below
        as $x\to\infty$; as $x\to-\infty$ the graph falls without bound.
        Check: $y(0)=3-2=1$. Root: $5^{-x}=\tfrac32$, so
        $x=-\dfrac{\ln1.5}{\ln5}\approx-0.25$.
  \item The base $\tfrac25$ lies between $0$ and $1$, so $(\tfrac25)^{x}$
        is a \emph{falling} exponential through $(0,1)$. Stretch vertically
        by $3$, shift up $1$. $D=\R$, $W=(1,\infty)$, horizontal asymptote
        $y=1$, approached from above as $x\to\infty$. $y(0)=1+3=4$.
        (Equivalently $3\cdot(\tfrac52)^{-x}+1$: the rising exponential
        $(\tfrac52)^{x}$ reflected in the $y$-axis.)
  \item Rewrite $\dfrac{12^{x}}{3^{x}}=\Bigl(\dfrac{12}{3}\Bigr)^{x}=4^{x}$,
        so $y=4^{x}-2$: the rising exponential $4^{x}$ shifted $2$ down.
        $D=\R$, $W=(-2,\infty)$, horizontal asymptote $y=-2$ as
        $x\to-\infty$. $y(0)=-1$; root $4^{x}=2$, i.e. $x=\tfrac12$.
  \item Rewrite $2^{2x}=4^{x}$, so $\dfrac{2^{2x}}{8^{x}}
        =\Bigl(\dfrac48\Bigr)^{x}=\bigl(\tfrac12\bigr)^{x}=2^{-x}$ and
        $y=4-2^{-x}$. From $2^{x}$: reflect in the $y$-axis, reflect in the
        $x$-axis, shift up $4$. $D=\R$, $W=(-\infty,4)$, horizontal
        asymptote $y=4$ as $x\to\infty$. $y(0)=3$; root $2^{-x}=4=2^{2}$,
        so $x=-2$.
  \item Factor out $-2$: $6-2x=-2(x-3)$, so $y=2\sqrt{-2(x-3)}+1$. From
        $\sqrt{x}$: reflect in the $y$-axis and compress horizontally by
        $2$ (giving $\sqrt{-2x}$), shift right $3$, stretch vertically by
        $2$, shift up $1$. Domain from $6-2x\ge0$: $D=(-\infty,3]$. The
        curve starts at $(3,1)$ and rises to the left without bound, so
        $W=[1,\infty)$; no asymptote. Check: $y(1)=2\sqrt4+1=5$.
  \item Factor out $-1$: $y=-\sqrt{-(x-5)}+2$. From $\sqrt{x}$: reflect in
        the $y$-axis, shift right $5$, reflect in the $x$-axis, shift up
        $2$. $D=(-\infty,5]$, starting point $(5,2)$, from which the curve
        falls to the left: $W=(-\infty,2]$. No asymptote. Root:
        $\sqrt{5-x}=2$, so $x=1$.
  \item Factor the inside: $2x+6=2(x+3)$, so $y=3\ln\bigl(2(x+3)\bigr)-1$.
        From $\ln x$: compress horizontally by $2$, shift left $3$, stretch
        vertically by $3$, shift down $1$. $D=(-3,\infty)$ from $2x+6>0$;
        vertical asymptote $x=-3$; $W=\R$. Second reading:
        $y=3\ln(x+3)+(3\ln2-1)$, i.e. $\ln(x+3)$ stretched by $3$ and
        shifted up by $3\ln2-1\approx1.08$. Root: $\ln(2x+6)=\tfrac13$,
        $x=\tfrac12\bigl(\mathrm{e}^{1/3}-6\bigr)\approx-2.30$.
  \item Factor out $-1$: $y=-2\ln\bigl(-(x-4)\bigr)+1$. From $\ln x$:
        reflect in the $y$-axis, shift right $4$, stretch vertically by
        $2$, reflect in the $x$-axis, shift up $1$. $D=(-\infty,4)$ from
        $4-x>0$; vertical asymptote $x=4$. As $x\to4^{-}$,
        $\ln(4-x)\to-\infty$, so $y\to+\infty$; the two reflections together
        leave an \emph{increasing} curve. $W=\R$. Root: $\ln(4-x)=\tfrac12$,
        $x=4-\sqrt{\mathrm{e}}\approx2.35$.
  \item Base $1/x^{2}$ (even power): stretch vertically by $2$, shift
        right $1$, down $3$. $D=\R\setminus\{1\}$, vertical asymptote $x=1$,
        horizontal asymptote $y=-3$. Since $\dfrac{2}{(x-1)^{2}}>0$ on both
        sides, both branches lie above $y=-3$: $W=(-3,\infty)$. Roots:
        $(x-1)^{2}=\tfrac23$, $x=1\pm\sqrt{2/3}\approx1\pm0.82$.
  \item Base $1/x^{3}$ (odd power): reflect in the $x$-axis, shift left
        $2$, up $1$. $D=\R\setminus\{-2\}$, asymptotes $x=-2$ and $y=1$.
        For $x>-2$ the cube $(x+2)^{3}$ is positive and $y<1$; for $x<-2$
        it is negative and $y>1$: the branches lie on opposite sides, and
        $W=\R\setminus\{1\}$. Root: $(x+2)^{3}=1$, $x=-1$; $y(0)=\tfrac78$.
  \item Factor the power: $(2x-6)^{2}=\bigl(2(x-3)\bigr)^{2}=4(x-3)^{2}$,
        so $y=\dfrac{1}{(x-3)^{2}}+1$. Base $1/x^{2}$: shift right $3$, up
        $1$. $D=\R\setminus\{3\}$, asymptotes $x=3$ and $y=1$, both branches
        above $y=1$: $W=(1,\infty)$. No roots.
\end{solutions}

\chapter{Inequalities in Two Variables: Regions in the Plane}

The algebra volume solved inequalities in one variable, where the answer is
a set of points on a line --- an interval, or a union of a few. With two variables
the same question has a two-dimensional answer: the solution set of
$2x+3y\le12$ is not a list of numbers but a \emph{region} of the plane, and
it is described by drawing it.

The mechanism is short. The corresponding \emph{equation} $2x+3y=12$ is a
curve, and a curve cuts the plane into pieces. On each piece the expression
$2x+3y-12$ cannot change sign without passing through zero, so its sign is
constant there. One test point per piece therefore settles the whole
region --- the same sign-table logic as in one variable, with curves in
place of points.

Two directions are worth drilling, and examinations favour the second. The
forward direction is: given inequalities, shade the region. The reverse is:
given a shaded region, write down the inequalities that describe it. The
reverse is harder and more useful, because it is what modelling actually
requires --- a feasible set is described in words or in a picture long before
anyone writes the constraints. It is also exactly the skill that the
graphical solution of a linear programme depends on.

\section{Core Rules}

\begin{rulebox}{The solution set is a region}
For an expression $F(x,y)$, the set $\{F(x,y)<0\}$ is bounded by the curve
$\{F(x,y)=0\}$. The recipe:
\begin{enumerate}[leftmargin=1.8em,itemsep=0.3ex]
  \item \textbf{Draw the boundary} $F(x,y)=0$. It is one of the curves of
        Chapter~2: a line, circle, ellipse, parabola or hyperbola.
  \item \textbf{Note which pieces it creates.} A line makes two half-planes;
        a circle makes an inside and an outside; a parabola makes an
        ``above'' and a ``below''.
  \item \textbf{Test one point per piece.} Substituting a single convenient
        point --- the origin whenever it is not on the boundary --- decides
        the sign for that whole piece.
  \item \textbf{Mark the boundary.} Solid if the inequality is
        $\le$ or $\ge$ (the boundary belongs to the set), dashed if it is
        $<$ or $>$ (it does not).
\end{enumerate}
\end{rulebox}

\begin{rulebox}{The standard regions}
\begin{itemize}[leftmargin=1.6em,itemsep=0.3ex]
  \item \textbf{Half-plane:} $ax+by\le c$. Draw the line, test the origin.
        With $c>0$ and the origin satisfying it, the region is the side
        containing the origin.
  \item \textbf{Disc and its complement:} $x^{2}+y^{2}\le r^{2}$ is the
        closed disc; $\ge$ gives everything outside. Shifted:
        $(x-x_{0})^{2}+(y-y_{0})^{2}\le r^{2}$.
  \item \textbf{Ellipse interior:} $\dfrac{x^{2}}{a^{2}}
        +\dfrac{y^{2}}{b^{2}}\le1$.
  \item \textbf{Above or below a parabola:} $y\ge x^{2}$ is the region
        \emph{inside} the parabola's opening; $y\le x^{2}$ is everything
        outside it. Rewriting an inequality as $y\ge\cdots$ or
        $y\le\cdots$ makes the side unambiguous: ``$y$ at least this'' means
        upwards.
  \item \textbf{Between hyperbola branches:} $xy\ge c$ with $c>0$ is the
        pair of regions beyond the two branches --- and on the positive
        quadrant alone it is the convex ``better set'' of Cobb--Douglas.
  \item \textbf{Strip:} $|x|\le a$ is the vertical strip $-a\le x\le a$;
        $|y|\le b$ a horizontal one.
\end{itemize}
\end{rulebox}

\begin{gallery}
\begin{tikzpicture}
  \begin{axis}[geoaxis, xmin=-1.4, xmax=7.4, ymin=-1.4, ymax=5.4,
      xtick={2,4,6}, ytick={2,4}, title={$2x+3y\le12$, $x,y\ge0$}]
    \addplot[draw=none, name path=lin, domain=0:6, samples=2]{(12-2*x)/3};
    \addplot[draw=none, name path=ax, domain=0:6, samples=2]{0};
    \region{lin}{ax}{}
    \addplot[domain=-1.2:7.2, samples=2]{(12-2*x)/3};
  \end{axis}
\end{tikzpicture}\hfill
\begin{tikzpicture}
  \begin{axis}[geoaxis, xmin=-3.4, xmax=3.4, ymin=-3.4, ymax=3.4,
      xtick={-2,2}, ytick={-2,2}, title={$x^{2}+y^{2}\le4$}]
    \addplot[draw=none, name path=up, domain=-2:2, samples=201]{sqrt(4-x^2)};
    \addplot[draw=none, name path=lo, domain=-2:2, samples=201]{-sqrt(4-x^2)};
    \region{up}{lo}{}
    \addplot[domain=0:360, samples=181]({2*cos(x)},{2*sin(x)});
  \end{axis}
\end{tikzpicture}\hfill
\begin{tikzpicture}
  \begin{axis}[geoaxis, xmin=-2.6, xmax=2.6, ymin=-0.8, ymax=4.6,
      xtick={-2,2}, ytick={2,4}, title={$x^{2}\le y\le 4$}]
    \addplot[draw=none, name path=par, domain=-2:2, samples=201]{x^2};
    \addplot[draw=none, name path=top, domain=-2:2, samples=2]{4};
    \region{par}{top}{}
    \addplot[domain=-2.15:2.15, samples=201]{x^2};
    \addplot[domain=-2.5:2.5, samples=2]{4};
  \end{axis}
\end{tikzpicture}
\gcap{Three regions: an intersection of three half-planes (a triangle), a
disc, and the region caught between a parabola and a line.}
\end{gallery}

\begin{rulebox}{Systems: intersect}
Several inequalities holding \emph{simultaneously} give the
\textbf{intersection} of their regions --- shade each one and keep only the
overlap. Two consequences:
\begin{itemize}[leftmargin=1.6em,itemsep=0.3ex]
  \item An intersection of half-planes is a \textbf{convex polygon}
        (possibly unbounded). This is precisely the feasible region of a
        linear programme, and the reason the optimum of an LP sits at a
        corner.
  \item An intersection of convex regions is always convex, so discs,
        ellipse interiors and half-planes may be combined freely and the
        result stays convex. A \emph{union}, by contrast, usually is not,
        and neither is the outside of a disc.
\end{itemize}
\end{rulebox}

\begin{rulebox}{The reverse direction: from picture to inequalities}
Given a shaded region, work boundary by boundary.
\begin{enumerate}[leftmargin=1.8em,itemsep=0.3ex]
  \item \textbf{Identify every boundary curve} and write its equation. A
        circle needs its centre and radius; a line two points or an
        intercept and a slope; a parabola its vertex and one further
        point.
  \item \textbf{Fix the direction of each} with a point that is clearly
        inside the shaded region: substitute it and keep the inequality
        sign that comes out true.
  \item \textbf{Join the conditions with ``and''}, one per boundary. A region
        bounded by three curves needs three inequalities.
  \item \textbf{Check a point just outside} each boundary: it must fail
        exactly one condition. If it fails none, a condition is missing.
\end{enumerate}
Read the boundary style too: a dashed boundary means a strict inequality.
\end{rulebox}

\begin{rulebox}{Which of these regions are convex}
One rule settles almost every case. For a function $f$ of one variable:
\begin{itemize}[leftmargin=1.6em,itemsep=0.3ex]
  \item $f$ \textbf{convex} $\Longrightarrow$ the region \emph{above} it,
        $\{y\ge f(x)\}$ (its \de{Epigraph}), is a convex set. The region
        \emph{below} it is then not convex.
  \item $f$ \textbf{concave} $\Longrightarrow$ the region \emph{below} it,
        $\{y\le f(x)\}$, is convex. The region above it is then not.
\end{itemize}
The proof is one line: if $y_{i}\ge f(x_{i})$ and $f$ is convex then
$ty_{1}+(1-t)y_{2}\ge tf(x_{1})+(1-t)f(x_{2})\ge f(tx_{1}+(1-t)x_{2})$.
So $y\ge x^{2}$ is convex and $y\le x^{2}$ is not, while $y\le4-x^{2}$
\emph{is} convex, because $4-x^{2}$ is concave. Add to this that
half-planes, discs and ellipse interiors are convex, that any
\textbf{intersection} of convex sets is convex, and that the
\emph{outside} of a disc and any \emph{union} of separated pieces is not,
and every region in this chapter is decided.
\end{rulebox}

\begin{rulebox}{Absolute values in two variables}
Split on the sign, exactly as in one variable, or recognise the standard
shapes:
\[
  |x|+|y|\le c \;\text{(a diamond, or square rotated }45^{\circ}),\qquad
  \max(|x|,|y|)\le c \;\text{(a square)} .
\]
For $|y|\le f(x)$ the region is squeezed between $y=f(x)$ and $y=-f(x)$,
and exists only where $f(x)\ge0$. For $y\ge|x|$ it is the wedge above both
lines $y=x$ and $y=-x$.
\end{rulebox}

\section{Worked Example}

\begin{worked}{Sketch $S=\{(x,y):x^{2}+y^{2}\le16,\ y\ge0,\ x\le2\}$, then decide whether $\{(x,y):y\ge\sqrt{x}\}$ is convex}
\wstep{Step 1. Take the constraints one at a time, and shade each.}
\begin{itemize}[leftmargin=1.6em,itemsep=0.2ex]
  \item $x^{2}+y^{2}\le16$: the closed disc of radius $4$ about the origin
        --- ``$\le$'' is the inside, because the origin satisfies it
        ($0\le16$).
  \item $y\ge0$: the upper half-plane, boundary included.
  \item $x\le2$: everything to the left of the vertical line $x=2$.
\end{itemize}
\wtrap{Decide inside-versus-outside by testing one point, not by staring at
the inequality sign. The origin is free and settles the disc immediately.
For a line not through the origin, any convenient point does.}

\wstep{Step 2. Intersect the three.}
The region is the part of the upper half-disc lying left of $x=2$: a
half-disc with a circular segment removed on the right. Its boundary
consists of three pieces --- the segment of the $x$-axis from $x=-4$ to
$x=2$, the vertical segment $x=2$ from $y=0$ up to $y=\sqrt{16-4}=2\sqrt3$,
and the arc of the circle from $(2,2\sqrt3)$ counter-clockwise round to
$(-4,0)$.

\wstep{Step 3. Find the corner points explicitly.}
Where $x=2$ meets the circle: $4+y^{2}=16$, so $y=2\sqrt3\approx3.46$ (the
negative root is cut off by $y\ge0$). Where $y=0$ meets the circle:
$x=\pm4$, of which $x=4$ is cut off by $x\le2$, leaving $(-4,0)$. And
$(2,0)$ where the two straight boundaries meet. Three corners:
$(-4,0)$, $(2,0)$, $(2,2\sqrt3)$.

\wstep{Step 4. Sketch, and check one interior point.}
\begin{center}
\begin{tikzpicture}
\begin{axis}[geoaxis, width=6.6cm, height=6.6cm,
             xmin=-5.2, xmax=5.2, ymin=-1.2, ymax=5.2,
             xtick={-4,-2,2,4}, ytick={2,3.464},
             yticklabels={$2$,$2\sqrt3$}]
  % the region: under the upper semicircle, above y=0, left of x=2
  \addplot[accent, very thick, domain=-4:2, samples=200,
           fill=regionfill, fill opacity=0.5] {sqrt(16-x^2)} \closedcycle;
  % the discarded part of the disc, for contrast
  \addplot[accent!35, thick, dashed, domain=2:4, samples=80]
           {sqrt(16-x^2)};
  \node[font=\scriptsize] at (axis cs:-1.2,1.4) {$S$};
  \node[font=\tiny, anchor=south west] at (axis cs:2.1,3.6) {$(2,2\sqrt3)$};
\end{axis}
\end{tikzpicture}
\end{center}
Test $(0,1)$: $0+1\le16$ \checkmark, $1\ge0$ \checkmark, $0\le2$
\checkmark. Inside, as drawn.

\wstep{Step 5. Note that $S$ is convex, and why you can say so instantly.}
Each of the three constraint sets is convex --- a disc, a half-plane, a
half-plane --- and \textbf{an intersection of convex sets is convex}. No
picture is needed for that conclusion, and no amount of picture-gazing
substitutes for it.

\wstep{Step 6. The second region: $\{y\ge\sqrt{x}\}$. Do not trust the eye.}
The boundary $y=\sqrt{x}$ is a smooth, gently bending curve, and the region
above it looks perfectly well behaved. Test the definition instead. Take the
two boundary points
\[
  P=(0,0),\qquad Q=(4,2),
\]
both in the set, since $0\ge\sqrt0$ and $2\ge\sqrt4$. Their midpoint is
$M=(2,1)$, and
\[
  \sqrt{2}\approx1.414>1 ,
\]
so $M$ violates $y\ge\sqrt{x}$ and lies \emph{outside}. A segment between
two points of the set leaves the set:
\[
  \boxed{\{(x,y):y\ge\sqrt{x}\}\ \text{is not convex.}}
\]

\wstep{Step 7. The principle behind it.}
$\sqrt{x}$ is \textbf{concave}. The set above the graph of a function --- its
epigraph --- is convex exactly when the function is \emph{convex}; the set
below --- the hypograph --- is convex exactly when the function is
\emph{concave}. So $\{y\ge\sqrt{x}\}$ fails while $\{y\le\sqrt{x}\}$
succeeds, and $\{y\ge x^{2}\}$ succeeds while $\{y\le x^{2}\}$ fails.
\wtrap{``Above'' is not a synonym for ``convex''. Match the direction of the
inequality to the curvature of the function: convex function with $\ge$, or
concave function with $\le$. Getting this backwards is the single most
common mistake in the chapter, and it recurs in Volume~7 where the feasible
set of a constrained problem must be convex for the first-order conditions
to be sufficient.}
\end{worked}

\section{Practice Problems}

Sketch the region described by each inequality in problems 1--12, and say
whether it is convex.

\begin{problems}
  \item $y\ge 2x-1$
  \item $x+y\le 3$
  \item $x^{2}+y^{2}\le 9$
  \item $x^{2}+y^{2}>4$
  \item $y\ge x^{2}$
  \item $y\le 4-x^{2}$
  \item $\dfrac{x^{2}}{9}+\dfrac{y^{2}}{4}\le1$
  \item $xy\ge 1$ on $x>0$
  \item $|x|\le 2$
  \item $|y|\le x$
  \item $(x-2)^{2}+(y+1)^{2}\le 9$
  \item $y>\dfrac1x$ on $x>0$
\end{problems}

In problems 13--20 sketch the region defined by the whole system.

\begin{problems}[resume]
  \item $x\ge0$, \;$y\ge0$, \;$2x+3y\le12$
  \item $x\ge0$, \;$y\ge0$, \;$x+y\le4$, \;$y\le2$
  \item $x^{2}+y^{2}\le16$, \;$y\ge0$
  \item $x^{2}+y^{2}\le16$, \;$y\ge0$, \;$x\le2$
  \item $x^{2}\le y$, \;$y\le x+2$
  \item $1\le x^{2}+y^{2}\le9$
  \item $x^{2}+y^{2}\le9$, \;$x\ge0$, \;$y\ge0$
  \item $y\ge x^{2}$, \;$y\le 4$, \;$x\ge0$
\end{problems}

Problems 21--28 refer to the eight shaded regions below. For each, give a
system of inequalities describing it.

\begin{gallery}
\begin{tikzpicture}
  \begin{axis}[geoaxis, xmin=-5.2, xmax=5.2, ymin=-5.2, ymax=5.2,
      xtick={-4,-2,2,4}, ytick={-4,-2,2,4}, title={(a)}]
    \addplot[draw=none, name path=up, domain=-4:2, samples=201]{sqrt(16-x^2)};
    \addplot[draw=none, name path=ax, domain=-4:2, samples=2]{0};
    \region{up}{ax}{}
    \addplot[domain=0:360, samples=181]({4*cos(x)},{4*sin(x)});
    \addplot[gray!60, dashed, thin, samples=2] coordinates {(2,-5)(2,5)};
  \end{axis}
\end{tikzpicture}\hfill
\begin{tikzpicture}
  \begin{axis}[geoaxis, xmin=-2.8, xmax=2.8, ymin=-0.9, ymax=5.1,
      xtick={-2,2}, ytick={2,4}, title={(b)}]
    \addplot[draw=none, name path=par, domain=-2:2, samples=201]{x^2};
    \addplot[draw=none, name path=top, domain=-2:2, samples=2]{4};
    \region{par}{top}{}
    \addplot[domain=-2.2:2.2, samples=201]{x^2};
    \addplot[domain=-2.6:2.6, samples=2]{4};
  \end{axis}
\end{tikzpicture}\hfill
\begin{tikzpicture}
  \begin{axis}[geoaxis, xmin=-3.9, xmax=3.9, ymin=-2.9, ymax=2.9,
      xtick={-2,2}, ytick={-2,2}, title={(c)}]
    \addplot[draw=none, name path=ell, domain=0:3, samples=201]{2*sqrt(1-x^2/9)};
    \addplot[draw=none, name path=ax, domain=0:3, samples=2]{0};
    \region{ell}{ax}{}
    \addplot[domain=0:360, samples=181]({3*cos(x)},{2*sin(x)});
  \end{axis}
\end{tikzpicture}
\end{gallery}

\begin{gallery}
\begin{tikzpicture}
  \begin{axis}[geoaxis, xmin=-2.4, xmax=3.4, ymin=-0.9, ymax=4.9,
      xtick={-2,2}, ytick={2,4}, title={(d)}]
    \addplot[draw=none, name path=par, domain=-1:2, samples=201]{x^2};
    \addplot[draw=none, name path=lin, domain=-1:2, samples=2]{x+2};
    \region{par}{lin}{}
    \addplot[domain=-2.1:2.2, samples=201]{x^2};
    \addplot[domain=-2.2:3.2, samples=2]{x+2};
  \end{axis}
\end{tikzpicture}\hfill
\begin{tikzpicture}
  \begin{axis}[geoaxis, xmin=-3.9, xmax=3.9, ymin=-3.9, ymax=3.9,
      xtick={-2,2}, ytick={-2,2}, title={(e)}]
    \addplot[draw=none, name path=o1, domain=-3:3, samples=201]{sqrt(9-x^2)};
    \addplot[draw=none, name path=o2, domain=-3:3, samples=201]{-sqrt(9-x^2)};
    \region{o1}{o2}{}
    \addplot[draw=none, name path=i1, domain=-1:1, samples=201]{sqrt(1-x^2)};
    \addplot[draw=none, name path=i2, domain=-1:1, samples=201]{-sqrt(1-x^2)};
    \addplot[white, forget plot, sharp plot] fill between[of=i1 and i2];
    \addplot[domain=0:360, samples=181]({3*cos(x)},{3*sin(x)});
    \addplot[domain=0:360, samples=181]({cos(x)},{sin(x)});
  \end{axis}
\end{tikzpicture}\hfill
\begin{tikzpicture}
  \begin{axis}[geoaxis, xmin=-1.4, xmax=7.4, ymin=-1.4, ymax=5.4,
      xtick={2,4,6}, ytick={2,4}, title={(f)}]
    \addplot[draw=none, name path=lin, domain=0:6, samples=2]{(12-2*x)/3};
    \addplot[draw=none, name path=ax, domain=0:6, samples=2]{0};
    \region{lin}{ax}{}
    \addplot[domain=-1.2:7.2, samples=2]{(12-2*x)/3};
  \end{axis}
\end{tikzpicture}
\end{gallery}

\begin{gallery}
\begin{tikzpicture}
  \begin{axis}[geoaxis, xmin=-1.2, xmax=5.4, ymin=-2.9, ymax=2.9,
      xtick={2,4}, ytick={-2,2}, title={(g)}]
    \addplot[draw=none, name path=up, domain=0:4, samples=201]{sqrt(x)};
    \addplot[draw=none, name path=lo, domain=0:4, samples=201]{-sqrt(x)};
    \region{up}{lo}{}
    \addplot[domain=0:4.6, samples=201]{sqrt(x)};
    \addplot[domain=0:4.6, samples=201]{-sqrt(x)};
    \addplot[gray!60, dashed, thin, samples=2] coordinates {(4,-2.8)(4,2.8)};
  \end{axis}
\end{tikzpicture}\hfill
\begin{tikzpicture}
  \begin{axis}[geoaxis, xmin=-3.4, xmax=3.4, ymin=-3.4, ymax=3.4,
      xtick={-2,2}, ytick={-2,2}, title={(h)}]
    \addplot[draw=none, name path=up, domain=-2:2, samples=101]{2-abs(x)};
    \addplot[draw=none, name path=lo, domain=-2:2, samples=101]{abs(x)-2};
    \region{up}{lo}{}
    \addplot[domain=-2:2, samples=101]{2-abs(x)};
    \addplot[domain=-2:2, samples=101]{abs(x)-2};
  \end{axis}
\end{tikzpicture}
\gcap{Regions (a)--(h) for problems 21--28. A dashed boundary means the
inequality there is strict.}
\end{gallery}

\begin{problems}[resume]
  \item Region (a).
  \item Region (b).
  \item Region (c).
  \item Region (d).
  \item Region (e).
  \item Region (f).
  \item Region (g).
  \item Region (h).
\end{problems}

\begin{problems}[resume]
  \item Which of the eight regions (a)--(h) are convex?
  \item Sketch $|x|+|y|\le 3$ and give its four vertices.
  \item Sketch $y\ge|x|$ and describe it in words.
  \item Sketch $|y|\le 4-x^{2}$ and give the $x$-values for which the
        region is non-empty.
  \item Write inequalities for the closed square with vertices
        $(\pm2,\pm2)$, using absolute values and then without them.
  \item A test point satisfies every inequality of a system but lies
        outside the intended region. What has gone wrong?
  \item Explain why the outside of a disc is never convex.
  \item Give a system of two inequalities whose solution set is empty, both
        of them satisfiable on their own.
  \item Show that the intersection of any two convex regions is convex.
  \item Is the union of the two half-planes $y\ge1$ and $y\le-1$ convex?
        Justify.
  \item Write inequalities for the set of points strictly inside the circle
        of radius $2$ about the origin and strictly above the line $y=x$.
  \item For which values of $c$ is the region $\{x^{2}+y^{2}\le c\}$ empty,
        a single point, and a disc?
\end{problems}

\begin{problems}[resume]
  \item \textbf{Budget set.} A consumer has income $60$ and faces prices
        $p_{x}=3$, $p_{y}=4$. Write the budget set as a system, sketch it,
        and give its three vertices.
  \item For the same consumer, write the set of bundles that are
        \emph{affordable and} give at least utility $12$ under
        $U(x,y)=xy$. Sketch it and say whether it is convex.
  \item \textbf{Production.} A firm can use at most $40$ hours of labour
        and $30$ units of capital; product A needs $2$ hours and $1$ unit,
        product B needs $1$ hour and $2$ units. Write the feasible set and
        sketch it.
  \item For the feasible set of the previous problem, give all its
        vertices, and say why a linear objective attains its maximum at one
        of them.
  \item \textbf{Tax wedge.} The set $\{(q,p): p\le 100-2q,\ p\ge 20+q,\
        q\ge0\}$ describes the quantities at which demand price exceeds
        supply price. Sketch it and find the largest $q$ in it.
\end{problems}

\section{Solutions}

\begin{solutions}
  \item The closed half-plane on or above the line $y=2x-1$. Testing
        $(0,0)$: $0\ge-1$ true, so the origin is in it. Convex.
  \item The closed half-plane on or below $x+y=3$ (through $(3,0)$ and
        $(0,3)$); the origin satisfies it. Convex.
  \item The closed disc of radius $3$ about the origin. Convex.
  \item Everything strictly outside the circle of radius $2$, boundary
        dashed and excluded. \textbf{Not} convex.
  \item The region on or above the parabola $y=x^{2}$ --- inside its
        opening. Convex.
  \item The region on or below the downward parabola $y=4-x^{2}$, which
        includes everything far below it. Since $4-x^{2}$ is
        \emph{concave}, the region below it is \textbf{convex}.
  \item The closed elliptical disc with semi-axes $3$ and $2$. Convex.
  \item The region on or beyond the hyperbola branch $y=1/x$ in the first
        quadrant. Convex (this is the quasiconcavity of $xy$).
  \item The closed vertical strip $-2\le x\le2$. Convex, being the
        intersection of two half-planes.
  \item Requires $x\ge0$ and $-x\le y\le x$: the wedge opening to the right
        between the lines $y=\pm x$. Convex.
  \item The closed disc of radius $3$ centred at $(2,-1)$. Convex.
  \item The region strictly above the hyperbola branch $y=1/x$ for $x>0$,
        boundary dashed. Convex.
  \item The triangle with vertices $(0,0)$, $(6,0)$, $(0,4)$ --- the
        intersection of three half-planes.
  \item The quadrilateral with vertices $(0,0)$, $(4,0)$, $(2,2)$,
        $(0,2)$: the triangle $x+y\le4$ in the first quadrant, with its top
        cut off by $y\le2$.
  \item The upper half-disc of radius $4$: the semicircular region above
        the $x$-axis.
  \item The same half-disc with everything right of $x=2$ removed --- the
        region shown in figure (a).
  \item The lens between the parabola and the line, from $x=-1$ to $x=2$
        (the boundaries meet where $x^{2}=x+2$, i.e. $x=-1,2$). Convex.
  \item The annulus between the circles of radius $1$ and $3$. Not convex.
  \item The quarter-disc of radius $3$ in the first quadrant. Convex.
  \item The region above the parabola, below $y=4$, right of the $y$-axis:
        the right half of figure (b).
  \item $x^{2}+y^{2}\le16$, \;$y\ge0$, \;$x<2$. The circle and the
        $x$-axis are solid, so those two are non-strict; the vertical
        boundary is dashed, so that one is strict.
  \item $y\ge x^{2}$ and $y\le4$. (Equivalently
        $x^{2}\le y\le 4$, which also forces $-2\le x\le2$ without needing
        a separate condition.)
  \item $\dfrac{x^{2}}{9}+\dfrac{y^{2}}{4}\le1$, \;$x\ge0$, \;$y\ge0$: the
        quarter of the elliptical disc in the first quadrant.
  \item $y\ge x^{2}$ and $y\le x+2$.
  \item $1\le x^{2}+y^{2}\le9$, that is $x^{2}+y^{2}\ge1$ \emph{and}
        $x^{2}+y^{2}\le9$.
  \item $2x+3y\le12$, \;$x\ge0$, \;$y\ge0$.
  \item $y^{2}\le x$ and $x<4$; equivalently
        $-\sqrt{x}\le y\le\sqrt{x}$ with $x<4$. The parabola opens to the
        right and is solid; the line $x=4$ is dashed.
  \item $|x|+|y|\le2$. Written without absolute values it is the four
        conditions $x+y\le2$, $x-y\le2$, $-x+y\le2$, $-x-y\le2$.
  \item All except (e). Each of (a), (b), (c), (d), (f), (g), (h) is an
        intersection of convex sets: discs and ellipse interiors,
        half-planes, the region above the convex $x^{2}$, the region
        \emph{right} of the parabola $x=y^{2}$ (the same rule with the axes
        exchanged, since $y\mapsto y^{2}$ is convex), and the diamond
        $|x|+|y|\le2$, which is an intersection of four half-planes.
        Only (e), the annulus, fails: it is the disc of radius $3$ with a
        hole, and the segment from $(-2,0)$ to $(2,0)$ runs straight
        through the hole.
  \item A diamond with vertices $(3,0)$, $(0,3)$, $(-3,0)$, $(0,-3)$: the
        square of side $3\sqrt2$ rotated by $45^{\circ}$.
  \item The wedge above both lines $y=x$ and $y=-x$ --- everything at or
        above the V of $y=|x|$, opening upwards with its vertex at the
        origin. Convex.
  \item Squeezed between $y=4-x^{2}$ and $y=-(4-x^{2})$, so it exists only
        where $4-x^{2}\ge0$, i.e. for $-2\le x\le2$. It is the
        lens-shaped region symmetric about both axes.
  \item With absolute values: $\max(|x|,|y|)\le2$, i.e. $|x|\le2$ and
        $|y|\le2$. Without: $-2\le x\le2$ and $-2\le y\le2$, four
        half-planes.
  \item The system describes a \emph{larger} set than intended --- a
        condition is missing. Every boundary of the intended region must
        contribute one inequality; a point outside that satisfies them all
        proves one boundary has been forgotten.
  \item Take two points diametrically opposite each other just outside the
        circle: the segment between them passes through the centre of the
        disc, which is not in the set. So the segment leaves the set and
        convexity fails.
  \item For instance $y\ge x+1$ and $y\le x-1$. Each is a satisfiable
        half-plane, but together they demand $x+1\le y\le x-1$, impossible
        since $x+1>x-1$. The two half-planes are on opposite sides of
        parallel lines.
  \item Let $C_{1},C_{2}$ be convex and take $P,Q\in C_{1}\cap C_{2}$.
        Since $P,Q\in C_{1}$ and $C_{1}$ is convex, the whole segment
        $[P,Q]$ lies in $C_{1}$; the same argument puts it in $C_{2}$.
        Hence $[P,Q]\subseteq C_{1}\cap C_{2}$, which is the definition of
        convexity. (The argument works for any number of sets.)
  \item No. The points $(0,1)$ and $(0,-1)$ both lie in the union, but the
        midpoint $(0,0)$ lies in neither half-plane. A union of disjoint
        convex sets is essentially never convex.
  \item $x^{2}+y^{2}<4$ \;and\; $y>x$. Both boundaries are dashed: the open
        half-disc lying above the diagonal. (It is convex, being an
        intersection of two convex sets.)
  \item Empty for $c<0$; the single point $(0,0)$ for $c=0$; the closed disc
        of radius $\sqrt{c}$ for $c>0$. The three cases are exactly the
        three things a level \emph{set} of $x^{2}+y^{2}$ can do, and they
        are why a sublevel set can be empty without the inequality being
        malformed.
  \item $3x+4y\le60$, \;$x\ge0$, \;$y\ge0$: the triangle with vertices
        $(0,0)$, $(20,0)$ and $(0,15)$.
  \item $3x+4y\le60$, \;$xy\ge12$, \;$x,y>0$: the part of the budget
        triangle lying beyond the indifference curve $y=12/x$. It is the
        intersection of a convex set (the triangle) with another convex set
        (the better set of a quasiconcave utility), hence \textbf{convex}.
  \item With $a$ units of A and $b$ of B: $2a+b\le40$ (labour),
        $a+2b\le30$ (capital), $a\ge0$, $b\ge0$. Four half-planes, so a
        convex polygon.
  \item Vertices: $(0,0)$, $(20,0)$, $(0,15)$, and the intersection of
        $2a+b=40$ with $a+2b=30$. Solving: $b=40-2a$ and
        $a+80-4a=30$, so $a=\tfrac{50}{3}$ and $b=\tfrac{20}{3}$. A linear
        objective has parallel straight level lines; pushing them across a
        convex polygon, the last point of contact is a vertex (or a whole
        edge, if the level lines happen to be parallel to it). Hence
        checking the vertices suffices --- the fundamental theorem of
        linear programming.
  \item The two lines meet where $100-2q=20+q$, i.e. $q=\tfrac{80}{3}$ with
        $p=\tfrac{140}{3}$. The region is the triangle with vertices
        $(0,20)$, $(0,100)$ and
        $\bigl(\tfrac{80}{3},\tfrac{140}{3}\bigr)$, so the largest feasible
        quantity is $q=\tfrac{80}{3}\approx26.7$ --- beyond it no price
        satisfies both buyer and seller, which is exactly the competitive
        equilibrium quantity.
\end{solutions}

\section{Additional Problems}

\begin{problems}
  \item Sketch $3x-2y>6$ and say which side of the line contains the
        origin.
  \item Sketch $y\le -x+5$ together with $x\ge1$.
  \item Sketch $(x+1)^{2}+(y-2)^{2}<4$ and give centre and radius.
  \item Sketch $x^{2}+y^{2}\ge 1$ and $x^{2}+y^{2}\le 4$ together.
  \item Sketch $y\ge x^{2}-4$.
  \item Sketch $y\le 2^{x}$.
  \item Sketch $y\ge\ln x$ and state the domain restriction it forces.
  \item Sketch $xy\le 4$ on the first quadrant.
  \item Sketch $|x-1|\le 2$.
  \item Sketch $|x|+|y|\ge 1$ and decide convexity.
  \item Sketch $y>x^{2}$ and $y<x+6$, and find the $x$-range of the
        region.
  \item Sketch $x^{2}+4y^{2}\le 16$ and $y\ge 0$.
  \item Write inequalities for the closed triangle with vertices $(0,0)$,
        $(4,0)$, $(0,3)$.
  \item Write inequalities for the closed disc of radius $2$ centred at
        $(1,1)$, intersected with the first quadrant.
  \item Write inequalities for the region above $y=x^{2}$ and below
        $y=2x+3$, and give the $x$-range.
  \item Write inequalities for the horizontal strip between $y=-1$ and
        $y=3$, once with an absolute value and once without.
  \item Write inequalities for everything outside the ellipse
        $x^{2}/4+y^{2}=1$ but inside the circle $x^{2}+y^{2}=9$.
  \item Is the region $\{y\ge x^{2}\}\cap\{y\le 4\}$ bounded? Give its
        extreme $x$- and $y$-values.
  \item Give a system of three inequalities whose solution set is a single
        point.
  \item Give a system of three inequalities whose solution set is an
        unbounded convex region with exactly two vertices.
  \item Show that $\{(x,y): y\ge \mathrm{e}^{x}\}$ is convex.
  \item Show that $\{(x,y): y\ge \sqrt{x},\ x\ge0\}$ is \emph{not}
        convex, and say which nearby set \emph{is}.
  \item \textbf{Two constraints, two goods.} Sketch
        $x+y\le 10$, \;$2x+y\le 16$, \;$x,y\ge0$, and give all vertices.
  \item \textbf{Diet problem.} A feed must supply at least $8$ units of
        protein and $12$ of energy. Grain gives $2$ protein and $2$ energy
        per kg; soy gives $1$ protein and $4$ energy per kg. Write the
        feasible set, sketch it, and say why it is unbounded.
  \item \textbf{Minimum-wage wedge.} Sketch
        $\{(L,w): w\ge 10,\ w\le 30-L,\ L\ge0\}$ and interpret the region
        and its largest $L$.
\end{problems}

\section{Additional Solutions}

\begin{solutions}
  \item The open half-plane strictly below the line $y=\tfrac32x-3$.
        Testing $(0,0)$: $0>6$ is false, so the origin is \emph{not} in the
        region; it lies on the other side.
  \item The intersection of the half-plane on or below $y=5-x$ with the
        half-plane $x\ge1$: an unbounded wedge opening down-right, with
        corner $(1,4)$.
  \item The \emph{open} disc of radius $2$ centred at $(-1,2)$; the
        boundary circle is dashed and excluded.
  \item The closed annulus $1\le x^{2}+y^{2}\le4$ between the two circles.
        Not convex.
  \item The region on or above the parabola $y=x^{2}-4$, whose vertex is
        $(0,-4)$ and which crosses the axis at $x=\pm2$. Convex.
  \item Everything on or below the increasing curve $y=2^{x}$, which
        passes through $(0,1)$ and hugs the $x$-axis to the left. Since
        $2^{x}$ is \emph{convex}, the region below it is \textbf{not}
        convex: $(-3,\tfrac18)$ and $(3,8)$ both satisfy the inequality,
        but their midpoint $(0,4.0625)$ does not, since $2^{0}=1<4.0625$.
  \item The region on or above $y=\ln x$; the logarithm's domain confines
        it to the open right half-plane $x>0$. Since $\ln x$ is
        \emph{concave}, the region above it is \textbf{not} convex: the
        boundary points $(0.1,\ln 0.1)$ and $(10,\ln 10)$ have midpoint
        $(5.05,0)$, and $0<\ln 5.05\approx1.62$, so the midpoint is
        outside.
  \item The part of the first quadrant on or below the hyperbola $y=4/x$
        --- the sublevel set $\{xy\le4\}$. \textbf{Not} convex: $(0.5,8)$
        and $(8,0.5)$ both lie on the boundary, but their midpoint
        $(4.25,4.25)$ has $xy\approx18.1>4$. (Equivalently: $4/x$ is convex
        on $x>0$, so the region below it is not.)
  \item $-2\le x-1\le2$, i.e. the closed vertical strip $-1\le x\le3$.
        Convex.
  \item Everything on or outside the diamond $|x|+|y|=1$. Not convex ---
        the segment from $(1,0)$ to $(-1,0)$ passes through the origin,
        where $|x|+|y|=0<1$.
  \item Strictly between the parabola and the line; they meet where
        $x^{2}=x+6$, i.e. $x=-2$ and $x=3$, so the region spans
        $-2<x<3$. Both boundaries dashed.
  \item $\dfrac{x^{2}}{16}+\dfrac{y^{2}}{4}\le1$ with $y\ge0$: the upper
        half of an elliptical disc with semi-axes $4$ and $2$.
  \item $x\ge0$, \;$y\ge0$, \;$3x+4y\le12$ (the line through $(4,0)$ and
        $(0,3)$).
  \item $(x-1)^{2}+(y-1)^{2}\le4$, \;$x\ge0$, \;$y\ge0$.
  \item $y\ge x^{2}$ and $y\le 2x+3$; the boundaries meet where
        $x^{2}=2x+3$, i.e. $x=-1$ and $x=3$, so $-1\le x\le3$.
  \item With an absolute value: $|y-1|\le2$. Without: $y\ge-1$ and
        $y\le3$.
  \item $\dfrac{x^{2}}{4}+y^{2}\ge1$ \;and\; $x^{2}+y^{2}\le9$: the disc of
        radius $3$ with the elliptical disc punched out. Not convex.
  \item Yes, bounded: $y$ ranges over $[0,4]$ and $x$ over $[-2,2]$, since
        $y\ge x^{2}$ together with $y\le4$ forces $x^{2}\le4$. The extreme
        points are $(\pm2,4)$ and $(0,0)$.
  \item For example $x\ge0$, \;$y\ge0$, \;$x+y\le0$: the first two force
        $x,y\ge0$ and the third then forces $x=y=0$, so the solution set is
        exactly $\{(0,0)\}$.
  \item $x\ge0$, \;$y\ge0$, \;$x+y\ge2$. The boundary consists of the
        ray $\{x=0,\ y\ge2\}$, the segment from $(0,2)$ to $(2,0)$, and the
        ray $\{y=0,\ x\ge2\}$, so there are exactly two corners, $(0,2)$
        and $(2,0)$. The region is unbounded (it contains every large
        point of the quadrant) and convex, being an intersection of three
        half-planes.
  \item $\mathrm{e}^{x}$ is a convex function, and the set lying \emph{on
        or above} the graph of a convex function --- its epigraph --- is a
        convex set: if $y_{1}\ge\mathrm{e}^{x_{1}}$ and
        $y_{2}\ge\mathrm{e}^{x_{2}}$ then for $t\in[0,1]$
        \[
          ty_{1}+(1-t)y_{2}\ \ge\ t\mathrm{e}^{x_{1}}+(1-t)\mathrm{e}^{x_{2}}
          \ \ge\ \mathrm{e}^{tx_{1}+(1-t)x_{2}},
        \]
        the last step being exactly convexity of $\mathrm{e}^{x}$.
  \item $\sqrt{x}$ is \emph{concave}, so the region \emph{above} it is
        not convex. Concretely, $(0,0)$ and $(4,2)$ both lie on the
        boundary and so belong to the set, but their midpoint $(2,1)$ does
        not, since $\sqrt2\approx1.414>1$. The set that \emph{is} convex is
        the one below the curve, $\{y\le\sqrt{x},\ x\ge0\}$ --- the
        hypograph of a concave function.
  \item The lines meet where $x+y=10$ and $2x+y=16$, i.e. $x=6$, $y=4$.
        Vertices: $(0,0)$, $(8,0)$, $(6,4)$, $(0,10)$.
  \item With $g$ kg of grain and $s$ kg of soy: $2g+s\ge8$,
        $2g+4s\ge12$, $g\ge0$, $s\ge0$. It is unbounded because the
        constraints are all ``at least'' --- buying arbitrarily much feed is
        always still feasible, so the region extends indefinitely up and to
        the right. A cost-minimising objective is nevertheless attained,
        at one of the two vertices on the lower-left boundary.
  \item The region is the triangle bounded below by $w=10$, above by
        $w=30-L$, and on the left by $L=0$; the two sloping boundaries
        meet at $L=20$, $w=10$. It shows the employment levels compatible
        with a wage floor of $10$ when labour demand is $w=30-L$: the
        largest feasible $L$ is $20$, the employment the floor permits.
        Any $L$ beyond that would require a wage below the minimum.
\end{solutions}

\chapter{Trigonometric Functions}

Trigonometry plays a small part in economics and a large part in the
prerequisites, for two reasons. The graphs of $\sin$, $\cos$ and $\tan$
belong to the catalogue of shapes one is expected to recognise, and they are
the standard test functions of calculus: many derivative and limit
exercises use them precisely because they oscillate and so expose careless
reasoning. Seasonality and business cycles supply the
occasional genuine application.

The chapter therefore covers what is actually needed and stops. The unit
circle and the handful of exact values; the three graphs with their periods,
amplitudes and asymptotes; the transformation form
$a\sin\bigl(b(x-c)\bigr)+d$, which is the machinery of \chref{functions}{3} applied to a
periodic function; the four identities that carry all the simplification
work; and how to write down \emph{every} solution of a trigonometric
equation rather than just the one a calculator returns.

\section{Core Rules}

\begin{rulebox}{Radians, and why}
An angle in \textbf{radians} is the arc length it cuts on the unit circle,
so a full turn is $2\pi$:
\[
  180^{\circ}=\pi,\qquad
  x_{\text{rad}}=x_{\deg}\cdot\frac{\pi}{180},\qquad
  x_{\deg}=x_{\text{rad}}\cdot\frac{180}{\pi} .
\]
Calculus uses radians exclusively, and for one reason: only in radians is
$\dfrac{\mathrm{d}}{\mathrm{d}x}\sin x=\cos x$. In degrees an unwanted
factor $\pi/180$ appears in every derivative.
\end{rulebox}

\begin{rulebox}{The unit circle definition}
For the point on the unit circle reached by turning through the angle $x$
from the positive $x$-axis, $\cos x$ is its first coordinate and $\sin x$
its second. Hence, immediately:
\[
  \sin^{2}x+\cos^{2}x=1,\qquad
  -1\le\sin x\le1,\qquad -1\le\cos x\le1,
\]
the first being Pythagoras. Also $\tan x=\dfrac{\sin x}{\cos x}$, the slope
of the ray, undefined where $\cos x=0$.

Signs by quadrant --- $\sin$ is the height, $\cos$ the width:
\[
  \begin{array}{c|cccc}
     & \text{I} & \text{II} & \text{III} & \text{IV}\\\hline
    \sin & + & + & - & -\\
    \cos & + & - & - & +\\
    \tan & + & - & + & -
  \end{array}
\]
\end{rulebox}

\begin{rulebox}{The values to know by heart}
\[
  \begin{array}{c|ccccc}
    x & 0 & \pi/6 & \pi/4 & \pi/3 & \pi/2\\\hline
    \sin x & 0 & \tfrac12 & \tfrac{\sqrt2}{2} & \tfrac{\sqrt3}{2} & 1\\[0.3ex]
    \cos x & 1 & \tfrac{\sqrt3}{2} & \tfrac{\sqrt2}{2} & \tfrac12 & 0\\[0.3ex]
    \tan x & 0 & \tfrac{1}{\sqrt3} & 1 & \sqrt3 & \text{---}
  \end{array}
\]
Everything else follows from three facts: $\sin$ is \textbf{odd}
($\sin(-x)=-\sin x$), $\cos$ is \textbf{even} ($\cos(-x)=\cos x$), and both
have period $2\pi$ while $\tan$ has period $\pi$. So
$\sin\tfrac{5\pi}{6}=\sin\tfrac{\pi}{6}=\tfrac12$ (second quadrant, $\sin$
positive) and $\cos\tfrac{4\pi}{3}=-\cos\tfrac{\pi}{3}=-\tfrac12$ (third
quadrant, $\cos$ negative).
\end{rulebox}

\begin{gallery}
\begin{tikzpicture}
  \begin{axis}[fnaxis, xmin=-0.6, xmax=7.0, ymin=-1.6, ymax=1.6,
      xtick={3.1416, 6.2832}, xticklabels={$\pi$, $2\pi$},
      ytick={-1,1}, title={$y=\sin x$}]
    \addplot[domain=-0.4:6.9, samples=201]{sin(deg(x))};
  \end{axis}
\end{tikzpicture}\hfill
\begin{tikzpicture}
  \begin{axis}[fnaxis, xmin=-0.6, xmax=7.0, ymin=-1.6, ymax=1.6,
      xtick={3.1416, 6.2832}, xticklabels={$\pi$, $2\pi$},
      ytick={-1,1}, title={$y=\cos x$}]
    \addplot[domain=-0.4:6.9, samples=201]{cos(deg(x))};
  \end{axis}
\end{tikzpicture}\hfill
\begin{tikzpicture}
  \begin{axis}[fnaxis, xmin=-3.4, xmax=3.4, ymin=-3.4, ymax=3.4,
      xtick={-1.5708, 1.5708}, xticklabels={$-\frac{\pi}{2}$,$\frac{\pi}{2}$},
      ytick={-2,2}, title={$y=\tan x$}]
    \addplot[domain=-1.32:1.32, samples=101]{tan(deg(x))};
    \addplot[domain=1.82:3.3, samples=101]{tan(deg(x))};
    \addplot[domain=-3.3:-1.82, samples=101]{tan(deg(x))};
    \addplot[gray!55, dashed, thin, samples=2, forget plot]
      coordinates {(1.5708,-3.3)(1.5708,3.3)};
    \addplot[gray!55, dashed, thin, samples=2, forget plot]
      coordinates {(-1.5708,-3.3)(-1.5708,3.3)};
  \end{axis}
\end{tikzpicture}
\gcap{$\sin$ and $\cos$ are the same wave shifted by $\pi/2$, with period
$2\pi$ and range $[-1,1]$. $\tan$ has period $\pi$, range $\R$, and a
vertical asymptote at every odd multiple of $\pi/2$.}
\end{gallery}

\begin{rulebox}{Reading $a\sin\bigl(b(x-c)\bigr)+d$}
The four constants are exactly the four moves of \chref{functions}{3}, and each has a
name here:
\begin{itemize}[leftmargin=1.6em,itemsep=0.3ex]
  \item $|a|$ is the \textbf{amplitude}: the wave runs from $d-|a|$ to
        $d+|a|$, so the range is $[d-|a|,\,d+|a|]$. Negative $a$ flips it.
  \item $b$ sets the \textbf{period}, $T=\dfrac{2\pi}{|b|}$ --- larger $b$
        means faster oscillation. (For $\tan$, $T=\pi/|b|$.)
  \item $c$ is the \textbf{phase shift}: the whole wave moves $c$ to the
        right. Factor before reading: $\sin(2x-\pi)=\sin\bigl(2(x-\pi/2)\bigr)$
        is a shift of $\pi/2$, not of $\pi$.
  \item $d$ is the \textbf{midline}, the horizontal line the wave
        oscillates about.
\end{itemize}
So $y=3\sin\bigl(2(x-\tfrac{\pi}{4})\bigr)+1$ has amplitude $3$, period
$\pi$, midline $y=1$, range $[-2,4]$, shifted $\pi/4$ right.
\end{rulebox}

\begin{rulebox}{The identities that do the work}
\[
  \sin^{2}x+\cos^{2}x=1,\qquad
  \tan x=\frac{\sin x}{\cos x},\qquad
  1+\tan^{2}x=\frac{1}{\cos^{2}x} .
\]
\[
  \sin(x\pm y)=\sin x\cos y\pm\cos x\sin y,\qquad
  \cos(x\pm y)=\cos x\cos y\mp\sin x\sin y .
\]
Setting $y=x$ in the addition formulas gives the \textbf{double-angle}
forms, which are the ones that appear in practice:
\[
  \sin 2x=2\sin x\cos x,\qquad
  \cos 2x=\cos^{2}x-\sin^{2}x=1-2\sin^{2}x=2\cos^{2}x-1 .
\]
Almost every simplification is one of: replace $1-\cos^{2}x$ by
$\sin^{2}x$, split a $\tan$ into $\sin/\cos$, or collapse a double angle.
\end{rulebox}

\begin{rulebox}{Solving trigonometric equations: give \emph{all} solutions}
A calculator returns one angle; periodicity supplies infinitely many.
\begin{itemize}[leftmargin=1.6em,itemsep=0.35ex]
  \item $\sin x=s$ with $|s|\le1$: two families,
        $x=x_{0}+2k\pi$ and $x=\pi-x_{0}+2k\pi$, where
        $x_{0}=\arcsin s$ and $k\in\mathbb{Z}$.
  \item $\cos x=s$: $x=\pm x_{0}+2k\pi$ with $x_{0}=\arccos s$.
  \item $\tan x=s$: one family, $x=\arctan s+k\pi$, since the period is
        $\pi$.
  \item If the argument is $bx$ rather than $x$, solve for $bx$ first and
        divide \emph{everything}, the $2k\pi$ included: $\sin 2x=1$ gives
        $2x=\tfrac{\pi}{2}+2k\pi$, hence
        $x=\tfrac{\pi}{4}+k\pi$.
\end{itemize}
\end{rulebox}

\begin{rulebox}{Inverse trigonometric functions}
None of the three is injective, so each is inverted only on a restricted
domain:
\[
  \arcsin:[-1,1]\to\Bigl[-\tfrac{\pi}{2},\tfrac{\pi}{2}\Bigr],\qquad
  \arccos:[-1,1]\to[0,\pi],\qquad
  \arctan:\R\to\Bigl(-\tfrac{\pi}{2},\tfrac{\pi}{2}\Bigr).
\]
So $\arcsin(\sin x)=x$ only for $x$ in that range: $\arcsin(\sin\pi)=0$,
not $\pi$. And $\arctan$ is bounded, with horizontal asymptotes
$y=\pm\pi/2$ --- the reason it appears as a saturating S-shaped curve.
\end{rulebox}

\section{Worked Example}

\begin{worked}{Solve $2\sin^{2}x-\sin x=0$ on $[0,2\pi]$ and in general, and read off the shape of $y=3\sin\!\bigl(2(x-\tfrac{\pi}{4})\bigr)+1$}
\wstep{Step 1. Factor. Do not divide.}
\[
  2\sin^{2}x-\sin x=0
  \quad\Longrightarrow\quad
  \sin x\,(2\sin x-1)=0 .
\]
\wtrap{Dividing both sides by $\sin x$ gives $2\sin x-1=0$ and loses every
solution with $\sin x=0$ --- three of the five on $[0,2\pi]$. You may divide
by a quantity only when you know it is non-zero, and $\sin x$ is exactly the
quantity in question. Factor and use the zero-product property instead.}

\wstep{Step 2. Solve each factor on the given interval.}
\emph{First factor:} $\sin x=0$ on $[0,2\pi]$ gives $x=0,\ \pi,\ 2\pi$.

\emph{Second factor:} $\sin x=\tfrac12$. The reference angle is
$\pi/6$, and the sine is positive in the first and second quadrants, so
\[
  x=\frac{\pi}{6}\qquad\text{and}\qquad x=\pi-\frac{\pi}{6}=\frac{5\pi}{6}.
\]
\wtrap{$\arcsin(\tfrac12)=\pi/6$ is one solution, not the solution. The
inverse sine returns a single value from $[-\pi/2,\pi/2]$; the second
solution in each period must be supplied by hand from the symmetry
$\sin(\pi-x)=\sin x$. For cosine the partner is $-x$ (or $2\pi-x$), for
tangent it is $x+\pi$.}

\wstep{Step 3. Collect the five solutions.}
\[
  \boxed{\;x\in\Bigl\{0,\ \tfrac{\pi}{6},\ \tfrac{5\pi}{6},\ \pi,\ 2\pi\Bigr\}\;}
\]
Check the smallest and one middle value: at $x=0$, $0-0=0$ \checkmark; at
$x=\pi/6$, $2\cdot\tfrac14-\tfrac12=0$ \checkmark.

\wstep{Step 4. Give the general solution, since ``all solutions'' usually means all of $\R$.}
\[
  x=k\pi,
  \qquad
  x=\frac{\pi}{6}+2k\pi,
  \qquad
  x=\frac{5\pi}{6}+2k\pi,
  \qquad k\in\mathbb{Z} .
\]
Note the two different periods: the zeros of $\sin$ repeat every $\pi$
(because $\sin(x+\pi)=-\sin x$, which is also zero), while the solutions of
$\sin x=\tfrac12$ repeat only every $2\pi$.
\wtrap{Adding $2k\pi$ to everything is safe but incomplete for the first
family; writing $k\pi$ for the second family is simply wrong. Work out the
period of each family separately from the equation it solves.}

\wstep{Step 5. Now the sinusoid: read the four parameters off $a\sin\bigl(b(x-c)\bigr)+d$.}
For $y=3\sin\bigl(2(x-\tfrac{\pi}{4})\bigr)+1$:
\begin{center}\footnotesize
\begin{tabular}{llll}
  \hline
  parameter & value & name & effect\\
  \hline
  $a$ & $3$ & amplitude & oscillates $3$ above and below the midline\\
  $b$ & $2$ & angular frequency & period $=2\pi/b=\pi$\\
  $c$ & $\pi/4$ & phase shift & shifted \emph{right} by $\pi/4$\\
  $d$ & $1$ & midline & centred on $y=1$, so range $[-2,4]$\\
  \hline
\end{tabular}
\end{center}

\wstep{Step 6. Sketch it from those four numbers.}
\begin{center}
\begin{tikzpicture}
\begin{axis}[fnaxis, width=8.2cm, height=4.4cm,
             domain=-0.4:3.6, samples=201,
             xmin=-0.5, xmax=3.7, ymin=-2.8, ymax=4.8,
             xtick={0.7854, 1.5708, 2.3562, 3.1416},
             xticklabels={$\frac{\pi}{4}$,$\frac{\pi}{2}$,$\frac{3\pi}{4}$,$\pi$},
             ytick={-2,1,4}]
  \addplot[accent!35, thick, dashed, samples=2, forget plot] {1};
  \addplot {3*sin(deg(2*(x-pi/4)))+1};
\end{axis}
\end{tikzpicture}
\end{center}
The curve crosses its midline $y=1$ going upward at $x=\pi/4$ --- that is
what the phase shift means --- peaks at $4$ a quarter-period later
($x=\pi/4+\pi/4=\pi/2$), and repeats after one full period $\pi$.
\wtrap{The period is $2\pi/b$, not $2\pi b$: a \emph{larger} $b$ gives a
\emph{shorter} period. And the phase shift is only readable once the inside
is factored as $b(x-c)$; from $\sin(2x-\tfrac{\pi}{2})$ the shift is
$\pi/4$, not $\pi/2$. It is the same trap as in \chref{functions}{3} on
transformations.}
\end{worked}

\section{Practice Problems}

Convert in problems 1--6.

\begin{problems}
  \item $60^{\circ}$ into radians.
  \item $270^{\circ}$ into radians.
  \item $\tfrac{3\pi}{4}$ into degrees.
  \item $\tfrac{5\pi}{3}$ into degrees.
  \item $1$ radian into degrees, to one decimal.
  \item Why does calculus insist on radians?
\end{problems}

Give the exact value in problems 7--16.

\begin{problems}[resume]
  \item $\sin\tfrac{\pi}{6}$
  \item $\cos\tfrac{\pi}{4}$
  \item $\tan\tfrac{\pi}{3}$
  \item $\sin\tfrac{5\pi}{6}$
  \item $\cos\tfrac{2\pi}{3}$
  \item $\tan\pi$
  \item $\sin\bigl(-\tfrac{\pi}{3}\bigr)$
  \item $\cos\tfrac{7\pi}{2}$
  \item $\sin\tfrac{9\pi}{4}$
  \item $\tan\tfrac{3\pi}{4}$
\end{problems}

In problems 17--24 give period, amplitude (where defined), range and any
asymptotes.

\begin{problems}[resume]
  \item $y=\sin x$
  \item $y=\cos x$
  \item $y=\tan x$
  \item $y=3\sin x$
  \item $y=\sin 2x$
  \item $y=\cos\tfrac{x}{2}$
  \item $y=2\sin x+1$
  \item $y=\tan 3x$
\end{problems}

\begin{problems}[resume]
  \item Describe $y=\sin\bigl(x-\tfrac{\pi}{2}\bigr)$ as a shift, and name
        the simpler function it equals.
  \item Give amplitude, period, midline and range of
        $y=3\sin\bigl(2(x-\tfrac{\pi}{4})\bigr)+1$.
  \item Give amplitude, period and phase shift of
        $y=\sin(2x-\pi)$ --- factor first.
  \item Give amplitude, period and range of $y=-2\cos 3x$.
  \item Write a sine function with amplitude $4$, period $\pi$ and midline
        $y=-1$.
  \item Write a cosine function with amplitude $2$, period $6\pi$ and no
        vertical shift.
  \item Monthly sales oscillate about $500$ units with a swing of $\pm80$
        and a twelve-month cycle, peaking in month $3$. Write $S(t)$.
  \item Why can no function of the form $a\sin(bx)+d$ ever be injective on
        $\R$?
\end{problems}

Simplify or verify in problems 33--38.

\begin{problems}[resume]
  \item $\dfrac{1-\cos^{2}x}{\sin x}$
  \item $\tan x\cdot\cos x$
  \item $\sin^{2}x-1$
  \item $\dfrac{\sin 2x}{2\cos x}$
  \item Verify $1+\tan^{2}x=\dfrac{1}{\cos^{2}x}$.
  \item Verify $\cos 2x=1-2\sin^{2}x$ from the addition formula.
\end{problems}

Give \emph{all} real solutions in problems 39--45.

\begin{problems}[resume]
  \item $\sin x=\tfrac12$
  \item $\cos x=0$
  \item $\tan x=1$
  \item $2\cos x+1=0$
  \item $\sin 2x=1$
  \item $\cos^{2}x=\tfrac14$
  \item $\sin x=\cos x$
\end{problems}

\section{Solutions}

\begin{solutions}
  \item $60\cdot\tfrac{\pi}{180}=\tfrac{\pi}{3}$.
  \item $270\cdot\tfrac{\pi}{180}=\tfrac{3\pi}{2}$.
  \item $\tfrac{3\pi}{4}\cdot\tfrac{180}{\pi}=135^{\circ}$.
  \item $\tfrac{5\pi}{3}\cdot\tfrac{180}{\pi}=300^{\circ}$.
  \item $\tfrac{180}{\pi}\approx57.3^{\circ}$.
  \item Because $\tfrac{\mathrm{d}}{\mathrm{d}x}\sin x=\cos x$ holds only
        in radians. In degrees the limit
        $\lim_{h\to0}\tfrac{\sin h}{h}$ equals $\tfrac{\pi}{180}$ instead
        of $1$, and that factor would appear in every derivative and integral
        of a trigonometric function.
  \item $\tfrac12$.
  \item $\tfrac{\sqrt2}{2}$.
  \item $\sqrt3$.
  \item Second quadrant, $\sin$ positive:
        $\sin\tfrac{5\pi}{6}=\sin\tfrac{\pi}{6}=\tfrac12$.
  \item Second quadrant, $\cos$ negative:
        $-\cos\tfrac{\pi}{3}=-\tfrac12$.
  \item $\tan\pi=\tfrac{\sin\pi}{\cos\pi}=\tfrac{0}{-1}=0$.
  \item $\sin$ is odd: $-\sin\tfrac{\pi}{3}=-\tfrac{\sqrt3}{2}$.
  \item $\tfrac{7\pi}{2}=\tfrac{3\pi}{2}+2\pi$, so
        $\cos\tfrac{7\pi}{2}=\cos\tfrac{3\pi}{2}=0$.
  \item $\tfrac{9\pi}{4}=\tfrac{\pi}{4}+2\pi$, so the value is
        $\sin\tfrac{\pi}{4}=\tfrac{\sqrt2}{2}$.
  \item Second quadrant, $\tan$ negative:
        $-\tan\tfrac{\pi}{4}=-1$.
  \item Period $2\pi$, amplitude $1$, range $[-1,1]$, no asymptotes.
  \item Period $2\pi$, amplitude $1$, range $[-1,1]$, no asymptotes.
  \item Period $\pi$, amplitude undefined (unbounded), range $\R$,
        vertical asymptotes at $x=\tfrac{\pi}{2}+k\pi$.
  \item Period $2\pi$, amplitude $3$, range $[-3,3]$.
  \item Period $\tfrac{2\pi}{2}=\pi$, amplitude $1$, range $[-1,1]$.
  \item Period $\tfrac{2\pi}{1/2}=4\pi$, amplitude $1$, range $[-1,1]$.
  \item Period $2\pi$, amplitude $2$, midline $y=1$, range $[-1,3]$.
  \item Period $\tfrac{\pi}{3}$, unbounded range $\R$, asymptotes where
        $3x=\tfrac{\pi}{2}+k\pi$, i.e. $x=\tfrac{\pi}{6}+\tfrac{k\pi}{3}$.
  \item A shift of $\tfrac{\pi}{2}$ to the right. It equals $-\cos x$,
        since $\sin\bigl(x-\tfrac{\pi}{2}\bigr)
        =\sin x\cos\tfrac{\pi}{2}-\cos x\sin\tfrac{\pi}{2}=-\cos x$.
  \item Amplitude $3$, period $\tfrac{2\pi}{2}=\pi$, midline $y=1$, range
        $[-2,4]$, shifted $\tfrac{\pi}{4}$ to the right.
  \item $\sin(2x-\pi)=\sin\bigl(2(x-\tfrac{\pi}{2})\bigr)$: amplitude $1$,
        period $\pi$, phase shift $\tfrac{\pi}{2}$ right. Reading the
        shift as $\pi$ from the unfactored form is the standard error.
  \item Amplitude $2$ (the minus sign reflects rather than changing the
        amplitude), period $\tfrac{2\pi}{3}$, range $[-2,2]$.
  \item Period $\pi$ needs $b=2$: $y=4\sin(2x)-1$.
  \item Period $6\pi$ needs $b=\tfrac13$:
        $y=2\cos\bigl(\tfrac{x}{3}\bigr)$.
  \item Midline $500$, amplitude $80$, period $12$ so
        $b=\tfrac{2\pi}{12}=\tfrac{\pi}{6}$, and a sine peaks a quarter
        period after its shift, so peaking at $t=3$ means $c=0$:
        \[
          S(t)=500+80\sin\!\left(\frac{\pi t}{6}\right).
        \]
        (Check: at $t=3$ the sine is $\sin\tfrac{\pi}{2}=1$, giving the
        peak $580$.)
  \item Because it is periodic: $f(x+T)=f(x)$ for $T=2\pi/|b|\neq0$, so two
        different inputs share a value. Injectivity fails everywhere, which
        is why the inverse functions need restricted domains.
  \item $\dfrac{\sin^{2}x}{\sin x}=\sin x$ (for $\sin x\neq0$).
  \item $\dfrac{\sin x}{\cos x}\cdot\cos x=\sin x$ (for
        $\cos x\neq0$).
  \item $-(1-\sin^{2}x)=-\cos^{2}x$.
  \item $\dfrac{2\sin x\cos x}{2\cos x}=\sin x$ (for $\cos x\neq0$).
  \item Divide $\sin^{2}x+\cos^{2}x=1$ by $\cos^{2}x$:
        $\tan^{2}x+1=\dfrac{1}{\cos^{2}x}$.
  \item $\cos2x=\cos^{2}x-\sin^{2}x$, and substituting
        $\cos^{2}x=1-\sin^{2}x$ gives $1-\sin^{2}x-\sin^{2}x
        =1-2\sin^{2}x$.
  \item $x=\tfrac{\pi}{6}+2k\pi$ or $x=\tfrac{5\pi}{6}+2k\pi$,
        $k\in\mathbb{Z}$.
  \item $x=\tfrac{\pi}{2}+k\pi$.
  \item $x=\tfrac{\pi}{4}+k\pi$.
  \item $\cos x=-\tfrac12$, so
        $x=\pm\tfrac{2\pi}{3}+2k\pi$.
  \item $2x=\tfrac{\pi}{2}+2k\pi$, hence
        $x=\tfrac{\pi}{4}+k\pi$.
  \item $\cos x=\pm\tfrac12$, giving
        $x=\pm\tfrac{\pi}{3}+2k\pi$ and $x=\pm\tfrac{2\pi}{3}+2k\pi$ ---
        more compactly $x=\tfrac{\pi}{3}+\tfrac{k\pi}{1}$ fails, so keep
        the four families, or write
        $x=\pm\tfrac{\pi}{3}+k\pi$.
  \item Dividing by $\cos x$ (legitimate, since $\cos x=0$ would force
        $\sin x=0$ too, impossible): $\tan x=1$, so
        $x=\tfrac{\pi}{4}+k\pi$.
\end{solutions}

\section{Additional Problems}

\begin{problems}
  \item $45^{\circ}$ and $150^{\circ}$ into radians.
  \item $\tfrac{7\pi}{6}$ and $\tfrac{11\pi}{6}$ into degrees.
  \item $\cos\tfrac{5\pi}{4}$
  \item $\sin\tfrac{4\pi}{3}$
  \item $\tan\tfrac{5\pi}{6}$
  \item $\sin\tfrac{11\pi}{6}$
  \item $\cos(-\tfrac{\pi}{6})$
  \item $\tan\tfrac{7\pi}{4}$
  \item Period, amplitude and range of $y=\tfrac12\cos 4x-2$.
  \item Period and asymptotes of $y=\tan\bigl(x-\tfrac{\pi}{4}\bigr)$.
  \item Amplitude, period, midline of
        $y=5\cos\bigl(\tfrac{\pi}{3}(x-2)\bigr)+3$.
  \item Write a cosine with maximum $7$, minimum $1$ and period $4$.
  \item Simplify $\dfrac{\cos^{2}x-1}{\cos x-1}$.
  \item Simplify $\sin x\cos x\cdot\dfrac{2}{\sin 2x}$.
  \item Simplify $\dfrac{1}{\cos^{2}x}-\tan^{2}x$.
  \item Verify $\sin(x+\pi)=-\sin x$.
  \item Verify $\cos\bigl(\tfrac{\pi}{2}-x\bigr)=\sin x$.
  \item Solve $\sin x=-\tfrac{\sqrt3}{2}$.
  \item Solve $\tan 2x=\sqrt3$.
  \item Solve $2\sin^{2}x-\sin x=0$.
  \item Solve $\cos 2x=\tfrac12$.
  \item Solve $\sin x=2$ and explain the answer.
  \item Give $\arcsin\tfrac12$, $\arccos(-1)$ and $\arctan(-1)$.
  \item Evaluate $\arcsin\bigl(\sin\tfrac{3\pi}{4}\bigr)$ and explain why
        the answer is not $\tfrac{3\pi}{4}$.
  \item Quarterly demand is $D(t)=200-40\cos\bigl(\tfrac{\pi t}{2}\bigr)$
        with $t$ in quarters. Give the period, the maximum and minimum
        demand, and the quarters in which each occurs.
\end{problems}

\section{Additional Solutions}

\begin{solutions}
  \item $\tfrac{\pi}{4}$ and $\tfrac{5\pi}{6}$.
  \item $210^{\circ}$ and $330^{\circ}$.
  \item Third quadrant, $\cos$ negative:
        $-\tfrac{\sqrt2}{2}$.
  \item Third quadrant, $\sin$ negative:
        $-\tfrac{\sqrt3}{2}$.
  \item Second quadrant, $\tan$ negative:
        $-\tfrac{1}{\sqrt3}$.
  \item Fourth quadrant, $\sin$ negative:
        $-\tfrac12$.
  \item $\cos$ is even: $\cos\tfrac{\pi}{6}=\tfrac{\sqrt3}{2}$.
  \item Fourth quadrant, $\tan$ negative: $-1$.
  \item Period $\tfrac{2\pi}{4}=\tfrac{\pi}{2}$, amplitude $\tfrac12$,
        midline $y=-2$, range $[-2.5,-1.5]$.
  \item Period $\pi$; asymptotes where
        $x-\tfrac{\pi}{4}=\tfrac{\pi}{2}+k\pi$, i.e.
        $x=\tfrac{3\pi}{4}+k\pi$.
  \item Amplitude $5$, period $\tfrac{2\pi}{\pi/3}=6$, midline $y=3$,
        shifted $2$ to the right; range $[-2,8]$.
  \item Midline $\tfrac{7+1}{2}=4$, amplitude $\tfrac{7-1}{2}=3$, period
        $4$ so $b=\tfrac{2\pi}{4}=\tfrac{\pi}{2}$:
        $y=3\cos\bigl(\tfrac{\pi x}{2}\bigr)+4$.
  \item $\dfrac{(\cos x-1)(\cos x+1)}{\cos x-1}=\cos x+1$ (for
        $\cos x\neq1$).
  \item $\dfrac{2\sin x\cos x}{\sin 2x}=\dfrac{\sin 2x}{\sin 2x}=1$.
  \item $\dfrac{1}{\cos^{2}x}-\tan^{2}x=1$ by the identity
        $1+\tan^{2}x=1/\cos^{2}x$.
  \item $\sin(x+\pi)=\sin x\cos\pi+\cos x\sin\pi=-\sin x$.
  \item $\cos\bigl(\tfrac{\pi}{2}-x\bigr)
        =\cos\tfrac{\pi}{2}\cos x+\sin\tfrac{\pi}{2}\sin x=\sin x$.
  \item $x=-\tfrac{\pi}{3}+2k\pi$ or
        $x=\pi+\tfrac{\pi}{3}+2k\pi=\tfrac{4\pi}{3}+2k\pi$.
  \item $2x=\tfrac{\pi}{3}+k\pi$, so
        $x=\tfrac{\pi}{6}+\tfrac{k\pi}{2}$.
  \item Factor: $\sin x\,(2\sin x-1)=0$, so $\sin x=0$ giving $x=k\pi$, or
        $\sin x=\tfrac12$ giving $x=\tfrac{\pi}{6}+2k\pi$ and
        $x=\tfrac{5\pi}{6}+2k\pi$.
  \item $2x=\pm\tfrac{\pi}{3}+2k\pi$, so
        $x=\pm\tfrac{\pi}{6}+k\pi$.
  \item No solution: $\sin$ never exceeds $1$, since it is a coordinate of
        a point on the unit circle. An equation can be perfectly
        well-formed and still have an empty solution set.
  \item $\tfrac{\pi}{6}$, \;$\pi$, \;$-\tfrac{\pi}{4}$.
  \item $\sin\tfrac{3\pi}{4}=\tfrac{\sqrt2}{2}$, and
        $\arcsin\tfrac{\sqrt2}{2}=\tfrac{\pi}{4}$. The answer is
        $\tfrac{\pi}{4}$, not $\tfrac{3\pi}{4}$, because $\arcsin$ returns
        values in $[-\tfrac{\pi}{2},\tfrac{\pi}{2}]$ only:
        $\arcsin(\sin x)=x$ holds just on that interval.
  \item Period $\tfrac{2\pi}{\pi/2}=4$ quarters, i.e. one year. Maximum
        $240$ where $\cos\bigl(\tfrac{\pi t}{2}\bigr)=-1$, i.e.
        $\tfrac{\pi t}{2}=\pi+2k\pi$ so $t=2,6,10,\dots$; minimum $160$
        where the cosine is $1$, i.e. $t=0,4,8,\dots$. Demand peaks in the
        third quarter of each year and troughs in the first.
\end{solutions}


\end{document}
"
\ifdefined\KindleEdition
  \documentclass[10pt,oneside]{book}
\else
  \documentclass[11pt,a4paper,twoside,openright]{book}
\fi
% Shared preamble for the Functions training sheet.
\usepackage[T1]{fontenc}
\usepackage[utf8]{inputenc}
\usepackage{lmodern}
\usepackage[a4paper,margin=2.4cm]{geometry}
\usepackage{amsmath,amssymb,amsthm}
\usepackage{microtype}
\usepackage{enumitem}
\usepackage{xcolor}
\usepackage[most]{tcolorbox}
\usepackage{fancyhdr}
\usepackage[hidelinks]{hyperref}

\definecolor{accent}{HTML}{1F4E79}
\definecolor{soft}{HTML}{F4F7FA}

% --- function graphs ------------------------------------------------------
\usepackage{pgfplots}
\pgfplotsset{compat=1.18}
\usetikzlibrary{arrows.meta}

% One consistent look for every function graph in these books.
\pgfplotsset{
  fnaxis/.style={
    width=5.9cm, height=4.6cm,
    axis lines=middle,
    axis line style={-{Stealth[length=4pt]}, gray!65!black, line width=0.3pt},
    tick style={gray!65!black},
    tick label style={font=\tiny, gray!50!black},
    label style={font=\scriptsize},
    xlabel={$x$}, ylabel={$y$},
    every axis x label/.style={at={(ticklabel* cs:1.0)}, anchor=west},
    every axis y label/.style={at={(ticklabel* cs:1.0)}, anchor=south east,
      xshift=-1pt},
    every axis plot/.append style={accent, very thick, smooth, samples=201},
    clip=true, enlargelimits=false,
    title style={font=\scriptsize, yshift=4pt},
  },
  fnaxis/.append style={/pgfplots/legend style={font=\tiny, draw=none,
    fill=none, at={(0.98,0.03)}, anchor=south east}},
}

% A row of small graphs, evenly spaced with \hfill between them, plus an
% optional caption line: \begin{gallery} <pic>\hfill<pic> \gcap{...} \end{gallery}
\newenvironment{gallery}
  {\par\medskip\noindent\ignorespaces}
  {\par\medskip\noindent}
\newcommand{\gcap}[1]{\par\smallskip{\centering\footnotesize #1\par}}

% --- plane regions, implicit curves, surfaces -----------------------------
\usepgfplotslibrary{fillbetween}
\definecolor{regionfill}{HTML}{9DB4D8}
\pgfplotsset{colormap={house}{color=(soft) color=(accent!45) color=(accent)}}

\pgfplotsset{
  % Squared, gridded axis for plane regions and implicitly given curves.
  geoaxis/.style={
    fnaxis,
    width=5.9cm, height=5.9cm,
    grid=both,
    grid style={gray!22, line width=0.2pt},
    minor tick num=1,
  },
  % Level-curve maps: many thin curves, no single curve emphasised.
  contouraxis/.style={
    geoaxis,
    every axis plot/.append style={accent!75, line width=0.5pt, smooth,
      samples=201},
  },
  % Surfaces: \addplot3[surf, domain=..., y domain=...] {expr};
  surfaxis/.style={
    width=6.0cm, height=5.0cm,
    view={130}{22},
    colormap name=house,
    ticks=none,
    xlabel={$x$}, ylabel={$y$}, zlabel={$f$},
    label style={font=\tiny, gray!50!black},
    title style={font=\scriptsize},
    every axis plot/.append style={samples=25, shader=faceted interp,
      draw=accent!35, line width=0.1pt},
  },
}

% Shorthand for a shaded region between two named paths.
% \region{<path A>}{<path B>}{<soft clip options>}
\newcommand{\region}[3]{%
  \addplot[regionfill, opacity=0.55, forget plot, sharp plot]
    fill between[of=#1 and #2, #3];}

% --- rule boxes -----------------------------------------------------------
\newtcolorbox{rulebox}[1]{%
  breakable, enhanced, colback=soft, colframe=accent,
  boxrule=0.4pt, left=7pt, right=7pt, top=6pt, bottom=6pt,
  fonttitle=\bfseries, coltitle=white, colbacktitle=accent,
  title={#1}}

% --- worked examples ------------------------------------------------------
% One fully narrated solution per chapter, carrying the chapter's central
% technique and the mistakes that technique invites:
%   \begin{worked}{Task} ... \wstep{Step 1.} ... \wtrap{...} ... \end{worked}
\newtcolorbox{worked}[1]{%
  breakable, enhanced, colback=white, colframe=accent!65,
  boxrule=0.4pt, left=7pt, right=7pt, top=6pt, bottom=6pt,
  fonttitle=\bfseries, coltitle=accent, colbacktitle=soft,
  title={#1}}

% Run-in step heading inside a worked example.
\newcommand{\wstep}[1]{\par\smallskip\noindent\textbf{#1}\enspace\ignorespaces}

% A pitfall, set off from the running solution so it can be found again.
\newcommand{\wtrap}[1]{%
  \par\smallskip
  {\leftskip=1.3em\noindent
   \llap{\textcolor{accent}{\rule[0.22ex]{0.9em}{0.6pt}}\hspace{0.4em}}%
   \textbf{Trap.}\enspace\itshape #1\par}%
  \smallskip}

% --- problem / solution lists --------------------------------------------
\newlist{problems}{enumerate}{1}
\setlist[problems]{label=\textbf{\arabic*.}, leftmargin=2.1em,
  itemsep=0.55ex, parsep=0pt, topsep=0.8ex}
\newlist{solutions}{enumerate}{1}
\setlist[solutions]{label=\textbf{\arabic*.}, leftmargin=2.1em,
  itemsep=1.0ex, parsep=0pt, topsep=0.8ex}

% --- running heads --------------------------------------------------------
\pagestyle{fancy}
\fancyhf{}
\fancyhead[LE,RO]{\thepage}
\fancyhead[RE]{\small\itshape\nouppercase{\leftmark}}
\fancyhead[LO]{\small\itshape\nouppercase{\rightmark}}
\setlength{\headheight}{14pt}
\renewcommand{\headrulewidth}{0.3pt}
\renewcommand{\chaptermark}[1]{\markboth{\chaptername\ \thechapter.\ #1}{}}
\renewcommand{\sectionmark}[1]{\markright{\thesection\ #1}}

% --- shorthands used throughout -------------------------------------------
\newcommand{\dd}[2]{\frac{\mathrm{d}#1}{\mathrm{d}#2}}
\newcommand{\pd}[2]{\frac{\partial #1}{\partial #2}}
\newcommand{\pdd}[3]{\frac{\partial^{2} #1}{\partial #2\,\partial #3}}
\newcommand{\R}{\mathbb{R}}
\newcommand{\grad}{\nabla}
\DeclareMathOperator{\sgn}{sgn}
\DeclareMathOperator{\Hess}{H}

% German gloss, used on first mention of a technical term
\newcommand{\de}[1]{\emph{#1}}

\renewcommand*{\thisvolume}{functions}

\title{\Huge Functions and Their Graphs\\[0.4em]
  \large A Refresher with Exercises and Worked Solutions\\[0.8em]
  \normalsize Domain \(\cdot\) Range \(\cdot\) Graphs \(\cdot\)
  Transformations \(\cdot\) Regions \(\cdot\) Trigonometry}
\author{}
\date{}

\begin{document}
\frontmatter
\maketitle

\chapter*{About this sheet}
\addcontentsline{toc}{chapter}{About this sheet}

The algebra volume manipulates \emph{expressions}. This one is about
\emph{functions}: an expression together with a decision about which numbers
may be substituted into it, a record of which numbers come out, and above
all a picture.

Everything that follows in the later volumes leans on this material and
rarely stops to explain it. A limit is taken at a point of the domain or at
its boundary; a derivative exists only where the function does; an
asymptote is read off a graph; a constrained optimum is the point where two
families of curves touch. A mathematics sequence for economics or business
takes this for granted: one is expected to know what the graphs of the usual
functions look like, and to be able to sketch curves defined by equations in
two variables. That is the brief of this volume.

Five chapters, in the order they build on each other. \textbf{Domain and
range} fixes what a function is and where it lives. \textbf{The standard
functions and their graphs} is the catalogue --- powers, roots, exponentials,
logarithms, rational functions with their poles and asymptotes --- together
with the curves that are not graphs at all: circles, ellipses, hyperbolas,
isoquants. \textbf{Transforming functions} covers shifting, stretching,
reflecting, composing and inverting, and closes with the question an
economist meets most often, which is the reverse one: given the shape you
need, which formula has it? \textbf{Inequalities in two variables} turns
curves into regions, the skill behind every feasible set. And
\textbf{trigonometric functions} supplies the periodic shapes that later material
assumes and the calculus chapters use as test cases.

The structure is the same in every chapter: \textbf{Core Rules} states the
rules compactly; \textbf{Practice Problems} drills them; \textbf{Solutions}
works every problem through in full; each chapter closes with
\emph{additional problems} for a second pass. Seventy problems per chapter,
every one with a worked solution.

Work the problems before reading the solutions, and draw the pictures by
hand. A curve you have sketched yourself is one you will recognise in an
examination.

\tableofcontents

\mainmatter
\chapter{Functions, Domain and Range}

Up to here every chapter has manipulated \emph{expressions}. A
\de{Funktion} is the next step: an expression together with a decision
about which numbers may be substituted into it, and a record of which
numbers come out. Those two sets, the \de{Definitionsbereich} (domain) and
the \de{Wertebereich} (range), are not decoration. A limit is taken at a
point of the domain or at its boundary; a derivative exists only where the
function does; an optimisation problem is solved over a domain; and a
demand function that returns a negative price has left the region where it
means anything.

Finding a maximal domain is pure algebra of the kind the previous chapters
drilled: a denominator must not vanish, an even root needs a non-negative
radicand, a logarithm needs a positive argument. Each condition is an
equation or an inequality, and the domain is what survives all of them at
once. Finding a range is slightly different work: either the function is
built from pieces whose ranges are known, or the equation $y=f(x)$ is
solved for $x$ and one asks which $y$ admit a solution.

\section{Core Rules}

\begin{rulebox}{Function, domain, range}
A function $f\colon D\to\R$ assigns to every $x\in D$ exactly one value
$f(x)\in\R$. Two sets belong to the description:
\begin{itemize}[leftmargin=1.6em,itemsep=0.3ex]
  \item the \textbf{domain} $D(f)$ \de{(Definitionsbereich)}: the $x$ that
        may be substituted;
  \item the \textbf{range} $W(f)=\{f(x): x\in D(f)\}$
        \de{(Wertebereich, Bildmenge)}: the values actually attained.
\end{itemize}
The \textbf{graph} of $f$ is the set $\{(x,f(x)) : x\in D(f)\}$ in the
plane. A curve is the graph of a function exactly when every vertical line
meets it at most once (\textbf{vertical line test}). Reading off a graph:
$D(f)$ is its shadow on the $x$-axis, $W(f)$ its shadow on the $y$-axis.
\end{rulebox}

\begin{gallery}
\begin{tikzpicture}
  \begin{axis}[fnaxis, xmin=-0.6, xmax=5.4, ymin=-0.6, ymax=3.4,
      xtick={1,3,5}, ytick={1,2,3}, title={$f(x)=\sqrt{x-1}$}]
    \addplot[domain=1:5]{sqrt(x-1)};
  \end{axis}
\end{tikzpicture}\hfill
\begin{tikzpicture}
  \begin{axis}[fnaxis, xmin=-3.4, xmax=3.4, ymin=-0.6, ymax=4.4,
      xtick={-2,2}, ytick={2,4}, title={$f(x)=x^{2}$}]
    \addplot[domain=-2.1:2.1]{x^2};
  \end{axis}
\end{tikzpicture}\hfill
\begin{tikzpicture}
  \begin{axis}[fnaxis, xmin=-3.4, xmax=3.4, ymin=-1.2, ymax=4.4,
      xtick={-2,2}, ytick={1,3}, title={$f(x)=\mathrm{e}^{x}$}]
    \addplot[domain=-3:1.45]{exp(x)};
  \end{axis}
\end{tikzpicture}
\gcap{$D=[1,\infty)$, $W=[0,\infty)$ \quad\textbullet\quad
$D=\R$, $W=[0,\infty)$ \quad\textbullet\quad $D=\R$, $W=(0,\infty)$.}
\end{gallery}

\begin{rulebox}{The maximal domain: three obstructions}
Unless a domain is prescribed, $D(f)$ is the \textbf{maximal domain}: every
real $x$ for which the formula produces a real number. Only three things
can go wrong, and each is an (in)equality to solve:
\begin{itemize}[leftmargin=1.6em,itemsep=0.3ex]
  \item \textbf{Denominator:} $\dfrac{\;\cdot\;}{g(x)}$ requires
        $g(x)\neq 0$. Solve $g(x)=0$ and delete those points.
  \item \textbf{Even root:} $\sqrt[2k]{g(x)}$ requires $g(x)\geq 0$.
        Odd roots ($\sqrt[3]{\;}$) are defined for all reals.
  \item \textbf{Logarithm:} $\ln g(x)$ requires $g(x)>0$ (strictly).
\end{itemize}
A power $g(x)^{r}$ with non-integer $r$ counts as a root: it requires
$g(x)\geq 0$ for $r>0$ and $g(x)>0$ for $r<0$. If several obstructions
occur, \textbf{intersect} the resulting sets. A root in a denominator
combines two conditions into one: $\dfrac{1}{\sqrt{g(x)}}$ needs
$g(x)>0$.
\end{rulebox}

\begin{rulebox}{Domains prescribed by context}
In economics the domain is usually narrower than the maximal domain and
comes from the model, not from the formula:
quantities and inputs satisfy $q\geq 0$, $K,L\geq 0$; a price must satisfy
$p\geq 0$; a share or probability lives in $[0,1]$; an average cost
$AC(q)=C(q)/q$ needs $q>0$. State the economic domain explicitly, because
the formula will happily return values outside it: the linear demand
$p=120-3q$ is defined for every real $q$, but only $q\in[0,40]$ keeps both
$q\geq0$ and $p\geq 0$.
\end{rulebox}

\begin{rulebox}{Finding the range}
Four routes, in increasing order of generality:
\begin{enumerate}[leftmargin=1.7em,itemsep=0.3ex]
  \item \textbf{Standard building blocks.}
        $x^{2}\colon[0,\infty)$, \;$\sqrt{x}\colon[0,\infty)$,
        \;$|x|\colon[0,\infty)$, \;$\mathrm{e}^{x}\colon(0,\infty)$,
        \;$\ln x\colon\R$, \;$1/x\colon\R\setminus\{0\}$,
        \;$x^{3}\colon\R$, \;$\sin x,\cos x\colon[-1,1]$.
        Then track what the outer operations do: adding $c$ shifts the
        range by $c$, multiplying by $a>0$ scales it, multiplying by $a<0$
        scales \emph{and} flips it.
  \item \textbf{Complete the square} for a quadratic:
        $a(x-h)^{2}+k$ has range $[k,\infty)$ if $a>0$ and
        $(-\infty,k]$ if $a<0$.
  \item \textbf{Monotonicity.} If $f$ is strictly increasing on $[a,b]$,
        then $W=[f(a),f(b)]$; open ends stay open, and an end at
        $\pm\infty$ is replaced by the limiting value.
  \item \textbf{Solve $y=f(x)$ for $x$} (the general method). The range is
        the set of $y$ for which the solution exists and lies in $D(f)$.
        For $f(x)=\dfrac{x}{x+1}$: \;$y(x+1)=x\Rightarrow x=\dfrac{y}{1-y}$,
        solvable for every $y\neq 1$, hence $W=\R\setminus\{1\}$.
\end{enumerate}
\end{rulebox}

\begin{rulebox}{Composite functions}
$(g\circ f)(x)=g\bigl(f(x)\bigr)$ is evaluated \textbf{inside out}. Its
domain carries two conditions at once:
\[
  D(g\circ f)=\{\,x\in D(f)\;:\;f(x)\in D(g)\,\}.
\]
So $x$ must be admissible for the inner function, and its \emph{output}
must be admissible for the outer one. For $g(u)=\ln u$ and $f(x)=x^{2}-1$:
$x^{2}-1>0$, hence $D=(-\infty,-1)\cup(1,\infty)$. Order matters:
$g\circ f\neq f\circ g$ in general, and the two domains usually differ.
\end{rulebox}

\begin{rulebox}{Symmetry, monotonicity, injectivity}
\begin{itemize}[leftmargin=1.6em,itemsep=0.3ex]
  \item $f$ is \textbf{even} \de{(gerade)} if $f(-x)=f(x)$ for all $x\in D$
        (graph symmetric about the $y$-axis), \textbf{odd}
        \de{(ungerade)} if $f(-x)=-f(x)$ (symmetric about the origin).
        Only even powers $\Rightarrow$ even; only odd powers
        $\Rightarrow$ odd; a mix is neither.
  \item $f$ is \textbf{strictly increasing} if $x_{1}<x_{2}$ implies
        $f(x_{1})<f(x_{2})$, strictly decreasing if the inequality
        reverses.
  \item $f$ is \textbf{injective} \de{(eineindeutig)} if $f(x_{1})=f(x_{2})$
        forces $x_{1}=x_{2}$: no horizontal line meets the graph twice.
        Strictly monotone $\Rightarrow$ injective $\Rightarrow$ invertible
        on its range (Chapter~3).
\end{itemize}
\end{rulebox}

\section{Worked Example}

\begin{worked}{Find the maximal domain of $f(x)=\dfrac{\ln(x+3)}{\sqrt{4-x^{2}}}$, and the range of $g(x)=\dfrac{x+3}{x-2}$}
\wstep{Step 1. List the obstructions, one per structure in the formula.}
Three of them appear here, which is why this example was chosen:
\begin{itemize}[leftmargin=1.6em,itemsep=0.2ex]
  \item the logarithm needs $x+3>0$;
  \item the even root needs $4-x^{2}\ge0$;
  \item the denominator needs $\sqrt{4-x^{2}}\neq0$.
\end{itemize}
\wtrap{The last two combine into the \emph{strict} inequality $4-x^{2}>0$.
A square root sitting in a denominator is not allowed to be zero, so the
usual ``$\ge$'' for an even root tightens to ``$>$''. Handling the root and
the denominator as separate conditions and then forgetting to intersect them
is the standard way to hand in $[-2,2]$ instead of the right answer.}

\wstep{Step 2. Solve each condition separately.}
\[
  x+3>0 \iff x>-3 ;
  \qquad
  4-x^{2}>0 \iff x^{2}<4 \iff -2<x<2 .
\]

\wstep{Step 3. Intersect --- never unite.}
The function is defined where \emph{all} conditions hold at once:
\[
  D_{f}=(-3,\infty)\cap(-2,2)=\boxed{(-2,2)} .
\]
\wtrap{``And'' means intersection. Each obstruction removes points; nothing
in the formula ever gives points back. If your answer is larger than one of
the individual conditions, you took a union by mistake.}

\wstep{Step 4. Check the boundary behaviour, which confirms the strictness.}
As $x\to2^{-}$ the denominator $\to0^{+}$ while $\ln5>0$, so $f\to+\infty$.
The endpoint really must be excluded. At $x=-2$ the same thing happens.

\wstep{Step 5. The range of $g$: solve $y=g(x)$ for $x$.}
The range is the set of $y$ for which the equation $y=g(x)$ has a solution.
So treat $y$ as given and solve:
\[
  y=\frac{x+3}{x-2}
  \;\Longrightarrow\;
  y(x-2)=x+3
  \;\Longrightarrow\;
  yx-x=2y+3
  \;\Longrightarrow\;
  x(y-1)=2y+3 .
\]

\wstep{Step 6. Read off which $y$ fail.}
\[
  x=\frac{2y+3}{y-1},
\]
which exists for every $y$ except $y=1$; and at $y=1$ the previous line reads
$0=5$, impossible. Hence
\[
  \boxed{\;\text{range}(g)=\R\setminus\{1\}\;}
\]

\wstep{Step 7. Confirm it against the graph.}
$g(x)=\frac{x+3}{x-2}=1+\frac{5}{x-2}$, so the graph is the hyperbola
$5/x$ shifted, with horizontal asymptote $y=1$. The value $1$ is approached
but never attained --- exactly what the algebra said. Two independent
routes agreeing is worth the extra line.
\wtrap{Do not assume the range of a rational function is all of $\R$, and do
not assume it is the same as the domain. The excluded input ($x=2$) and the
excluded output ($y=1$) are different numbers, read off different parts of
the formula.}
\end{worked}

\section{Practice Problems}

State the maximal domain $D(f)\subseteq\R$ in problems 1--18, using
interval notation.

\begin{problems}
  \item $f(x)=\dfrac{x+3}{x-2}$
  \item $f(x)=\sqrt{x-5}$
  \item $f(x)=\sqrt{9-x^{2}}$
  \item $f(x)=\ln(x-4)$
  \item $f(x)=\ln(x^{2}-4)$
  \item $f(x)=\dfrac{\sqrt{x}}{x-1}$
  \item $f(x)=\dfrac{1}{\sqrt{x-3}}$
  \item $f(x)=\sqrt{x^{2}-x-6}$
  \item $f(x)=\dfrac{\ln x}{x-2}$
  \item $f(x)=\mathrm{e}^{1/x}$
  \item $f(x)=\sqrt{x+2}+\sqrt{4-x}$
  \item $f(x)=\dfrac{x^{2}-1}{x^{2}-3x+2}$
  \item $f(x)=\ln(\ln x)$
  \item $f(x)=x^{3/4}$
  \item $f(x)=\dfrac{1}{1-\mathrm{e}^{x}}$
  \item $f(x)=\sqrt{\ln x}$
  \item $f(x)=\sqrt[3]{x-7}$
  \item $f(x)=\dfrac{\sqrt{4-x}}{x^{2}-x}$
\end{problems}

\begin{problems}[resume]
  \item A firm has cost $C(q)=50+2q$. State the economic domain of $C$
        and of the average cost $AC(q)=C(q)/q$.
  \item For the demand schedule $p(q)=120-3q$, state the largest interval
        of quantities on which both $q\geq0$ and $p\geq0$ hold.
  \item State the domain of the Cobb--Douglas product
        $Y(K,L)=K^{0.3}L^{0.7}$ as a subset of the $(K,L)$-plane, and say
        why the exponents force the restriction.
  \item Utility is $U(x)=\ln x$. State the economic domain and explain what
        happens to $U$ as $x\to 0^{+}$.
\end{problems}

Determine the range $W(f)$ in problems 23--36. The domain is the maximal
one unless stated otherwise.

\begin{problems}[resume]
  \item $f(x)=x^{2}+4$
  \item $f(x)=x^{2}-6x+11$
  \item $f(x)=-2x^{2}+8x-3$
  \item $f(x)=\mathrm{e}^{x}-3$
  \item $f(x)=5-\mathrm{e}^{x}$
  \item $f(x)=\ln(x+2)$
  \item $f(x)=\dfrac{1}{x-2}+3$
  \item $f(x)=\sqrt{x-1}$
  \item $f(x)=\sqrt{16-x^{2}}$
  \item $f(x)=|x-3|$
  \item $f(x)=\dfrac{2x+1}{x-3}$
  \item $f(x)=\dfrac{x^{2}}{x^{2}+1}$
  \item $f(x)=x^{2}$ on $D=[-1,3]$
  \item $f(x)=2x-5$ on $D=(1,4]$
\end{problems}

\begin{problems}[resume]
  \item With $f(u)=\sqrt{u}$ and $g(x)=x-4$: give $(f\circ g)(x)$ and
        $(g\circ f)(x)$ with both domains.
  \item With $f(u)=1/u$ and $g(x)=x+2$: give $(f\circ g)(x)$ and its
        domain.
  \item With $f(u)=\ln u$ and $g(x)=9-x^{2}$: give $(f\circ g)(x)$ and its
        domain.
  \item With $f(u)=\sqrt{u}$ and $g(x)=1-x^{2}$: give $(f\circ g)(x)$, its
        domain and its range.
  \item Classify as even, odd, or neither:
        (a) $x^{4}-3x^{2}$, \;(b) $x^{3}-x$, \;(c) $x^{2}+x$,
        \;(d) $\dfrac{1}{x}$.
  \item Which of $x^{3}$, $x^{2}$, $\mathrm{e}^{x}$, $|x|$ are injective on
        $\R$?
  \item Give domain and range of
        $f(x)=\begin{cases} x^{2}, & x<1,\\ 2x-1, & x\geq 1.\end{cases}$
  \item Does $x^{2}+y^{2}=25$ define $y$ as a function of $x$? Justify with
        the vertical line test, and give the two functions obtained by
        solving for $y$, with their domains.
  \item A graph consists of the segment from $(-2,1)$ to $(0,5)$ together
        with the horizontal ray from $(0,5)$ to the right. Read off
        $D$ and $W$.
\end{problems}

\section{Solutions}

\begin{solutions}
  \item Denominator $x-2\neq0$: \;$D=\R\setminus\{2\}
        =(-\infty,2)\cup(2,\infty)$.
  \item $x-5\geq0$: \;$D=[5,\infty)$.
  \item $9-x^{2}\geq0\iff x^{2}\leq9\iff|x|\leq3$: \;$D=[-3,3]$.
  \item $x-4>0$: \;$D=(4,\infty)$.
  \item $x^{2}-4>0\iff|x|>2$: \;$D=(-\infty,-2)\cup(2,\infty)$.
  \item Two conditions: $x\geq0$ for the root and $x\neq1$ for the
        denominator: \;$D=[0,1)\cup(1,\infty)$.
  \item The root sits in the denominator, so $x-3>0$ (not merely
        $\geq0$): \;$D=(3,\infty)$.
  \item $x^{2}-x-6=(x-3)(x+2)\geq0$; the parabola opens upwards, so the
        product is non-negative outside the roots:
        \;$D=(-\infty,-2]\cup[3,\infty)$.
  \item $\ln x$ needs $x>0$, the denominator needs $x\neq2$:
        \;$D=(0,2)\cup(2,\infty)$.
  \item Only $x\neq0$; the exponential itself is unrestricted:
        \;$D=\R\setminus\{0\}$.
  \item $x+2\geq0$ and $4-x\geq0$, i.e. $x\geq-2$ \emph{and} $x\leq4$:
        \;$D=[-2,4]$.
  \item $x^{2}-3x+2=(x-1)(x-2)$: \;$D=\R\setminus\{1,2\}$. (The numerator
        also vanishes at $x=1$, but a cancelling factor does not repair the
        domain: $x=1$ stays excluded.)
  \item Inner: $x>0$; outer: $\ln x>0\iff x>1$. The stronger condition
        wins: \;$D=(1,\infty)$.
  \item Exponent $3/4$ is not an integer, so the base must satisfy
        $x\geq0$: \;$D=[0,\infty)$.
  \item $1-\mathrm{e}^{x}=0\iff \mathrm{e}^{x}=1\iff x=0$:
        \;$D=\R\setminus\{0\}$.
  \item $\ln x$ needs $x>0$, and $\ln x\geq0$ needs $x\geq1$; the
        intersection is \;$D=[1,\infty)$.
  \item An odd root accepts every real radicand: \;$D=\R$.
  \item Root: $4-x\geq0\iff x\leq4$. Denominator: $x^{2}-x=x(x-1)\neq0$,
        so $x\neq0$ and $x\neq1$. Together:
        \;$D=(-\infty,0)\cup(0,1)\cup(1,4]$.
  \item $C$: production cannot be negative, so $D(C)=[0,\infty)$; at $q=0$
        the value $C(0)=50$ is the fixed cost. $AC$: dividing by $q$
        excludes the origin, so $D(AC)=(0,\infty)$.
  \item $p\geq0\iff 120-3q\geq0\iff q\leq40$; with $q\geq0$ this gives
        $q\in[0,40]$.
  \item $D=\{(K,L): K\geq0,\;L\geq0\}$, the non-negative quadrant. The
        exponents $0.3$ and $0.7$ are not integers, so $K^{0.3}$ and
        $L^{0.7}$ require non-negative bases; negative input quantities
        would be meaningless anyway.
  \item $D=(0,\infty)$: consumption must be strictly positive, since
        $\ln 0$ is undefined. As $x\to0^{+}$, $U(x)\to-\infty$: the
        logarithmic utility punishes zero consumption infinitely hard,
        which is exactly why interior solutions appear in Cobb--Douglas
        problems.
  \item $x^{2}\geq0$, so $x^{2}+4\geq4$: \;$W=[4,\infty)$.
  \item $x^{2}-6x+11=(x-3)^{2}+2$: \;$W=[2,\infty)$.
  \item $-2x^{2}+8x-3=-2(x^{2}-4x)-3=-2(x-2)^{2}+5$; the parabola opens
        downwards: \;$W=(-\infty,5]$.
  \item $\mathrm{e}^{x}\in(0,\infty)$, shifted down by $3$:
        \;$W=(-3,\infty)$.
  \item $\mathrm{e}^{x}\in(0,\infty)$, negated to $(-\infty,0)$, shifted up
        by $5$: \;$W=(-\infty,5)$.
  \item $\ln$ attains every real value, and $x+2$ runs through all of
        $(0,\infty)$ on $D=(-2,\infty)$: \;$W=\R$.
  \item $1/(x-2)$ attains every value except $0$; adding $3$ shifts this:
        \;$W=\R\setminus\{3\}$.
  \item A square root returns exactly the non-negative numbers, and
        $x-1$ runs through $[0,\infty)$: \;$W=[0,\infty)$.
  \item $D=[-4,4]$, on which $16-x^{2}$ runs through $[0,16]$; the root of
        that is \;$W=[0,4]$.
  \item An absolute value is non-negative and attains every non-negative
        value (at $x=3\pm w$): \;$W=[0,\infty)$.
  \item Solve $y=\dfrac{2x+1}{x-3}$: \;$y(x-3)=2x+1\Rightarrow
        x(y-2)=3y+1\Rightarrow x=\dfrac{3y+1}{y-2}$, which exists exactly
        for $y\neq2$: \;$W=\R\setminus\{2\}$. ($y=2$ is the horizontal
        asymptote.)
  \item Solve $y=\dfrac{x^{2}}{x^{2}+1}$: \;$y(x^{2}+1)=x^{2}\Rightarrow
        x^{2}=\dfrac{y}{1-y}$. A real $x$ exists iff
        $\dfrac{y}{1-y}\geq0$, i.e. $0\leq y<1$: \;$W=[0,1)$.
  \item On $[-1,3]$ the square is not monotone: it falls to $0$ at $x=0$
        and rises to $9$ at $x=3$. Minimum $0$, maximum $9$:
        \;$W=[0,9]$.
  \item $f$ is strictly increasing, so the endpoints map to the endpoints:
        $f(1)=-3$ (excluded, as $1\notin D$) and $f(4)=3$ (included):
        \;$W=(-3,3]$.
  \item $(f\circ g)(x)=\sqrt{x-4}$ with $D=[4,\infty)$;
        $(g\circ f)(x)=\sqrt{x}-4$ with $D=[0,\infty)$. Same ingredients,
        different functions and different domains.
  \item $(f\circ g)(x)=\dfrac{1}{x+2}$, requiring $x+2\neq0$:
        \;$D=\R\setminus\{-2\}$.
  \item $(f\circ g)(x)=\ln(9-x^{2})$, requiring $9-x^{2}>0\iff|x|<3$:
        \;$D=(-3,3)$.
  \item $(f\circ g)(x)=\sqrt{1-x^{2}}$, requiring $1-x^{2}\geq0$:
        \;$D=[-1,1]$. On it $1-x^{2}$ runs through $[0,1]$, so
        $W=[0,1]$ (the upper unit semicircle).
  \item (a) even (only even powers); (b) odd (only odd powers);
        (c) neither: $f(-x)=x^{2}-x$ equals neither $f(x)$ nor $-f(x)$;
        (d) odd: $1/(-x)=-1/x$.
  \item $x^{3}$ and $\mathrm{e}^{x}$ are injective (both strictly
        increasing). $x^{2}$ and $|x|$ are not: $f(-1)=f(1)$ in both cases.
  \item $D=\R$. First piece: $x$ runs through $(-\infty,1)$, so $x^{2}$
        attains every value in $[0,\infty)$. Second piece: $2x-1$ on
        $[1,\infty)$ gives $[1,\infty)$. Union: \;$W=[0,\infty)$.
  \item No. The vertical line $x=0$ meets the circle at $(0,5)$ and
        $(0,-5)$, so the assignment is not unique. Solving gives two
        functions, $y=\sqrt{25-x^{2}}$ (upper semicircle) and
        $y=-\sqrt{25-x^{2}}$ (lower), each with domain $[-5,5]$ and ranges
        $[0,5]$ and $[-5,0]$.
  \item $D=[-2,\infty)$: the segment covers $x\in[-2,0]$, the ray covers
        $x\in[0,\infty)$. On the segment $y$ climbs from $1$ to $5$, on the
        ray $y$ stays at $5$: \;$W=[1,5]$.
\end{solutions}

\section{Additional Problems}

\begin{problems}
  \item $f(x)=\dfrac{3x}{2x+7}$: maximal domain.
  \item $f(x)=\sqrt{2x+8}$: maximal domain and range.
  \item $f(x)=\ln(5-2x)$: maximal domain.
  \item $f(x)=\dfrac{1}{\ln x}$: maximal domain.
  \item $f(x)=\dfrac{x-1}{\sqrt{x^{2}-4}}$: maximal domain.
  \item $f(x)=\sqrt{\dfrac{x-1}{x+2}}$: maximal domain.
  \item $f(x)=\ln\!\bigl(x^{2}+1\bigr)$: maximal domain and range.
  \item $f(x)=\mathrm{e}^{-x^{2}}$: maximal domain and range.
  \item $f(x)=\dfrac{1}{x^{2}-5x+6}$: maximal domain.
  \item $f(x)=\sqrt{4-x}+\dfrac{1}{x+1}$: maximal domain.
  \item $f(x)=(x-3)^{-1/2}$: maximal domain.
  \item $f(x)=\dfrac{|x|}{x}$: maximal domain and range.
  \item $f(x)=3-\dfrac{2}{x+1}$: range.
  \item $f(x)=\dfrac{x-1}{x+1}$: range.
  \item $f(x)=4-\sqrt{x+9}$: range.
  \item $f(x)=x^{2}-4x+7$ on $D=[0,3]$: range.
  \item $f(x)=\mathrm{e}^{2x}+1$ on $D=[0,\ln 2]$: range.
  \item $f(x)=\ln x$ on $D=[1,\mathrm{e}^{3}]$: range.
  \item With $f(u)=u^{2}$ and $g(x)=\sqrt{x-3}$: give $(f\circ g)(x)$ and
        its domain, and explain why the simplified formula alone would give
        the wrong answer.
  \item With $f(u)=\dfrac{1}{u-1}$ and $g(x)=\mathrm{e}^{x}$: give
        $(f\circ g)(x)$ and its domain.
  \item Revenue is $R(q)=q\,(200-4q)$. State the economic domain on which
        both $q\ge 0$ and the price stay non-negative, and the range of $R$
        on it.
  \item A production function is $Q(L)=10\sqrt{L}$ with $L\in[0,100]$.
        Give the range, and the average product $Q(L)/L$ with its domain.
  \item Classify as even, odd, or neither: (a) $\mathrm{e}^{-x^{2}}$,
        \;(b) $x\,|x|$, \;(c) $\ln x$, \;(d) $\dfrac{x^{3}}{x^{2}+1}$.
  \item Show that $f(x)=\dfrac{2x-1}{x+4}$ is injective on its domain by
        solving $f(x_{1})=f(x_{2})$.
  \item Give domain and range of
        $f(x)=\begin{cases} -x, & x<0,\\ \sqrt{x}, & 0\le x\le 4,\\
        2, & x>4.\end{cases}$
\end{problems}

\section{Additional Solutions}

\begin{solutions}
  \item $2x+7\neq0\iff x\neq-\tfrac{7}{2}$:
        \;$D=\R\setminus\{-\tfrac{7}{2}\}$.
  \item $2x+8\geq0\iff x\geq-4$: \;$D=[-4,\infty)$, and since $2x+8$ then
        runs through $[0,\infty)$, \;$W=[0,\infty)$.
  \item $5-2x>0\iff x<\tfrac{5}{2}$: \;$D=(-\infty,\tfrac{5}{2})$.
  \item $x>0$ for the logarithm and $\ln x\neq0\iff x\neq1$:
        \;$D=(0,1)\cup(1,\infty)$.
  \item The root is in the denominator, so $x^{2}-4>0\iff|x|>2$:
        \;$D=(-\infty,-2)\cup(2,\infty)$.
  \item Need $\dfrac{x-1}{x+2}\geq0$ with $x\neq-2$. The quotient is
        non-negative outside the interval between the two critical points
        $-2$ and $1$: \;$D=(-\infty,-2)\cup[1,\infty)$. (At $x=1$ the
        quotient is $0$, which is allowed; $x=-2$ is not.)
  \item $x^{2}+1\geq1>0$ always, so $D=\R$; the argument runs through
        $[1,\infty)$, hence \;$W=[0,\infty)$.
  \item $D=\R$. The exponent $-x^{2}$ runs through $(-\infty,0]$, and
        $\mathrm{e}^{u}$ maps that to $(0,1]$: \;$W=(0,1]$ (the bell curve,
        maximum $1$ at $x=0$).
  \item $x^{2}-5x+6=(x-2)(x-3)$: \;$D=\R\setminus\{2,3\}$.
  \item $4-x\geq0\iff x\leq4$, and $x\neq-1$:
        \;$D=(-\infty,-1)\cup(-1,4]$.
  \item A negative non-integer exponent means a root in a denominator:
        $x-3>0$, \;$D=(3,\infty)$.
  \item $D=\R\setminus\{0\}$. The quotient equals $1$ for $x>0$ and $-1$
        for $x<0$: \;$W=\{-1,1\}$, a range consisting of two points.
  \item $2/(x+1)$ attains every value except $0$; negating and adding $3$
        keeps exactly one value out of reach: \;$W=\R\setminus\{3\}$.
  \item $y=\dfrac{x-1}{x+1}\Rightarrow y(x+1)=x-1\Rightarrow
        x(y-1)=-1-y\Rightarrow x=\dfrac{-(1+y)}{y-1}=\dfrac{1+y}{1-y}$,
        which exists for $y\neq1$: \;$W=\R\setminus\{1\}$.
  \item $D=[-9,\infty)$; $\sqrt{x+9}$ runs through $[0,\infty)$, and
        $4-\sqrt{x+9}$ therefore through \;$W=(-\infty,4]$.
  \item $x^{2}-4x+7=(x-2)^{2}+3$. On $[0,3]$ the vertex $x=2$ lies inside,
        giving the minimum $3$; the endpoints give $f(0)=7$ and $f(3)=4$,
        so the maximum is $7$: \;$W=[3,7]$.
  \item Strictly increasing: $f(0)=2$ and
        $f(\ln 2)=\mathrm{e}^{2\ln 2}+1=4+1=5$: \;$W=[2,5]$.
  \item Strictly increasing: $\ln 1=0$ and $\ln \mathrm{e}^{3}=3$:
        \;$W=[0,3]$.
  \item $(f\circ g)(x)=\bigl(\sqrt{x-3}\bigr)^{2}=x-3$, but only where the
        inner function is defined: \;$D=[3,\infty)$. The simplified formula
        $x-3$ suggests $D=\R$; the domain is inherited from the inner
        function and is not restored by cancellation.
  \item $(f\circ g)(x)=\dfrac{1}{\mathrm{e}^{x}-1}$, requiring
        $\mathrm{e}^{x}\neq1\iff x\neq0$: \;$D=\R\setminus\{0\}$.
  \item Price $200-4q\ge0\iff q\le50$, so the domain is $[0,50]$.
        $R(q)=200q-4q^{2}=-4(q-25)^{2}+2500$, a downward parabola with
        vertex at $q=25$; on $[0,50]$ it runs from $0$ up to $2500$ and
        back to $0$: \;$W=[0,2500]$.
  \item $\sqrt{L}$ runs through $[0,10]$ on $[0,100]$, so
        $W=[0,100]$. The average product is
        $Q(L)/L=10/\sqrt{L}$, which needs $L>0$: domain $(0,100]$.
  \item (a) even: $(-x)^{2}=x^{2}$; (b) odd:
        $(-x)|-x|=-x|x|$; (c) neither, and in fact the test cannot even be
        applied: the domain $(0,\infty)$ is not symmetric about $0$;
        (d) odd: numerator odd, denominator even.
  \item $\dfrac{2x_{1}-1}{x_{1}+4}=\dfrac{2x_{2}-1}{x_{2}+4}$
        $\Rightarrow(2x_{1}-1)(x_{2}+4)=(2x_{2}-1)(x_{1}+4)$
        $\Rightarrow 2x_{1}x_{2}+8x_{1}-x_{2}-4
        =2x_{1}x_{2}+8x_{2}-x_{1}-4$
        $\Rightarrow 9x_{1}=9x_{2}\Rightarrow x_{1}=x_{2}$. Injective on
        $\R\setminus\{-4\}$.
  \item $D=\R$. First piece: $-x$ on $(-\infty,0)$ gives $(0,\infty)$.
        Second: $\sqrt{x}$ on $[0,4]$ gives $[0,2]$. Third: the constant
        $2$. The union is \;$W=[0,\infty)$.
\end{solutions}

\chapter{The Standard Functions and Their Graphs}

A first course in mathematics for economists assumes that the graphs of the
usual functions are known by sight: which way a cubic runs off at its ends,
where a rational function has a pole, that $\ln x$ never turns around, that
an isoquant is a hyperbola. That knowledge is not decoration either. Every
optimisation problem later on begins with a picture, every limit is checked
against one, and a wrong sketch produces a wrong answer faster than any
arithmetic slip.

This chapter is the catalogue, and the drill that makes it stick. First the
standard shapes; then the features that can be read straight off a formula
without any calculus: intercepts, roots, symmetry, poles and asymptotes,
and the behaviour at the far ends of the axis. The last block leaves the
world of $y=f(x)$ and sketches curves given by an \emph{equation} in $x$
and $y$: circles, ellipses, hyperbolas, budget lines and indifference
curves.

\section{Core Rules}

\begin{rulebox}{The power functions}
Whole-number powers alternate: even powers are symmetric about the
$y$-axis and non-negative, odd powers are symmetric about the origin and
run from $-\infty$ to $+\infty$. Negative powers have a pole at $0$; the
root $\sqrt{x}$ is the upper half of a sideways parabola.
\end{rulebox}

\begin{gallery}
\begin{tikzpicture}
  \begin{axis}[fnaxis, xmin=-2.6, xmax=2.6, ymin=-2.6, ymax=2.6,
      xtick={-2,2}, ytick={-2,2}, title={$y=x$}]
    \addplot[domain=-2.4:2.4]{x};
  \end{axis}
\end{tikzpicture}\hfill
\begin{tikzpicture}
  \begin{axis}[fnaxis, xmin=-2.6, xmax=2.6, ymin=-1.0, ymax=4.4,
      xtick={-2,2}, ytick={2,4}, title={$y=x^{2}$}]
    \addplot[domain=-2.1:2.1]{x^2};
  \end{axis}
\end{tikzpicture}\hfill
\begin{tikzpicture}
  \begin{axis}[fnaxis, xmin=-2.6, xmax=2.6, ymin=-4.4, ymax=4.4,
      xtick={-2,2}, ytick={-4,4}, title={$y=x^{3}$}]
    \addplot[domain=-1.62:1.62]{x^3};
  \end{axis}
\end{tikzpicture}
\end{gallery}

\begin{gallery}
\begin{tikzpicture}
  \begin{axis}[fnaxis, xmin=-0.6, xmax=4.6, ymin=-0.8, ymax=2.6,
      xtick={2,4}, ytick={1,2}, title={$y=\sqrt{x}$}]
    \addplot[domain=0:4.4]{sqrt(x)};
  \end{axis}
\end{tikzpicture}\hfill
\begin{tikzpicture}
  \begin{axis}[fnaxis, xmin=-3.4, xmax=3.4, ymin=-3.4, ymax=3.4,
      xtick={-2,2}, ytick={-2,2}, title={$y=1/x$}]
    \addplot[domain=-3.2:-0.31]{1/x};
    \addplot[domain=0.31:3.2]{1/x};
  \end{axis}
\end{tikzpicture}\hfill
\begin{tikzpicture}
  \begin{axis}[fnaxis, xmin=-3.4, xmax=3.4, ymin=-0.6, ymax=4.4,
      xtick={-2,2}, ytick={2,4}, title={$y=1/x^{2}$}]
    \addplot[domain=-3.2:-0.48]{1/x^2};
    \addplot[domain=0.48:3.2]{1/x^2};
  \end{axis}
\end{tikzpicture}
\end{gallery}

\begin{rulebox}{Exponential, logarithm, absolute value}
$\mathrm{e}^{x}$ is positive everywhere, strictly increasing, and flattens
onto the asymptote $y=0$ as $x\to-\infty$; it passes through $(0,1)$.
$\ln x$ is its mirror image at the line $y=x$: defined only for $x>0$,
strictly increasing, passing through $(1,0)$, with the $y$-axis as a
vertical asymptote. A general exponential $a^{x}$ rises for $a>1$ and falls
for $0<a<1$, always through $(0,1)$; note $\bigl(\tfrac12\bigr)^{x}
=2^{-x}$, so the decreasing case is the increasing one reflected in the
$y$-axis. The absolute value $|x|$ is a V with its vertex at the origin.
\end{rulebox}

\begin{gallery}
\begin{tikzpicture}
  \begin{axis}[fnaxis, xmin=-3.4, xmax=2.4, ymin=-0.8, ymax=4.4,
      xtick={-2,2}, ytick={1,3}, title={$y=\mathrm{e}^{x}$}]
    \addplot[domain=-3.2:1.45]{exp(x)};
  \end{axis}
\end{tikzpicture}\hfill
\begin{tikzpicture}
  \begin{axis}[fnaxis, xmin=-0.8, xmax=5.4, ymin=-3.4, ymax=2.4,
      xtick={1,3,5}, ytick={-2,1}, title={$y=\ln x$}]
    \addplot[domain=0.045:5.2]{ln(x)};
  \end{axis}
\end{tikzpicture}\hfill
\begin{tikzpicture}
  \begin{axis}[fnaxis, xmin=-3.4, xmax=3.4, ymin=-0.8, ymax=3.4,
      xtick={-2,2}, ytick={1,3}, title={$y=|x|$}]
    \addplot[domain=-3:3]{abs(x)};
  \end{axis}
\end{tikzpicture}
\end{gallery}

\begin{rulebox}{What can be read off a formula without calculus}
\begin{itemize}[leftmargin=1.6em,itemsep=0.3ex]
  \item \textbf{$y$-intercept:} evaluate $f(0)$ (if $0\in D$).
  \item \textbf{Roots} \de{(Nullstellen)}: solve $f(x)=0$. For a product,
        set each factor to zero. A factor $(x-a)^{k}$ with $k$ odd makes
        the graph \emph{cross} the axis at $a$; with $k$ even it only
        \emph{touches} it.
  \item \textbf{Sign:} between consecutive roots and poles the sign cannot
        change. Pick one test point per interval; that fixes the whole
        picture.
  \item \textbf{Symmetry:} $f(-x)=f(x)$ gives mirror symmetry in the
        $y$-axis, $f(-x)=-f(x)$ point symmetry in the origin.
  \item \textbf{End behaviour of a polynomial:} only the leading term
        $a_{n}x^{n}$ matters as $x\to\pm\infty$. Even $n$: both ends go the
        same way (up if $a_{n}>0$). Odd $n$: the ends go opposite ways
        (down--left, up--right if $a_{n}>0$).
\end{itemize}
\end{rulebox}

\begin{rulebox}{Rational functions: poles and asymptotes}
Write $f(x)=\dfrac{p(x)}{q(x)}$ with $p,q$ polynomials and
\textbf{cancel common factors first}.
\begin{itemize}[leftmargin=1.6em,itemsep=0.3ex]
  \item A zero of $q$ that does \emph{not} cancel is a \textbf{pole}
        \de{(Polstelle)}: the vertical line $x=a$ is a
        \textbf{vertical asymptote}, and $|f|\to\infty$ there.
  \item A zero of $q$ that \emph{does} cancel is a \textbf{hole}: the graph
        is that of the reduced function, minus the single point $x=a$.
  \item \textbf{Horizontal asymptote}, by comparing degrees:
        $\deg p<\deg q\Rightarrow y=0$; \;$\deg p=\deg q\Rightarrow
        y=\dfrac{\text{leading coeff. of }p}{\text{leading coeff. of }q}$;
        \;$\deg p>\deg q\Rightarrow$ none.
  \item If $\deg p=\deg q+1$, polynomial division gives
        $f(x)=(\text{linear})+\dfrac{\text{remainder}}{q(x)}$, and the
        linear part is an \textbf{oblique asymptote}
        \de{(schiefe Asymptote)}, because the remainder term dies out.
\end{itemize}
\end{rulebox}

\begin{gallery}
\begin{tikzpicture}
  \begin{axis}[fnaxis, xmin=-2.4, xmax=6.4, ymin=-4.4, ymax=6.4,
      xtick={2,4,6}, ytick={-2,2,4}, title={$y=\dfrac{2x+1}{x-3}$}]
    \addplot[domain=-2.2:2.55]{(2*x+1)/(x-3)};
    \addplot[domain=3.42:6.2]{(2*x+1)/(x-3)};
    \addplot[gray, dashed, thin, forget plot, samples=2, domain=-2.2:6.2]{2};
    \addplot[gray, dashed, thin, forget plot, samples=2] coordinates {(3,-4.2)(3,6.2)};
  \end{axis}
\end{tikzpicture}\hfill
\begin{tikzpicture}
  \begin{axis}[fnaxis, xmin=-4.4, xmax=4.4, ymin=-6.4, ymax=6.4,
      xtick={-2,2}, ytick={-4,4}, title={$y=\dfrac{x^{2}+1}{x}$}]
    \addplot[domain=-4.2:-0.17]{(x^2+1)/x};
    \addplot[domain=0.17:4.2]{(x^2+1)/x};
    \addplot[gray, dashed, thin, forget plot, samples=2, domain=-4.2:4.2]{x};
  \end{axis}
\end{tikzpicture}\hfill
\begin{tikzpicture}
  \begin{axis}[fnaxis, xmin=-3, xmax=42, ymin=-3, ymax=32,
      xtick={10,20,30,40}, ytick={5,15,25},
      title={$AC(q)=\dfrac{100+5q}{q}$}]
    \addplot[domain=3.9:40]{(100+5*x)/x};
    \addplot[gray, dashed, thin, forget plot, samples=2, domain=-2.8:40]{5};
  \end{axis}
\end{tikzpicture}
\gcap{Horizontal asymptote $y=2$ and pole $x=3$ \quad\textbullet\quad
oblique asymptote $y=x$ \quad\textbullet\quad average cost falling towards
the asymptote $AC=5$.}
\end{gallery}

\begin{rulebox}{Curves given by an equation in two variables}
Not every curve is a graph $y=f(x)$. The standard ones:
\begin{itemize}[leftmargin=1.6em,itemsep=0.3ex]
  \item \textbf{Line:} $ax+by=c$. Sketch it from the two intercepts
        $x=c/a$ and $y=c/b$; solved for $y$ the slope is $-a/b$. This is
        the budget line $p_{x}x+p_{y}y=M$, with slope $-p_{x}/p_{y}$.
  \item \textbf{Circle:} $(x-x_{0})^{2}+(y-y_{0})^{2}=r^{2}$, centre
        $(x_{0},y_{0})$, radius $r$. A general
        $x^{2}+y^{2}+Dx+Ey+F=0$ is brought to this form by
        \textbf{completing the square} in $x$ and in $y$.
  \item \textbf{Ellipse:} $\dfrac{x^{2}}{a^{2}}+\dfrac{y^{2}}{b^{2}}=1$,
        crossing the axes at $(\pm a,0)$ and $(0,\pm b)$.
  \item \textbf{Hyperbola:} $xy=c$ with $c>0$ lies in the first and third
        quadrant with both axes as asymptotes. Solved for $y$ it is
        $y=c/x$. Every Cobb--Douglas isoquant $K^{\alpha}L^{\beta}=\bar Q$
        and every indifference curve $xy=\bar U$ has this shape.
\end{itemize}
To sketch, either solve for $y$, remembering that a square root splits the
curve into \textbf{two branches} $y=\pm\sqrt{\cdots}$, or read off centre,
radius and intercepts directly.
\end{rulebox}

\section{Worked Example}

\begin{worked}{Sketch $f(x)=\dfrac{2x+1}{x-3}$ and $g(x)=(x-1)^{2}(x+3)$ from the formulas alone, and say what $h(x)=\dfrac{x^{2}-1}{x-1}$ really is}
\wstep{Step 1. The rational function: find the poles.}
The denominator of $f$ vanishes at $x=3$ and the numerator does not
($2\cdot3+1=7\neq0$), so $x=3$ is a genuine \textbf{vertical asymptote}.

\wstep{Step 2. Find the horizontal asymptote by comparing degrees.}
Numerator and denominator both have degree $1$, so the limit at $\pm\infty$
is the ratio of leading coefficients:
\[
  \lim_{x\to\pm\infty}\frac{2x+1}{x-3}=\frac{2}{1}=2 ,
\]
a horizontal asymptote $y=2$. Equivalently, dividing out gives
$f(x)=2+\dfrac{7}{x-3}$, which shows the whole graph at once: the hyperbola
$7/x$ shifted right $3$ and up $2$.

\wstep{Step 3. Intercepts and signs.}
$f(x)=0$ when $2x+1=0$, i.e.\ $x=-\tfrac12$; and $f(0)=\tfrac{1}{-3}=-\tfrac13$.
The sign changes at $x=-\tfrac12$ (zero) and at $x=3$ (pole), giving
$+,-,+$ on the three intervals.

\wstep{Step 4. Approach the pole from both sides.}
As $x\to3^{+}$, $x-3\to0^{+}$ and the numerator is near $7>0$, so
$f\to+\infty$; as $x\to3^{-}$, $f\to-\infty$. Now the sketch is determined:
two branches, one below-left rising to the asymptote $y=2$ from below, one
above-right falling to it from above.
\wtrap{Test the sign on \emph{each side} of a pole rather than assuming
symmetry. A pole of odd order flips the sign, one of even order does not,
and drawing both branches on the same side of the horizontal asymptote is
the usual consequence of guessing.}

\wstep{Step 5. The polynomial: read the roots with their multiplicities.}
$g(x)=(x-1)^{2}(x+3)$ has a root of multiplicity $2$ at $x=1$ and a simple
root at $x=-3$. Even multiplicity means the graph \textbf{touches} the axis
and turns back; odd multiplicity means it \textbf{crosses}. So the curve
crosses at $-3$ and bounces off the axis at $1$.

\wstep{Step 6. End behaviour is decided by the leading term alone.}
Expanding just the leading term: $x^{2}\cdot x=x^{3}$, degree $3$ with a
positive coefficient, so $g\to-\infty$ as $x\to-\infty$ and $g\to+\infty$ as
$x\to+\infty$. With $g(0)=1\cdot3=3>0$, the shape is now forced: up through
$-3$, over a local maximum, down to touch the axis at $1$, then away
upward.
\wtrap{Degree and leading sign come from multiplying the leading terms of
every factor, counted with multiplicity. $(x-1)^{2}(x+3)$ is cubic, not
quadratic; miscounting the degree reverses one of the two ends.}

\wstep{Step 7. The impostor: $h(x)=\dfrac{x^{2}-1}{x-1}$.}
Here the numerator \emph{also} vanishes at $x=1$. Factor before concluding
anything:
\[
  h(x)=\frac{(x-1)(x+1)}{x-1}=x+1\qquad\text{for }x\neq1 .
\]
So $h$ has no asymptote at all. Its graph is the straight line $y=x+1$ with a
single point punched out at $(1,2)$ --- a \textbf{removable discontinuity},
a hole, not a pole.
\wtrap{A zero denominator does not by itself make an asymptote. Always
factor numerator and denominator first: a shared factor gives a hole, an
unshared one gives a pole. This is the same distinction that decides whether
a limit in Volume~4 exists or diverges.}
\end{worked}

\section{Practice Problems}

For each polynomial in problems 1--8, give the roots (with multiplicity
where it matters), the $y$-intercept, the symmetry, and the end behaviour
as $x\to\pm\infty$.

\begin{problems}
  \item $f(x)=x^{3}-4x$
  \item $f(x)=x^{4}-5x^{2}+4$
  \item $f(x)=-2x^{3}+6x$
  \item $f(x)=(x-1)^{2}(x+3)$
  \item $f(x)=x^{2}-6x+8$
  \item $f(x)=-x^{2}+4x-4$
  \item $f(x)=x^{5}-x$
  \item A polynomial satisfies $f(x)\to-\infty$ both as $x\to+\infty$ and
        as $x\to-\infty$. What does that say about its degree and its
        leading coefficient?
\end{problems}

In problems 9--20, give all vertical asymptotes, holes, and horizontal or
oblique asymptotes.

\begin{problems}[resume]
  \item $f(x)=\dfrac{1}{x-3}$
  \item $f(x)=\dfrac{2x+1}{x-3}$
  \item $f(x)=\dfrac{x^{2}-1}{x-1}$
  \item $f(x)=\dfrac{3x}{x^{2}-4}$
  \item $f(x)=\dfrac{x^{2}+1}{x}$
  \item $f(x)=\dfrac{x^{2}-4}{x^{2}-9}$
  \item $f(x)=\dfrac{x^{2}+3x+2}{x+1}$
  \item $f(x)=\dfrac{x^{2}-x-6}{x-3}$
  \item $f(x)=\dfrac{2x^{2}+3}{x-1}$
  \item $f(x)=\dfrac{5}{x^{2}+1}$
  \item $AC(q)=\dfrac{100+5q}{q}$ \;(average cost, $q>0$)
  \item $f(x)=\dfrac{x+3}{x^{2}+x-6}$
\end{problems}

\begin{problems}[resume]
  \item For $y=2^{x}$: give the $y$-intercept, the asymptote, the
        monotonicity and the range.
  \item Explain why the graph of $y=\bigl(\tfrac12\bigr)^{x}$ is the graph
        of $y=2^{x}$ reflected in the $y$-axis.
  \item For $y=\ln x$: give the $x$-intercept, the asymptote, the domain
        and the range.
  \item Sketch $y=\mathrm{e}^{-x}$ and state its asymptote and range.
  \item Both $\ln x$ and $\log_{10}x$ pass through $(1,0)$ and rise. What
        exactly distinguishes their graphs?
  \item The adjustment path $y=1-\mathrm{e}^{-x}$, $x\ge0$, is used for
        saturation effects. Give $y(0)$, the asymptote, and say in words
        what the curve does.
\end{problems}

\begin{problems}[resume]
  \item Match each graph (A)--(F) below to one of the formulas
        $\;\ln x$, \;$|x|$, \;$\sqrt{x}$, \;$x^{3}-x$,
        \;$\mathrm{e}^{-x}$, \;$1/x^{2}$.
\end{problems}

\begin{gallery}
\begin{tikzpicture}
  \begin{axis}[fnaxis, xmin=-2.6, xmax=2.6, ymin=-2.2, ymax=2.6,
      xtick={-2,2}, ytick={-2,2}, title={(A)}]
    \addplot[domain=-1.85:1.85]{x^3-x};
  \end{axis}
\end{tikzpicture}\hfill
\begin{tikzpicture}
  \begin{axis}[fnaxis, xmin=-0.7, xmax=4.6, ymin=-0.8, ymax=2.6,
      xtick={2,4}, ytick={1,2}, title={(B)}]
    \addplot[domain=0:4.4]{sqrt(x)};
  \end{axis}
\end{tikzpicture}\hfill
\begin{tikzpicture}
  \begin{axis}[fnaxis, xmin=-3.4, xmax=3.4, ymin=-0.6, ymax=4.4,
      xtick={-2,2}, ytick={2,4}, title={(C)}]
    \addplot[domain=-3.2:-0.48]{1/x^2};
    \addplot[domain=0.48:3.2]{1/x^2};
  \end{axis}
\end{tikzpicture}
\end{gallery}

\begin{gallery}
\begin{tikzpicture}
  \begin{axis}[fnaxis, xmin=-2.4, xmax=3.4, ymin=-0.8, ymax=4.4,
      xtick={-2,2}, ytick={1,3}, title={(D)}]
    \addplot[domain=-1.45:3.2]{exp(-x)};
  \end{axis}
\end{tikzpicture}\hfill
\begin{tikzpicture}
  \begin{axis}[fnaxis, xmin=-3.4, xmax=3.4, ymin=-0.8, ymax=3.4,
      xtick={-2,2}, ytick={1,3}, title={(E)}]
    \addplot[domain=-3:3]{abs(x)};
  \end{axis}
\end{tikzpicture}\hfill
\begin{tikzpicture}
  \begin{axis}[fnaxis, xmin=-0.8, xmax=5.4, ymin=-3.4, ymax=2.4,
      xtick={1,3,5}, ytick={-2,1}, title={(F)}]
    \addplot[domain=0.045:5.2]{ln(x)};
  \end{axis}
\end{tikzpicture}
\end{gallery}

\begin{problems}[resume]
  \item Sketch $f(x)=x^{2}-4x+3$: give the vertex, both intercepts and the
        direction of opening.
  \item Sketch $f(x)=\dfrac{x+1}{x-1}$: give both asymptotes and both
        intercepts.
  \item Sketch $f(x)=\sqrt{x+4}$: give the domain, the range and the
        starting point of the curve.
  \item Sketch $f(x)=|x|-2$: give the vertex and the roots.
  \item Sketch $f(x)=x^{3}-3x^{2}$: give the roots with multiplicity and
        the end behaviour.
  \item Classify as even, odd, or neither: (a) $x^{4}+2$, \;(b) $x^{3}+x$,
        \;(c) $1/x^{2}$, \;(d) $\mathrm{e}^{x}$.
  \item Which of these have a horizontal asymptote, and what is it?
        (a) $\dfrac{3x+1}{x-2}$, \;(b) $\dfrac{x^{2}}{x+1}$,
        \;(c) $\mathrm{e}^{-x}$ as $x\to+\infty$, \;(d) $\ln x$.
  \item On which intervals is $f(x)=\dfrac{(x-1)(x+2)}{x-3}$ positive?
  \item Describe the graph of $f(x)=\sqrt{x^{2}-1}$: domain, symmetry,
        range, and the two points where it meets the $x$-axis.
\end{problems}

In problems 37--45, identify the curve, and sketch it.

\begin{problems}[resume]
  \item $x^{2}+y^{2}=16$. Also solve for $y$ and give the two branches with
        their domains.
  \item $(x-2)^{2}+(y+1)^{2}=9$. Give centre and radius.
  \item $\dfrac{x^{2}}{9}+\dfrac{y^{2}}{4}=1$. Give the four axis
        intercepts.
  \item $xy=6$. Solve for $y$; in which quadrants does the curve lie?
  \item $y^{2}=x$. Is $y$ a function of $x$? Is $x$ a function of $y$?
  \item $2x+3y=12$. Give both intercepts and the slope.
  \item $x^{2}+y^{2}-4x+6y+9=0$. Complete the square; give centre and
        radius.
  \item The isoquant $K^{0.5}L^{0.5}=10$. Solve for $L$ and describe the
        shape for $K>0$.
  \item The budget line $4x+2y=40$ and the indifference curve $xy=50$.
        Give the slope of the budget line and $y$ as a function of $x$ on
        the indifference curve.
\end{problems}

\section{Solutions}

\begin{solutions}
  \item $x^{3}-4x=x(x^{2}-4)=x(x-2)(x+2)$: roots $-2,0,2$, all simple, all
        crossings; $f(0)=0$; odd (only odd powers); degree $3$ with
        positive leading coefficient, so $f\to-\infty$ as $x\to-\infty$ and
        $f\to+\infty$ as $x\to+\infty$.
  \item $x^{4}-5x^{2}+4=(x^{2}-1)(x^{2}-4)$: roots $\pm1,\pm2$;
        $f(0)=4$; even; degree $4$, positive leading coefficient, so both
        ends go to $+\infty$ (a W-shape).
  \item $-2x^{3}+6x=-2x(x^{2}-3)$: roots $0,\pm\sqrt3$; $f(0)=0$; odd;
        degree $3$ with \emph{negative} leading coefficient, so
        $f\to+\infty$ as $x\to-\infty$ and $f\to-\infty$ as $x\to+\infty$.
  \item Roots $x=1$ (double, the graph touches and turns back) and $x=-3$
        (simple, a crossing); $f(0)=(-1)^{2}\cdot3=3$; neither even nor
        odd; degree $3$, leading coefficient $+1$: down on the left, up on
        the right.
  \item $(x-2)(x-4)$: roots $2$ and $4$; $f(0)=8$; vertex at
        $x=3$, $f(3)=-1$; neither even nor odd; parabola opening upwards,
        both ends $+\infty$.
  \item $-(x-2)^{2}$: a double root at $x=2$ (the vertex touches the axis);
        $f(0)=-4$; opens downwards, both ends $-\infty$.
  \item $x^{5}-x=x(x^{4}-1)=x(x^{2}-1)(x^{2}+1)$: roots $-1,0,1$;
        $f(0)=0$; odd; degree $5$, positive leading coefficient: down on
        the left, up on the right.
  \item Both ends behave alike, so the degree is \textbf{even}; both go
        \emph{down}, so the leading coefficient is \textbf{negative}.
  \item Vertical asymptote $x=3$; $\deg p=0<1=\deg q$, so the horizontal
        asymptote is $y=0$. No holes.
  \item Vertical asymptote $x=3$; equal degrees, so
        $y=\tfrac{2}{1}=2$ is the horizontal asymptote.
  \item $\dfrac{(x-1)(x+1)}{x-1}=x+1$ for $x\neq1$: no pole but a
        \textbf{hole} at $(1,2)$. The graph is the line $y=x+1$ with that
        one point punched out; no horizontal asymptote.
  \item $x^{2}-4=(x-2)(x+2)$: poles at $x=\pm2$; $\deg p=1<2=\deg q$, so
        $y=0$.
  \item Pole at $x=0$. Degrees differ by one, so divide:
        $\dfrac{x^{2}+1}{x}=x+\dfrac{1}{x}$, giving the oblique asymptote
        $y=x$.
  \item Poles at $x=\pm3$; equal degrees with leading coefficients $1$ and
        $1$, so the horizontal asymptote is $y=1$.
  \item $\dfrac{(x+1)(x+2)}{x+1}=x+2$ for $x\neq-1$: a hole at $(-1,1)$,
        no pole, no horizontal asymptote.
  \item $\dfrac{(x-3)(x+2)}{x-3}=x+2$ for $x\neq3$: a hole at $(3,5)$.
  \item Pole at $x=1$. Division: $2x^{2}+3=(x-1)(2x+2)+5$, so
        $f(x)=2x+2+\dfrac{5}{x-1}$ and the oblique asymptote is $y=2x+2$.
  \item $x^{2}+1$ never vanishes: no poles at all. $\deg p=0<2$, so
        $y=0$; the graph is a bell-shaped hump with maximum $5$ at $x=0$,
        even.
  \item $AC(q)=\dfrac{100}{q}+5$: pole at $q=0$ and horizontal asymptote
        $AC=5$. Average cost falls towards the marginal cost $5$ as the
        fixed cost $100$ is spread over more units.
  \item $x^{2}+x-6=(x+3)(x-2)$, so
        $f(x)=\dfrac{x+3}{(x+3)(x-2)}=\dfrac{1}{x-2}$ for $x\neq-3$:
        a hole at $\bigl(-3,-\tfrac15\bigr)$, a pole at $x=2$, and the
        horizontal asymptote $y=0$.
  \item $y$-intercept $2^{0}=1$; asymptote $y=0$ (approached as
        $x\to-\infty$); strictly increasing; range $(0,\infty)$.
  \item Because $\bigl(\tfrac12\bigr)^{x}=(2^{-1})^{x}=2^{-x}$: the
        argument $x$ has been replaced by $-x$, and replacing $x$ by $-x$
        mirrors any graph in the $y$-axis.
  \item $x$-intercept $(1,0)$; the $y$-axis $x=0$ is a vertical asymptote
        (as $x\to0^{+}$, $\ln x\to-\infty$); $D=(0,\infty)$, $W=\R$.
  \item The decreasing exponential: it falls from $+\infty$ through
        $(0,1)$ towards the asymptote $y=0$ as $x\to+\infty$;
        $W=(0,\infty)$. It is $\mathrm{e}^{x}$ mirrored in the $y$-axis.
  \item Only the vertical scale: $\log_{10}x=\dfrac{\ln x}{\ln 10}$, so the
        common logarithm is the natural one compressed by the constant
        factor $\ln 10\approx2.303$. Same domain, same asymptote, same
        intercept, same shape.
  \item $y(0)=1-1=0$; as $x\to\infty$, $\mathrm{e}^{-x}\to0$, so the
        asymptote is $y=1$. The curve starts at the origin, rises steeply
        at first and then ever more slowly, approaching the saturation
        level $1$ without ever reaching it.
  \item (A) $x^{3}-x$ \;(odd, three roots); \;(B) $\sqrt{x}$;
        \;(C) $1/x^{2}$ \;(even, positive, pole at $0$);
        \;(D) $\mathrm{e}^{-x}$; \;(E) $|x|$; \;(F) $\ln x$.
  \item $x^{2}-4x+3=(x-1)(x-3)=(x-2)^{2}-1$: opens upwards, vertex
        $(2,-1)$, roots $1$ and $3$, $y$-intercept $3$.
  \item $\dfrac{x+1}{x-1}$: pole $x=1$, horizontal asymptote $y=1$ (equal
        degrees); root $x=-1$; $y$-intercept $f(0)=-1$. Two branches, one
        in each half-plane of the pole.
  \item $D=[-4,\infty)$, $W=[0,\infty)$; the curve starts at $(-4,0)$ and
        rises, ever more slowly: the $\sqrt{\;}$ shape moved four units to
        the left.
  \item A V with vertex $(0,-2)$; roots where $|x|=2$, i.e. $x=\pm2$;
        $y$-intercept $-2$.
  \item $x^{3}-3x^{2}=x^{2}(x-3)$: a double root at $x=0$ (touch) and a
        simple root at $x=3$ (crossing); degree $3$, positive leading
        coefficient: down on the left, up on the right.
  \item (a) even; (b) odd; (c) even, since $1/(-x)^{2}=1/x^{2}$;
        (d) neither.
  \item (a) yes, $y=3$ (equal degrees); (b) no: the numerator degree
        exceeds the denominator degree by one, so there is an oblique
        asymptote $y=x-1$ instead; (c) yes, $y=0$; (d) no: $\ln x$ grows
        without bound, however slowly.
  \item Critical points $-2,1,3$. Test one point per interval:
        $x=-3$ gives $\frac{(-4)(-1)}{-6}<0$; $x=0$ gives
        $\frac{(-1)(2)}{-3}>0$; $x=2$ gives $\frac{(1)(4)}{-1}<0$;
        $x=4$ gives $\frac{(3)(6)}{1}>0$. Hence $f>0$ on
        $(-2,1)\cup(3,\infty)$.
  \item $D=(-\infty,-1]\cup[1,\infty)$ (from $x^{2}-1\ge0$); even, so
        symmetric about the $y$-axis; $W=[0,\infty)$; it meets the axis at
        $(-1,0)$ and $(1,0)$ and rises from there, approaching the lines
        $y=\pm x$ from below for large $|x|$.
  \item Circle, centre $(0,0)$, radius $4$. Solving:
        $y=\pm\sqrt{16-x^{2}}$, two branches, each with domain $[-4,4]$;
        the plus sign is the upper semicircle, the minus sign the lower.
  \item Circle with centre $(2,-1)$ and radius $3$.
  \item Ellipse with $a=3$ and $b=2$: it crosses the axes at
        $(\pm3,0)$ and $(0,\pm2)$.
  \item $y=6/x$, a hyperbola with both axes as asymptotes. Since the
        product is positive, $x$ and $y$ share a sign: the curve lies in
        the first and third quadrants.
  \item A parabola opening to the right. $y$ is \emph{not} a function of
        $x$: the line $x=4$ meets the curve at $y=\pm2$. But $x=y^{2}$
        \emph{is} a function of $y$.
  \item Intercepts $(6,0)$ and $(0,4)$; solved for $y$,
        $y=4-\tfrac23 x$, so the slope is $-\tfrac23$.
  \item $(x^{2}-4x+4)+(y^{2}+6y+9)=-9+4+9$, i.e.
        $(x-2)^{2}+(y+3)^{2}=4$: centre $(2,-3)$, radius $2$.
  \item $K^{0.5}L^{0.5}=10\Rightarrow (KL)^{0.5}=10\Rightarrow KL=100
        \Rightarrow L=\dfrac{100}{K}$. A hyperbola in the first quadrant,
        falling and convex, with both axes as asymptotes: the familiar
        isoquant, along which more capital buys ever less labour saving.
  \item Budget line: $y=20-2x$, slope $-2$ (the price ratio
        $-p_{x}/p_{y}=-4/2$), intercepts $(10,0)$ and $(0,20)$.
        Indifference curve: $y=50/x$, a hyperbola in the first quadrant.
\end{solutions}

\section{Additional Problems}

\begin{problems}
  \item Asymptotes of $f(x)=\dfrac{4x-1}{2x+6}$.
  \item Asymptotes and holes of $f(x)=\dfrac{x^{2}-9}{x+3}$.
  \item Asymptotes and symmetry of $f(x)=\dfrac{x^{3}}{x^{2}-1}$.
  \item Asymptotes and holes of $f(x)=\dfrac{x^{2}-2x}{x^{2}-4}$.
  \item Asymptotes of $f(x)=\dfrac{6}{x^{2}-x-2}$.
  \item Asymptotes of $f(x)=\dfrac{x-4}{x^{2}+4}$.
  \item Roots, symmetry and end behaviour of $f(x)=x^{4}-4x^{2}$.
  \item Roots, sign and end behaviour of $f(x)=(x+2)(x-1)(x-3)$.
  \item Symmetry and asymptote of $f(x)=\dfrac{x}{x^{2}+1}$.
  \item Where do $y=3^{x}$ and $y=3^{-x}$ intersect, and what asymptote do
        they share?
  \item The graphs of $y=\mathrm{e}^{x}$ and $y=\ln x$ are mirror images.
        In which line, and which points correspond to each other?
  \item Domain, range, symmetry and asymptotes of $y=1/x^{2}$.
  \item Domain, symmetry and shape of $y=\sqrt[3]{x}$.
  \item Vertex and roots of $y=|x-2|$.
  \item For the cost function $C(q)=q^{3}-6q^{2}+15q$, tabulate
        $C(0),C(1),\dots,C(5)$ and describe the shape of the curve.
  \item Tabulate $f(x)=x^{2}-2x-3$ for $x=-2,\dots,4$ and sketch it from
        the table; give vertex and roots.
  \item Identify $9x^{2}+4y^{2}=36$ and give its axis intercepts.
  \item Identify $x^{2}-y^{2}=9$ and solve it for $y$.
  \item What curve is $(x+1)^{2}+(y-2)^{2}=0$?
  \item Identify $x^{2}+y^{2}+2x-8=0$: centre and radius.
  \item Compare $y=\sqrt{9-x^{2}}$ with $x^{2}+y^{2}=9$: which is a
        function, and what are its domain and range?
  \item The isocost line $5K+2L=100$: give both intercepts and $L$ as a
        function of $K$.
  \item For the indifference curve $xy=36$, tabulate $y$ at
        $x=1,2,3,4,6,9,12$ and describe the shape.
  \item The demand curve $p=100/q$ has the property that revenue $pq$ is
        constant. Give that constant, the asymptotes, and the economic
        reading of each.
  \item A curve passes through the origin, rises throughout, is concave,
        and flattens onto the horizontal line $y=4$. Which of
        $y=4-\mathrm{e}^{-x}$, \;$y=4(1-\mathrm{e}^{-x})$,
        \;$y=4\mathrm{e}^{-x}$ fits, and why do the other two fail?
\end{problems}

\section{Additional Solutions}

\begin{solutions}
  \item Pole where $2x+6=0$, i.e. $x=-3$; equal degrees, so the horizontal
        asymptote is $y=\tfrac42=2$.
  \item $\dfrac{(x-3)(x+3)}{x+3}=x-3$ for $x\neq-3$: a hole at
        $(-3,-6)$, no pole, no horizontal asymptote (the graph is a line).
  \item Poles at $x=\pm1$. Division: $x^{3}=x(x^{2}-1)+x$, so
        $f(x)=x+\dfrac{x}{x^{2}-1}$ and the oblique asymptote is $y=x$.
        Odd: numerator odd, denominator even.
  \item $\dfrac{x(x-2)}{(x-2)(x+2)}=\dfrac{x}{x+2}$ for $x\neq2$: a hole at
        $\bigl(2,\tfrac12\bigr)$, a pole at $x=-2$, horizontal asymptote
        $y=1$.
  \item $x^{2}-x-2=(x-2)(x+1)$: poles at $x=2$ and $x=-1$; horizontal
        asymptote $y=0$.
  \item $x^{2}+4>0$ always: no poles. Numerator degree smaller, so $y=0$.
  \item $x^{4}-4x^{2}=x^{2}(x^{2}-4)=x^{2}(x-2)(x+2)$: double root at $0$
        (touch), simple roots at $\pm2$; even; degree $4$ with positive
        leading coefficient, so both ends rise to $+\infty$.
  \item Roots $-2,1,3$, all simple. Degree $3$, positive leading
        coefficient: down on the left, up on the right. Signs alternate
        between the roots: negative on $(-\infty,-2)$, positive on
        $(-2,1)$, negative on $(1,3)$, positive on $(3,\infty)$.
  \item Odd (odd numerator over even denominator); no poles since
        $x^{2}+1>0$; horizontal asymptote $y=0$ at both ends. The graph is
        a hump above the axis on the right and its mirror image below on
        the left.
  \item $3^{x}=3^{-x}\iff x=-x\iff x=0$: they meet at $(0,1)$. Both have
        the horizontal asymptote $y=0$, the first as $x\to-\infty$, the
        second as $x\to+\infty$; each is the other reflected in the
        $y$-axis.
  \item In the line $y=x$. The point $(a,\mathrm{e}^{a})$ on the
        exponential corresponds to $(\mathrm{e}^{a},a)$ on the logarithm;
        in particular $(0,1)\leftrightarrow(1,0)$ and
        $(1,\mathrm{e})\leftrightarrow(\mathrm{e},1)$. The two functions
        are inverse to each other.
  \item $D=\R\setminus\{0\}$, $W=(0,\infty)$; even, hence symmetric about
        the $y$-axis; vertical asymptote $x=0$ and horizontal asymptote
        $y=0$. Both branches lie above the axis.
  \item $D=\R$: an odd root accepts negative radicands. It is odd, since
        $\sqrt[3]{-x}=-\sqrt[3]{x}$, strictly increasing, steep near the
        origin and flattening out: the graph of $x^{3}$ mirrored in
        $y=x$.
  \item A V with vertex $(2,0)$; a single (double) root at $x=2$;
        $y$-intercept $|0-2|=2$.
  \item $C(0)=0$, $C(1)=10$, $C(2)=14$, $C(3)=18$, $C(4)=28$, $C(5)=50$.
        Increasing throughout; the increments $10,4,4,10,22$ first fall and
        then rise, so the curve is concave up to about $q=2$ and convex
        afterwards: the S-shaped total cost curve of the textbooks, with
        its inflection at $q=2$.
  \item $f(-2)=5$, $f(-1)=0$, $f(0)=-3$, $f(1)=-4$, $f(2)=-3$, $f(3)=0$,
        $f(4)=5$. The table is symmetric about $x=1$: an upward parabola
        with vertex $(1,-4)$ and roots $-1$ and $3$.
  \item Divide by $36$: $\dfrac{x^{2}}{4}+\dfrac{y^{2}}{9}=1$, an ellipse
        with $a=2$, $b=3$, crossing the axes at $(\pm2,0)$ and
        $(0,\pm3)$.
  \item A hyperbola opening left and right. Solved for $y$:
        $y=\pm\sqrt{x^{2}-9}$, defined for $|x|\ge3$, two branches.
  \item A sum of two squares is zero only if both vanish: the ``curve'' is
        the single point $(-1,2)$.
  \item $(x+1)^{2}-1+y^{2}=8$, so $(x+1)^{2}+y^{2}=9$: a circle with centre
        $(-1,0)$ and radius $3$.
  \item $y=\sqrt{9-x^{2}}$ is a function, the upper semicircle, with
        $D=[-3,3]$ and $W=[0,3]$. The full equation $x^{2}+y^{2}=9$ is not
        a function of $x$: it also contains the lower half.
  \item Intercepts $(20,0)$ on the $K$-axis and $(0,50)$ on the $L$-axis;
        $L=50-2.5K$, a falling line with slope $-2.5$, the ratio of the
        input prices.
  \item $y=36/x$: $36,18,12,9,6,4,3$. A falling, convex hyperbola branch in
        the first quadrant, approaching both axes without touching them:
        halving $x$ doubles $y$.
  \item $pq=100$. The asymptotes are $p=0$ and $q=0$: as the quantity
        grows, the price falls towards zero without reaching it, and no
        finite price clears an arbitrarily small quantity. Revenue stays at
        $100$ everywhere, which is what unit elasticity means.
  \item $y=4(1-\mathrm{e}^{-x})$. It gives $y(0)=0$ and rises to the
        asymptote $4$. \;$y=4-\mathrm{e}^{-x}$ has the right asymptote but
        starts at $y(0)=3$, not at the origin; \;$y=4\mathrm{e}^{-x}$ falls
        from $4$ towards $0$ instead of rising.
\end{solutions}

\chapter{Transforming Functions}

Knowing the catalogue of standard shapes is worth much more once it is
combined with the second half of the skill: recognising that an unfamiliar
formula is a familiar function in disguise. The graph of
$y=-2\mathrm{e}^{-(x-3)}+5$ never has to be plotted point by point. It is
the exponential, mirrored twice, stretched, and moved, and all of this is
read off the formula without any arithmetic.

Three related skills belong here. \textbf{Transforming}: what the constants
in $a\,f\bigl(b(x-c)\bigr)+d$ do to a graph, with the standing trap that
everything \emph{inside} the function acts in the opposite direction and by
the reciprocal factor. \textbf{Composing and decomposing}: building a
function out of simpler ones, and taking one apart again, which is exactly
the reading that makes the chain rule mechanical later. And
\textbf{inverting}: swapping the roles of $x$ and $y$, the operation that
turns an exponential into a logarithm and a demand curve into an inverse
demand curve.

The chapter ends with the reverse question, the one an economist meets
most often: not ``what does this formula look like'', but ``which formula has
the shape I need?'' Given a vertex, an asymptote, a saturation level or two
data points, determine the parameters.

\section{Core Rules}

\begin{rulebox}{The four elementary moves}
Start from a known graph $y=f(x)$ and read the general form
\[
  y=a\,f\bigl(b(x-c)\bigr)+d .
\]
\begin{itemize}[leftmargin=1.6em,itemsep=0.3ex]
  \item \textbf{$d$: vertical shift.} $f(x)+d$ moves the graph \emph{up}
        by $d$ (down if $d<0$). Outside the function, and in the direction
        you expect.
  \item \textbf{$c$: horizontal shift.} $f(x-c)$ moves the graph
        \emph{right} by $c$. Inside the function, and \emph{against}
        the sign: $f(x+2)$ moves two units to the \emph{left}.
  \item \textbf{$a$: vertical scaling.} $a\,f(x)$ stretches the graph
        vertically by the factor $|a|$ ($0<|a|<1$ compresses), and
        \textbf{reflects it in the $x$-axis} if $a<0$. Roots stay put;
        every $y$-value is multiplied by $a$.
  \item \textbf{$b$: horizontal scaling.} $f(bx)$ \emph{compresses} the
        graph horizontally by the factor $|b|$ --- again the reciprocal of
        what the number suggests --- and \textbf{reflects it in the
        $y$-axis} if $b<0$. The $y$-intercept stays put.
\end{itemize}
Order matters: inside operations act on $x$ before $f$ is applied, outside
operations act on the result. Apply them as
\emph{horizontal scaling and shifting first, then $f$, then vertical
scaling and shifting}. And factor before reading: $f(2x-6)=f\bigl(2(x-3)\bigr)$
is a compression by $2$ together with a shift of $3$ to the right, not a
shift of $6$.
\end{rulebox}

\begin{gallery}
\begin{tikzpicture}
  \begin{axis}[fnaxis, xmin=-2.6, xmax=4.6, ymin=-1.2, ymax=4.6,
      xtick={-2,2,4}, ytick={2,4}, title={$y=(x-2)^{2}+1$}]
    \addplot[gray!60, dashed, thin, samples=101, domain=-2.1:2.1]{x^2};
    \addplot[domain=0.15:3.85]{(x-2)^2+1};
  \end{axis}
\end{tikzpicture}\hfill
\begin{tikzpicture}
  \begin{axis}[fnaxis, xmin=-2.6, xmax=2.6, ymin=-1.2, ymax=4.6,
      xtick={-2,2}, ytick={2,4}, title={$y=2x^{2}$}]
    \addplot[gray!60, dashed, thin, samples=101, domain=-2.1:2.1]{x^2};
    \addplot[domain=-1.5:1.5]{2*x^2};
  \end{axis}
\end{tikzpicture}\hfill
\begin{tikzpicture}
  \begin{axis}[fnaxis, xmin=-3.4, xmax=3.4, ymin=-1.2, ymax=4.6,
      xtick={-2,2}, ytick={2,4}, title={$y=\mathrm{e}^{-x}$}]
    \addplot[gray!60, dashed, thin, domain=-3.2:1.5]{exp(x)};
    \addplot[domain=-1.5:3.2]{exp(-x)};
  \end{axis}
\end{tikzpicture}
\gcap{Shift right and up \quad\textbullet\quad vertical stretch
\quad\textbullet\quad reflection in the $y$-axis. The original is dashed.}
\end{gallery}

\begin{rulebox}{Rewrite first, then read: other bases, negative inner factors, higher powers}
The four moves can be read off only once the formula has the shape
$a\,f\bigl(b(x-c)\bigr)+d$ for a base function you know. Four rewritings
get it there.
\begin{itemize}[leftmargin=1.6em,itemsep=0.3ex]
  \item \textbf{Exponentials with any base.} For $q>0$, $q\neq1$, the graph
        of $q^{x}$ passes through $(0,1)$ and has the asymptote $y=0$; it
        rises if $q>1$ and falls if $0<q<1$. Since $q^{-x}=(1/q)^{x}$, a
        negative exponent is a reflection in the $y$-axis, turning a rising
        exponential into a falling one. Collect everything into a single
        base before sketching:
        \[
          \frac{a^{x}}{b^{x}}=\Bigl(\frac{a}{b}\Bigr)^{x},\qquad
          a^{2x}=\bigl(a^{2}\bigr)^{x},\qquad
          \text{e.g. }\ \frac{3^{2x}}{2^{x}}=\Bigl(\frac{9}{2}\Bigr)^{x}.
        \]
        Then $y=c\,q^{x}+d$ has domain $\R$, the horizontal asymptote
        $y=d$, and range $(d,\infty)$ if $c>0$, $(-\infty,d)$ if $c<0$.
  \item \textbf{A negative factor inside.} In $\sqrt{k-mx}$ with $m>0$,
        factor out $-m$, not $m$:
        \[
          \sqrt{k-mx}=\sqrt{-m\Bigl(x-\frac{k}{m}\Bigr)} .
        \]
        The factor $-m$ is a reflection in the $y$-axis together with a
        horizontal compression by $m$; then comes a shift \emph{right} by
        $k/m$. The domain is $x\le k/m$: the curve starts at $x=k/m$ and
        runs to the \emph{left}. The same factoring handles $\ln(k-mx)$,
        with domain $x<k/m$ and the asymptote $x=k/m$.
  \item \textbf{Logarithms with an inner factor.}
        $c\ln(ax+b)+d=c\ln\bigl(a(x+\tfrac{b}{a})\bigr)+d$: domain
        $ax+b>0$, vertical asymptote $x=-b/a$, range $\R$. For $a>0$ there
        is a second reading, $c\ln a+c\ln\bigl(x+\tfrac{b}{a}\bigr)+d$: the
        horizontal compression by $a$ is the same as a vertical shift by
        $c\ln a$.
  \item \textbf{Higher powers in a denominator.} $y=\dfrac{c}{(x-a)^{n}}+b$
        has the asymptotes $x=a$ and $y=b$. For \emph{even} $n$ the
        denominator is positive on both sides of the pole, so both branches
        lie on the \emph{same} side of $y=b$ (above it if $c>0$, below if
        $c<0$), and the range is $(b,\infty)$ or $(-\infty,b)$. For
        \emph{odd} $n$ the denominator changes sign at $x=a$, the branches
        lie on \emph{opposite} sides, as for $1/x$, and the range is
        $\R\setminus\{b\}$.
\end{itemize}
Factor powers too: $(2x-6)^{2}=4(x-3)^{2}$, so
$\dfrac{8}{(2x-6)^{2}}=\dfrac{2}{(x-3)^{2}}$, with the pole at $x=3$.
\end{rulebox}

\begin{gallery}
\begin{tikzpicture}
  \begin{axis}[fnaxis, xmin=-3.4, xmax=3.4, ymin=-1.2, ymax=4.6,
      xtick={-2,2}, ytick={2,4}, title={$y=2^{-x}=\bigl(\tfrac12\bigr)^{x}$}]
    \addplot[gray!60, dashed, thin, domain=-3.2:2.1]{2^x};
    \addplot[domain=-2.1:3.2]{2^(-x)};
  \end{axis}
\end{tikzpicture}\hfill
\begin{tikzpicture}
  \begin{axis}[fnaxis, xmin=-3.4, xmax=4.4, ymin=-1.2, ymax=4.6,
      xtick={-2,2,4}, ytick={2,4}, title={$y=\sqrt{4-2x}$}]
    \addplot[gray!60, dashed, thin, samples=101, domain=0:4.2]{sqrt(x)};
    \addplot[samples=201, domain=-3.2:2]{sqrt(4-2*x)};
  \end{axis}
\end{tikzpicture}\hfill
\begin{tikzpicture}
  \begin{axis}[fnaxis, xmin=-2.4, xmax=3.4, ymin=-3.2, ymax=4.6,
      xtick={-2,2}, ytick={-2,2,4},
      title={$\tfrac{1}{(x-1)^{2}}$ and $\tfrac{1}{(x-1)^{3}}$}]
    \addplot[samples=101, domain=-2.3:0.53]{1/(x-1)^2};
    \addplot[samples=101, domain=1.47:3.3]{1/(x-1)^2};
    \addplot[gray!60, dashed, thin, samples=101, domain=-2.3:0.31]{1/(x-1)^3};
    \addplot[gray!60, dashed, thin, samples=101, domain=1.62:3.3]{1/(x-1)^3};
  \end{axis}
\end{tikzpicture}
\gcap{A negative exponent reflects in the $y$-axis \quad\textbullet\quad
$\sqrt{-2(x-2)}$ starts at $x=2$ and runs left \quad\textbullet\quad even
power (solid): both branches above; odd power (dashed): opposite sides.}
\end{gallery}

\begin{rulebox}{Absolute values: two different operations}
\begin{itemize}[leftmargin=1.6em,itemsep=0.3ex]
  \item $y=|f(x)|$ acts on the \textbf{output}: everything below the
        $x$-axis is folded \emph{up}, everything above stays. Roots become
        points where the graph touches the axis and turns.
  \item $y=f(|x|)$ acts on the \textbf{input}: the part of the graph with
        $x\geq0$ is kept and \emph{mirrored} to the left, and whatever the
        function did for $x<0$ is discarded. The result is always even.
\end{itemize}
The two coincide only by accident. Note also $|ax+b|=|a|\,
\bigl|x+\tfrac{b}{a}\bigr|$, which turns any such expression into a shifted
and stretched V.
\end{rulebox}

\begin{gallery}
\begin{tikzpicture}
  \begin{axis}[fnaxis, xmin=-2.6, xmax=2.6, ymin=-2.6, ymax=2.6,
      xtick={-2,2}, ytick={-2,2}, title={$y=f(x)=x^{3}-3x$}]
    \addplot[domain=-2.1:2.1]{x^3-3*x};
  \end{axis}
\end{tikzpicture}\hfill
\begin{tikzpicture}
  \begin{axis}[fnaxis, xmin=-2.6, xmax=2.6, ymin=-2.6, ymax=2.6,
      xtick={-2,2}, ytick={2}, title={$y=|f(x)|$}]
    \addplot[domain=-2.1:2.1]{abs(x^3-3*x)};
  \end{axis}
\end{tikzpicture}\hfill
\begin{tikzpicture}
  \begin{axis}[fnaxis, xmin=-2.6, xmax=2.6, ymin=-2.6, ymax=2.6,
      xtick={-2,2}, ytick={-2,2}, title={$y=f(|x|)$}]
    \addplot[domain=-2.1:2.1]{abs(x)^3-3*abs(x)};
  \end{axis}
\end{tikzpicture}
\end{gallery}

\begin{rulebox}{Composing and decomposing}
$(f\circ g)(x)=f\bigl(g(x)\bigr)$: the \textbf{inner} function $g$ is
applied first. To \emph{decompose} a formula, ask what you would compute
first with a calculator; that is the inner function:
\[
  \begin{aligned}
    \sqrt{3x+1}\;&:\quad g(x)=3x+1, \quad f(u)=\sqrt{u},\\
    \mathrm{e}^{x^{2}}\;&:\quad g(x)=x^{2}, \quad f(u)=\mathrm{e}^{u},\\
    (\ln x)^{3}\;&:\quad g(x)=\ln x, \quad f(u)=u^{3}.
  \end{aligned}
\]
Decompositions are not unique, but the ``innermost bracket first'' reading
is the one the chain rule wants. Note that all four elementary moves are
themselves compositions with linear functions: $f(b(x-c))$ is $f$ composed
with a linear inner function.
\end{rulebox}

\begin{rulebox}{Inverse functions}
$f^{-1}$ undoes $f$: \;$f^{-1}\bigl(f(x)\bigr)=x$ and
$f\bigl(f^{-1}(y)\bigr)=y$. It exists exactly when $f$ is
\textbf{injective}; if it is not, restrict the domain first (as with
$x^{2}$ on $[0,\infty)$). Three consequences and one recipe:
\begin{itemize}[leftmargin=1.6em,itemsep=0.3ex]
  \item $D(f^{-1})=W(f)$ and $W(f^{-1})=D(f)$: domain and range swap.
  \item The graph of $f^{-1}$ is the graph of $f$ \textbf{mirrored in the
        line $y=x$}; the point $(p,q)$ becomes $(q,p)$.
  \item $f^{-1}$ is \emph{not} $1/f$.
  \item \textbf{Recipe:} write $y=f(x)$, solve for $x$, then rename the
        variables. \;$y=\mathrm{e}^{2x}\Rightarrow \ln y=2x\Rightarrow
        x=\tfrac12\ln y$, so $f^{-1}(x)=\tfrac12\ln x$ with
        $D(f^{-1})=(0,\infty)$.
\end{itemize}
This is the operation behind every \emph{inverse demand} function: the
demand $q(p)$ solved for $p$ gives $p(q)$, the same curve read along the
other axis.
\end{rulebox}

\begin{gallery}
\begin{tikzpicture}
  \begin{axis}[fnaxis, xmin=-3.2, xmax=4.4, ymin=-3.2, ymax=4.4,
      xtick={-2,2,4}, ytick={-2,2,4},
      title={$\mathrm{e}^{x}$ and $\ln x$}]
    \addplot[domain=-3:1.45]{exp(x)};
    \addplot[domain=0.05:4.2]{ln(x)};
    \addplot[gray!60, dashed, thin, samples=2, domain=-3:4.2]{x};
  \end{axis}
\end{tikzpicture}\hfill
\begin{tikzpicture}
  \begin{axis}[fnaxis, xmin=-0.8, xmax=4.4, ymin=-0.8, ymax=4.4,
      xtick={2,4}, ytick={2,4},
      title={$x^{2}$ on $[0,\infty)$ and $\sqrt{x}$}]
    \addplot[domain=0:2.05]{x^2};
    \addplot[domain=0:4.2]{sqrt(x)};
    \addplot[gray!60, dashed, thin, samples=2, domain=0:4.2]{x};
  \end{axis}
\end{tikzpicture}
\gcap{Each pair is symmetric about the dashed line $y=x$.}
\end{gallery}

\begin{rulebox}{Reverse engineering: choosing the parameters}
Given the shape you want, pick the family and fit the constants.
\begin{itemize}[leftmargin=1.6em,itemsep=0.3ex]
  \item \textbf{Prescribed vertex:} use $y=a(x-h)^{2}+k$ with the vertex
        $(h,k)$, and determine $a$ from one further point.
  \item \textbf{Prescribed asymptotes:} $y=\dfrac{k}{x-c}+d$ has the
        vertical asymptote $x=c$ and the horizontal one $y=d$; a rational
        $\dfrac{ax+b}{x-c}$ has $y=a$ as horizontal asymptote.
  \item \textbf{Prescribed saturation:} $y=S\bigl(1-\mathrm{e}^{-kx}\bigr)$
        starts at $0$ and rises to the ceiling $S$; $y=S\mathrm{e}^{-kx}$
        starts at $S$ and decays to $0$.
  \item \textbf{Two data points, exponential family:} $y=A\,q^{x}$ with
        $A=y(0)$, and $q$ from the ratio of the two values. Halving every
        $T$ units means $q=2^{-1/T}$.
  \item \textbf{Economic shifts:} a change in fixed cost shifts $C(q)$
        \emph{vertically}; a per-unit tax $t$ adds $tq$ and therefore
        changes the \emph{slope}, not the intercept; measuring the input in
        thousands instead of units rescales $x$ horizontally.
\end{itemize}
\end{rulebox}

\section{Worked Example}

\begin{worked}{Describe the moves taking $y=f(x)$ to $y=-3f(2x-4)+1$, apply them to $f(x)=\sqrt{x}$, and invert $g(x)=2\ln(x-1)+3$}
\wstep{Step 1. Factor the inside before saying anything about the shift.}
The inside is $2x-4$, and it must be written as
\[
  2x-4=2(x-2)
\]
before it can be read. In this form the horizontal instructions are
unambiguous: \textbf{compress by a factor $2$}, and \textbf{shift right by
$2$} --- not by $4$.
\wtrap{This is the chapter's central trap. $f(2x-4)$ shifts by $2$, because
what matters is the value of $x$ that makes the argument zero. Reading the
shift straight off the $-4$ is the most common error with
transformations.}

\wstep{Step 2. Separate inside from outside, and get the directions right.}
\begin{center}\footnotesize
\begin{tabular}{lll}
  \hline
  where & operation & effect on the graph\\
  \hline
  inside, $2(\cdot)$ & multiply $x$ by $2$ & horizontal compression by $2$ \\
  inside, $(\cdot-2)$ & subtract $2$ from $x$ & shift \emph{right} $2$ \\
  outside, $-3(\cdot)$ & multiply by $-3$ & vertical stretch $3$, flip in the $x$-axis\\
  outside, $+1$ & add $1$ & shift \emph{up} $1$ \\
  \hline
\end{tabular}
\end{center}
Inside changes act \textbf{opposite} to how they read (subtracting moves
right, multiplying compresses); outside changes act \textbf{as} they read.
The reason is that inside operations act on the input before $f$ sees it,
so to keep the output the same you must compensate.

\wstep{Step 3. Apply the moves in the right order.}
Inside: shift and scale, in the order that the factored form dictates ---
compress first, then shift, or equivalently shift the compressed picture by
$2$. Outside: stretch and reflect first, then translate. Applying the
vertical shift before the stretch would multiply the shift by $-3$ as well,
landing at $-3f(\cdot)-3$ instead.

\wstep{Step 4. Apply it to $f(x)=\sqrt{x}$.}
\[
  y=-3\sqrt{2x-4}+1=-3\sqrt{2(x-2)}+1 .
\]
The domain moves with the graph: $\sqrt{\cdot}$ needs $2x-4\ge0$, so
$x\ge2$. The starting point $(0,0)$ of $\sqrt{x}$ goes to $(2,1)$: $x=2$
because the shift is right by~$2$, and $y=1$ because $-3\cdot0+1=1$. From
there the curve falls (the reflection), so its range is $(-\infty,1]$.

\wstep{Step 5. Track one more point, as a check.}
The point $(4,2)$ on $y=\sqrt{x}$: solve $2x-4=4$ for the new abscissa,
giving $x=4$; the new ordinate is $-3\cdot2+1=-5$. So $(4,2)\mapsto(4,-5)$,
and substituting $x=4$ into the final formula gives
$-3\sqrt{4}+1=-5$. \checkmark

\wstep{Step 6. Now the inverse of $g(x)=2\ln(x-1)+3$: undo the operations in reverse order.}
Forward, $g$ does: subtract $1$, take $\ln$, double, add $3$. Backwards:
subtract $3$, halve, exponentiate, add $1$. In symbols, set $y=g(x)$ and
solve for $x$:
\[
  y-3=2\ln(x-1)
  \;\Longrightarrow\;
  \frac{y-3}{2}=\ln(x-1)
  \;\Longrightarrow\;
  \mathrm{e}^{(y-3)/2}=x-1
  \;\Longrightarrow\;
  x=1+\mathrm{e}^{(y-3)/2} .
\]
Renaming the variable,
\[
  \boxed{\;g^{-1}(x)=1+\mathrm{e}^{(x-3)/2}\;}
\]

\wstep{Step 7. State the domain and range, which swap.}
$g$ has domain $(1,\infty)$ and range $\R$. Therefore $g^{-1}$ has domain
$\R$ and range $(1,\infty)$ --- visible in the formula, since
$\mathrm{e}^{(x-3)/2}>0$ always, so $g^{-1}>1$ always.
\wtrap{The domain of the inverse is the \emph{range} of the original. This is
not bookkeeping: it is what makes $g^{-1}$ a function at all. And check
$g^{-1}(g(x))=x$ on one value --- $g(1+\mathrm{e}^{0})=g(2)=2\ln1+3=3$ and
$g^{-1}(3)=1+\mathrm{e}^{0}=2$. \checkmark}
\end{worked}

\section{Practice Problems}

In problems 1--10, describe in words how the graph of $y=f(x)$ has to be
moved, scaled or reflected.

\begin{problems}
  \item $y=f(x-3)$
  \item $y=f(x)+2$
  \item $y=-f(x)$
  \item $y=f(-x)$
  \item $y=2f(x)$
  \item $y=f(2x)$
  \item $y=f\!\left(\dfrac{x}{2}\right)+1$
  \item $y=-3f(x-1)+4$
  \item $y=f(x+2)-5$
  \item $y=\tfrac12 f(-x)$
\end{problems}

In problems 11--18, write down the formula.

\begin{problems}[resume]
  \item $y=x^{2}$ shifted $3$ to the right and $1$ up.
  \item $y=\sqrt{x}$ reflected in the $x$-axis and shifted $4$ up.
  \item $y=\mathrm{e}^{x}$ compressed horizontally by the factor $2$.
  \item $y=\ln x$ shifted $5$ to the left.
  \item $y=1/x$ shifted $2$ to the right and $3$ down.
  \item $y=|x|$ shifted $1$ to the left and stretched vertically by $3$.
  \item $y=x^{3}$ reflected in the $y$-axis and shifted $2$ up.
  \item A downward parabola with vertex $(-1,5)$, twice as steep as
        $y=x^{2}$.
\end{problems}

In problems 19--24, name the base function and the transformations, and
give the features asked for.

\begin{problems}[resume]
  \item $y=2(x+1)^{2}-8$: vertex, roots, range.
  \item $y=-\sqrt{x-2}+3$: domain, starting point, range.
  \item $y=3\mathrm{e}^{x+1}$: asymptote, value at $x=-1$.
  \item $y=\ln(2x)$: give \emph{two} readings of this graph, one
        horizontal and one vertical.
  \item $y=\dfrac{1}{x-4}+2$: both asymptotes.
  \item $y=-|x+3|+5$: vertex, roots, direction of opening.
\end{problems}

\begin{problems}[resume]
  \item For $f(x)=x^{2}-2x-3$, describe the graph of $y=|f(x)|$: which
        part moves, and what happens at the roots?
  \item Sketch $y=|2x-6|$ by first writing it as a stretched, shifted V.
  \item For $f(x)=\sqrt{x}$, describe $y=f(|x|)$: give its domain and its
        symmetry.
  \item With $f(u)=u^{2}$ and $g(x)=x-5$: give $(f\circ g)(x)$ and
        $(g\circ f)(x)$, and say which transformation of $u^{2}$ each one
        is.
  \item With $f(u)=\sqrt{u}$ and $g(x)=x^{2}+9$: give $(f\circ g)(x)$ and
        its domain.
  \item Decompose $h(x)=\sqrt{3x+1}$ into $f\circ g$.
  \item Decompose $h(x)=\mathrm{e}^{x^{2}}$ into $f\circ g$.
  \item Decompose $h(x)=(\ln x)^{3}$ into $f\circ g$.
  \item Decompose $h(x)=\dfrac{1}{x^{2}+1}$ into $f\circ g$.
\end{problems}

In problems 34--39, find the inverse together with its domain and range.

\begin{problems}[resume]
  \item $f(x)=3x-7$
  \item $f(x)=\dfrac{x+2}{x-1}$
  \item $f(x)=x^{2}$ on $D=[0,\infty)$
  \item $f(x)=\ln(x-1)$
  \item $f(x)=\sqrt{x+4}$
  \item $f(x)=\dfrac{2x+1}{x-3}$
  \item Verify for $f(x)=3x-7$ that $f\bigl(f^{-1}(x)\bigr)=x$ and
        $f^{-1}\bigl(f(x)\bigr)=x$.
  \item The demand function is $q(p)=50-2p$. Give the inverse demand
        function $p(q)$, and state on which interval of quantities it is
        economically meaningful.
\end{problems}

\begin{problems}[resume]
  \item Determine $a$ so that the parabola with vertex $(2,-3)$ passes
        through $(0,1)$, and write the function.
  \item An exponential $y=A\,q^{x}$ passes through $(0,5)$ and $(2,45)$.
        Determine $A$ and $q$.
  \item Write a rational function with vertical asymptote $x=2$,
        horizontal asymptote $y=3$, and a root at the origin.
  \item Total cost is $C(q)=5q+100$. Write the new cost function and name
        the transformation when (a) the fixed cost rises by $40$;
        (b) a tax of $2$ per unit is levied.
\end{problems}

\section{Solutions}

\begin{solutions}
  \item Shifted $3$ units to the \emph{right} (inside, against the sign).
  \item Shifted $2$ units up.
  \item Reflected in the $x$-axis: every $y$-value changes sign.
  \item Reflected in the $y$-axis: left and right are swapped.
  \item Stretched vertically by the factor $2$; roots and $x$-intercepts
        stay where they are.
  \item Compressed horizontally by the factor $2$: what happened at $x=4$
        now happens at $x=2$. The $y$-intercept is unchanged.
  \item Stretched horizontally by the factor $2$ (the inner factor
        $\tfrac12$ is a compression by $\tfrac12$, i.e. a stretch), then
        moved $1$ up.
  \item Right $1$, stretched vertically by $3$, reflected in the $x$-axis,
        then up $4$ --- inside first, outside afterwards.
  \item Left $2$ and down $5$.
  \item Reflected in the $y$-axis and compressed vertically to half its
        height.
  \item $y=(x-3)^{2}+1$.
  \item $y=-\sqrt{x}+4$.
  \item $y=\mathrm{e}^{2x}$ (a factor $2$ \emph{inside} compresses).
  \item $y=\ln(x+5)$, with the vertical asymptote moved to $x=-5$.
  \item $y=\dfrac{1}{x-2}-3$, asymptotes $x=2$ and $y=-3$.
  \item $y=3|x+1|$.
  \item $y=(-x)^{3}+2=-x^{3}+2$.
  \item $y=-2(x+1)^{2}+5$.
  \item Base $x^{2}$: left $1$, stretched by $2$, down $8$. Vertex
        $(-1,-8)$; roots from $2(x+1)^{2}=8\iff(x+1)^{2}=4\iff x=1$ or
        $x=-3$; range $[-8,\infty)$.
  \item Base $\sqrt{x}$: right $2$, reflected in the $x$-axis, up $3$.
        $D=[2,\infty)$, the curve starts at $(2,3)$ and falls,
        $W=(-\infty,3]$.
  \item Base $\mathrm{e}^{x}$: left $1$, stretched by $3$. Asymptote
        $y=0$; $y(-1)=3\mathrm{e}^{0}=3$. (Equivalently
        $3\mathrm{e}\cdot\mathrm{e}^{x}$: a horizontal shift of an
        exponential is also a vertical stretch.)
  \item Horizontally: $\ln x$ compressed by the factor $2$. Vertically:
        since $\ln(2x)=\ln 2+\ln x$, the same graph is $\ln x$ shifted up
        by $\ln2\approx0.69$. Both readings are correct --- a peculiarity
        of the logarithm.
  \item Base $1/x$: right $4$, up $2$. Vertical asymptote $x=4$,
        horizontal asymptote $y=2$.
  \item Base $|x|$: left $3$, reflected (so it opens downwards), up $5$.
        Vertex $(-3,5)$; roots where $|x+3|=5$, i.e. $x=2$ and $x=-8$.
  \item $f(x)=(x-3)(x+1)$ is negative exactly on $(-1,3)$, and that arc is
        folded up above the axis; outside it the graph is unchanged. At the
        roots $x=-1$ and $x=3$ the curve now touches the axis and turns
        instead of crossing it, and the old minimum $f(1)=-4$ becomes a
        local maximum of height $4$.
  \item $|2x-6|=2|x-3|$: a V with vertex $(3,0)$, twice as steep as
        $|x|$, i.e. with slopes $-2$ and $+2$; $y$-intercept $6$.
  \item $y=\sqrt{|x|}$: $D=\R$ now, since $|x|\geq0$ always. The right half
        is unchanged and mirrored to the left, so the function is even,
        with range $[0,\infty)$ and a sharp point at the origin.
  \item $(f\circ g)(x)=(x-5)^{2}$: the parabola shifted $5$ to the right.
        $(g\circ f)(x)=x^{2}-5$: the parabola shifted $5$ down. Inner
        shifts move horizontally, outer shifts vertically.
  \item $(f\circ g)(x)=\sqrt{x^{2}+9}$, and since $x^{2}+9\geq9>0$ for
        every $x$, $D=\R$. (Its range is $[3,\infty)$.)
  \item $g(x)=3x+1$, $f(u)=\sqrt{u}$.
  \item $g(x)=x^{2}$, $f(u)=\mathrm{e}^{u}$.
  \item $g(x)=\ln x$, $f(u)=u^{3}$.
  \item $g(x)=x^{2}+1$, $f(u)=\dfrac{1}{u}$.
  \item $y=3x-7\Rightarrow x=\dfrac{y+7}{3}$, so
        $f^{-1}(x)=\dfrac{x+7}{3}$, with $D=W=\R$.
  \item $y=\dfrac{x+2}{x-1}\Rightarrow y(x-1)=x+2\Rightarrow x(y-1)=y+2
        \Rightarrow x=\dfrac{y+2}{y-1}$. Hence
        $f^{-1}(x)=\dfrac{x+2}{x-1}=f(x)$: the function is its own
        inverse, and its graph is symmetric about $y=x$. Domain and range
        are both $\R\setminus\{1\}$.
  \item $f^{-1}(x)=\sqrt{x}$, with $D(f^{-1})=[0,\infty)=W(f)$ and
        $W(f^{-1})=[0,\infty)=D(f)$. The restriction to $[0,\infty)$ was
        needed to make $f$ injective in the first place.
  \item $y=\ln(x-1)\Rightarrow \mathrm{e}^{y}=x-1\Rightarrow
        x=\mathrm{e}^{y}+1$: $f^{-1}(x)=\mathrm{e}^{x}+1$, with $D=\R$ and
        $W=(1,\infty)$ --- exactly the domain $(1,\infty)$ of $f$.
  \item $y=\sqrt{x+4}\Rightarrow y^{2}=x+4\Rightarrow x=y^{2}-4$, and only
        $y\geq0$ occurs: $f^{-1}(x)=x^{2}-4$ with $D=[0,\infty)$ and
        $W=[-4,\infty)$. Without the restriction $D=[0,\infty)$ the formula
        $x^{2}-4$ would not be the inverse.
  \item $y=\dfrac{2x+1}{x-3}\Rightarrow y(x-3)=2x+1\Rightarrow
        x(y-2)=3y+1\Rightarrow x=\dfrac{3y+1}{y-2}$. So
        $f^{-1}(x)=\dfrac{3x+1}{x-2}$, with $D=\R\setminus\{2\}$ (the
        range of $f$) and $W=\R\setminus\{3\}$ (the domain of $f$). The two
        asymptotes have swapped roles.
  \item $f\bigl(f^{-1}(x)\bigr)=3\cdot\dfrac{x+7}{3}-7=x+7-7=x$, and
        $f^{-1}\bigl(f(x)\bigr)=\dfrac{(3x-7)+7}{3}=\dfrac{3x}{3}=x$.
  \item $q=50-2p\Rightarrow 2p=50-q\Rightarrow p(q)=25-\dfrac{q}{2}$. It
        is meaningful for $q\in[0,50]$: at $q=50$ the price has fallen to
        $0$, and larger quantities would require a negative price.
  \item $y=a(x-2)^{2}-3$ with $1=a\cdot4-3$, so $a=1$ and
        $y=(x-2)^{2}-3$.
  \item $A=y(0)=5$. Then $45=5q^{2}\Rightarrow q^{2}=9\Rightarrow q=3$
        (positive base), so $y=5\cdot3^{x}$.
  \item $y=\dfrac{3x}{x-2}$: the denominator gives the pole $x=2$, equal
        degrees give the horizontal asymptote $y=3$, and the numerator
        vanishes at $x=0$, so the graph passes through the origin.
  \item (a) $C(q)=5q+140$: a pure vertical shift by $40$; the curve keeps
        its slope. (b) $C(q)=7q+100$: the tax adds $2q$, so the
        \emph{slope} changes from $5$ to $7$ while the intercept stays.
        A vertical shift and a change of slope are different operations,
        and only the first one is a translation.
\end{solutions}

\section{Additional Problems}

\begin{problems}
  \item $y=(x+4)^{2}-2$: base function, transformations, vertex, range.
  \item $y=-2\sqrt{x}+1$: transformations, domain, range.
  \item $y=\mathrm{e}^{-(x-1)}$: transformations, asymptote, value at
        $x=1$.
  \item $y=\ln(x+3)$: transformation, vertical asymptote, domain.
  \item $y=\dfrac{3}{x+2}-1$: transformations and both asymptotes.
  \item $y=\tfrac12|x-6|$: transformations and vertex.
  \item Write: $y=1/x$ reflected in the $x$-axis and shifted $3$ up.
  \item Write: $y=\mathrm{e}^{x}$ compressed horizontally by the factor
        $3$ and shifted $1$ down.
  \item Write: $y=\sqrt{x}$ stretched vertically by $5$ and shifted $1$ to
        the right.
  \item Write: a downward parabola with vertex $(3,-4)$, twice as steep as
        $y=x^{2}$.
  \item Describe the graph of $y=|x^{2}-1|$: which arc is folded, and what
        does the graph do at $x=\pm1$?
  \item Describe the graph of $y=|\ln x|$: domain, roots, and behaviour as
        $x\to0^{+}$.
  \item For $f(x)=x^{3}-3x$, give the formula for $f(|x|)$ and state its
        symmetry.
  \item With $f(u)=u^{2}$ and $g(x)=x+3$: give both composites and say why
        they differ.
  \item Decompose $h(x)=\sqrt{1-x^{2}}$ into $f\circ g$ and give its
        domain.
  \item Decompose $h(x)=\mathrm{e}^{\sqrt{x}}$ into $f\circ g$.
  \item Decompose $h(x)=(2x+5)^{10}$ into $f\circ g$.
  \item Inverse of $f(x)=5-2x$.
  \item Inverse of $f(x)=x^{3}+1$.
  \item Inverse of $f(x)=\dfrac{1}{x-2}$, with domain and range.
  \item Inverse of $f(x)=10^{x}$, with domain and range.
  \item Inverse of $f(x)=\ln(3x)$, with domain and range.
  \item Why does $f(x)=x^{2}$ have no inverse on $\R$, and which two
        restrictions repair it?
  \item A quantity is $200$ at $x=0$ and halves every $5$ units. Write it
        both as $A\,q^{x}$ and in the form $A\,\mathrm{e}^{-kx}$.
  \item Cost is $C(q)=3q+40$. Fixed cost rises to $60$ and the unit cost
        falls to $2.5$. Write the new function and name both changes as
        transformations of the graph.
\end{problems}

In problems 26--36, rewrite the formula first (single base, factored
inside, factored power), then give the chain of moves from the base
function, the maximal domain, the range and every asymptote.

\begin{problems}[resume]
  \item $y=3-2\cdot5^{-x}$
  \item $y=1+3\cdot\bigl(\tfrac25\bigr)^{x}$
  \item $y=\dfrac{12^{x}}{3^{x}}-2$
  \item $y=4-\dfrac{2^{2x}}{8^{x}}$
  \item $y=2\sqrt{6-2x}+1$
  \item $y=2-\sqrt{5-x}$
  \item $y=3\ln(2x+6)-1$
  \item $y=1-2\ln(4-x)$
  \item $y=\dfrac{2}{(x-1)^{2}}-3$
  \item $y=1-\dfrac{1}{(x+2)^{3}}$
  \item $y=\dfrac{4}{(2x-6)^{2}}+1$
\end{problems}

\section{Additional Solutions}

\begin{solutions}
  \item Base $x^{2}$, shifted $4$ to the left and $2$ down; vertex
        $(-4,-2)$; range $[-2,\infty)$.
  \item Base $\sqrt{x}$: stretched by $2$, reflected in the $x$-axis, then
        up $1$. $D=[0,\infty)$, and since $-2\sqrt{x}\in(-\infty,0]$,
        $W=(-\infty,1]$.
  \item $\mathrm{e}^{-(x-1)}$: reflect $\mathrm{e}^{x}$ in the $y$-axis,
        then shift $1$ to the right (the inner expression is
        $-(x-1)$). Asymptote $y=0$; $y(1)=\mathrm{e}^{0}=1$.
  \item $\ln x$ shifted $3$ to the left; the asymptote moves with it to
        $x=-3$; $D=(-3,\infty)$.
  \item Base $1/x$: stretched by $3$, left $2$, down $1$. Vertical
        asymptote $x=-2$, horizontal asymptote $y=-1$.
  \item Base $|x|$: right $6$ and compressed vertically to half the height;
        vertex $(6,0)$, slopes $\pm\tfrac12$.
  \item $y=-\dfrac{1}{x}+3$.
  \item $y=\mathrm{e}^{3x}-1$.
  \item $y=5\sqrt{x-1}$.
  \item $y=-2(x-3)^{2}-4$.
  \item The arc between the roots, where $x^{2}-1<0$, i.e. over
        $(-1,1)$, is folded up; it becomes a hump of height $1$ at $x=0$.
        At $x=\pm1$ the graph touches the axis and turns, forming two
        corners.
  \item $D=(0,\infty)$; the single root is $x=1$; on $(0,1)$ the logarithm
        is negative and is folded up, so the graph falls from $+\infty$ to
        $0$ there and rises again afterwards, with a corner at $(1,0)$. As
        $x\to0^{+}$ it goes to $+\infty$.
  \item $f(|x|)=|x|^{3}-3|x|$. It is even: the right half of the original
        cubic is kept and mirrored into the left half-plane, so the old
        left branch is discarded.
  \item $(f\circ g)(x)=(x+3)^{2}$, the parabola shifted $3$ to the left;
        $(g\circ f)(x)=x^{2}+3$, the parabola shifted $3$ up. Composition
        is not commutative: an inner shift moves the graph horizontally, an
        outer one vertically.
  \item $g(x)=1-x^{2}$, $f(u)=\sqrt{u}$; $D=[-1,1]$ from $1-x^{2}\geq0$.
  \item $g(x)=\sqrt{x}$, $f(u)=\mathrm{e}^{u}$; $D=[0,\infty)$.
  \item $g(x)=2x+5$, $f(u)=u^{10}$.
  \item $y=5-2x\Rightarrow x=\dfrac{5-y}{2}$, so
        $f^{-1}(x)=\dfrac{5-x}{2}$.
  \item $y=x^{3}+1\Rightarrow x=\sqrt[3]{y-1}$, so
        $f^{-1}(x)=\sqrt[3]{x-1}$, with $D=W=\R$.
  \item $y=\dfrac{1}{x-2}\Rightarrow x-2=\dfrac1y\Rightarrow
        x=2+\dfrac1y$: $f^{-1}(x)=2+\dfrac{1}{x}$, with
        $D=\R\setminus\{0\}$ and $W=\R\setminus\{2\}$.
  \item $f^{-1}(x)=\log_{10}x$, with $D=(0,\infty)$ and $W=\R$.
  \item $y=\ln(3x)\Rightarrow \mathrm{e}^{y}=3x\Rightarrow
        x=\dfrac{\mathrm{e}^{y}}{3}$: $f^{-1}(x)=\dfrac{\mathrm{e}^{x}}{3}$,
        with $D=\R$ and $W=(0,\infty)$.
  \item It is not injective: $f(-2)=f(2)=4$, so a horizontal line meets the
        graph twice and the inverse assignment would not be unique.
        Restricting to $[0,\infty)$ gives the inverse $\sqrt{x}$;
        restricting to $(-\infty,0]$ gives $-\sqrt{x}$.
  \item Halving every $5$ units means $q=2^{-1/5}$, so
        $y=200\cdot2^{-x/5}$. Since $2^{-x/5}
        =\mathrm{e}^{-(\ln2/5)x}$, the same function is
        $y=200\,\mathrm{e}^{-kx}$ with
        $k=\dfrac{\ln2}{5}\approx0.1386$.
  \item $C(q)=2.5q+60$. Raising the fixed cost from $40$ to $60$ shifts the
        line $20$ units up; lowering the unit cost from $3$ to $2.5$
        flattens it, a change of slope rather than a translation. Only the
        first is a shift of the graph.
  \item Rewrite $5^{-x}=\bigl(\tfrac15\bigr)^{x}$. From $y=5^{x}$:
        reflect in the $y$-axis (the graph now falls), stretch vertically by
        $2$, reflect in the $x$-axis (it now rises again), shift up $3$.
        $D=\R$. Since $2\cdot5^{-x}$ takes every value in $(0,\infty)$,
        $-2\cdot5^{-x}$ takes every value in $(-\infty,0)$ and
        $W=(-\infty,3)$. Horizontal asymptote $y=3$, approached from below
        as $x\to\infty$; as $x\to-\infty$ the graph falls without bound.
        Check: $y(0)=3-2=1$. Root: $5^{-x}=\tfrac32$, so
        $x=-\dfrac{\ln1.5}{\ln5}\approx-0.25$.
  \item The base $\tfrac25$ lies between $0$ and $1$, so $(\tfrac25)^{x}$
        is a \emph{falling} exponential through $(0,1)$. Stretch vertically
        by $3$, shift up $1$. $D=\R$, $W=(1,\infty)$, horizontal asymptote
        $y=1$, approached from above as $x\to\infty$. $y(0)=1+3=4$.
        (Equivalently $3\cdot(\tfrac52)^{-x}+1$: the rising exponential
        $(\tfrac52)^{x}$ reflected in the $y$-axis.)
  \item Rewrite $\dfrac{12^{x}}{3^{x}}=\Bigl(\dfrac{12}{3}\Bigr)^{x}=4^{x}$,
        so $y=4^{x}-2$: the rising exponential $4^{x}$ shifted $2$ down.
        $D=\R$, $W=(-2,\infty)$, horizontal asymptote $y=-2$ as
        $x\to-\infty$. $y(0)=-1$; root $4^{x}=2$, i.e. $x=\tfrac12$.
  \item Rewrite $2^{2x}=4^{x}$, so $\dfrac{2^{2x}}{8^{x}}
        =\Bigl(\dfrac48\Bigr)^{x}=\bigl(\tfrac12\bigr)^{x}=2^{-x}$ and
        $y=4-2^{-x}$. From $2^{x}$: reflect in the $y$-axis, reflect in the
        $x$-axis, shift up $4$. $D=\R$, $W=(-\infty,4)$, horizontal
        asymptote $y=4$ as $x\to\infty$. $y(0)=3$; root $2^{-x}=4=2^{2}$,
        so $x=-2$.
  \item Factor out $-2$: $6-2x=-2(x-3)$, so $y=2\sqrt{-2(x-3)}+1$. From
        $\sqrt{x}$: reflect in the $y$-axis and compress horizontally by
        $2$ (giving $\sqrt{-2x}$), shift right $3$, stretch vertically by
        $2$, shift up $1$. Domain from $6-2x\ge0$: $D=(-\infty,3]$. The
        curve starts at $(3,1)$ and rises to the left without bound, so
        $W=[1,\infty)$; no asymptote. Check: $y(1)=2\sqrt4+1=5$.
  \item Factor out $-1$: $y=-\sqrt{-(x-5)}+2$. From $\sqrt{x}$: reflect in
        the $y$-axis, shift right $5$, reflect in the $x$-axis, shift up
        $2$. $D=(-\infty,5]$, starting point $(5,2)$, from which the curve
        falls to the left: $W=(-\infty,2]$. No asymptote. Root:
        $\sqrt{5-x}=2$, so $x=1$.
  \item Factor the inside: $2x+6=2(x+3)$, so $y=3\ln\bigl(2(x+3)\bigr)-1$.
        From $\ln x$: compress horizontally by $2$, shift left $3$, stretch
        vertically by $3$, shift down $1$. $D=(-3,\infty)$ from $2x+6>0$;
        vertical asymptote $x=-3$; $W=\R$. Second reading:
        $y=3\ln(x+3)+(3\ln2-1)$, i.e. $\ln(x+3)$ stretched by $3$ and
        shifted up by $3\ln2-1\approx1.08$. Root: $\ln(2x+6)=\tfrac13$,
        $x=\tfrac12\bigl(\mathrm{e}^{1/3}-6\bigr)\approx-2.30$.
  \item Factor out $-1$: $y=-2\ln\bigl(-(x-4)\bigr)+1$. From $\ln x$:
        reflect in the $y$-axis, shift right $4$, stretch vertically by
        $2$, reflect in the $x$-axis, shift up $1$. $D=(-\infty,4)$ from
        $4-x>0$; vertical asymptote $x=4$. As $x\to4^{-}$,
        $\ln(4-x)\to-\infty$, so $y\to+\infty$; the two reflections together
        leave an \emph{increasing} curve. $W=\R$. Root: $\ln(4-x)=\tfrac12$,
        $x=4-\sqrt{\mathrm{e}}\approx2.35$.
  \item Base $1/x^{2}$ (even power): stretch vertically by $2$, shift
        right $1$, down $3$. $D=\R\setminus\{1\}$, vertical asymptote $x=1$,
        horizontal asymptote $y=-3$. Since $\dfrac{2}{(x-1)^{2}}>0$ on both
        sides, both branches lie above $y=-3$: $W=(-3,\infty)$. Roots:
        $(x-1)^{2}=\tfrac23$, $x=1\pm\sqrt{2/3}\approx1\pm0.82$.
  \item Base $1/x^{3}$ (odd power): reflect in the $x$-axis, shift left
        $2$, up $1$. $D=\R\setminus\{-2\}$, asymptotes $x=-2$ and $y=1$.
        For $x>-2$ the cube $(x+2)^{3}$ is positive and $y<1$; for $x<-2$
        it is negative and $y>1$: the branches lie on opposite sides, and
        $W=\R\setminus\{1\}$. Root: $(x+2)^{3}=1$, $x=-1$; $y(0)=\tfrac78$.
  \item Factor the power: $(2x-6)^{2}=\bigl(2(x-3)\bigr)^{2}=4(x-3)^{2}$,
        so $y=\dfrac{1}{(x-3)^{2}}+1$. Base $1/x^{2}$: shift right $3$, up
        $1$. $D=\R\setminus\{3\}$, asymptotes $x=3$ and $y=1$, both branches
        above $y=1$: $W=(1,\infty)$. No roots.
\end{solutions}

\chapter{Inequalities in Two Variables: Regions in the Plane}

The algebra volume solved inequalities in one variable, where the answer is
a set of points on a line --- an interval, or a union of a few. With two variables
the same question has a two-dimensional answer: the solution set of
$2x+3y\le12$ is not a list of numbers but a \emph{region} of the plane, and
it is described by drawing it.

The mechanism is short. The corresponding \emph{equation} $2x+3y=12$ is a
curve, and a curve cuts the plane into pieces. On each piece the expression
$2x+3y-12$ cannot change sign without passing through zero, so its sign is
constant there. One test point per piece therefore settles the whole
region --- the same sign-table logic as in one variable, with curves in
place of points.

Two directions are worth drilling, and examinations favour the second. The
forward direction is: given inequalities, shade the region. The reverse is:
given a shaded region, write down the inequalities that describe it. The
reverse is harder and more useful, because it is what modelling actually
requires --- a feasible set is described in words or in a picture long before
anyone writes the constraints. It is also exactly the skill that the
graphical solution of a linear programme depends on.

\section{Core Rules}

\begin{rulebox}{The solution set is a region}
For an expression $F(x,y)$, the set $\{F(x,y)<0\}$ is bounded by the curve
$\{F(x,y)=0\}$. The recipe:
\begin{enumerate}[leftmargin=1.8em,itemsep=0.3ex]
  \item \textbf{Draw the boundary} $F(x,y)=0$. It is one of the curves of
        Chapter~2: a line, circle, ellipse, parabola or hyperbola.
  \item \textbf{Note which pieces it creates.} A line makes two half-planes;
        a circle makes an inside and an outside; a parabola makes an
        ``above'' and a ``below''.
  \item \textbf{Test one point per piece.} Substituting a single convenient
        point --- the origin whenever it is not on the boundary --- decides
        the sign for that whole piece.
  \item \textbf{Mark the boundary.} Solid if the inequality is
        $\le$ or $\ge$ (the boundary belongs to the set), dashed if it is
        $<$ or $>$ (it does not).
\end{enumerate}
\end{rulebox}

\begin{rulebox}{The standard regions}
\begin{itemize}[leftmargin=1.6em,itemsep=0.3ex]
  \item \textbf{Half-plane:} $ax+by\le c$. Draw the line, test the origin.
        With $c>0$ and the origin satisfying it, the region is the side
        containing the origin.
  \item \textbf{Disc and its complement:} $x^{2}+y^{2}\le r^{2}$ is the
        closed disc; $\ge$ gives everything outside. Shifted:
        $(x-x_{0})^{2}+(y-y_{0})^{2}\le r^{2}$.
  \item \textbf{Ellipse interior:} $\dfrac{x^{2}}{a^{2}}
        +\dfrac{y^{2}}{b^{2}}\le1$.
  \item \textbf{Above or below a parabola:} $y\ge x^{2}$ is the region
        \emph{inside} the parabola's opening; $y\le x^{2}$ is everything
        outside it. Rewriting an inequality as $y\ge\cdots$ or
        $y\le\cdots$ makes the side unambiguous: ``$y$ at least this'' means
        upwards.
  \item \textbf{Between hyperbola branches:} $xy\ge c$ with $c>0$ is the
        pair of regions beyond the two branches --- and on the positive
        quadrant alone it is the convex ``better set'' of Cobb--Douglas.
  \item \textbf{Strip:} $|x|\le a$ is the vertical strip $-a\le x\le a$;
        $|y|\le b$ a horizontal one.
\end{itemize}
\end{rulebox}

\begin{gallery}
\begin{tikzpicture}
  \begin{axis}[geoaxis, xmin=-1.4, xmax=7.4, ymin=-1.4, ymax=5.4,
      xtick={2,4,6}, ytick={2,4}, title={$2x+3y\le12$, $x,y\ge0$}]
    \addplot[draw=none, name path=lin, domain=0:6, samples=2]{(12-2*x)/3};
    \addplot[draw=none, name path=ax, domain=0:6, samples=2]{0};
    \region{lin}{ax}{}
    \addplot[domain=-1.2:7.2, samples=2]{(12-2*x)/3};
  \end{axis}
\end{tikzpicture}\hfill
\begin{tikzpicture}
  \begin{axis}[geoaxis, xmin=-3.4, xmax=3.4, ymin=-3.4, ymax=3.4,
      xtick={-2,2}, ytick={-2,2}, title={$x^{2}+y^{2}\le4$}]
    \addplot[draw=none, name path=up, domain=-2:2, samples=201]{sqrt(4-x^2)};
    \addplot[draw=none, name path=lo, domain=-2:2, samples=201]{-sqrt(4-x^2)};
    \region{up}{lo}{}
    \addplot[domain=0:360, samples=181]({2*cos(x)},{2*sin(x)});
  \end{axis}
\end{tikzpicture}\hfill
\begin{tikzpicture}
  \begin{axis}[geoaxis, xmin=-2.6, xmax=2.6, ymin=-0.8, ymax=4.6,
      xtick={-2,2}, ytick={2,4}, title={$x^{2}\le y\le 4$}]
    \addplot[draw=none, name path=par, domain=-2:2, samples=201]{x^2};
    \addplot[draw=none, name path=top, domain=-2:2, samples=2]{4};
    \region{par}{top}{}
    \addplot[domain=-2.15:2.15, samples=201]{x^2};
    \addplot[domain=-2.5:2.5, samples=2]{4};
  \end{axis}
\end{tikzpicture}
\gcap{Three regions: an intersection of three half-planes (a triangle), a
disc, and the region caught between a parabola and a line.}
\end{gallery}

\begin{rulebox}{Systems: intersect}
Several inequalities holding \emph{simultaneously} give the
\textbf{intersection} of their regions --- shade each one and keep only the
overlap. Two consequences:
\begin{itemize}[leftmargin=1.6em,itemsep=0.3ex]
  \item An intersection of half-planes is a \textbf{convex polygon}
        (possibly unbounded). This is precisely the feasible region of a
        linear programme, and the reason the optimum of an LP sits at a
        corner.
  \item An intersection of convex regions is always convex, so discs,
        ellipse interiors and half-planes may be combined freely and the
        result stays convex. A \emph{union}, by contrast, usually is not,
        and neither is the outside of a disc.
\end{itemize}
\end{rulebox}

\begin{rulebox}{The reverse direction: from picture to inequalities}
Given a shaded region, work boundary by boundary.
\begin{enumerate}[leftmargin=1.8em,itemsep=0.3ex]
  \item \textbf{Identify every boundary curve} and write its equation. A
        circle needs its centre and radius; a line two points or an
        intercept and a slope; a parabola its vertex and one further
        point.
  \item \textbf{Fix the direction of each} with a point that is clearly
        inside the shaded region: substitute it and keep the inequality
        sign that comes out true.
  \item \textbf{Join the conditions with ``and''}, one per boundary. A region
        bounded by three curves needs three inequalities.
  \item \textbf{Check a point just outside} each boundary: it must fail
        exactly one condition. If it fails none, a condition is missing.
\end{enumerate}
Read the boundary style too: a dashed boundary means a strict inequality.
\end{rulebox}

\begin{rulebox}{Which of these regions are convex}
One rule settles almost every case. For a function $f$ of one variable:
\begin{itemize}[leftmargin=1.6em,itemsep=0.3ex]
  \item $f$ \textbf{convex} $\Longrightarrow$ the region \emph{above} it,
        $\{y\ge f(x)\}$ (its \de{Epigraph}), is a convex set. The region
        \emph{below} it is then not convex.
  \item $f$ \textbf{concave} $\Longrightarrow$ the region \emph{below} it,
        $\{y\le f(x)\}$, is convex. The region above it is then not.
\end{itemize}
The proof is one line: if $y_{i}\ge f(x_{i})$ and $f$ is convex then
$ty_{1}+(1-t)y_{2}\ge tf(x_{1})+(1-t)f(x_{2})\ge f(tx_{1}+(1-t)x_{2})$.
So $y\ge x^{2}$ is convex and $y\le x^{2}$ is not, while $y\le4-x^{2}$
\emph{is} convex, because $4-x^{2}$ is concave. Add to this that
half-planes, discs and ellipse interiors are convex, that any
\textbf{intersection} of convex sets is convex, and that the
\emph{outside} of a disc and any \emph{union} of separated pieces is not,
and every region in this chapter is decided.
\end{rulebox}

\begin{rulebox}{Absolute values in two variables}
Split on the sign, exactly as in one variable, or recognise the standard
shapes:
\[
  |x|+|y|\le c \;\text{(a diamond, or square rotated }45^{\circ}),\qquad
  \max(|x|,|y|)\le c \;\text{(a square)} .
\]
For $|y|\le f(x)$ the region is squeezed between $y=f(x)$ and $y=-f(x)$,
and exists only where $f(x)\ge0$. For $y\ge|x|$ it is the wedge above both
lines $y=x$ and $y=-x$.
\end{rulebox}

\section{Worked Example}

\begin{worked}{Sketch $S=\{(x,y):x^{2}+y^{2}\le16,\ y\ge0,\ x\le2\}$, then decide whether $\{(x,y):y\ge\sqrt{x}\}$ is convex}
\wstep{Step 1. Take the constraints one at a time, and shade each.}
\begin{itemize}[leftmargin=1.6em,itemsep=0.2ex]
  \item $x^{2}+y^{2}\le16$: the closed disc of radius $4$ about the origin
        --- ``$\le$'' is the inside, because the origin satisfies it
        ($0\le16$).
  \item $y\ge0$: the upper half-plane, boundary included.
  \item $x\le2$: everything to the left of the vertical line $x=2$.
\end{itemize}
\wtrap{Decide inside-versus-outside by testing one point, not by staring at
the inequality sign. The origin is free and settles the disc immediately.
For a line not through the origin, any convenient point does.}

\wstep{Step 2. Intersect the three.}
The region is the part of the upper half-disc lying left of $x=2$: a
half-disc with a circular segment removed on the right. Its boundary
consists of three pieces --- the segment of the $x$-axis from $x=-4$ to
$x=2$, the vertical segment $x=2$ from $y=0$ up to $y=\sqrt{16-4}=2\sqrt3$,
and the arc of the circle from $(2,2\sqrt3)$ counter-clockwise round to
$(-4,0)$.

\wstep{Step 3. Find the corner points explicitly.}
Where $x=2$ meets the circle: $4+y^{2}=16$, so $y=2\sqrt3\approx3.46$ (the
negative root is cut off by $y\ge0$). Where $y=0$ meets the circle:
$x=\pm4$, of which $x=4$ is cut off by $x\le2$, leaving $(-4,0)$. And
$(2,0)$ where the two straight boundaries meet. Three corners:
$(-4,0)$, $(2,0)$, $(2,2\sqrt3)$.

\wstep{Step 4. Sketch, and check one interior point.}
\begin{center}
\begin{tikzpicture}
\begin{axis}[geoaxis, width=6.6cm, height=6.6cm,
             xmin=-5.2, xmax=5.2, ymin=-1.2, ymax=5.2,
             xtick={-4,-2,2,4}, ytick={2,3.464},
             yticklabels={$2$,$2\sqrt3$}]
  % the region: under the upper semicircle, above y=0, left of x=2
  \addplot[accent, very thick, domain=-4:2, samples=200,
           fill=regionfill, fill opacity=0.5] {sqrt(16-x^2)} \closedcycle;
  % the discarded part of the disc, for contrast
  \addplot[accent!35, thick, dashed, domain=2:4, samples=80]
           {sqrt(16-x^2)};
  \node[font=\scriptsize] at (axis cs:-1.2,1.4) {$S$};
  \node[font=\tiny, anchor=south west] at (axis cs:2.1,3.6) {$(2,2\sqrt3)$};
\end{axis}
\end{tikzpicture}
\end{center}
Test $(0,1)$: $0+1\le16$ \checkmark, $1\ge0$ \checkmark, $0\le2$
\checkmark. Inside, as drawn.

\wstep{Step 5. Note that $S$ is convex, and why you can say so instantly.}
Each of the three constraint sets is convex --- a disc, a half-plane, a
half-plane --- and \textbf{an intersection of convex sets is convex}. No
picture is needed for that conclusion, and no amount of picture-gazing
substitutes for it.

\wstep{Step 6. The second region: $\{y\ge\sqrt{x}\}$. Do not trust the eye.}
The boundary $y=\sqrt{x}$ is a smooth, gently bending curve, and the region
above it looks perfectly well behaved. Test the definition instead. Take the
two boundary points
\[
  P=(0,0),\qquad Q=(4,2),
\]
both in the set, since $0\ge\sqrt0$ and $2\ge\sqrt4$. Their midpoint is
$M=(2,1)$, and
\[
  \sqrt{2}\approx1.414>1 ,
\]
so $M$ violates $y\ge\sqrt{x}$ and lies \emph{outside}. A segment between
two points of the set leaves the set:
\[
  \boxed{\{(x,y):y\ge\sqrt{x}\}\ \text{is not convex.}}
\]

\wstep{Step 7. The principle behind it.}
$\sqrt{x}$ is \textbf{concave}. The set above the graph of a function --- its
epigraph --- is convex exactly when the function is \emph{convex}; the set
below --- the hypograph --- is convex exactly when the function is
\emph{concave}. So $\{y\ge\sqrt{x}\}$ fails while $\{y\le\sqrt{x}\}$
succeeds, and $\{y\ge x^{2}\}$ succeeds while $\{y\le x^{2}\}$ fails.
\wtrap{``Above'' is not a synonym for ``convex''. Match the direction of the
inequality to the curvature of the function: convex function with $\ge$, or
concave function with $\le$. Getting this backwards is the single most
common mistake in the chapter, and it recurs in Volume~7 where the feasible
set of a constrained problem must be convex for the first-order conditions
to be sufficient.}
\end{worked}

\section{Practice Problems}

Sketch the region described by each inequality in problems 1--12, and say
whether it is convex.

\begin{problems}
  \item $y\ge 2x-1$
  \item $x+y\le 3$
  \item $x^{2}+y^{2}\le 9$
  \item $x^{2}+y^{2}>4$
  \item $y\ge x^{2}$
  \item $y\le 4-x^{2}$
  \item $\dfrac{x^{2}}{9}+\dfrac{y^{2}}{4}\le1$
  \item $xy\ge 1$ on $x>0$
  \item $|x|\le 2$
  \item $|y|\le x$
  \item $(x-2)^{2}+(y+1)^{2}\le 9$
  \item $y>\dfrac1x$ on $x>0$
\end{problems}

In problems 13--20 sketch the region defined by the whole system.

\begin{problems}[resume]
  \item $x\ge0$, \;$y\ge0$, \;$2x+3y\le12$
  \item $x\ge0$, \;$y\ge0$, \;$x+y\le4$, \;$y\le2$
  \item $x^{2}+y^{2}\le16$, \;$y\ge0$
  \item $x^{2}+y^{2}\le16$, \;$y\ge0$, \;$x\le2$
  \item $x^{2}\le y$, \;$y\le x+2$
  \item $1\le x^{2}+y^{2}\le9$
  \item $x^{2}+y^{2}\le9$, \;$x\ge0$, \;$y\ge0$
  \item $y\ge x^{2}$, \;$y\le 4$, \;$x\ge0$
\end{problems}

Problems 21--28 refer to the eight shaded regions below. For each, give a
system of inequalities describing it.

\begin{gallery}
\begin{tikzpicture}
  \begin{axis}[geoaxis, xmin=-5.2, xmax=5.2, ymin=-5.2, ymax=5.2,
      xtick={-4,-2,2,4}, ytick={-4,-2,2,4}, title={(a)}]
    \addplot[draw=none, name path=up, domain=-4:2, samples=201]{sqrt(16-x^2)};
    \addplot[draw=none, name path=ax, domain=-4:2, samples=2]{0};
    \region{up}{ax}{}
    \addplot[domain=0:360, samples=181]({4*cos(x)},{4*sin(x)});
    \addplot[gray!60, dashed, thin, samples=2] coordinates {(2,-5)(2,5)};
  \end{axis}
\end{tikzpicture}\hfill
\begin{tikzpicture}
  \begin{axis}[geoaxis, xmin=-2.8, xmax=2.8, ymin=-0.9, ymax=5.1,
      xtick={-2,2}, ytick={2,4}, title={(b)}]
    \addplot[draw=none, name path=par, domain=-2:2, samples=201]{x^2};
    \addplot[draw=none, name path=top, domain=-2:2, samples=2]{4};
    \region{par}{top}{}
    \addplot[domain=-2.2:2.2, samples=201]{x^2};
    \addplot[domain=-2.6:2.6, samples=2]{4};
  \end{axis}
\end{tikzpicture}\hfill
\begin{tikzpicture}
  \begin{axis}[geoaxis, xmin=-3.9, xmax=3.9, ymin=-2.9, ymax=2.9,
      xtick={-2,2}, ytick={-2,2}, title={(c)}]
    \addplot[draw=none, name path=ell, domain=0:3, samples=201]{2*sqrt(1-x^2/9)};
    \addplot[draw=none, name path=ax, domain=0:3, samples=2]{0};
    \region{ell}{ax}{}
    \addplot[domain=0:360, samples=181]({3*cos(x)},{2*sin(x)});
  \end{axis}
\end{tikzpicture}
\end{gallery}

\begin{gallery}
\begin{tikzpicture}
  \begin{axis}[geoaxis, xmin=-2.4, xmax=3.4, ymin=-0.9, ymax=4.9,
      xtick={-2,2}, ytick={2,4}, title={(d)}]
    \addplot[draw=none, name path=par, domain=-1:2, samples=201]{x^2};
    \addplot[draw=none, name path=lin, domain=-1:2, samples=2]{x+2};
    \region{par}{lin}{}
    \addplot[domain=-2.1:2.2, samples=201]{x^2};
    \addplot[domain=-2.2:3.2, samples=2]{x+2};
  \end{axis}
\end{tikzpicture}\hfill
\begin{tikzpicture}
  \begin{axis}[geoaxis, xmin=-3.9, xmax=3.9, ymin=-3.9, ymax=3.9,
      xtick={-2,2}, ytick={-2,2}, title={(e)}]
    \addplot[draw=none, name path=o1, domain=-3:3, samples=201]{sqrt(9-x^2)};
    \addplot[draw=none, name path=o2, domain=-3:3, samples=201]{-sqrt(9-x^2)};
    \region{o1}{o2}{}
    \addplot[draw=none, name path=i1, domain=-1:1, samples=201]{sqrt(1-x^2)};
    \addplot[draw=none, name path=i2, domain=-1:1, samples=201]{-sqrt(1-x^2)};
    \addplot[white, forget plot, sharp plot] fill between[of=i1 and i2];
    \addplot[domain=0:360, samples=181]({3*cos(x)},{3*sin(x)});
    \addplot[domain=0:360, samples=181]({cos(x)},{sin(x)});
  \end{axis}
\end{tikzpicture}\hfill
\begin{tikzpicture}
  \begin{axis}[geoaxis, xmin=-1.4, xmax=7.4, ymin=-1.4, ymax=5.4,
      xtick={2,4,6}, ytick={2,4}, title={(f)}]
    \addplot[draw=none, name path=lin, domain=0:6, samples=2]{(12-2*x)/3};
    \addplot[draw=none, name path=ax, domain=0:6, samples=2]{0};
    \region{lin}{ax}{}
    \addplot[domain=-1.2:7.2, samples=2]{(12-2*x)/3};
  \end{axis}
\end{tikzpicture}
\end{gallery}

\begin{gallery}
\begin{tikzpicture}
  \begin{axis}[geoaxis, xmin=-1.2, xmax=5.4, ymin=-2.9, ymax=2.9,
      xtick={2,4}, ytick={-2,2}, title={(g)}]
    \addplot[draw=none, name path=up, domain=0:4, samples=201]{sqrt(x)};
    \addplot[draw=none, name path=lo, domain=0:4, samples=201]{-sqrt(x)};
    \region{up}{lo}{}
    \addplot[domain=0:4.6, samples=201]{sqrt(x)};
    \addplot[domain=0:4.6, samples=201]{-sqrt(x)};
    \addplot[gray!60, dashed, thin, samples=2] coordinates {(4,-2.8)(4,2.8)};
  \end{axis}
\end{tikzpicture}\hfill
\begin{tikzpicture}
  \begin{axis}[geoaxis, xmin=-3.4, xmax=3.4, ymin=-3.4, ymax=3.4,
      xtick={-2,2}, ytick={-2,2}, title={(h)}]
    \addplot[draw=none, name path=up, domain=-2:2, samples=101]{2-abs(x)};
    \addplot[draw=none, name path=lo, domain=-2:2, samples=101]{abs(x)-2};
    \region{up}{lo}{}
    \addplot[domain=-2:2, samples=101]{2-abs(x)};
    \addplot[domain=-2:2, samples=101]{abs(x)-2};
  \end{axis}
\end{tikzpicture}
\gcap{Regions (a)--(h) for problems 21--28. A dashed boundary means the
inequality there is strict.}
\end{gallery}

\begin{problems}[resume]
  \item Region (a).
  \item Region (b).
  \item Region (c).
  \item Region (d).
  \item Region (e).
  \item Region (f).
  \item Region (g).
  \item Region (h).
\end{problems}

\begin{problems}[resume]
  \item Which of the eight regions (a)--(h) are convex?
  \item Sketch $|x|+|y|\le 3$ and give its four vertices.
  \item Sketch $y\ge|x|$ and describe it in words.
  \item Sketch $|y|\le 4-x^{2}$ and give the $x$-values for which the
        region is non-empty.
  \item Write inequalities for the closed square with vertices
        $(\pm2,\pm2)$, using absolute values and then without them.
  \item A test point satisfies every inequality of a system but lies
        outside the intended region. What has gone wrong?
  \item Explain why the outside of a disc is never convex.
  \item Give a system of two inequalities whose solution set is empty, both
        of them satisfiable on their own.
  \item Show that the intersection of any two convex regions is convex.
  \item Is the union of the two half-planes $y\ge1$ and $y\le-1$ convex?
        Justify.
  \item Write inequalities for the set of points strictly inside the circle
        of radius $2$ about the origin and strictly above the line $y=x$.
  \item For which values of $c$ is the region $\{x^{2}+y^{2}\le c\}$ empty,
        a single point, and a disc?
\end{problems}

\begin{problems}[resume]
  \item \textbf{Budget set.} A consumer has income $60$ and faces prices
        $p_{x}=3$, $p_{y}=4$. Write the budget set as a system, sketch it,
        and give its three vertices.
  \item For the same consumer, write the set of bundles that are
        \emph{affordable and} give at least utility $12$ under
        $U(x,y)=xy$. Sketch it and say whether it is convex.
  \item \textbf{Production.} A firm can use at most $40$ hours of labour
        and $30$ units of capital; product A needs $2$ hours and $1$ unit,
        product B needs $1$ hour and $2$ units. Write the feasible set and
        sketch it.
  \item For the feasible set of the previous problem, give all its
        vertices, and say why a linear objective attains its maximum at one
        of them.
  \item \textbf{Tax wedge.} The set $\{(q,p): p\le 100-2q,\ p\ge 20+q,\
        q\ge0\}$ describes the quantities at which demand price exceeds
        supply price. Sketch it and find the largest $q$ in it.
\end{problems}

\section{Solutions}

\begin{solutions}
  \item The closed half-plane on or above the line $y=2x-1$. Testing
        $(0,0)$: $0\ge-1$ true, so the origin is in it. Convex.
  \item The closed half-plane on or below $x+y=3$ (through $(3,0)$ and
        $(0,3)$); the origin satisfies it. Convex.
  \item The closed disc of radius $3$ about the origin. Convex.
  \item Everything strictly outside the circle of radius $2$, boundary
        dashed and excluded. \textbf{Not} convex.
  \item The region on or above the parabola $y=x^{2}$ --- inside its
        opening. Convex.
  \item The region on or below the downward parabola $y=4-x^{2}$, which
        includes everything far below it. Since $4-x^{2}$ is
        \emph{concave}, the region below it is \textbf{convex}.
  \item The closed elliptical disc with semi-axes $3$ and $2$. Convex.
  \item The region on or beyond the hyperbola branch $y=1/x$ in the first
        quadrant. Convex (this is the quasiconcavity of $xy$).
  \item The closed vertical strip $-2\le x\le2$. Convex, being the
        intersection of two half-planes.
  \item Requires $x\ge0$ and $-x\le y\le x$: the wedge opening to the right
        between the lines $y=\pm x$. Convex.
  \item The closed disc of radius $3$ centred at $(2,-1)$. Convex.
  \item The region strictly above the hyperbola branch $y=1/x$ for $x>0$,
        boundary dashed. Convex.
  \item The triangle with vertices $(0,0)$, $(6,0)$, $(0,4)$ --- the
        intersection of three half-planes.
  \item The quadrilateral with vertices $(0,0)$, $(4,0)$, $(2,2)$,
        $(0,2)$: the triangle $x+y\le4$ in the first quadrant, with its top
        cut off by $y\le2$.
  \item The upper half-disc of radius $4$: the semicircular region above
        the $x$-axis.
  \item The same half-disc with everything right of $x=2$ removed --- the
        region shown in figure (a).
  \item The lens between the parabola and the line, from $x=-1$ to $x=2$
        (the boundaries meet where $x^{2}=x+2$, i.e. $x=-1,2$). Convex.
  \item The annulus between the circles of radius $1$ and $3$. Not convex.
  \item The quarter-disc of radius $3$ in the first quadrant. Convex.
  \item The region above the parabola, below $y=4$, right of the $y$-axis:
        the right half of figure (b).
  \item $x^{2}+y^{2}\le16$, \;$y\ge0$, \;$x<2$. The circle and the
        $x$-axis are solid, so those two are non-strict; the vertical
        boundary is dashed, so that one is strict.
  \item $y\ge x^{2}$ and $y\le4$. (Equivalently
        $x^{2}\le y\le 4$, which also forces $-2\le x\le2$ without needing
        a separate condition.)
  \item $\dfrac{x^{2}}{9}+\dfrac{y^{2}}{4}\le1$, \;$x\ge0$, \;$y\ge0$: the
        quarter of the elliptical disc in the first quadrant.
  \item $y\ge x^{2}$ and $y\le x+2$.
  \item $1\le x^{2}+y^{2}\le9$, that is $x^{2}+y^{2}\ge1$ \emph{and}
        $x^{2}+y^{2}\le9$.
  \item $2x+3y\le12$, \;$x\ge0$, \;$y\ge0$.
  \item $y^{2}\le x$ and $x<4$; equivalently
        $-\sqrt{x}\le y\le\sqrt{x}$ with $x<4$. The parabola opens to the
        right and is solid; the line $x=4$ is dashed.
  \item $|x|+|y|\le2$. Written without absolute values it is the four
        conditions $x+y\le2$, $x-y\le2$, $-x+y\le2$, $-x-y\le2$.
  \item All except (e). Each of (a), (b), (c), (d), (f), (g), (h) is an
        intersection of convex sets: discs and ellipse interiors,
        half-planes, the region above the convex $x^{2}$, the region
        \emph{right} of the parabola $x=y^{2}$ (the same rule with the axes
        exchanged, since $y\mapsto y^{2}$ is convex), and the diamond
        $|x|+|y|\le2$, which is an intersection of four half-planes.
        Only (e), the annulus, fails: it is the disc of radius $3$ with a
        hole, and the segment from $(-2,0)$ to $(2,0)$ runs straight
        through the hole.
  \item A diamond with vertices $(3,0)$, $(0,3)$, $(-3,0)$, $(0,-3)$: the
        square of side $3\sqrt2$ rotated by $45^{\circ}$.
  \item The wedge above both lines $y=x$ and $y=-x$ --- everything at or
        above the V of $y=|x|$, opening upwards with its vertex at the
        origin. Convex.
  \item Squeezed between $y=4-x^{2}$ and $y=-(4-x^{2})$, so it exists only
        where $4-x^{2}\ge0$, i.e. for $-2\le x\le2$. It is the
        lens-shaped region symmetric about both axes.
  \item With absolute values: $\max(|x|,|y|)\le2$, i.e. $|x|\le2$ and
        $|y|\le2$. Without: $-2\le x\le2$ and $-2\le y\le2$, four
        half-planes.
  \item The system describes a \emph{larger} set than intended --- a
        condition is missing. Every boundary of the intended region must
        contribute one inequality; a point outside that satisfies them all
        proves one boundary has been forgotten.
  \item Take two points diametrically opposite each other just outside the
        circle: the segment between them passes through the centre of the
        disc, which is not in the set. So the segment leaves the set and
        convexity fails.
  \item For instance $y\ge x+1$ and $y\le x-1$. Each is a satisfiable
        half-plane, but together they demand $x+1\le y\le x-1$, impossible
        since $x+1>x-1$. The two half-planes are on opposite sides of
        parallel lines.
  \item Let $C_{1},C_{2}$ be convex and take $P,Q\in C_{1}\cap C_{2}$.
        Since $P,Q\in C_{1}$ and $C_{1}$ is convex, the whole segment
        $[P,Q]$ lies in $C_{1}$; the same argument puts it in $C_{2}$.
        Hence $[P,Q]\subseteq C_{1}\cap C_{2}$, which is the definition of
        convexity. (The argument works for any number of sets.)
  \item No. The points $(0,1)$ and $(0,-1)$ both lie in the union, but the
        midpoint $(0,0)$ lies in neither half-plane. A union of disjoint
        convex sets is essentially never convex.
  \item $x^{2}+y^{2}<4$ \;and\; $y>x$. Both boundaries are dashed: the open
        half-disc lying above the diagonal. (It is convex, being an
        intersection of two convex sets.)
  \item Empty for $c<0$; the single point $(0,0)$ for $c=0$; the closed disc
        of radius $\sqrt{c}$ for $c>0$. The three cases are exactly the
        three things a level \emph{set} of $x^{2}+y^{2}$ can do, and they
        are why a sublevel set can be empty without the inequality being
        malformed.
  \item $3x+4y\le60$, \;$x\ge0$, \;$y\ge0$: the triangle with vertices
        $(0,0)$, $(20,0)$ and $(0,15)$.
  \item $3x+4y\le60$, \;$xy\ge12$, \;$x,y>0$: the part of the budget
        triangle lying beyond the indifference curve $y=12/x$. It is the
        intersection of a convex set (the triangle) with another convex set
        (the better set of a quasiconcave utility), hence \textbf{convex}.
  \item With $a$ units of A and $b$ of B: $2a+b\le40$ (labour),
        $a+2b\le30$ (capital), $a\ge0$, $b\ge0$. Four half-planes, so a
        convex polygon.
  \item Vertices: $(0,0)$, $(20,0)$, $(0,15)$, and the intersection of
        $2a+b=40$ with $a+2b=30$. Solving: $b=40-2a$ and
        $a+80-4a=30$, so $a=\tfrac{50}{3}$ and $b=\tfrac{20}{3}$. A linear
        objective has parallel straight level lines; pushing them across a
        convex polygon, the last point of contact is a vertex (or a whole
        edge, if the level lines happen to be parallel to it). Hence
        checking the vertices suffices --- the fundamental theorem of
        linear programming.
  \item The two lines meet where $100-2q=20+q$, i.e. $q=\tfrac{80}{3}$ with
        $p=\tfrac{140}{3}$. The region is the triangle with vertices
        $(0,20)$, $(0,100)$ and
        $\bigl(\tfrac{80}{3},\tfrac{140}{3}\bigr)$, so the largest feasible
        quantity is $q=\tfrac{80}{3}\approx26.7$ --- beyond it no price
        satisfies both buyer and seller, which is exactly the competitive
        equilibrium quantity.
\end{solutions}

\section{Additional Problems}

\begin{problems}
  \item Sketch $3x-2y>6$ and say which side of the line contains the
        origin.
  \item Sketch $y\le -x+5$ together with $x\ge1$.
  \item Sketch $(x+1)^{2}+(y-2)^{2}<4$ and give centre and radius.
  \item Sketch $x^{2}+y^{2}\ge 1$ and $x^{2}+y^{2}\le 4$ together.
  \item Sketch $y\ge x^{2}-4$.
  \item Sketch $y\le 2^{x}$.
  \item Sketch $y\ge\ln x$ and state the domain restriction it forces.
  \item Sketch $xy\le 4$ on the first quadrant.
  \item Sketch $|x-1|\le 2$.
  \item Sketch $|x|+|y|\ge 1$ and decide convexity.
  \item Sketch $y>x^{2}$ and $y<x+6$, and find the $x$-range of the
        region.
  \item Sketch $x^{2}+4y^{2}\le 16$ and $y\ge 0$.
  \item Write inequalities for the closed triangle with vertices $(0,0)$,
        $(4,0)$, $(0,3)$.
  \item Write inequalities for the closed disc of radius $2$ centred at
        $(1,1)$, intersected with the first quadrant.
  \item Write inequalities for the region above $y=x^{2}$ and below
        $y=2x+3$, and give the $x$-range.
  \item Write inequalities for the horizontal strip between $y=-1$ and
        $y=3$, once with an absolute value and once without.
  \item Write inequalities for everything outside the ellipse
        $x^{2}/4+y^{2}=1$ but inside the circle $x^{2}+y^{2}=9$.
  \item Is the region $\{y\ge x^{2}\}\cap\{y\le 4\}$ bounded? Give its
        extreme $x$- and $y$-values.
  \item Give a system of three inequalities whose solution set is a single
        point.
  \item Give a system of three inequalities whose solution set is an
        unbounded convex region with exactly two vertices.
  \item Show that $\{(x,y): y\ge \mathrm{e}^{x}\}$ is convex.
  \item Show that $\{(x,y): y\ge \sqrt{x},\ x\ge0\}$ is \emph{not}
        convex, and say which nearby set \emph{is}.
  \item \textbf{Two constraints, two goods.} Sketch
        $x+y\le 10$, \;$2x+y\le 16$, \;$x,y\ge0$, and give all vertices.
  \item \textbf{Diet problem.} A feed must supply at least $8$ units of
        protein and $12$ of energy. Grain gives $2$ protein and $2$ energy
        per kg; soy gives $1$ protein and $4$ energy per kg. Write the
        feasible set, sketch it, and say why it is unbounded.
  \item \textbf{Minimum-wage wedge.} Sketch
        $\{(L,w): w\ge 10,\ w\le 30-L,\ L\ge0\}$ and interpret the region
        and its largest $L$.
\end{problems}

\section{Additional Solutions}

\begin{solutions}
  \item The open half-plane strictly below the line $y=\tfrac32x-3$.
        Testing $(0,0)$: $0>6$ is false, so the origin is \emph{not} in the
        region; it lies on the other side.
  \item The intersection of the half-plane on or below $y=5-x$ with the
        half-plane $x\ge1$: an unbounded wedge opening down-right, with
        corner $(1,4)$.
  \item The \emph{open} disc of radius $2$ centred at $(-1,2)$; the
        boundary circle is dashed and excluded.
  \item The closed annulus $1\le x^{2}+y^{2}\le4$ between the two circles.
        Not convex.
  \item The region on or above the parabola $y=x^{2}-4$, whose vertex is
        $(0,-4)$ and which crosses the axis at $x=\pm2$. Convex.
  \item Everything on or below the increasing curve $y=2^{x}$, which
        passes through $(0,1)$ and hugs the $x$-axis to the left. Since
        $2^{x}$ is \emph{convex}, the region below it is \textbf{not}
        convex: $(-3,\tfrac18)$ and $(3,8)$ both satisfy the inequality,
        but their midpoint $(0,4.0625)$ does not, since $2^{0}=1<4.0625$.
  \item The region on or above $y=\ln x$; the logarithm's domain confines
        it to the open right half-plane $x>0$. Since $\ln x$ is
        \emph{concave}, the region above it is \textbf{not} convex: the
        boundary points $(0.1,\ln 0.1)$ and $(10,\ln 10)$ have midpoint
        $(5.05,0)$, and $0<\ln 5.05\approx1.62$, so the midpoint is
        outside.
  \item The part of the first quadrant on or below the hyperbola $y=4/x$
        --- the sublevel set $\{xy\le4\}$. \textbf{Not} convex: $(0.5,8)$
        and $(8,0.5)$ both lie on the boundary, but their midpoint
        $(4.25,4.25)$ has $xy\approx18.1>4$. (Equivalently: $4/x$ is convex
        on $x>0$, so the region below it is not.)
  \item $-2\le x-1\le2$, i.e. the closed vertical strip $-1\le x\le3$.
        Convex.
  \item Everything on or outside the diamond $|x|+|y|=1$. Not convex ---
        the segment from $(1,0)$ to $(-1,0)$ passes through the origin,
        where $|x|+|y|=0<1$.
  \item Strictly between the parabola and the line; they meet where
        $x^{2}=x+6$, i.e. $x=-2$ and $x=3$, so the region spans
        $-2<x<3$. Both boundaries dashed.
  \item $\dfrac{x^{2}}{16}+\dfrac{y^{2}}{4}\le1$ with $y\ge0$: the upper
        half of an elliptical disc with semi-axes $4$ and $2$.
  \item $x\ge0$, \;$y\ge0$, \;$3x+4y\le12$ (the line through $(4,0)$ and
        $(0,3)$).
  \item $(x-1)^{2}+(y-1)^{2}\le4$, \;$x\ge0$, \;$y\ge0$.
  \item $y\ge x^{2}$ and $y\le 2x+3$; the boundaries meet where
        $x^{2}=2x+3$, i.e. $x=-1$ and $x=3$, so $-1\le x\le3$.
  \item With an absolute value: $|y-1|\le2$. Without: $y\ge-1$ and
        $y\le3$.
  \item $\dfrac{x^{2}}{4}+y^{2}\ge1$ \;and\; $x^{2}+y^{2}\le9$: the disc of
        radius $3$ with the elliptical disc punched out. Not convex.
  \item Yes, bounded: $y$ ranges over $[0,4]$ and $x$ over $[-2,2]$, since
        $y\ge x^{2}$ together with $y\le4$ forces $x^{2}\le4$. The extreme
        points are $(\pm2,4)$ and $(0,0)$.
  \item For example $x\ge0$, \;$y\ge0$, \;$x+y\le0$: the first two force
        $x,y\ge0$ and the third then forces $x=y=0$, so the solution set is
        exactly $\{(0,0)\}$.
  \item $x\ge0$, \;$y\ge0$, \;$x+y\ge2$. The boundary consists of the
        ray $\{x=0,\ y\ge2\}$, the segment from $(0,2)$ to $(2,0)$, and the
        ray $\{y=0,\ x\ge2\}$, so there are exactly two corners, $(0,2)$
        and $(2,0)$. The region is unbounded (it contains every large
        point of the quadrant) and convex, being an intersection of three
        half-planes.
  \item $\mathrm{e}^{x}$ is a convex function, and the set lying \emph{on
        or above} the graph of a convex function --- its epigraph --- is a
        convex set: if $y_{1}\ge\mathrm{e}^{x_{1}}$ and
        $y_{2}\ge\mathrm{e}^{x_{2}}$ then for $t\in[0,1]$
        \[
          ty_{1}+(1-t)y_{2}\ \ge\ t\mathrm{e}^{x_{1}}+(1-t)\mathrm{e}^{x_{2}}
          \ \ge\ \mathrm{e}^{tx_{1}+(1-t)x_{2}},
        \]
        the last step being exactly convexity of $\mathrm{e}^{x}$.
  \item $\sqrt{x}$ is \emph{concave}, so the region \emph{above} it is
        not convex. Concretely, $(0,0)$ and $(4,2)$ both lie on the
        boundary and so belong to the set, but their midpoint $(2,1)$ does
        not, since $\sqrt2\approx1.414>1$. The set that \emph{is} convex is
        the one below the curve, $\{y\le\sqrt{x},\ x\ge0\}$ --- the
        hypograph of a concave function.
  \item The lines meet where $x+y=10$ and $2x+y=16$, i.e. $x=6$, $y=4$.
        Vertices: $(0,0)$, $(8,0)$, $(6,4)$, $(0,10)$.
  \item With $g$ kg of grain and $s$ kg of soy: $2g+s\ge8$,
        $2g+4s\ge12$, $g\ge0$, $s\ge0$. It is unbounded because the
        constraints are all ``at least'' --- buying arbitrarily much feed is
        always still feasible, so the region extends indefinitely up and to
        the right. A cost-minimising objective is nevertheless attained,
        at one of the two vertices on the lower-left boundary.
  \item The region is the triangle bounded below by $w=10$, above by
        $w=30-L$, and on the left by $L=0$; the two sloping boundaries
        meet at $L=20$, $w=10$. It shows the employment levels compatible
        with a wage floor of $10$ when labour demand is $w=30-L$: the
        largest feasible $L$ is $20$, the employment the floor permits.
        Any $L$ beyond that would require a wage below the minimum.
\end{solutions}

\chapter{Trigonometric Functions}

Trigonometry plays a small part in economics and a large part in the
prerequisites, for two reasons. The graphs of $\sin$, $\cos$ and $\tan$
belong to the catalogue of shapes one is expected to recognise, and they are
the standard test functions of calculus: many derivative and limit
exercises use them precisely because they oscillate and so expose careless
reasoning. Seasonality and business cycles supply the
occasional genuine application.

The chapter therefore covers what is actually needed and stops. The unit
circle and the handful of exact values; the three graphs with their periods,
amplitudes and asymptotes; the transformation form
$a\sin\bigl(b(x-c)\bigr)+d$, which is the machinery of \chref{functions}{3} applied to a
periodic function; the four identities that carry all the simplification
work; and how to write down \emph{every} solution of a trigonometric
equation rather than just the one a calculator returns.

\section{Core Rules}

\begin{rulebox}{Radians, and why}
An angle in \textbf{radians} is the arc length it cuts on the unit circle,
so a full turn is $2\pi$:
\[
  180^{\circ}=\pi,\qquad
  x_{\text{rad}}=x_{\deg}\cdot\frac{\pi}{180},\qquad
  x_{\deg}=x_{\text{rad}}\cdot\frac{180}{\pi} .
\]
Calculus uses radians exclusively, and for one reason: only in radians is
$\dfrac{\mathrm{d}}{\mathrm{d}x}\sin x=\cos x$. In degrees an unwanted
factor $\pi/180$ appears in every derivative.
\end{rulebox}

\begin{rulebox}{The unit circle definition}
For the point on the unit circle reached by turning through the angle $x$
from the positive $x$-axis, $\cos x$ is its first coordinate and $\sin x$
its second. Hence, immediately:
\[
  \sin^{2}x+\cos^{2}x=1,\qquad
  -1\le\sin x\le1,\qquad -1\le\cos x\le1,
\]
the first being Pythagoras. Also $\tan x=\dfrac{\sin x}{\cos x}$, the slope
of the ray, undefined where $\cos x=0$.

Signs by quadrant --- $\sin$ is the height, $\cos$ the width:
\[
  \begin{array}{c|cccc}
     & \text{I} & \text{II} & \text{III} & \text{IV}\\\hline
    \sin & + & + & - & -\\
    \cos & + & - & - & +\\
    \tan & + & - & + & -
  \end{array}
\]
\end{rulebox}

\begin{rulebox}{The values to know by heart}
\[
  \begin{array}{c|ccccc}
    x & 0 & \pi/6 & \pi/4 & \pi/3 & \pi/2\\\hline
    \sin x & 0 & \tfrac12 & \tfrac{\sqrt2}{2} & \tfrac{\sqrt3}{2} & 1\\[0.3ex]
    \cos x & 1 & \tfrac{\sqrt3}{2} & \tfrac{\sqrt2}{2} & \tfrac12 & 0\\[0.3ex]
    \tan x & 0 & \tfrac{1}{\sqrt3} & 1 & \sqrt3 & \text{---}
  \end{array}
\]
Everything else follows from three facts: $\sin$ is \textbf{odd}
($\sin(-x)=-\sin x$), $\cos$ is \textbf{even} ($\cos(-x)=\cos x$), and both
have period $2\pi$ while $\tan$ has period $\pi$. So
$\sin\tfrac{5\pi}{6}=\sin\tfrac{\pi}{6}=\tfrac12$ (second quadrant, $\sin$
positive) and $\cos\tfrac{4\pi}{3}=-\cos\tfrac{\pi}{3}=-\tfrac12$ (third
quadrant, $\cos$ negative).
\end{rulebox}

\begin{gallery}
\begin{tikzpicture}
  \begin{axis}[fnaxis, xmin=-0.6, xmax=7.0, ymin=-1.6, ymax=1.6,
      xtick={3.1416, 6.2832}, xticklabels={$\pi$, $2\pi$},
      ytick={-1,1}, title={$y=\sin x$}]
    \addplot[domain=-0.4:6.9, samples=201]{sin(deg(x))};
  \end{axis}
\end{tikzpicture}\hfill
\begin{tikzpicture}
  \begin{axis}[fnaxis, xmin=-0.6, xmax=7.0, ymin=-1.6, ymax=1.6,
      xtick={3.1416, 6.2832}, xticklabels={$\pi$, $2\pi$},
      ytick={-1,1}, title={$y=\cos x$}]
    \addplot[domain=-0.4:6.9, samples=201]{cos(deg(x))};
  \end{axis}
\end{tikzpicture}\hfill
\begin{tikzpicture}
  \begin{axis}[fnaxis, xmin=-3.4, xmax=3.4, ymin=-3.4, ymax=3.4,
      xtick={-1.5708, 1.5708}, xticklabels={$-\frac{\pi}{2}$,$\frac{\pi}{2}$},
      ytick={-2,2}, title={$y=\tan x$}]
    \addplot[domain=-1.32:1.32, samples=101]{tan(deg(x))};
    \addplot[domain=1.82:3.3, samples=101]{tan(deg(x))};
    \addplot[domain=-3.3:-1.82, samples=101]{tan(deg(x))};
    \addplot[gray!55, dashed, thin, samples=2, forget plot]
      coordinates {(1.5708,-3.3)(1.5708,3.3)};
    \addplot[gray!55, dashed, thin, samples=2, forget plot]
      coordinates {(-1.5708,-3.3)(-1.5708,3.3)};
  \end{axis}
\end{tikzpicture}
\gcap{$\sin$ and $\cos$ are the same wave shifted by $\pi/2$, with period
$2\pi$ and range $[-1,1]$. $\tan$ has period $\pi$, range $\R$, and a
vertical asymptote at every odd multiple of $\pi/2$.}
\end{gallery}

\begin{rulebox}{Reading $a\sin\bigl(b(x-c)\bigr)+d$}
The four constants are exactly the four moves of \chref{functions}{3}, and each has a
name here:
\begin{itemize}[leftmargin=1.6em,itemsep=0.3ex]
  \item $|a|$ is the \textbf{amplitude}: the wave runs from $d-|a|$ to
        $d+|a|$, so the range is $[d-|a|,\,d+|a|]$. Negative $a$ flips it.
  \item $b$ sets the \textbf{period}, $T=\dfrac{2\pi}{|b|}$ --- larger $b$
        means faster oscillation. (For $\tan$, $T=\pi/|b|$.)
  \item $c$ is the \textbf{phase shift}: the whole wave moves $c$ to the
        right. Factor before reading: $\sin(2x-\pi)=\sin\bigl(2(x-\pi/2)\bigr)$
        is a shift of $\pi/2$, not of $\pi$.
  \item $d$ is the \textbf{midline}, the horizontal line the wave
        oscillates about.
\end{itemize}
So $y=3\sin\bigl(2(x-\tfrac{\pi}{4})\bigr)+1$ has amplitude $3$, period
$\pi$, midline $y=1$, range $[-2,4]$, shifted $\pi/4$ right.
\end{rulebox}

\begin{rulebox}{The identities that do the work}
\[
  \sin^{2}x+\cos^{2}x=1,\qquad
  \tan x=\frac{\sin x}{\cos x},\qquad
  1+\tan^{2}x=\frac{1}{\cos^{2}x} .
\]
\[
  \sin(x\pm y)=\sin x\cos y\pm\cos x\sin y,\qquad
  \cos(x\pm y)=\cos x\cos y\mp\sin x\sin y .
\]
Setting $y=x$ in the addition formulas gives the \textbf{double-angle}
forms, which are the ones that appear in practice:
\[
  \sin 2x=2\sin x\cos x,\qquad
  \cos 2x=\cos^{2}x-\sin^{2}x=1-2\sin^{2}x=2\cos^{2}x-1 .
\]
Almost every simplification is one of: replace $1-\cos^{2}x$ by
$\sin^{2}x$, split a $\tan$ into $\sin/\cos$, or collapse a double angle.
\end{rulebox}

\begin{rulebox}{Solving trigonometric equations: give \emph{all} solutions}
A calculator returns one angle; periodicity supplies infinitely many.
\begin{itemize}[leftmargin=1.6em,itemsep=0.35ex]
  \item $\sin x=s$ with $|s|\le1$: two families,
        $x=x_{0}+2k\pi$ and $x=\pi-x_{0}+2k\pi$, where
        $x_{0}=\arcsin s$ and $k\in\mathbb{Z}$.
  \item $\cos x=s$: $x=\pm x_{0}+2k\pi$ with $x_{0}=\arccos s$.
  \item $\tan x=s$: one family, $x=\arctan s+k\pi$, since the period is
        $\pi$.
  \item If the argument is $bx$ rather than $x$, solve for $bx$ first and
        divide \emph{everything}, the $2k\pi$ included: $\sin 2x=1$ gives
        $2x=\tfrac{\pi}{2}+2k\pi$, hence
        $x=\tfrac{\pi}{4}+k\pi$.
\end{itemize}
\end{rulebox}

\begin{rulebox}{Inverse trigonometric functions}
None of the three is injective, so each is inverted only on a restricted
domain:
\[
  \arcsin:[-1,1]\to\Bigl[-\tfrac{\pi}{2},\tfrac{\pi}{2}\Bigr],\qquad
  \arccos:[-1,1]\to[0,\pi],\qquad
  \arctan:\R\to\Bigl(-\tfrac{\pi}{2},\tfrac{\pi}{2}\Bigr).
\]
So $\arcsin(\sin x)=x$ only for $x$ in that range: $\arcsin(\sin\pi)=0$,
not $\pi$. And $\arctan$ is bounded, with horizontal asymptotes
$y=\pm\pi/2$ --- the reason it appears as a saturating S-shaped curve.
\end{rulebox}

\section{Worked Example}

\begin{worked}{Solve $2\sin^{2}x-\sin x=0$ on $[0,2\pi]$ and in general, and read off the shape of $y=3\sin\!\bigl(2(x-\tfrac{\pi}{4})\bigr)+1$}
\wstep{Step 1. Factor. Do not divide.}
\[
  2\sin^{2}x-\sin x=0
  \quad\Longrightarrow\quad
  \sin x\,(2\sin x-1)=0 .
\]
\wtrap{Dividing both sides by $\sin x$ gives $2\sin x-1=0$ and loses every
solution with $\sin x=0$ --- three of the five on $[0,2\pi]$. You may divide
by a quantity only when you know it is non-zero, and $\sin x$ is exactly the
quantity in question. Factor and use the zero-product property instead.}

\wstep{Step 2. Solve each factor on the given interval.}
\emph{First factor:} $\sin x=0$ on $[0,2\pi]$ gives $x=0,\ \pi,\ 2\pi$.

\emph{Second factor:} $\sin x=\tfrac12$. The reference angle is
$\pi/6$, and the sine is positive in the first and second quadrants, so
\[
  x=\frac{\pi}{6}\qquad\text{and}\qquad x=\pi-\frac{\pi}{6}=\frac{5\pi}{6}.
\]
\wtrap{$\arcsin(\tfrac12)=\pi/6$ is one solution, not the solution. The
inverse sine returns a single value from $[-\pi/2,\pi/2]$; the second
solution in each period must be supplied by hand from the symmetry
$\sin(\pi-x)=\sin x$. For cosine the partner is $-x$ (or $2\pi-x$), for
tangent it is $x+\pi$.}

\wstep{Step 3. Collect the five solutions.}
\[
  \boxed{\;x\in\Bigl\{0,\ \tfrac{\pi}{6},\ \tfrac{5\pi}{6},\ \pi,\ 2\pi\Bigr\}\;}
\]
Check the smallest and one middle value: at $x=0$, $0-0=0$ \checkmark; at
$x=\pi/6$, $2\cdot\tfrac14-\tfrac12=0$ \checkmark.

\wstep{Step 4. Give the general solution, since ``all solutions'' usually means all of $\R$.}
\[
  x=k\pi,
  \qquad
  x=\frac{\pi}{6}+2k\pi,
  \qquad
  x=\frac{5\pi}{6}+2k\pi,
  \qquad k\in\mathbb{Z} .
\]
Note the two different periods: the zeros of $\sin$ repeat every $\pi$
(because $\sin(x+\pi)=-\sin x$, which is also zero), while the solutions of
$\sin x=\tfrac12$ repeat only every $2\pi$.
\wtrap{Adding $2k\pi$ to everything is safe but incomplete for the first
family; writing $k\pi$ for the second family is simply wrong. Work out the
period of each family separately from the equation it solves.}

\wstep{Step 5. Now the sinusoid: read the four parameters off $a\sin\bigl(b(x-c)\bigr)+d$.}
For $y=3\sin\bigl(2(x-\tfrac{\pi}{4})\bigr)+1$:
\begin{center}\footnotesize
\begin{tabular}{llll}
  \hline
  parameter & value & name & effect\\
  \hline
  $a$ & $3$ & amplitude & oscillates $3$ above and below the midline\\
  $b$ & $2$ & angular frequency & period $=2\pi/b=\pi$\\
  $c$ & $\pi/4$ & phase shift & shifted \emph{right} by $\pi/4$\\
  $d$ & $1$ & midline & centred on $y=1$, so range $[-2,4]$\\
  \hline
\end{tabular}
\end{center}

\wstep{Step 6. Sketch it from those four numbers.}
\begin{center}
\begin{tikzpicture}
\begin{axis}[fnaxis, width=8.2cm, height=4.4cm,
             domain=-0.4:3.6, samples=201,
             xmin=-0.5, xmax=3.7, ymin=-2.8, ymax=4.8,
             xtick={0.7854, 1.5708, 2.3562, 3.1416},
             xticklabels={$\frac{\pi}{4}$,$\frac{\pi}{2}$,$\frac{3\pi}{4}$,$\pi$},
             ytick={-2,1,4}]
  \addplot[accent!35, thick, dashed, samples=2, forget plot] {1};
  \addplot {3*sin(deg(2*(x-pi/4)))+1};
\end{axis}
\end{tikzpicture}
\end{center}
The curve crosses its midline $y=1$ going upward at $x=\pi/4$ --- that is
what the phase shift means --- peaks at $4$ a quarter-period later
($x=\pi/4+\pi/4=\pi/2$), and repeats after one full period $\pi$.
\wtrap{The period is $2\pi/b$, not $2\pi b$: a \emph{larger} $b$ gives a
\emph{shorter} period. And the phase shift is only readable once the inside
is factored as $b(x-c)$; from $\sin(2x-\tfrac{\pi}{2})$ the shift is
$\pi/4$, not $\pi/2$. It is the same trap as in \chref{functions}{3} on
transformations.}
\end{worked}

\section{Practice Problems}

Convert in problems 1--6.

\begin{problems}
  \item $60^{\circ}$ into radians.
  \item $270^{\circ}$ into radians.
  \item $\tfrac{3\pi}{4}$ into degrees.
  \item $\tfrac{5\pi}{3}$ into degrees.
  \item $1$ radian into degrees, to one decimal.
  \item Why does calculus insist on radians?
\end{problems}

Give the exact value in problems 7--16.

\begin{problems}[resume]
  \item $\sin\tfrac{\pi}{6}$
  \item $\cos\tfrac{\pi}{4}$
  \item $\tan\tfrac{\pi}{3}$
  \item $\sin\tfrac{5\pi}{6}$
  \item $\cos\tfrac{2\pi}{3}$
  \item $\tan\pi$
  \item $\sin\bigl(-\tfrac{\pi}{3}\bigr)$
  \item $\cos\tfrac{7\pi}{2}$
  \item $\sin\tfrac{9\pi}{4}$
  \item $\tan\tfrac{3\pi}{4}$
\end{problems}

In problems 17--24 give period, amplitude (where defined), range and any
asymptotes.

\begin{problems}[resume]
  \item $y=\sin x$
  \item $y=\cos x$
  \item $y=\tan x$
  \item $y=3\sin x$
  \item $y=\sin 2x$
  \item $y=\cos\tfrac{x}{2}$
  \item $y=2\sin x+1$
  \item $y=\tan 3x$
\end{problems}

\begin{problems}[resume]
  \item Describe $y=\sin\bigl(x-\tfrac{\pi}{2}\bigr)$ as a shift, and name
        the simpler function it equals.
  \item Give amplitude, period, midline and range of
        $y=3\sin\bigl(2(x-\tfrac{\pi}{4})\bigr)+1$.
  \item Give amplitude, period and phase shift of
        $y=\sin(2x-\pi)$ --- factor first.
  \item Give amplitude, period and range of $y=-2\cos 3x$.
  \item Write a sine function with amplitude $4$, period $\pi$ and midline
        $y=-1$.
  \item Write a cosine function with amplitude $2$, period $6\pi$ and no
        vertical shift.
  \item Monthly sales oscillate about $500$ units with a swing of $\pm80$
        and a twelve-month cycle, peaking in month $3$. Write $S(t)$.
  \item Why can no function of the form $a\sin(bx)+d$ ever be injective on
        $\R$?
\end{problems}

Simplify or verify in problems 33--38.

\begin{problems}[resume]
  \item $\dfrac{1-\cos^{2}x}{\sin x}$
  \item $\tan x\cdot\cos x$
  \item $\sin^{2}x-1$
  \item $\dfrac{\sin 2x}{2\cos x}$
  \item Verify $1+\tan^{2}x=\dfrac{1}{\cos^{2}x}$.
  \item Verify $\cos 2x=1-2\sin^{2}x$ from the addition formula.
\end{problems}

Give \emph{all} real solutions in problems 39--45.

\begin{problems}[resume]
  \item $\sin x=\tfrac12$
  \item $\cos x=0$
  \item $\tan x=1$
  \item $2\cos x+1=0$
  \item $\sin 2x=1$
  \item $\cos^{2}x=\tfrac14$
  \item $\sin x=\cos x$
\end{problems}

\section{Solutions}

\begin{solutions}
  \item $60\cdot\tfrac{\pi}{180}=\tfrac{\pi}{3}$.
  \item $270\cdot\tfrac{\pi}{180}=\tfrac{3\pi}{2}$.
  \item $\tfrac{3\pi}{4}\cdot\tfrac{180}{\pi}=135^{\circ}$.
  \item $\tfrac{5\pi}{3}\cdot\tfrac{180}{\pi}=300^{\circ}$.
  \item $\tfrac{180}{\pi}\approx57.3^{\circ}$.
  \item Because $\tfrac{\mathrm{d}}{\mathrm{d}x}\sin x=\cos x$ holds only
        in radians. In degrees the limit
        $\lim_{h\to0}\tfrac{\sin h}{h}$ equals $\tfrac{\pi}{180}$ instead
        of $1$, and that factor would appear in every derivative and integral
        of a trigonometric function.
  \item $\tfrac12$.
  \item $\tfrac{\sqrt2}{2}$.
  \item $\sqrt3$.
  \item Second quadrant, $\sin$ positive:
        $\sin\tfrac{5\pi}{6}=\sin\tfrac{\pi}{6}=\tfrac12$.
  \item Second quadrant, $\cos$ negative:
        $-\cos\tfrac{\pi}{3}=-\tfrac12$.
  \item $\tan\pi=\tfrac{\sin\pi}{\cos\pi}=\tfrac{0}{-1}=0$.
  \item $\sin$ is odd: $-\sin\tfrac{\pi}{3}=-\tfrac{\sqrt3}{2}$.
  \item $\tfrac{7\pi}{2}=\tfrac{3\pi}{2}+2\pi$, so
        $\cos\tfrac{7\pi}{2}=\cos\tfrac{3\pi}{2}=0$.
  \item $\tfrac{9\pi}{4}=\tfrac{\pi}{4}+2\pi$, so the value is
        $\sin\tfrac{\pi}{4}=\tfrac{\sqrt2}{2}$.
  \item Second quadrant, $\tan$ negative:
        $-\tan\tfrac{\pi}{4}=-1$.
  \item Period $2\pi$, amplitude $1$, range $[-1,1]$, no asymptotes.
  \item Period $2\pi$, amplitude $1$, range $[-1,1]$, no asymptotes.
  \item Period $\pi$, amplitude undefined (unbounded), range $\R$,
        vertical asymptotes at $x=\tfrac{\pi}{2}+k\pi$.
  \item Period $2\pi$, amplitude $3$, range $[-3,3]$.
  \item Period $\tfrac{2\pi}{2}=\pi$, amplitude $1$, range $[-1,1]$.
  \item Period $\tfrac{2\pi}{1/2}=4\pi$, amplitude $1$, range $[-1,1]$.
  \item Period $2\pi$, amplitude $2$, midline $y=1$, range $[-1,3]$.
  \item Period $\tfrac{\pi}{3}$, unbounded range $\R$, asymptotes where
        $3x=\tfrac{\pi}{2}+k\pi$, i.e. $x=\tfrac{\pi}{6}+\tfrac{k\pi}{3}$.
  \item A shift of $\tfrac{\pi}{2}$ to the right. It equals $-\cos x$,
        since $\sin\bigl(x-\tfrac{\pi}{2}\bigr)
        =\sin x\cos\tfrac{\pi}{2}-\cos x\sin\tfrac{\pi}{2}=-\cos x$.
  \item Amplitude $3$, period $\tfrac{2\pi}{2}=\pi$, midline $y=1$, range
        $[-2,4]$, shifted $\tfrac{\pi}{4}$ to the right.
  \item $\sin(2x-\pi)=\sin\bigl(2(x-\tfrac{\pi}{2})\bigr)$: amplitude $1$,
        period $\pi$, phase shift $\tfrac{\pi}{2}$ right. Reading the
        shift as $\pi$ from the unfactored form is the standard error.
  \item Amplitude $2$ (the minus sign reflects rather than changing the
        amplitude), period $\tfrac{2\pi}{3}$, range $[-2,2]$.
  \item Period $\pi$ needs $b=2$: $y=4\sin(2x)-1$.
  \item Period $6\pi$ needs $b=\tfrac13$:
        $y=2\cos\bigl(\tfrac{x}{3}\bigr)$.
  \item Midline $500$, amplitude $80$, period $12$ so
        $b=\tfrac{2\pi}{12}=\tfrac{\pi}{6}$, and a sine peaks a quarter
        period after its shift, so peaking at $t=3$ means $c=0$:
        \[
          S(t)=500+80\sin\!\left(\frac{\pi t}{6}\right).
        \]
        (Check: at $t=3$ the sine is $\sin\tfrac{\pi}{2}=1$, giving the
        peak $580$.)
  \item Because it is periodic: $f(x+T)=f(x)$ for $T=2\pi/|b|\neq0$, so two
        different inputs share a value. Injectivity fails everywhere, which
        is why the inverse functions need restricted domains.
  \item $\dfrac{\sin^{2}x}{\sin x}=\sin x$ (for $\sin x\neq0$).
  \item $\dfrac{\sin x}{\cos x}\cdot\cos x=\sin x$ (for
        $\cos x\neq0$).
  \item $-(1-\sin^{2}x)=-\cos^{2}x$.
  \item $\dfrac{2\sin x\cos x}{2\cos x}=\sin x$ (for $\cos x\neq0$).
  \item Divide $\sin^{2}x+\cos^{2}x=1$ by $\cos^{2}x$:
        $\tan^{2}x+1=\dfrac{1}{\cos^{2}x}$.
  \item $\cos2x=\cos^{2}x-\sin^{2}x$, and substituting
        $\cos^{2}x=1-\sin^{2}x$ gives $1-\sin^{2}x-\sin^{2}x
        =1-2\sin^{2}x$.
  \item $x=\tfrac{\pi}{6}+2k\pi$ or $x=\tfrac{5\pi}{6}+2k\pi$,
        $k\in\mathbb{Z}$.
  \item $x=\tfrac{\pi}{2}+k\pi$.
  \item $x=\tfrac{\pi}{4}+k\pi$.
  \item $\cos x=-\tfrac12$, so
        $x=\pm\tfrac{2\pi}{3}+2k\pi$.
  \item $2x=\tfrac{\pi}{2}+2k\pi$, hence
        $x=\tfrac{\pi}{4}+k\pi$.
  \item $\cos x=\pm\tfrac12$, giving
        $x=\pm\tfrac{\pi}{3}+2k\pi$ and $x=\pm\tfrac{2\pi}{3}+2k\pi$ ---
        more compactly $x=\tfrac{\pi}{3}+\tfrac{k\pi}{1}$ fails, so keep
        the four families, or write
        $x=\pm\tfrac{\pi}{3}+k\pi$.
  \item Dividing by $\cos x$ (legitimate, since $\cos x=0$ would force
        $\sin x=0$ too, impossible): $\tan x=1$, so
        $x=\tfrac{\pi}{4}+k\pi$.
\end{solutions}

\section{Additional Problems}

\begin{problems}
  \item $45^{\circ}$ and $150^{\circ}$ into radians.
  \item $\tfrac{7\pi}{6}$ and $\tfrac{11\pi}{6}$ into degrees.
  \item $\cos\tfrac{5\pi}{4}$
  \item $\sin\tfrac{4\pi}{3}$
  \item $\tan\tfrac{5\pi}{6}$
  \item $\sin\tfrac{11\pi}{6}$
  \item $\cos(-\tfrac{\pi}{6})$
  \item $\tan\tfrac{7\pi}{4}$
  \item Period, amplitude and range of $y=\tfrac12\cos 4x-2$.
  \item Period and asymptotes of $y=\tan\bigl(x-\tfrac{\pi}{4}\bigr)$.
  \item Amplitude, period, midline of
        $y=5\cos\bigl(\tfrac{\pi}{3}(x-2)\bigr)+3$.
  \item Write a cosine with maximum $7$, minimum $1$ and period $4$.
  \item Simplify $\dfrac{\cos^{2}x-1}{\cos x-1}$.
  \item Simplify $\sin x\cos x\cdot\dfrac{2}{\sin 2x}$.
  \item Simplify $\dfrac{1}{\cos^{2}x}-\tan^{2}x$.
  \item Verify $\sin(x+\pi)=-\sin x$.
  \item Verify $\cos\bigl(\tfrac{\pi}{2}-x\bigr)=\sin x$.
  \item Solve $\sin x=-\tfrac{\sqrt3}{2}$.
  \item Solve $\tan 2x=\sqrt3$.
  \item Solve $2\sin^{2}x-\sin x=0$.
  \item Solve $\cos 2x=\tfrac12$.
  \item Solve $\sin x=2$ and explain the answer.
  \item Give $\arcsin\tfrac12$, $\arccos(-1)$ and $\arctan(-1)$.
  \item Evaluate $\arcsin\bigl(\sin\tfrac{3\pi}{4}\bigr)$ and explain why
        the answer is not $\tfrac{3\pi}{4}$.
  \item Quarterly demand is $D(t)=200-40\cos\bigl(\tfrac{\pi t}{2}\bigr)$
        with $t$ in quarters. Give the period, the maximum and minimum
        demand, and the quarters in which each occurs.
\end{problems}

\section{Additional Solutions}

\begin{solutions}
  \item $\tfrac{\pi}{4}$ and $\tfrac{5\pi}{6}$.
  \item $210^{\circ}$ and $330^{\circ}$.
  \item Third quadrant, $\cos$ negative:
        $-\tfrac{\sqrt2}{2}$.
  \item Third quadrant, $\sin$ negative:
        $-\tfrac{\sqrt3}{2}$.
  \item Second quadrant, $\tan$ negative:
        $-\tfrac{1}{\sqrt3}$.
  \item Fourth quadrant, $\sin$ negative:
        $-\tfrac12$.
  \item $\cos$ is even: $\cos\tfrac{\pi}{6}=\tfrac{\sqrt3}{2}$.
  \item Fourth quadrant, $\tan$ negative: $-1$.
  \item Period $\tfrac{2\pi}{4}=\tfrac{\pi}{2}$, amplitude $\tfrac12$,
        midline $y=-2$, range $[-2.5,-1.5]$.
  \item Period $\pi$; asymptotes where
        $x-\tfrac{\pi}{4}=\tfrac{\pi}{2}+k\pi$, i.e.
        $x=\tfrac{3\pi}{4}+k\pi$.
  \item Amplitude $5$, period $\tfrac{2\pi}{\pi/3}=6$, midline $y=3$,
        shifted $2$ to the right; range $[-2,8]$.
  \item Midline $\tfrac{7+1}{2}=4$, amplitude $\tfrac{7-1}{2}=3$, period
        $4$ so $b=\tfrac{2\pi}{4}=\tfrac{\pi}{2}$:
        $y=3\cos\bigl(\tfrac{\pi x}{2}\bigr)+4$.
  \item $\dfrac{(\cos x-1)(\cos x+1)}{\cos x-1}=\cos x+1$ (for
        $\cos x\neq1$).
  \item $\dfrac{2\sin x\cos x}{\sin 2x}=\dfrac{\sin 2x}{\sin 2x}=1$.
  \item $\dfrac{1}{\cos^{2}x}-\tan^{2}x=1$ by the identity
        $1+\tan^{2}x=1/\cos^{2}x$.
  \item $\sin(x+\pi)=\sin x\cos\pi+\cos x\sin\pi=-\sin x$.
  \item $\cos\bigl(\tfrac{\pi}{2}-x\bigr)
        =\cos\tfrac{\pi}{2}\cos x+\sin\tfrac{\pi}{2}\sin x=\sin x$.
  \item $x=-\tfrac{\pi}{3}+2k\pi$ or
        $x=\pi+\tfrac{\pi}{3}+2k\pi=\tfrac{4\pi}{3}+2k\pi$.
  \item $2x=\tfrac{\pi}{3}+k\pi$, so
        $x=\tfrac{\pi}{6}+\tfrac{k\pi}{2}$.
  \item Factor: $\sin x\,(2\sin x-1)=0$, so $\sin x=0$ giving $x=k\pi$, or
        $\sin x=\tfrac12$ giving $x=\tfrac{\pi}{6}+2k\pi$ and
        $x=\tfrac{5\pi}{6}+2k\pi$.
  \item $2x=\pm\tfrac{\pi}{3}+2k\pi$, so
        $x=\pm\tfrac{\pi}{6}+k\pi$.
  \item No solution: $\sin$ never exceeds $1$, since it is a coordinate of
        a point on the unit circle. An equation can be perfectly
        well-formed and still have an empty solution set.
  \item $\tfrac{\pi}{6}$, \;$\pi$, \;$-\tfrac{\pi}{4}$.
  \item $\sin\tfrac{3\pi}{4}=\tfrac{\sqrt2}{2}$, and
        $\arcsin\tfrac{\sqrt2}{2}=\tfrac{\pi}{4}$. The answer is
        $\tfrac{\pi}{4}$, not $\tfrac{3\pi}{4}$, because $\arcsin$ returns
        values in $[-\tfrac{\pi}{2},\tfrac{\pi}{2}]$ only:
        $\arcsin(\sin x)=x$ holds just on that interval.
  \item Period $\tfrac{2\pi}{\pi/2}=4$ quarters, i.e. one year. Maximum
        $240$ where $\cos\bigl(\tfrac{\pi t}{2}\bigr)=-1$, i.e.
        $\tfrac{\pi t}{2}=\pi+2k\pi$ so $t=2,6,10,\dots$; minimum $160$
        where the cosine is $1$, i.e. $t=0,4,8,\dots$. Demand peaks in the
        third quarter of each year and troughs in the first.
\end{solutions}


\end{document}
"
\ifdefined\KindleEdition
  \documentclass[10pt,oneside]{book}
\else
  \documentclass[11pt,a4paper,twoside,openright]{book}
\fi
% Shared preamble for the Functions training sheet.
\usepackage[T1]{fontenc}
\usepackage[utf8]{inputenc}
\usepackage{lmodern}
\usepackage[a4paper,margin=2.4cm]{geometry}
\usepackage{amsmath,amssymb,amsthm}
\usepackage{microtype}
\usepackage{enumitem}
\usepackage{xcolor}
\usepackage[most]{tcolorbox}
\usepackage{fancyhdr}
\usepackage[hidelinks]{hyperref}

\definecolor{accent}{HTML}{1F4E79}
\definecolor{soft}{HTML}{F4F7FA}

% --- function graphs ------------------------------------------------------
\usepackage{pgfplots}
\pgfplotsset{compat=1.18}
\usetikzlibrary{arrows.meta}

% One consistent look for every function graph in these books.
\pgfplotsset{
  fnaxis/.style={
    width=5.9cm, height=4.6cm,
    axis lines=middle,
    axis line style={-{Stealth[length=4pt]}, gray!65!black, line width=0.3pt},
    tick style={gray!65!black},
    tick label style={font=\tiny, gray!50!black},
    label style={font=\scriptsize},
    xlabel={$x$}, ylabel={$y$},
    every axis x label/.style={at={(ticklabel* cs:1.0)}, anchor=west},
    every axis y label/.style={at={(ticklabel* cs:1.0)}, anchor=south east,
      xshift=-1pt},
    every axis plot/.append style={accent, very thick, smooth, samples=201},
    clip=true, enlargelimits=false,
    title style={font=\scriptsize, yshift=4pt},
  },
  fnaxis/.append style={/pgfplots/legend style={font=\tiny, draw=none,
    fill=none, at={(0.98,0.03)}, anchor=south east}},
}

% A row of small graphs, evenly spaced with \hfill between them, plus an
% optional caption line: \begin{gallery} <pic>\hfill<pic> \gcap{...} \end{gallery}
\newenvironment{gallery}
  {\par\medskip\noindent\ignorespaces}
  {\par\medskip\noindent}
\newcommand{\gcap}[1]{\par\smallskip{\centering\footnotesize #1\par}}

% --- plane regions, implicit curves, surfaces -----------------------------
\usepgfplotslibrary{fillbetween}
\definecolor{regionfill}{HTML}{9DB4D8}
\pgfplotsset{colormap={house}{color=(soft) color=(accent!45) color=(accent)}}

\pgfplotsset{
  % Squared, gridded axis for plane regions and implicitly given curves.
  geoaxis/.style={
    fnaxis,
    width=5.9cm, height=5.9cm,
    grid=both,
    grid style={gray!22, line width=0.2pt},
    minor tick num=1,
  },
  % Level-curve maps: many thin curves, no single curve emphasised.
  contouraxis/.style={
    geoaxis,
    every axis plot/.append style={accent!75, line width=0.5pt, smooth,
      samples=201},
  },
  % Surfaces: \addplot3[surf, domain=..., y domain=...] {expr};
  surfaxis/.style={
    width=6.0cm, height=5.0cm,
    view={130}{22},
    colormap name=house,
    ticks=none,
    xlabel={$x$}, ylabel={$y$}, zlabel={$f$},
    label style={font=\tiny, gray!50!black},
    title style={font=\scriptsize},
    every axis plot/.append style={samples=25, shader=faceted interp,
      draw=accent!35, line width=0.1pt},
  },
}

% Shorthand for a shaded region between two named paths.
% \region{<path A>}{<path B>}{<soft clip options>}
\newcommand{\region}[3]{%
  \addplot[regionfill, opacity=0.55, forget plot, sharp plot]
    fill between[of=#1 and #2, #3];}

% --- rule boxes -----------------------------------------------------------
\newtcolorbox{rulebox}[1]{%
  breakable, enhanced, colback=soft, colframe=accent,
  boxrule=0.4pt, left=7pt, right=7pt, top=6pt, bottom=6pt,
  fonttitle=\bfseries, coltitle=white, colbacktitle=accent,
  title={#1}}

% --- worked examples ------------------------------------------------------
% One fully narrated solution per chapter, carrying the chapter's central
% technique and the mistakes that technique invites:
%   \begin{worked}{Task} ... \wstep{Step 1.} ... \wtrap{...} ... \end{worked}
\newtcolorbox{worked}[1]{%
  breakable, enhanced, colback=white, colframe=accent!65,
  boxrule=0.4pt, left=7pt, right=7pt, top=6pt, bottom=6pt,
  fonttitle=\bfseries, coltitle=accent, colbacktitle=soft,
  title={#1}}

% Run-in step heading inside a worked example.
\newcommand{\wstep}[1]{\par\smallskip\noindent\textbf{#1}\enspace\ignorespaces}

% A pitfall, set off from the running solution so it can be found again.
\newcommand{\wtrap}[1]{%
  \par\smallskip
  {\leftskip=1.3em\noindent
   \llap{\textcolor{accent}{\rule[0.22ex]{0.9em}{0.6pt}}\hspace{0.4em}}%
   \textbf{Trap.}\enspace\itshape #1\par}%
  \smallskip}

% --- problem / solution lists --------------------------------------------
\newlist{problems}{enumerate}{1}
\setlist[problems]{label=\textbf{\arabic*.}, leftmargin=2.1em,
  itemsep=0.55ex, parsep=0pt, topsep=0.8ex}
\newlist{solutions}{enumerate}{1}
\setlist[solutions]{label=\textbf{\arabic*.}, leftmargin=2.1em,
  itemsep=1.0ex, parsep=0pt, topsep=0.8ex}

% --- running heads --------------------------------------------------------
\pagestyle{fancy}
\fancyhf{}
\fancyhead[LE,RO]{\thepage}
\fancyhead[RE]{\small\itshape\nouppercase{\leftmark}}
\fancyhead[LO]{\small\itshape\nouppercase{\rightmark}}
\setlength{\headheight}{14pt}
\renewcommand{\headrulewidth}{0.3pt}
\renewcommand{\chaptermark}[1]{\markboth{\chaptername\ \thechapter.\ #1}{}}
\renewcommand{\sectionmark}[1]{\markright{\thesection\ #1}}

% --- shorthands used throughout -------------------------------------------
\newcommand{\dd}[2]{\frac{\mathrm{d}#1}{\mathrm{d}#2}}
\newcommand{\pd}[2]{\frac{\partial #1}{\partial #2}}
\newcommand{\pdd}[3]{\frac{\partial^{2} #1}{\partial #2\,\partial #3}}
\newcommand{\R}{\mathbb{R}}
\newcommand{\grad}{\nabla}
\DeclareMathOperator{\sgn}{sgn}
\DeclareMathOperator{\Hess}{H}

% German gloss, used on first mention of a technical term
\newcommand{\de}[1]{\emph{#1}}

\renewcommand*{\thisvolume}{functions}

\title{\Huge Functions and Their Graphs\\[0.4em]
  \large A Refresher with Exercises and Worked Solutions\\[0.8em]
  \normalsize Domain \(\cdot\) Range \(\cdot\) Graphs \(\cdot\)
  Transformations \(\cdot\) Regions \(\cdot\) Trigonometry}
\author{}
\date{}

\begin{document}
\frontmatter
\maketitle

\chapter*{About this sheet}
\addcontentsline{toc}{chapter}{About this sheet}

The algebra volume manipulates \emph{expressions}. This one is about
\emph{functions}: an expression together with a decision about which numbers
may be substituted into it, a record of which numbers come out, and above
all a picture.

Everything that follows in the later volumes leans on this material and
rarely stops to explain it. A limit is taken at a point of the domain or at
its boundary; a derivative exists only where the function does; an
asymptote is read off a graph; a constrained optimum is the point where two
families of curves touch. A mathematics sequence for economics or business
takes this for granted: one is expected to know what the graphs of the usual
functions look like, and to be able to sketch curves defined by equations in
two variables. That is the brief of this volume.

Five chapters, in the order they build on each other. \textbf{Domain and
range} fixes what a function is and where it lives. \textbf{The standard
functions and their graphs} is the catalogue --- powers, roots, exponentials,
logarithms, rational functions with their poles and asymptotes --- together
with the curves that are not graphs at all: circles, ellipses, hyperbolas,
isoquants. \textbf{Transforming functions} covers shifting, stretching,
reflecting, composing and inverting, and closes with the question an
economist meets most often, which is the reverse one: given the shape you
need, which formula has it? \textbf{Inequalities in two variables} turns
curves into regions, the skill behind every feasible set. And
\textbf{trigonometric functions} supplies the periodic shapes that later material
assumes and the calculus chapters use as test cases.

The structure is the same in every chapter: \textbf{Core Rules} states the
rules compactly; \textbf{Practice Problems} drills them; \textbf{Solutions}
works every problem through in full; each chapter closes with
\emph{additional problems} for a second pass. Seventy problems per chapter,
every one with a worked solution.

Work the problems before reading the solutions, and draw the pictures by
hand. A curve you have sketched yourself is one you will recognise in an
examination.

\tableofcontents

\mainmatter
\chapter{Functions, Domain and Range}

Up to here every chapter has manipulated \emph{expressions}. A
\de{Funktion} is the next step: an expression together with a decision
about which numbers may be substituted into it, and a record of which
numbers come out. Those two sets, the \de{Definitionsbereich} (domain) and
the \de{Wertebereich} (range), are not decoration. A limit is taken at a
point of the domain or at its boundary; a derivative exists only where the
function does; an optimisation problem is solved over a domain; and a
demand function that returns a negative price has left the region where it
means anything.

Finding a maximal domain is pure algebra of the kind the previous chapters
drilled: a denominator must not vanish, an even root needs a non-negative
radicand, a logarithm needs a positive argument. Each condition is an
equation or an inequality, and the domain is what survives all of them at
once. Finding a range is slightly different work: either the function is
built from pieces whose ranges are known, or the equation $y=f(x)$ is
solved for $x$ and one asks which $y$ admit a solution.

\section{Core Rules}

\begin{rulebox}{Function, domain, range}
A function $f\colon D\to\R$ assigns to every $x\in D$ exactly one value
$f(x)\in\R$. Two sets belong to the description:
\begin{itemize}[leftmargin=1.6em,itemsep=0.3ex]
  \item the \textbf{domain} $D(f)$ \de{(Definitionsbereich)}: the $x$ that
        may be substituted;
  \item the \textbf{range} $W(f)=\{f(x): x\in D(f)\}$
        \de{(Wertebereich, Bildmenge)}: the values actually attained.
\end{itemize}
The \textbf{graph} of $f$ is the set $\{(x,f(x)) : x\in D(f)\}$ in the
plane. A curve is the graph of a function exactly when every vertical line
meets it at most once (\textbf{vertical line test}). Reading off a graph:
$D(f)$ is its shadow on the $x$-axis, $W(f)$ its shadow on the $y$-axis.
\end{rulebox}

\begin{gallery}
\begin{tikzpicture}
  \begin{axis}[fnaxis, xmin=-0.6, xmax=5.4, ymin=-0.6, ymax=3.4,
      xtick={1,3,5}, ytick={1,2,3}, title={$f(x)=\sqrt{x-1}$}]
    \addplot[domain=1:5]{sqrt(x-1)};
  \end{axis}
\end{tikzpicture}\hfill
\begin{tikzpicture}
  \begin{axis}[fnaxis, xmin=-3.4, xmax=3.4, ymin=-0.6, ymax=4.4,
      xtick={-2,2}, ytick={2,4}, title={$f(x)=x^{2}$}]
    \addplot[domain=-2.1:2.1]{x^2};
  \end{axis}
\end{tikzpicture}\hfill
\begin{tikzpicture}
  \begin{axis}[fnaxis, xmin=-3.4, xmax=3.4, ymin=-1.2, ymax=4.4,
      xtick={-2,2}, ytick={1,3}, title={$f(x)=\mathrm{e}^{x}$}]
    \addplot[domain=-3:1.45]{exp(x)};
  \end{axis}
\end{tikzpicture}
\gcap{$D=[1,\infty)$, $W=[0,\infty)$ \quad\textbullet\quad
$D=\R$, $W=[0,\infty)$ \quad\textbullet\quad $D=\R$, $W=(0,\infty)$.}
\end{gallery}

\begin{rulebox}{The maximal domain: three obstructions}
Unless a domain is prescribed, $D(f)$ is the \textbf{maximal domain}: every
real $x$ for which the formula produces a real number. Only three things
can go wrong, and each is an (in)equality to solve:
\begin{itemize}[leftmargin=1.6em,itemsep=0.3ex]
  \item \textbf{Denominator:} $\dfrac{\;\cdot\;}{g(x)}$ requires
        $g(x)\neq 0$. Solve $g(x)=0$ and delete those points.
  \item \textbf{Even root:} $\sqrt[2k]{g(x)}$ requires $g(x)\geq 0$.
        Odd roots ($\sqrt[3]{\;}$) are defined for all reals.
  \item \textbf{Logarithm:} $\ln g(x)$ requires $g(x)>0$ (strictly).
\end{itemize}
A power $g(x)^{r}$ with non-integer $r$ counts as a root: it requires
$g(x)\geq 0$ for $r>0$ and $g(x)>0$ for $r<0$. If several obstructions
occur, \textbf{intersect} the resulting sets. A root in a denominator
combines two conditions into one: $\dfrac{1}{\sqrt{g(x)}}$ needs
$g(x)>0$.
\end{rulebox}

\begin{rulebox}{Domains prescribed by context}
In economics the domain is usually narrower than the maximal domain and
comes from the model, not from the formula:
quantities and inputs satisfy $q\geq 0$, $K,L\geq 0$; a price must satisfy
$p\geq 0$; a share or probability lives in $[0,1]$; an average cost
$AC(q)=C(q)/q$ needs $q>0$. State the economic domain explicitly, because
the formula will happily return values outside it: the linear demand
$p=120-3q$ is defined for every real $q$, but only $q\in[0,40]$ keeps both
$q\geq0$ and $p\geq 0$.
\end{rulebox}

\begin{rulebox}{Finding the range}
Four routes, in increasing order of generality:
\begin{enumerate}[leftmargin=1.7em,itemsep=0.3ex]
  \item \textbf{Standard building blocks.}
        $x^{2}\colon[0,\infty)$, \;$\sqrt{x}\colon[0,\infty)$,
        \;$|x|\colon[0,\infty)$, \;$\mathrm{e}^{x}\colon(0,\infty)$,
        \;$\ln x\colon\R$, \;$1/x\colon\R\setminus\{0\}$,
        \;$x^{3}\colon\R$, \;$\sin x,\cos x\colon[-1,1]$.
        Then track what the outer operations do: adding $c$ shifts the
        range by $c$, multiplying by $a>0$ scales it, multiplying by $a<0$
        scales \emph{and} flips it.
  \item \textbf{Complete the square} for a quadratic:
        $a(x-h)^{2}+k$ has range $[k,\infty)$ if $a>0$ and
        $(-\infty,k]$ if $a<0$.
  \item \textbf{Monotonicity.} If $f$ is strictly increasing on $[a,b]$,
        then $W=[f(a),f(b)]$; open ends stay open, and an end at
        $\pm\infty$ is replaced by the limiting value.
  \item \textbf{Solve $y=f(x)$ for $x$} (the general method). The range is
        the set of $y$ for which the solution exists and lies in $D(f)$.
        For $f(x)=\dfrac{x}{x+1}$: \;$y(x+1)=x\Rightarrow x=\dfrac{y}{1-y}$,
        solvable for every $y\neq 1$, hence $W=\R\setminus\{1\}$.
\end{enumerate}
\end{rulebox}

\begin{rulebox}{Composite functions}
$(g\circ f)(x)=g\bigl(f(x)\bigr)$ is evaluated \textbf{inside out}. Its
domain carries two conditions at once:
\[
  D(g\circ f)=\{\,x\in D(f)\;:\;f(x)\in D(g)\,\}.
\]
So $x$ must be admissible for the inner function, and its \emph{output}
must be admissible for the outer one. For $g(u)=\ln u$ and $f(x)=x^{2}-1$:
$x^{2}-1>0$, hence $D=(-\infty,-1)\cup(1,\infty)$. Order matters:
$g\circ f\neq f\circ g$ in general, and the two domains usually differ.
\end{rulebox}

\begin{rulebox}{Symmetry, monotonicity, injectivity}
\begin{itemize}[leftmargin=1.6em,itemsep=0.3ex]
  \item $f$ is \textbf{even} \de{(gerade)} if $f(-x)=f(x)$ for all $x\in D$
        (graph symmetric about the $y$-axis), \textbf{odd}
        \de{(ungerade)} if $f(-x)=-f(x)$ (symmetric about the origin).
        Only even powers $\Rightarrow$ even; only odd powers
        $\Rightarrow$ odd; a mix is neither.
  \item $f$ is \textbf{strictly increasing} if $x_{1}<x_{2}$ implies
        $f(x_{1})<f(x_{2})$, strictly decreasing if the inequality
        reverses.
  \item $f$ is \textbf{injective} \de{(eineindeutig)} if $f(x_{1})=f(x_{2})$
        forces $x_{1}=x_{2}$: no horizontal line meets the graph twice.
        Strictly monotone $\Rightarrow$ injective $\Rightarrow$ invertible
        on its range (Chapter~3).
\end{itemize}
\end{rulebox}

\section{Worked Example}

\begin{worked}{Find the maximal domain of $f(x)=\dfrac{\ln(x+3)}{\sqrt{4-x^{2}}}$, and the range of $g(x)=\dfrac{x+3}{x-2}$}
\wstep{Step 1. List the obstructions, one per structure in the formula.}
Three of them appear here, which is why this example was chosen:
\begin{itemize}[leftmargin=1.6em,itemsep=0.2ex]
  \item the logarithm needs $x+3>0$;
  \item the even root needs $4-x^{2}\ge0$;
  \item the denominator needs $\sqrt{4-x^{2}}\neq0$.
\end{itemize}
\wtrap{The last two combine into the \emph{strict} inequality $4-x^{2}>0$.
A square root sitting in a denominator is not allowed to be zero, so the
usual ``$\ge$'' for an even root tightens to ``$>$''. Handling the root and
the denominator as separate conditions and then forgetting to intersect them
is the standard way to hand in $[-2,2]$ instead of the right answer.}

\wstep{Step 2. Solve each condition separately.}
\[
  x+3>0 \iff x>-3 ;
  \qquad
  4-x^{2}>0 \iff x^{2}<4 \iff -2<x<2 .
\]

\wstep{Step 3. Intersect --- never unite.}
The function is defined where \emph{all} conditions hold at once:
\[
  D_{f}=(-3,\infty)\cap(-2,2)=\boxed{(-2,2)} .
\]
\wtrap{``And'' means intersection. Each obstruction removes points; nothing
in the formula ever gives points back. If your answer is larger than one of
the individual conditions, you took a union by mistake.}

\wstep{Step 4. Check the boundary behaviour, which confirms the strictness.}
As $x\to2^{-}$ the denominator $\to0^{+}$ while $\ln5>0$, so $f\to+\infty$.
The endpoint really must be excluded. At $x=-2$ the same thing happens.

\wstep{Step 5. The range of $g$: solve $y=g(x)$ for $x$.}
The range is the set of $y$ for which the equation $y=g(x)$ has a solution.
So treat $y$ as given and solve:
\[
  y=\frac{x+3}{x-2}
  \;\Longrightarrow\;
  y(x-2)=x+3
  \;\Longrightarrow\;
  yx-x=2y+3
  \;\Longrightarrow\;
  x(y-1)=2y+3 .
\]

\wstep{Step 6. Read off which $y$ fail.}
\[
  x=\frac{2y+3}{y-1},
\]
which exists for every $y$ except $y=1$; and at $y=1$ the previous line reads
$0=5$, impossible. Hence
\[
  \boxed{\;\text{range}(g)=\R\setminus\{1\}\;}
\]

\wstep{Step 7. Confirm it against the graph.}
$g(x)=\frac{x+3}{x-2}=1+\frac{5}{x-2}$, so the graph is the hyperbola
$5/x$ shifted, with horizontal asymptote $y=1$. The value $1$ is approached
but never attained --- exactly what the algebra said. Two independent
routes agreeing is worth the extra line.
\wtrap{Do not assume the range of a rational function is all of $\R$, and do
not assume it is the same as the domain. The excluded input ($x=2$) and the
excluded output ($y=1$) are different numbers, read off different parts of
the formula.}
\end{worked}

\section{Practice Problems}

State the maximal domain $D(f)\subseteq\R$ in problems 1--18, using
interval notation.

\begin{problems}
  \item $f(x)=\dfrac{x+3}{x-2}$
  \item $f(x)=\sqrt{x-5}$
  \item $f(x)=\sqrt{9-x^{2}}$
  \item $f(x)=\ln(x-4)$
  \item $f(x)=\ln(x^{2}-4)$
  \item $f(x)=\dfrac{\sqrt{x}}{x-1}$
  \item $f(x)=\dfrac{1}{\sqrt{x-3}}$
  \item $f(x)=\sqrt{x^{2}-x-6}$
  \item $f(x)=\dfrac{\ln x}{x-2}$
  \item $f(x)=\mathrm{e}^{1/x}$
  \item $f(x)=\sqrt{x+2}+\sqrt{4-x}$
  \item $f(x)=\dfrac{x^{2}-1}{x^{2}-3x+2}$
  \item $f(x)=\ln(\ln x)$
  \item $f(x)=x^{3/4}$
  \item $f(x)=\dfrac{1}{1-\mathrm{e}^{x}}$
  \item $f(x)=\sqrt{\ln x}$
  \item $f(x)=\sqrt[3]{x-7}$
  \item $f(x)=\dfrac{\sqrt{4-x}}{x^{2}-x}$
\end{problems}

\begin{problems}[resume]
  \item A firm has cost $C(q)=50+2q$. State the economic domain of $C$
        and of the average cost $AC(q)=C(q)/q$.
  \item For the demand schedule $p(q)=120-3q$, state the largest interval
        of quantities on which both $q\geq0$ and $p\geq0$ hold.
  \item State the domain of the Cobb--Douglas product
        $Y(K,L)=K^{0.3}L^{0.7}$ as a subset of the $(K,L)$-plane, and say
        why the exponents force the restriction.
  \item Utility is $U(x)=\ln x$. State the economic domain and explain what
        happens to $U$ as $x\to 0^{+}$.
\end{problems}

Determine the range $W(f)$ in problems 23--36. The domain is the maximal
one unless stated otherwise.

\begin{problems}[resume]
  \item $f(x)=x^{2}+4$
  \item $f(x)=x^{2}-6x+11$
  \item $f(x)=-2x^{2}+8x-3$
  \item $f(x)=\mathrm{e}^{x}-3$
  \item $f(x)=5-\mathrm{e}^{x}$
  \item $f(x)=\ln(x+2)$
  \item $f(x)=\dfrac{1}{x-2}+3$
  \item $f(x)=\sqrt{x-1}$
  \item $f(x)=\sqrt{16-x^{2}}$
  \item $f(x)=|x-3|$
  \item $f(x)=\dfrac{2x+1}{x-3}$
  \item $f(x)=\dfrac{x^{2}}{x^{2}+1}$
  \item $f(x)=x^{2}$ on $D=[-1,3]$
  \item $f(x)=2x-5$ on $D=(1,4]$
\end{problems}

\begin{problems}[resume]
  \item With $f(u)=\sqrt{u}$ and $g(x)=x-4$: give $(f\circ g)(x)$ and
        $(g\circ f)(x)$ with both domains.
  \item With $f(u)=1/u$ and $g(x)=x+2$: give $(f\circ g)(x)$ and its
        domain.
  \item With $f(u)=\ln u$ and $g(x)=9-x^{2}$: give $(f\circ g)(x)$ and its
        domain.
  \item With $f(u)=\sqrt{u}$ and $g(x)=1-x^{2}$: give $(f\circ g)(x)$, its
        domain and its range.
  \item Classify as even, odd, or neither:
        (a) $x^{4}-3x^{2}$, \;(b) $x^{3}-x$, \;(c) $x^{2}+x$,
        \;(d) $\dfrac{1}{x}$.
  \item Which of $x^{3}$, $x^{2}$, $\mathrm{e}^{x}$, $|x|$ are injective on
        $\R$?
  \item Give domain and range of
        $f(x)=\begin{cases} x^{2}, & x<1,\\ 2x-1, & x\geq 1.\end{cases}$
  \item Does $x^{2}+y^{2}=25$ define $y$ as a function of $x$? Justify with
        the vertical line test, and give the two functions obtained by
        solving for $y$, with their domains.
  \item A graph consists of the segment from $(-2,1)$ to $(0,5)$ together
        with the horizontal ray from $(0,5)$ to the right. Read off
        $D$ and $W$.
\end{problems}

\section{Solutions}

\begin{solutions}
  \item Denominator $x-2\neq0$: \;$D=\R\setminus\{2\}
        =(-\infty,2)\cup(2,\infty)$.
  \item $x-5\geq0$: \;$D=[5,\infty)$.
  \item $9-x^{2}\geq0\iff x^{2}\leq9\iff|x|\leq3$: \;$D=[-3,3]$.
  \item $x-4>0$: \;$D=(4,\infty)$.
  \item $x^{2}-4>0\iff|x|>2$: \;$D=(-\infty,-2)\cup(2,\infty)$.
  \item Two conditions: $x\geq0$ for the root and $x\neq1$ for the
        denominator: \;$D=[0,1)\cup(1,\infty)$.
  \item The root sits in the denominator, so $x-3>0$ (not merely
        $\geq0$): \;$D=(3,\infty)$.
  \item $x^{2}-x-6=(x-3)(x+2)\geq0$; the parabola opens upwards, so the
        product is non-negative outside the roots:
        \;$D=(-\infty,-2]\cup[3,\infty)$.
  \item $\ln x$ needs $x>0$, the denominator needs $x\neq2$:
        \;$D=(0,2)\cup(2,\infty)$.
  \item Only $x\neq0$; the exponential itself is unrestricted:
        \;$D=\R\setminus\{0\}$.
  \item $x+2\geq0$ and $4-x\geq0$, i.e. $x\geq-2$ \emph{and} $x\leq4$:
        \;$D=[-2,4]$.
  \item $x^{2}-3x+2=(x-1)(x-2)$: \;$D=\R\setminus\{1,2\}$. (The numerator
        also vanishes at $x=1$, but a cancelling factor does not repair the
        domain: $x=1$ stays excluded.)
  \item Inner: $x>0$; outer: $\ln x>0\iff x>1$. The stronger condition
        wins: \;$D=(1,\infty)$.
  \item Exponent $3/4$ is not an integer, so the base must satisfy
        $x\geq0$: \;$D=[0,\infty)$.
  \item $1-\mathrm{e}^{x}=0\iff \mathrm{e}^{x}=1\iff x=0$:
        \;$D=\R\setminus\{0\}$.
  \item $\ln x$ needs $x>0$, and $\ln x\geq0$ needs $x\geq1$; the
        intersection is \;$D=[1,\infty)$.
  \item An odd root accepts every real radicand: \;$D=\R$.
  \item Root: $4-x\geq0\iff x\leq4$. Denominator: $x^{2}-x=x(x-1)\neq0$,
        so $x\neq0$ and $x\neq1$. Together:
        \;$D=(-\infty,0)\cup(0,1)\cup(1,4]$.
  \item $C$: production cannot be negative, so $D(C)=[0,\infty)$; at $q=0$
        the value $C(0)=50$ is the fixed cost. $AC$: dividing by $q$
        excludes the origin, so $D(AC)=(0,\infty)$.
  \item $p\geq0\iff 120-3q\geq0\iff q\leq40$; with $q\geq0$ this gives
        $q\in[0,40]$.
  \item $D=\{(K,L): K\geq0,\;L\geq0\}$, the non-negative quadrant. The
        exponents $0.3$ and $0.7$ are not integers, so $K^{0.3}$ and
        $L^{0.7}$ require non-negative bases; negative input quantities
        would be meaningless anyway.
  \item $D=(0,\infty)$: consumption must be strictly positive, since
        $\ln 0$ is undefined. As $x\to0^{+}$, $U(x)\to-\infty$: the
        logarithmic utility punishes zero consumption infinitely hard,
        which is exactly why interior solutions appear in Cobb--Douglas
        problems.
  \item $x^{2}\geq0$, so $x^{2}+4\geq4$: \;$W=[4,\infty)$.
  \item $x^{2}-6x+11=(x-3)^{2}+2$: \;$W=[2,\infty)$.
  \item $-2x^{2}+8x-3=-2(x^{2}-4x)-3=-2(x-2)^{2}+5$; the parabola opens
        downwards: \;$W=(-\infty,5]$.
  \item $\mathrm{e}^{x}\in(0,\infty)$, shifted down by $3$:
        \;$W=(-3,\infty)$.
  \item $\mathrm{e}^{x}\in(0,\infty)$, negated to $(-\infty,0)$, shifted up
        by $5$: \;$W=(-\infty,5)$.
  \item $\ln$ attains every real value, and $x+2$ runs through all of
        $(0,\infty)$ on $D=(-2,\infty)$: \;$W=\R$.
  \item $1/(x-2)$ attains every value except $0$; adding $3$ shifts this:
        \;$W=\R\setminus\{3\}$.
  \item A square root returns exactly the non-negative numbers, and
        $x-1$ runs through $[0,\infty)$: \;$W=[0,\infty)$.
  \item $D=[-4,4]$, on which $16-x^{2}$ runs through $[0,16]$; the root of
        that is \;$W=[0,4]$.
  \item An absolute value is non-negative and attains every non-negative
        value (at $x=3\pm w$): \;$W=[0,\infty)$.
  \item Solve $y=\dfrac{2x+1}{x-3}$: \;$y(x-3)=2x+1\Rightarrow
        x(y-2)=3y+1\Rightarrow x=\dfrac{3y+1}{y-2}$, which exists exactly
        for $y\neq2$: \;$W=\R\setminus\{2\}$. ($y=2$ is the horizontal
        asymptote.)
  \item Solve $y=\dfrac{x^{2}}{x^{2}+1}$: \;$y(x^{2}+1)=x^{2}\Rightarrow
        x^{2}=\dfrac{y}{1-y}$. A real $x$ exists iff
        $\dfrac{y}{1-y}\geq0$, i.e. $0\leq y<1$: \;$W=[0,1)$.
  \item On $[-1,3]$ the square is not monotone: it falls to $0$ at $x=0$
        and rises to $9$ at $x=3$. Minimum $0$, maximum $9$:
        \;$W=[0,9]$.
  \item $f$ is strictly increasing, so the endpoints map to the endpoints:
        $f(1)=-3$ (excluded, as $1\notin D$) and $f(4)=3$ (included):
        \;$W=(-3,3]$.
  \item $(f\circ g)(x)=\sqrt{x-4}$ with $D=[4,\infty)$;
        $(g\circ f)(x)=\sqrt{x}-4$ with $D=[0,\infty)$. Same ingredients,
        different functions and different domains.
  \item $(f\circ g)(x)=\dfrac{1}{x+2}$, requiring $x+2\neq0$:
        \;$D=\R\setminus\{-2\}$.
  \item $(f\circ g)(x)=\ln(9-x^{2})$, requiring $9-x^{2}>0\iff|x|<3$:
        \;$D=(-3,3)$.
  \item $(f\circ g)(x)=\sqrt{1-x^{2}}$, requiring $1-x^{2}\geq0$:
        \;$D=[-1,1]$. On it $1-x^{2}$ runs through $[0,1]$, so
        $W=[0,1]$ (the upper unit semicircle).
  \item (a) even (only even powers); (b) odd (only odd powers);
        (c) neither: $f(-x)=x^{2}-x$ equals neither $f(x)$ nor $-f(x)$;
        (d) odd: $1/(-x)=-1/x$.
  \item $x^{3}$ and $\mathrm{e}^{x}$ are injective (both strictly
        increasing). $x^{2}$ and $|x|$ are not: $f(-1)=f(1)$ in both cases.
  \item $D=\R$. First piece: $x$ runs through $(-\infty,1)$, so $x^{2}$
        attains every value in $[0,\infty)$. Second piece: $2x-1$ on
        $[1,\infty)$ gives $[1,\infty)$. Union: \;$W=[0,\infty)$.
  \item No. The vertical line $x=0$ meets the circle at $(0,5)$ and
        $(0,-5)$, so the assignment is not unique. Solving gives two
        functions, $y=\sqrt{25-x^{2}}$ (upper semicircle) and
        $y=-\sqrt{25-x^{2}}$ (lower), each with domain $[-5,5]$ and ranges
        $[0,5]$ and $[-5,0]$.
  \item $D=[-2,\infty)$: the segment covers $x\in[-2,0]$, the ray covers
        $x\in[0,\infty)$. On the segment $y$ climbs from $1$ to $5$, on the
        ray $y$ stays at $5$: \;$W=[1,5]$.
\end{solutions}

\section{Additional Problems}

\begin{problems}
  \item $f(x)=\dfrac{3x}{2x+7}$: maximal domain.
  \item $f(x)=\sqrt{2x+8}$: maximal domain and range.
  \item $f(x)=\ln(5-2x)$: maximal domain.
  \item $f(x)=\dfrac{1}{\ln x}$: maximal domain.
  \item $f(x)=\dfrac{x-1}{\sqrt{x^{2}-4}}$: maximal domain.
  \item $f(x)=\sqrt{\dfrac{x-1}{x+2}}$: maximal domain.
  \item $f(x)=\ln\!\bigl(x^{2}+1\bigr)$: maximal domain and range.
  \item $f(x)=\mathrm{e}^{-x^{2}}$: maximal domain and range.
  \item $f(x)=\dfrac{1}{x^{2}-5x+6}$: maximal domain.
  \item $f(x)=\sqrt{4-x}+\dfrac{1}{x+1}$: maximal domain.
  \item $f(x)=(x-3)^{-1/2}$: maximal domain.
  \item $f(x)=\dfrac{|x|}{x}$: maximal domain and range.
  \item $f(x)=3-\dfrac{2}{x+1}$: range.
  \item $f(x)=\dfrac{x-1}{x+1}$: range.
  \item $f(x)=4-\sqrt{x+9}$: range.
  \item $f(x)=x^{2}-4x+7$ on $D=[0,3]$: range.
  \item $f(x)=\mathrm{e}^{2x}+1$ on $D=[0,\ln 2]$: range.
  \item $f(x)=\ln x$ on $D=[1,\mathrm{e}^{3}]$: range.
  \item With $f(u)=u^{2}$ and $g(x)=\sqrt{x-3}$: give $(f\circ g)(x)$ and
        its domain, and explain why the simplified formula alone would give
        the wrong answer.
  \item With $f(u)=\dfrac{1}{u-1}$ and $g(x)=\mathrm{e}^{x}$: give
        $(f\circ g)(x)$ and its domain.
  \item Revenue is $R(q)=q\,(200-4q)$. State the economic domain on which
        both $q\ge 0$ and the price stay non-negative, and the range of $R$
        on it.
  \item A production function is $Q(L)=10\sqrt{L}$ with $L\in[0,100]$.
        Give the range, and the average product $Q(L)/L$ with its domain.
  \item Classify as even, odd, or neither: (a) $\mathrm{e}^{-x^{2}}$,
        \;(b) $x\,|x|$, \;(c) $\ln x$, \;(d) $\dfrac{x^{3}}{x^{2}+1}$.
  \item Show that $f(x)=\dfrac{2x-1}{x+4}$ is injective on its domain by
        solving $f(x_{1})=f(x_{2})$.
  \item Give domain and range of
        $f(x)=\begin{cases} -x, & x<0,\\ \sqrt{x}, & 0\le x\le 4,\\
        2, & x>4.\end{cases}$
\end{problems}

\section{Additional Solutions}

\begin{solutions}
  \item $2x+7\neq0\iff x\neq-\tfrac{7}{2}$:
        \;$D=\R\setminus\{-\tfrac{7}{2}\}$.
  \item $2x+8\geq0\iff x\geq-4$: \;$D=[-4,\infty)$, and since $2x+8$ then
        runs through $[0,\infty)$, \;$W=[0,\infty)$.
  \item $5-2x>0\iff x<\tfrac{5}{2}$: \;$D=(-\infty,\tfrac{5}{2})$.
  \item $x>0$ for the logarithm and $\ln x\neq0\iff x\neq1$:
        \;$D=(0,1)\cup(1,\infty)$.
  \item The root is in the denominator, so $x^{2}-4>0\iff|x|>2$:
        \;$D=(-\infty,-2)\cup(2,\infty)$.
  \item Need $\dfrac{x-1}{x+2}\geq0$ with $x\neq-2$. The quotient is
        non-negative outside the interval between the two critical points
        $-2$ and $1$: \;$D=(-\infty,-2)\cup[1,\infty)$. (At $x=1$ the
        quotient is $0$, which is allowed; $x=-2$ is not.)
  \item $x^{2}+1\geq1>0$ always, so $D=\R$; the argument runs through
        $[1,\infty)$, hence \;$W=[0,\infty)$.
  \item $D=\R$. The exponent $-x^{2}$ runs through $(-\infty,0]$, and
        $\mathrm{e}^{u}$ maps that to $(0,1]$: \;$W=(0,1]$ (the bell curve,
        maximum $1$ at $x=0$).
  \item $x^{2}-5x+6=(x-2)(x-3)$: \;$D=\R\setminus\{2,3\}$.
  \item $4-x\geq0\iff x\leq4$, and $x\neq-1$:
        \;$D=(-\infty,-1)\cup(-1,4]$.
  \item A negative non-integer exponent means a root in a denominator:
        $x-3>0$, \;$D=(3,\infty)$.
  \item $D=\R\setminus\{0\}$. The quotient equals $1$ for $x>0$ and $-1$
        for $x<0$: \;$W=\{-1,1\}$, a range consisting of two points.
  \item $2/(x+1)$ attains every value except $0$; negating and adding $3$
        keeps exactly one value out of reach: \;$W=\R\setminus\{3\}$.
  \item $y=\dfrac{x-1}{x+1}\Rightarrow y(x+1)=x-1\Rightarrow
        x(y-1)=-1-y\Rightarrow x=\dfrac{-(1+y)}{y-1}=\dfrac{1+y}{1-y}$,
        which exists for $y\neq1$: \;$W=\R\setminus\{1\}$.
  \item $D=[-9,\infty)$; $\sqrt{x+9}$ runs through $[0,\infty)$, and
        $4-\sqrt{x+9}$ therefore through \;$W=(-\infty,4]$.
  \item $x^{2}-4x+7=(x-2)^{2}+3$. On $[0,3]$ the vertex $x=2$ lies inside,
        giving the minimum $3$; the endpoints give $f(0)=7$ and $f(3)=4$,
        so the maximum is $7$: \;$W=[3,7]$.
  \item Strictly increasing: $f(0)=2$ and
        $f(\ln 2)=\mathrm{e}^{2\ln 2}+1=4+1=5$: \;$W=[2,5]$.
  \item Strictly increasing: $\ln 1=0$ and $\ln \mathrm{e}^{3}=3$:
        \;$W=[0,3]$.
  \item $(f\circ g)(x)=\bigl(\sqrt{x-3}\bigr)^{2}=x-3$, but only where the
        inner function is defined: \;$D=[3,\infty)$. The simplified formula
        $x-3$ suggests $D=\R$; the domain is inherited from the inner
        function and is not restored by cancellation.
  \item $(f\circ g)(x)=\dfrac{1}{\mathrm{e}^{x}-1}$, requiring
        $\mathrm{e}^{x}\neq1\iff x\neq0$: \;$D=\R\setminus\{0\}$.
  \item Price $200-4q\ge0\iff q\le50$, so the domain is $[0,50]$.
        $R(q)=200q-4q^{2}=-4(q-25)^{2}+2500$, a downward parabola with
        vertex at $q=25$; on $[0,50]$ it runs from $0$ up to $2500$ and
        back to $0$: \;$W=[0,2500]$.
  \item $\sqrt{L}$ runs through $[0,10]$ on $[0,100]$, so
        $W=[0,100]$. The average product is
        $Q(L)/L=10/\sqrt{L}$, which needs $L>0$: domain $(0,100]$.
  \item (a) even: $(-x)^{2}=x^{2}$; (b) odd:
        $(-x)|-x|=-x|x|$; (c) neither, and in fact the test cannot even be
        applied: the domain $(0,\infty)$ is not symmetric about $0$;
        (d) odd: numerator odd, denominator even.
  \item $\dfrac{2x_{1}-1}{x_{1}+4}=\dfrac{2x_{2}-1}{x_{2}+4}$
        $\Rightarrow(2x_{1}-1)(x_{2}+4)=(2x_{2}-1)(x_{1}+4)$
        $\Rightarrow 2x_{1}x_{2}+8x_{1}-x_{2}-4
        =2x_{1}x_{2}+8x_{2}-x_{1}-4$
        $\Rightarrow 9x_{1}=9x_{2}\Rightarrow x_{1}=x_{2}$. Injective on
        $\R\setminus\{-4\}$.
  \item $D=\R$. First piece: $-x$ on $(-\infty,0)$ gives $(0,\infty)$.
        Second: $\sqrt{x}$ on $[0,4]$ gives $[0,2]$. Third: the constant
        $2$. The union is \;$W=[0,\infty)$.
\end{solutions}

\chapter{The Standard Functions and Their Graphs}

A first course in mathematics for economists assumes that the graphs of the
usual functions are known by sight: which way a cubic runs off at its ends,
where a rational function has a pole, that $\ln x$ never turns around, that
an isoquant is a hyperbola. That knowledge is not decoration either. Every
optimisation problem later on begins with a picture, every limit is checked
against one, and a wrong sketch produces a wrong answer faster than any
arithmetic slip.

This chapter is the catalogue, and the drill that makes it stick. First the
standard shapes; then the features that can be read straight off a formula
without any calculus: intercepts, roots, symmetry, poles and asymptotes,
and the behaviour at the far ends of the axis. The last block leaves the
world of $y=f(x)$ and sketches curves given by an \emph{equation} in $x$
and $y$: circles, ellipses, hyperbolas, budget lines and indifference
curves.

\section{Core Rules}

\begin{rulebox}{The power functions}
Whole-number powers alternate: even powers are symmetric about the
$y$-axis and non-negative, odd powers are symmetric about the origin and
run from $-\infty$ to $+\infty$. Negative powers have a pole at $0$; the
root $\sqrt{x}$ is the upper half of a sideways parabola.
\end{rulebox}

\begin{gallery}
\begin{tikzpicture}
  \begin{axis}[fnaxis, xmin=-2.6, xmax=2.6, ymin=-2.6, ymax=2.6,
      xtick={-2,2}, ytick={-2,2}, title={$y=x$}]
    \addplot[domain=-2.4:2.4]{x};
  \end{axis}
\end{tikzpicture}\hfill
\begin{tikzpicture}
  \begin{axis}[fnaxis, xmin=-2.6, xmax=2.6, ymin=-1.0, ymax=4.4,
      xtick={-2,2}, ytick={2,4}, title={$y=x^{2}$}]
    \addplot[domain=-2.1:2.1]{x^2};
  \end{axis}
\end{tikzpicture}\hfill
\begin{tikzpicture}
  \begin{axis}[fnaxis, xmin=-2.6, xmax=2.6, ymin=-4.4, ymax=4.4,
      xtick={-2,2}, ytick={-4,4}, title={$y=x^{3}$}]
    \addplot[domain=-1.62:1.62]{x^3};
  \end{axis}
\end{tikzpicture}
\end{gallery}

\begin{gallery}
\begin{tikzpicture}
  \begin{axis}[fnaxis, xmin=-0.6, xmax=4.6, ymin=-0.8, ymax=2.6,
      xtick={2,4}, ytick={1,2}, title={$y=\sqrt{x}$}]
    \addplot[domain=0:4.4]{sqrt(x)};
  \end{axis}
\end{tikzpicture}\hfill
\begin{tikzpicture}
  \begin{axis}[fnaxis, xmin=-3.4, xmax=3.4, ymin=-3.4, ymax=3.4,
      xtick={-2,2}, ytick={-2,2}, title={$y=1/x$}]
    \addplot[domain=-3.2:-0.31]{1/x};
    \addplot[domain=0.31:3.2]{1/x};
  \end{axis}
\end{tikzpicture}\hfill
\begin{tikzpicture}
  \begin{axis}[fnaxis, xmin=-3.4, xmax=3.4, ymin=-0.6, ymax=4.4,
      xtick={-2,2}, ytick={2,4}, title={$y=1/x^{2}$}]
    \addplot[domain=-3.2:-0.48]{1/x^2};
    \addplot[domain=0.48:3.2]{1/x^2};
  \end{axis}
\end{tikzpicture}
\end{gallery}

\begin{rulebox}{Exponential, logarithm, absolute value}
$\mathrm{e}^{x}$ is positive everywhere, strictly increasing, and flattens
onto the asymptote $y=0$ as $x\to-\infty$; it passes through $(0,1)$.
$\ln x$ is its mirror image at the line $y=x$: defined only for $x>0$,
strictly increasing, passing through $(1,0)$, with the $y$-axis as a
vertical asymptote. A general exponential $a^{x}$ rises for $a>1$ and falls
for $0<a<1$, always through $(0,1)$; note $\bigl(\tfrac12\bigr)^{x}
=2^{-x}$, so the decreasing case is the increasing one reflected in the
$y$-axis. The absolute value $|x|$ is a V with its vertex at the origin.
\end{rulebox}

\begin{gallery}
\begin{tikzpicture}
  \begin{axis}[fnaxis, xmin=-3.4, xmax=2.4, ymin=-0.8, ymax=4.4,
      xtick={-2,2}, ytick={1,3}, title={$y=\mathrm{e}^{x}$}]
    \addplot[domain=-3.2:1.45]{exp(x)};
  \end{axis}
\end{tikzpicture}\hfill
\begin{tikzpicture}
  \begin{axis}[fnaxis, xmin=-0.8, xmax=5.4, ymin=-3.4, ymax=2.4,
      xtick={1,3,5}, ytick={-2,1}, title={$y=\ln x$}]
    \addplot[domain=0.045:5.2]{ln(x)};
  \end{axis}
\end{tikzpicture}\hfill
\begin{tikzpicture}
  \begin{axis}[fnaxis, xmin=-3.4, xmax=3.4, ymin=-0.8, ymax=3.4,
      xtick={-2,2}, ytick={1,3}, title={$y=|x|$}]
    \addplot[domain=-3:3]{abs(x)};
  \end{axis}
\end{tikzpicture}
\end{gallery}

\begin{rulebox}{What can be read off a formula without calculus}
\begin{itemize}[leftmargin=1.6em,itemsep=0.3ex]
  \item \textbf{$y$-intercept:} evaluate $f(0)$ (if $0\in D$).
  \item \textbf{Roots} \de{(Nullstellen)}: solve $f(x)=0$. For a product,
        set each factor to zero. A factor $(x-a)^{k}$ with $k$ odd makes
        the graph \emph{cross} the axis at $a$; with $k$ even it only
        \emph{touches} it.
  \item \textbf{Sign:} between consecutive roots and poles the sign cannot
        change. Pick one test point per interval; that fixes the whole
        picture.
  \item \textbf{Symmetry:} $f(-x)=f(x)$ gives mirror symmetry in the
        $y$-axis, $f(-x)=-f(x)$ point symmetry in the origin.
  \item \textbf{End behaviour of a polynomial:} only the leading term
        $a_{n}x^{n}$ matters as $x\to\pm\infty$. Even $n$: both ends go the
        same way (up if $a_{n}>0$). Odd $n$: the ends go opposite ways
        (down--left, up--right if $a_{n}>0$).
\end{itemize}
\end{rulebox}

\begin{rulebox}{Rational functions: poles and asymptotes}
Write $f(x)=\dfrac{p(x)}{q(x)}$ with $p,q$ polynomials and
\textbf{cancel common factors first}.
\begin{itemize}[leftmargin=1.6em,itemsep=0.3ex]
  \item A zero of $q$ that does \emph{not} cancel is a \textbf{pole}
        \de{(Polstelle)}: the vertical line $x=a$ is a
        \textbf{vertical asymptote}, and $|f|\to\infty$ there.
  \item A zero of $q$ that \emph{does} cancel is a \textbf{hole}: the graph
        is that of the reduced function, minus the single point $x=a$.
  \item \textbf{Horizontal asymptote}, by comparing degrees:
        $\deg p<\deg q\Rightarrow y=0$; \;$\deg p=\deg q\Rightarrow
        y=\dfrac{\text{leading coeff. of }p}{\text{leading coeff. of }q}$;
        \;$\deg p>\deg q\Rightarrow$ none.
  \item If $\deg p=\deg q+1$, polynomial division gives
        $f(x)=(\text{linear})+\dfrac{\text{remainder}}{q(x)}$, and the
        linear part is an \textbf{oblique asymptote}
        \de{(schiefe Asymptote)}, because the remainder term dies out.
\end{itemize}
\end{rulebox}

\begin{gallery}
\begin{tikzpicture}
  \begin{axis}[fnaxis, xmin=-2.4, xmax=6.4, ymin=-4.4, ymax=6.4,
      xtick={2,4,6}, ytick={-2,2,4}, title={$y=\dfrac{2x+1}{x-3}$}]
    \addplot[domain=-2.2:2.55]{(2*x+1)/(x-3)};
    \addplot[domain=3.42:6.2]{(2*x+1)/(x-3)};
    \addplot[gray, dashed, thin, forget plot, samples=2, domain=-2.2:6.2]{2};
    \addplot[gray, dashed, thin, forget plot, samples=2] coordinates {(3,-4.2)(3,6.2)};
  \end{axis}
\end{tikzpicture}\hfill
\begin{tikzpicture}
  \begin{axis}[fnaxis, xmin=-4.4, xmax=4.4, ymin=-6.4, ymax=6.4,
      xtick={-2,2}, ytick={-4,4}, title={$y=\dfrac{x^{2}+1}{x}$}]
    \addplot[domain=-4.2:-0.17]{(x^2+1)/x};
    \addplot[domain=0.17:4.2]{(x^2+1)/x};
    \addplot[gray, dashed, thin, forget plot, samples=2, domain=-4.2:4.2]{x};
  \end{axis}
\end{tikzpicture}\hfill
\begin{tikzpicture}
  \begin{axis}[fnaxis, xmin=-3, xmax=42, ymin=-3, ymax=32,
      xtick={10,20,30,40}, ytick={5,15,25},
      title={$AC(q)=\dfrac{100+5q}{q}$}]
    \addplot[domain=3.9:40]{(100+5*x)/x};
    \addplot[gray, dashed, thin, forget plot, samples=2, domain=-2.8:40]{5};
  \end{axis}
\end{tikzpicture}
\gcap{Horizontal asymptote $y=2$ and pole $x=3$ \quad\textbullet\quad
oblique asymptote $y=x$ \quad\textbullet\quad average cost falling towards
the asymptote $AC=5$.}
\end{gallery}

\begin{rulebox}{Curves given by an equation in two variables}
Not every curve is a graph $y=f(x)$. The standard ones:
\begin{itemize}[leftmargin=1.6em,itemsep=0.3ex]
  \item \textbf{Line:} $ax+by=c$. Sketch it from the two intercepts
        $x=c/a$ and $y=c/b$; solved for $y$ the slope is $-a/b$. This is
        the budget line $p_{x}x+p_{y}y=M$, with slope $-p_{x}/p_{y}$.
  \item \textbf{Circle:} $(x-x_{0})^{2}+(y-y_{0})^{2}=r^{2}$, centre
        $(x_{0},y_{0})$, radius $r$. A general
        $x^{2}+y^{2}+Dx+Ey+F=0$ is brought to this form by
        \textbf{completing the square} in $x$ and in $y$.
  \item \textbf{Ellipse:} $\dfrac{x^{2}}{a^{2}}+\dfrac{y^{2}}{b^{2}}=1$,
        crossing the axes at $(\pm a,0)$ and $(0,\pm b)$.
  \item \textbf{Hyperbola:} $xy=c$ with $c>0$ lies in the first and third
        quadrant with both axes as asymptotes. Solved for $y$ it is
        $y=c/x$. Every Cobb--Douglas isoquant $K^{\alpha}L^{\beta}=\bar Q$
        and every indifference curve $xy=\bar U$ has this shape.
\end{itemize}
To sketch, either solve for $y$, remembering that a square root splits the
curve into \textbf{two branches} $y=\pm\sqrt{\cdots}$, or read off centre,
radius and intercepts directly.
\end{rulebox}

\section{Worked Example}

\begin{worked}{Sketch $f(x)=\dfrac{2x+1}{x-3}$ and $g(x)=(x-1)^{2}(x+3)$ from the formulas alone, and say what $h(x)=\dfrac{x^{2}-1}{x-1}$ really is}
\wstep{Step 1. The rational function: find the poles.}
The denominator of $f$ vanishes at $x=3$ and the numerator does not
($2\cdot3+1=7\neq0$), so $x=3$ is a genuine \textbf{vertical asymptote}.

\wstep{Step 2. Find the horizontal asymptote by comparing degrees.}
Numerator and denominator both have degree $1$, so the limit at $\pm\infty$
is the ratio of leading coefficients:
\[
  \lim_{x\to\pm\infty}\frac{2x+1}{x-3}=\frac{2}{1}=2 ,
\]
a horizontal asymptote $y=2$. Equivalently, dividing out gives
$f(x)=2+\dfrac{7}{x-3}$, which shows the whole graph at once: the hyperbola
$7/x$ shifted right $3$ and up $2$.

\wstep{Step 3. Intercepts and signs.}
$f(x)=0$ when $2x+1=0$, i.e.\ $x=-\tfrac12$; and $f(0)=\tfrac{1}{-3}=-\tfrac13$.
The sign changes at $x=-\tfrac12$ (zero) and at $x=3$ (pole), giving
$+,-,+$ on the three intervals.

\wstep{Step 4. Approach the pole from both sides.}
As $x\to3^{+}$, $x-3\to0^{+}$ and the numerator is near $7>0$, so
$f\to+\infty$; as $x\to3^{-}$, $f\to-\infty$. Now the sketch is determined:
two branches, one below-left rising to the asymptote $y=2$ from below, one
above-right falling to it from above.
\wtrap{Test the sign on \emph{each side} of a pole rather than assuming
symmetry. A pole of odd order flips the sign, one of even order does not,
and drawing both branches on the same side of the horizontal asymptote is
the usual consequence of guessing.}

\wstep{Step 5. The polynomial: read the roots with their multiplicities.}
$g(x)=(x-1)^{2}(x+3)$ has a root of multiplicity $2$ at $x=1$ and a simple
root at $x=-3$. Even multiplicity means the graph \textbf{touches} the axis
and turns back; odd multiplicity means it \textbf{crosses}. So the curve
crosses at $-3$ and bounces off the axis at $1$.

\wstep{Step 6. End behaviour is decided by the leading term alone.}
Expanding just the leading term: $x^{2}\cdot x=x^{3}$, degree $3$ with a
positive coefficient, so $g\to-\infty$ as $x\to-\infty$ and $g\to+\infty$ as
$x\to+\infty$. With $g(0)=1\cdot3=3>0$, the shape is now forced: up through
$-3$, over a local maximum, down to touch the axis at $1$, then away
upward.
\wtrap{Degree and leading sign come from multiplying the leading terms of
every factor, counted with multiplicity. $(x-1)^{2}(x+3)$ is cubic, not
quadratic; miscounting the degree reverses one of the two ends.}

\wstep{Step 7. The impostor: $h(x)=\dfrac{x^{2}-1}{x-1}$.}
Here the numerator \emph{also} vanishes at $x=1$. Factor before concluding
anything:
\[
  h(x)=\frac{(x-1)(x+1)}{x-1}=x+1\qquad\text{for }x\neq1 .
\]
So $h$ has no asymptote at all. Its graph is the straight line $y=x+1$ with a
single point punched out at $(1,2)$ --- a \textbf{removable discontinuity},
a hole, not a pole.
\wtrap{A zero denominator does not by itself make an asymptote. Always
factor numerator and denominator first: a shared factor gives a hole, an
unshared one gives a pole. This is the same distinction that decides whether
a limit in Volume~4 exists or diverges.}
\end{worked}

\section{Practice Problems}

For each polynomial in problems 1--8, give the roots (with multiplicity
where it matters), the $y$-intercept, the symmetry, and the end behaviour
as $x\to\pm\infty$.

\begin{problems}
  \item $f(x)=x^{3}-4x$
  \item $f(x)=x^{4}-5x^{2}+4$
  \item $f(x)=-2x^{3}+6x$
  \item $f(x)=(x-1)^{2}(x+3)$
  \item $f(x)=x^{2}-6x+8$
  \item $f(x)=-x^{2}+4x-4$
  \item $f(x)=x^{5}-x$
  \item A polynomial satisfies $f(x)\to-\infty$ both as $x\to+\infty$ and
        as $x\to-\infty$. What does that say about its degree and its
        leading coefficient?
\end{problems}

In problems 9--20, give all vertical asymptotes, holes, and horizontal or
oblique asymptotes.

\begin{problems}[resume]
  \item $f(x)=\dfrac{1}{x-3}$
  \item $f(x)=\dfrac{2x+1}{x-3}$
  \item $f(x)=\dfrac{x^{2}-1}{x-1}$
  \item $f(x)=\dfrac{3x}{x^{2}-4}$
  \item $f(x)=\dfrac{x^{2}+1}{x}$
  \item $f(x)=\dfrac{x^{2}-4}{x^{2}-9}$
  \item $f(x)=\dfrac{x^{2}+3x+2}{x+1}$
  \item $f(x)=\dfrac{x^{2}-x-6}{x-3}$
  \item $f(x)=\dfrac{2x^{2}+3}{x-1}$
  \item $f(x)=\dfrac{5}{x^{2}+1}$
  \item $AC(q)=\dfrac{100+5q}{q}$ \;(average cost, $q>0$)
  \item $f(x)=\dfrac{x+3}{x^{2}+x-6}$
\end{problems}

\begin{problems}[resume]
  \item For $y=2^{x}$: give the $y$-intercept, the asymptote, the
        monotonicity and the range.
  \item Explain why the graph of $y=\bigl(\tfrac12\bigr)^{x}$ is the graph
        of $y=2^{x}$ reflected in the $y$-axis.
  \item For $y=\ln x$: give the $x$-intercept, the asymptote, the domain
        and the range.
  \item Sketch $y=\mathrm{e}^{-x}$ and state its asymptote and range.
  \item Both $\ln x$ and $\log_{10}x$ pass through $(1,0)$ and rise. What
        exactly distinguishes their graphs?
  \item The adjustment path $y=1-\mathrm{e}^{-x}$, $x\ge0$, is used for
        saturation effects. Give $y(0)$, the asymptote, and say in words
        what the curve does.
\end{problems}

\begin{problems}[resume]
  \item Match each graph (A)--(F) below to one of the formulas
        $\;\ln x$, \;$|x|$, \;$\sqrt{x}$, \;$x^{3}-x$,
        \;$\mathrm{e}^{-x}$, \;$1/x^{2}$.
\end{problems}

\begin{gallery}
\begin{tikzpicture}
  \begin{axis}[fnaxis, xmin=-2.6, xmax=2.6, ymin=-2.2, ymax=2.6,
      xtick={-2,2}, ytick={-2,2}, title={(A)}]
    \addplot[domain=-1.85:1.85]{x^3-x};
  \end{axis}
\end{tikzpicture}\hfill
\begin{tikzpicture}
  \begin{axis}[fnaxis, xmin=-0.7, xmax=4.6, ymin=-0.8, ymax=2.6,
      xtick={2,4}, ytick={1,2}, title={(B)}]
    \addplot[domain=0:4.4]{sqrt(x)};
  \end{axis}
\end{tikzpicture}\hfill
\begin{tikzpicture}
  \begin{axis}[fnaxis, xmin=-3.4, xmax=3.4, ymin=-0.6, ymax=4.4,
      xtick={-2,2}, ytick={2,4}, title={(C)}]
    \addplot[domain=-3.2:-0.48]{1/x^2};
    \addplot[domain=0.48:3.2]{1/x^2};
  \end{axis}
\end{tikzpicture}
\end{gallery}

\begin{gallery}
\begin{tikzpicture}
  \begin{axis}[fnaxis, xmin=-2.4, xmax=3.4, ymin=-0.8, ymax=4.4,
      xtick={-2,2}, ytick={1,3}, title={(D)}]
    \addplot[domain=-1.45:3.2]{exp(-x)};
  \end{axis}
\end{tikzpicture}\hfill
\begin{tikzpicture}
  \begin{axis}[fnaxis, xmin=-3.4, xmax=3.4, ymin=-0.8, ymax=3.4,
      xtick={-2,2}, ytick={1,3}, title={(E)}]
    \addplot[domain=-3:3]{abs(x)};
  \end{axis}
\end{tikzpicture}\hfill
\begin{tikzpicture}
  \begin{axis}[fnaxis, xmin=-0.8, xmax=5.4, ymin=-3.4, ymax=2.4,
      xtick={1,3,5}, ytick={-2,1}, title={(F)}]
    \addplot[domain=0.045:5.2]{ln(x)};
  \end{axis}
\end{tikzpicture}
\end{gallery}

\begin{problems}[resume]
  \item Sketch $f(x)=x^{2}-4x+3$: give the vertex, both intercepts and the
        direction of opening.
  \item Sketch $f(x)=\dfrac{x+1}{x-1}$: give both asymptotes and both
        intercepts.
  \item Sketch $f(x)=\sqrt{x+4}$: give the domain, the range and the
        starting point of the curve.
  \item Sketch $f(x)=|x|-2$: give the vertex and the roots.
  \item Sketch $f(x)=x^{3}-3x^{2}$: give the roots with multiplicity and
        the end behaviour.
  \item Classify as even, odd, or neither: (a) $x^{4}+2$, \;(b) $x^{3}+x$,
        \;(c) $1/x^{2}$, \;(d) $\mathrm{e}^{x}$.
  \item Which of these have a horizontal asymptote, and what is it?
        (a) $\dfrac{3x+1}{x-2}$, \;(b) $\dfrac{x^{2}}{x+1}$,
        \;(c) $\mathrm{e}^{-x}$ as $x\to+\infty$, \;(d) $\ln x$.
  \item On which intervals is $f(x)=\dfrac{(x-1)(x+2)}{x-3}$ positive?
  \item Describe the graph of $f(x)=\sqrt{x^{2}-1}$: domain, symmetry,
        range, and the two points where it meets the $x$-axis.
\end{problems}

In problems 37--45, identify the curve, and sketch it.

\begin{problems}[resume]
  \item $x^{2}+y^{2}=16$. Also solve for $y$ and give the two branches with
        their domains.
  \item $(x-2)^{2}+(y+1)^{2}=9$. Give centre and radius.
  \item $\dfrac{x^{2}}{9}+\dfrac{y^{2}}{4}=1$. Give the four axis
        intercepts.
  \item $xy=6$. Solve for $y$; in which quadrants does the curve lie?
  \item $y^{2}=x$. Is $y$ a function of $x$? Is $x$ a function of $y$?
  \item $2x+3y=12$. Give both intercepts and the slope.
  \item $x^{2}+y^{2}-4x+6y+9=0$. Complete the square; give centre and
        radius.
  \item The isoquant $K^{0.5}L^{0.5}=10$. Solve for $L$ and describe the
        shape for $K>0$.
  \item The budget line $4x+2y=40$ and the indifference curve $xy=50$.
        Give the slope of the budget line and $y$ as a function of $x$ on
        the indifference curve.
\end{problems}

\section{Solutions}

\begin{solutions}
  \item $x^{3}-4x=x(x^{2}-4)=x(x-2)(x+2)$: roots $-2,0,2$, all simple, all
        crossings; $f(0)=0$; odd (only odd powers); degree $3$ with
        positive leading coefficient, so $f\to-\infty$ as $x\to-\infty$ and
        $f\to+\infty$ as $x\to+\infty$.
  \item $x^{4}-5x^{2}+4=(x^{2}-1)(x^{2}-4)$: roots $\pm1,\pm2$;
        $f(0)=4$; even; degree $4$, positive leading coefficient, so both
        ends go to $+\infty$ (a W-shape).
  \item $-2x^{3}+6x=-2x(x^{2}-3)$: roots $0,\pm\sqrt3$; $f(0)=0$; odd;
        degree $3$ with \emph{negative} leading coefficient, so
        $f\to+\infty$ as $x\to-\infty$ and $f\to-\infty$ as $x\to+\infty$.
  \item Roots $x=1$ (double, the graph touches and turns back) and $x=-3$
        (simple, a crossing); $f(0)=(-1)^{2}\cdot3=3$; neither even nor
        odd; degree $3$, leading coefficient $+1$: down on the left, up on
        the right.
  \item $(x-2)(x-4)$: roots $2$ and $4$; $f(0)=8$; vertex at
        $x=3$, $f(3)=-1$; neither even nor odd; parabola opening upwards,
        both ends $+\infty$.
  \item $-(x-2)^{2}$: a double root at $x=2$ (the vertex touches the axis);
        $f(0)=-4$; opens downwards, both ends $-\infty$.
  \item $x^{5}-x=x(x^{4}-1)=x(x^{2}-1)(x^{2}+1)$: roots $-1,0,1$;
        $f(0)=0$; odd; degree $5$, positive leading coefficient: down on
        the left, up on the right.
  \item Both ends behave alike, so the degree is \textbf{even}; both go
        \emph{down}, so the leading coefficient is \textbf{negative}.
  \item Vertical asymptote $x=3$; $\deg p=0<1=\deg q$, so the horizontal
        asymptote is $y=0$. No holes.
  \item Vertical asymptote $x=3$; equal degrees, so
        $y=\tfrac{2}{1}=2$ is the horizontal asymptote.
  \item $\dfrac{(x-1)(x+1)}{x-1}=x+1$ for $x\neq1$: no pole but a
        \textbf{hole} at $(1,2)$. The graph is the line $y=x+1$ with that
        one point punched out; no horizontal asymptote.
  \item $x^{2}-4=(x-2)(x+2)$: poles at $x=\pm2$; $\deg p=1<2=\deg q$, so
        $y=0$.
  \item Pole at $x=0$. Degrees differ by one, so divide:
        $\dfrac{x^{2}+1}{x}=x+\dfrac{1}{x}$, giving the oblique asymptote
        $y=x$.
  \item Poles at $x=\pm3$; equal degrees with leading coefficients $1$ and
        $1$, so the horizontal asymptote is $y=1$.
  \item $\dfrac{(x+1)(x+2)}{x+1}=x+2$ for $x\neq-1$: a hole at $(-1,1)$,
        no pole, no horizontal asymptote.
  \item $\dfrac{(x-3)(x+2)}{x-3}=x+2$ for $x\neq3$: a hole at $(3,5)$.
  \item Pole at $x=1$. Division: $2x^{2}+3=(x-1)(2x+2)+5$, so
        $f(x)=2x+2+\dfrac{5}{x-1}$ and the oblique asymptote is $y=2x+2$.
  \item $x^{2}+1$ never vanishes: no poles at all. $\deg p=0<2$, so
        $y=0$; the graph is a bell-shaped hump with maximum $5$ at $x=0$,
        even.
  \item $AC(q)=\dfrac{100}{q}+5$: pole at $q=0$ and horizontal asymptote
        $AC=5$. Average cost falls towards the marginal cost $5$ as the
        fixed cost $100$ is spread over more units.
  \item $x^{2}+x-6=(x+3)(x-2)$, so
        $f(x)=\dfrac{x+3}{(x+3)(x-2)}=\dfrac{1}{x-2}$ for $x\neq-3$:
        a hole at $\bigl(-3,-\tfrac15\bigr)$, a pole at $x=2$, and the
        horizontal asymptote $y=0$.
  \item $y$-intercept $2^{0}=1$; asymptote $y=0$ (approached as
        $x\to-\infty$); strictly increasing; range $(0,\infty)$.
  \item Because $\bigl(\tfrac12\bigr)^{x}=(2^{-1})^{x}=2^{-x}$: the
        argument $x$ has been replaced by $-x$, and replacing $x$ by $-x$
        mirrors any graph in the $y$-axis.
  \item $x$-intercept $(1,0)$; the $y$-axis $x=0$ is a vertical asymptote
        (as $x\to0^{+}$, $\ln x\to-\infty$); $D=(0,\infty)$, $W=\R$.
  \item The decreasing exponential: it falls from $+\infty$ through
        $(0,1)$ towards the asymptote $y=0$ as $x\to+\infty$;
        $W=(0,\infty)$. It is $\mathrm{e}^{x}$ mirrored in the $y$-axis.
  \item Only the vertical scale: $\log_{10}x=\dfrac{\ln x}{\ln 10}$, so the
        common logarithm is the natural one compressed by the constant
        factor $\ln 10\approx2.303$. Same domain, same asymptote, same
        intercept, same shape.
  \item $y(0)=1-1=0$; as $x\to\infty$, $\mathrm{e}^{-x}\to0$, so the
        asymptote is $y=1$. The curve starts at the origin, rises steeply
        at first and then ever more slowly, approaching the saturation
        level $1$ without ever reaching it.
  \item (A) $x^{3}-x$ \;(odd, three roots); \;(B) $\sqrt{x}$;
        \;(C) $1/x^{2}$ \;(even, positive, pole at $0$);
        \;(D) $\mathrm{e}^{-x}$; \;(E) $|x|$; \;(F) $\ln x$.
  \item $x^{2}-4x+3=(x-1)(x-3)=(x-2)^{2}-1$: opens upwards, vertex
        $(2,-1)$, roots $1$ and $3$, $y$-intercept $3$.
  \item $\dfrac{x+1}{x-1}$: pole $x=1$, horizontal asymptote $y=1$ (equal
        degrees); root $x=-1$; $y$-intercept $f(0)=-1$. Two branches, one
        in each half-plane of the pole.
  \item $D=[-4,\infty)$, $W=[0,\infty)$; the curve starts at $(-4,0)$ and
        rises, ever more slowly: the $\sqrt{\;}$ shape moved four units to
        the left.
  \item A V with vertex $(0,-2)$; roots where $|x|=2$, i.e. $x=\pm2$;
        $y$-intercept $-2$.
  \item $x^{3}-3x^{2}=x^{2}(x-3)$: a double root at $x=0$ (touch) and a
        simple root at $x=3$ (crossing); degree $3$, positive leading
        coefficient: down on the left, up on the right.
  \item (a) even; (b) odd; (c) even, since $1/(-x)^{2}=1/x^{2}$;
        (d) neither.
  \item (a) yes, $y=3$ (equal degrees); (b) no: the numerator degree
        exceeds the denominator degree by one, so there is an oblique
        asymptote $y=x-1$ instead; (c) yes, $y=0$; (d) no: $\ln x$ grows
        without bound, however slowly.
  \item Critical points $-2,1,3$. Test one point per interval:
        $x=-3$ gives $\frac{(-4)(-1)}{-6}<0$; $x=0$ gives
        $\frac{(-1)(2)}{-3}>0$; $x=2$ gives $\frac{(1)(4)}{-1}<0$;
        $x=4$ gives $\frac{(3)(6)}{1}>0$. Hence $f>0$ on
        $(-2,1)\cup(3,\infty)$.
  \item $D=(-\infty,-1]\cup[1,\infty)$ (from $x^{2}-1\ge0$); even, so
        symmetric about the $y$-axis; $W=[0,\infty)$; it meets the axis at
        $(-1,0)$ and $(1,0)$ and rises from there, approaching the lines
        $y=\pm x$ from below for large $|x|$.
  \item Circle, centre $(0,0)$, radius $4$. Solving:
        $y=\pm\sqrt{16-x^{2}}$, two branches, each with domain $[-4,4]$;
        the plus sign is the upper semicircle, the minus sign the lower.
  \item Circle with centre $(2,-1)$ and radius $3$.
  \item Ellipse with $a=3$ and $b=2$: it crosses the axes at
        $(\pm3,0)$ and $(0,\pm2)$.
  \item $y=6/x$, a hyperbola with both axes as asymptotes. Since the
        product is positive, $x$ and $y$ share a sign: the curve lies in
        the first and third quadrants.
  \item A parabola opening to the right. $y$ is \emph{not} a function of
        $x$: the line $x=4$ meets the curve at $y=\pm2$. But $x=y^{2}$
        \emph{is} a function of $y$.
  \item Intercepts $(6,0)$ and $(0,4)$; solved for $y$,
        $y=4-\tfrac23 x$, so the slope is $-\tfrac23$.
  \item $(x^{2}-4x+4)+(y^{2}+6y+9)=-9+4+9$, i.e.
        $(x-2)^{2}+(y+3)^{2}=4$: centre $(2,-3)$, radius $2$.
  \item $K^{0.5}L^{0.5}=10\Rightarrow (KL)^{0.5}=10\Rightarrow KL=100
        \Rightarrow L=\dfrac{100}{K}$. A hyperbola in the first quadrant,
        falling and convex, with both axes as asymptotes: the familiar
        isoquant, along which more capital buys ever less labour saving.
  \item Budget line: $y=20-2x$, slope $-2$ (the price ratio
        $-p_{x}/p_{y}=-4/2$), intercepts $(10,0)$ and $(0,20)$.
        Indifference curve: $y=50/x$, a hyperbola in the first quadrant.
\end{solutions}

\section{Additional Problems}

\begin{problems}
  \item Asymptotes of $f(x)=\dfrac{4x-1}{2x+6}$.
  \item Asymptotes and holes of $f(x)=\dfrac{x^{2}-9}{x+3}$.
  \item Asymptotes and symmetry of $f(x)=\dfrac{x^{3}}{x^{2}-1}$.
  \item Asymptotes and holes of $f(x)=\dfrac{x^{2}-2x}{x^{2}-4}$.
  \item Asymptotes of $f(x)=\dfrac{6}{x^{2}-x-2}$.
  \item Asymptotes of $f(x)=\dfrac{x-4}{x^{2}+4}$.
  \item Roots, symmetry and end behaviour of $f(x)=x^{4}-4x^{2}$.
  \item Roots, sign and end behaviour of $f(x)=(x+2)(x-1)(x-3)$.
  \item Symmetry and asymptote of $f(x)=\dfrac{x}{x^{2}+1}$.
  \item Where do $y=3^{x}$ and $y=3^{-x}$ intersect, and what asymptote do
        they share?
  \item The graphs of $y=\mathrm{e}^{x}$ and $y=\ln x$ are mirror images.
        In which line, and which points correspond to each other?
  \item Domain, range, symmetry and asymptotes of $y=1/x^{2}$.
  \item Domain, symmetry and shape of $y=\sqrt[3]{x}$.
  \item Vertex and roots of $y=|x-2|$.
  \item For the cost function $C(q)=q^{3}-6q^{2}+15q$, tabulate
        $C(0),C(1),\dots,C(5)$ and describe the shape of the curve.
  \item Tabulate $f(x)=x^{2}-2x-3$ for $x=-2,\dots,4$ and sketch it from
        the table; give vertex and roots.
  \item Identify $9x^{2}+4y^{2}=36$ and give its axis intercepts.
  \item Identify $x^{2}-y^{2}=9$ and solve it for $y$.
  \item What curve is $(x+1)^{2}+(y-2)^{2}=0$?
  \item Identify $x^{2}+y^{2}+2x-8=0$: centre and radius.
  \item Compare $y=\sqrt{9-x^{2}}$ with $x^{2}+y^{2}=9$: which is a
        function, and what are its domain and range?
  \item The isocost line $5K+2L=100$: give both intercepts and $L$ as a
        function of $K$.
  \item For the indifference curve $xy=36$, tabulate $y$ at
        $x=1,2,3,4,6,9,12$ and describe the shape.
  \item The demand curve $p=100/q$ has the property that revenue $pq$ is
        constant. Give that constant, the asymptotes, and the economic
        reading of each.
  \item A curve passes through the origin, rises throughout, is concave,
        and flattens onto the horizontal line $y=4$. Which of
        $y=4-\mathrm{e}^{-x}$, \;$y=4(1-\mathrm{e}^{-x})$,
        \;$y=4\mathrm{e}^{-x}$ fits, and why do the other two fail?
\end{problems}

\section{Additional Solutions}

\begin{solutions}
  \item Pole where $2x+6=0$, i.e. $x=-3$; equal degrees, so the horizontal
        asymptote is $y=\tfrac42=2$.
  \item $\dfrac{(x-3)(x+3)}{x+3}=x-3$ for $x\neq-3$: a hole at
        $(-3,-6)$, no pole, no horizontal asymptote (the graph is a line).
  \item Poles at $x=\pm1$. Division: $x^{3}=x(x^{2}-1)+x$, so
        $f(x)=x+\dfrac{x}{x^{2}-1}$ and the oblique asymptote is $y=x$.
        Odd: numerator odd, denominator even.
  \item $\dfrac{x(x-2)}{(x-2)(x+2)}=\dfrac{x}{x+2}$ for $x\neq2$: a hole at
        $\bigl(2,\tfrac12\bigr)$, a pole at $x=-2$, horizontal asymptote
        $y=1$.
  \item $x^{2}-x-2=(x-2)(x+1)$: poles at $x=2$ and $x=-1$; horizontal
        asymptote $y=0$.
  \item $x^{2}+4>0$ always: no poles. Numerator degree smaller, so $y=0$.
  \item $x^{4}-4x^{2}=x^{2}(x^{2}-4)=x^{2}(x-2)(x+2)$: double root at $0$
        (touch), simple roots at $\pm2$; even; degree $4$ with positive
        leading coefficient, so both ends rise to $+\infty$.
  \item Roots $-2,1,3$, all simple. Degree $3$, positive leading
        coefficient: down on the left, up on the right. Signs alternate
        between the roots: negative on $(-\infty,-2)$, positive on
        $(-2,1)$, negative on $(1,3)$, positive on $(3,\infty)$.
  \item Odd (odd numerator over even denominator); no poles since
        $x^{2}+1>0$; horizontal asymptote $y=0$ at both ends. The graph is
        a hump above the axis on the right and its mirror image below on
        the left.
  \item $3^{x}=3^{-x}\iff x=-x\iff x=0$: they meet at $(0,1)$. Both have
        the horizontal asymptote $y=0$, the first as $x\to-\infty$, the
        second as $x\to+\infty$; each is the other reflected in the
        $y$-axis.
  \item In the line $y=x$. The point $(a,\mathrm{e}^{a})$ on the
        exponential corresponds to $(\mathrm{e}^{a},a)$ on the logarithm;
        in particular $(0,1)\leftrightarrow(1,0)$ and
        $(1,\mathrm{e})\leftrightarrow(\mathrm{e},1)$. The two functions
        are inverse to each other.
  \item $D=\R\setminus\{0\}$, $W=(0,\infty)$; even, hence symmetric about
        the $y$-axis; vertical asymptote $x=0$ and horizontal asymptote
        $y=0$. Both branches lie above the axis.
  \item $D=\R$: an odd root accepts negative radicands. It is odd, since
        $\sqrt[3]{-x}=-\sqrt[3]{x}$, strictly increasing, steep near the
        origin and flattening out: the graph of $x^{3}$ mirrored in
        $y=x$.
  \item A V with vertex $(2,0)$; a single (double) root at $x=2$;
        $y$-intercept $|0-2|=2$.
  \item $C(0)=0$, $C(1)=10$, $C(2)=14$, $C(3)=18$, $C(4)=28$, $C(5)=50$.
        Increasing throughout; the increments $10,4,4,10,22$ first fall and
        then rise, so the curve is concave up to about $q=2$ and convex
        afterwards: the S-shaped total cost curve of the textbooks, with
        its inflection at $q=2$.
  \item $f(-2)=5$, $f(-1)=0$, $f(0)=-3$, $f(1)=-4$, $f(2)=-3$, $f(3)=0$,
        $f(4)=5$. The table is symmetric about $x=1$: an upward parabola
        with vertex $(1,-4)$ and roots $-1$ and $3$.
  \item Divide by $36$: $\dfrac{x^{2}}{4}+\dfrac{y^{2}}{9}=1$, an ellipse
        with $a=2$, $b=3$, crossing the axes at $(\pm2,0)$ and
        $(0,\pm3)$.
  \item A hyperbola opening left and right. Solved for $y$:
        $y=\pm\sqrt{x^{2}-9}$, defined for $|x|\ge3$, two branches.
  \item A sum of two squares is zero only if both vanish: the ``curve'' is
        the single point $(-1,2)$.
  \item $(x+1)^{2}-1+y^{2}=8$, so $(x+1)^{2}+y^{2}=9$: a circle with centre
        $(-1,0)$ and radius $3$.
  \item $y=\sqrt{9-x^{2}}$ is a function, the upper semicircle, with
        $D=[-3,3]$ and $W=[0,3]$. The full equation $x^{2}+y^{2}=9$ is not
        a function of $x$: it also contains the lower half.
  \item Intercepts $(20,0)$ on the $K$-axis and $(0,50)$ on the $L$-axis;
        $L=50-2.5K$, a falling line with slope $-2.5$, the ratio of the
        input prices.
  \item $y=36/x$: $36,18,12,9,6,4,3$. A falling, convex hyperbola branch in
        the first quadrant, approaching both axes without touching them:
        halving $x$ doubles $y$.
  \item $pq=100$. The asymptotes are $p=0$ and $q=0$: as the quantity
        grows, the price falls towards zero without reaching it, and no
        finite price clears an arbitrarily small quantity. Revenue stays at
        $100$ everywhere, which is what unit elasticity means.
  \item $y=4(1-\mathrm{e}^{-x})$. It gives $y(0)=0$ and rises to the
        asymptote $4$. \;$y=4-\mathrm{e}^{-x}$ has the right asymptote but
        starts at $y(0)=3$, not at the origin; \;$y=4\mathrm{e}^{-x}$ falls
        from $4$ towards $0$ instead of rising.
\end{solutions}

\chapter{Transforming Functions}

Knowing the catalogue of standard shapes is worth much more once it is
combined with the second half of the skill: recognising that an unfamiliar
formula is a familiar function in disguise. The graph of
$y=-2\mathrm{e}^{-(x-3)}+5$ never has to be plotted point by point. It is
the exponential, mirrored twice, stretched, and moved, and all of this is
read off the formula without any arithmetic.

Three related skills belong here. \textbf{Transforming}: what the constants
in $a\,f\bigl(b(x-c)\bigr)+d$ do to a graph, with the standing trap that
everything \emph{inside} the function acts in the opposite direction and by
the reciprocal factor. \textbf{Composing and decomposing}: building a
function out of simpler ones, and taking one apart again, which is exactly
the reading that makes the chain rule mechanical later. And
\textbf{inverting}: swapping the roles of $x$ and $y$, the operation that
turns an exponential into a logarithm and a demand curve into an inverse
demand curve.

The chapter ends with the reverse question, the one an economist meets
most often: not ``what does this formula look like'', but ``which formula has
the shape I need?'' Given a vertex, an asymptote, a saturation level or two
data points, determine the parameters.

\section{Core Rules}

\begin{rulebox}{The four elementary moves}
Start from a known graph $y=f(x)$ and read the general form
\[
  y=a\,f\bigl(b(x-c)\bigr)+d .
\]
\begin{itemize}[leftmargin=1.6em,itemsep=0.3ex]
  \item \textbf{$d$: vertical shift.} $f(x)+d$ moves the graph \emph{up}
        by $d$ (down if $d<0$). Outside the function, and in the direction
        you expect.
  \item \textbf{$c$: horizontal shift.} $f(x-c)$ moves the graph
        \emph{right} by $c$. Inside the function, and \emph{against}
        the sign: $f(x+2)$ moves two units to the \emph{left}.
  \item \textbf{$a$: vertical scaling.} $a\,f(x)$ stretches the graph
        vertically by the factor $|a|$ ($0<|a|<1$ compresses), and
        \textbf{reflects it in the $x$-axis} if $a<0$. Roots stay put;
        every $y$-value is multiplied by $a$.
  \item \textbf{$b$: horizontal scaling.} $f(bx)$ \emph{compresses} the
        graph horizontally by the factor $|b|$ --- again the reciprocal of
        what the number suggests --- and \textbf{reflects it in the
        $y$-axis} if $b<0$. The $y$-intercept stays put.
\end{itemize}
Order matters: inside operations act on $x$ before $f$ is applied, outside
operations act on the result. Apply them as
\emph{horizontal scaling and shifting first, then $f$, then vertical
scaling and shifting}. And factor before reading: $f(2x-6)=f\bigl(2(x-3)\bigr)$
is a compression by $2$ together with a shift of $3$ to the right, not a
shift of $6$.
\end{rulebox}

\begin{gallery}
\begin{tikzpicture}
  \begin{axis}[fnaxis, xmin=-2.6, xmax=4.6, ymin=-1.2, ymax=4.6,
      xtick={-2,2,4}, ytick={2,4}, title={$y=(x-2)^{2}+1$}]
    \addplot[gray!60, dashed, thin, samples=101, domain=-2.1:2.1]{x^2};
    \addplot[domain=0.15:3.85]{(x-2)^2+1};
  \end{axis}
\end{tikzpicture}\hfill
\begin{tikzpicture}
  \begin{axis}[fnaxis, xmin=-2.6, xmax=2.6, ymin=-1.2, ymax=4.6,
      xtick={-2,2}, ytick={2,4}, title={$y=2x^{2}$}]
    \addplot[gray!60, dashed, thin, samples=101, domain=-2.1:2.1]{x^2};
    \addplot[domain=-1.5:1.5]{2*x^2};
  \end{axis}
\end{tikzpicture}\hfill
\begin{tikzpicture}
  \begin{axis}[fnaxis, xmin=-3.4, xmax=3.4, ymin=-1.2, ymax=4.6,
      xtick={-2,2}, ytick={2,4}, title={$y=\mathrm{e}^{-x}$}]
    \addplot[gray!60, dashed, thin, domain=-3.2:1.5]{exp(x)};
    \addplot[domain=-1.5:3.2]{exp(-x)};
  \end{axis}
\end{tikzpicture}
\gcap{Shift right and up \quad\textbullet\quad vertical stretch
\quad\textbullet\quad reflection in the $y$-axis. The original is dashed.}
\end{gallery}

\begin{rulebox}{Rewrite first, then read: other bases, negative inner factors, higher powers}
The four moves can be read off only once the formula has the shape
$a\,f\bigl(b(x-c)\bigr)+d$ for a base function you know. Four rewritings
get it there.
\begin{itemize}[leftmargin=1.6em,itemsep=0.3ex]
  \item \textbf{Exponentials with any base.} For $q>0$, $q\neq1$, the graph
        of $q^{x}$ passes through $(0,1)$ and has the asymptote $y=0$; it
        rises if $q>1$ and falls if $0<q<1$. Since $q^{-x}=(1/q)^{x}$, a
        negative exponent is a reflection in the $y$-axis, turning a rising
        exponential into a falling one. Collect everything into a single
        base before sketching:
        \[
          \frac{a^{x}}{b^{x}}=\Bigl(\frac{a}{b}\Bigr)^{x},\qquad
          a^{2x}=\bigl(a^{2}\bigr)^{x},\qquad
          \text{e.g. }\ \frac{3^{2x}}{2^{x}}=\Bigl(\frac{9}{2}\Bigr)^{x}.
        \]
        Then $y=c\,q^{x}+d$ has domain $\R$, the horizontal asymptote
        $y=d$, and range $(d,\infty)$ if $c>0$, $(-\infty,d)$ if $c<0$.
  \item \textbf{A negative factor inside.} In $\sqrt{k-mx}$ with $m>0$,
        factor out $-m$, not $m$:
        \[
          \sqrt{k-mx}=\sqrt{-m\Bigl(x-\frac{k}{m}\Bigr)} .
        \]
        The factor $-m$ is a reflection in the $y$-axis together with a
        horizontal compression by $m$; then comes a shift \emph{right} by
        $k/m$. The domain is $x\le k/m$: the curve starts at $x=k/m$ and
        runs to the \emph{left}. The same factoring handles $\ln(k-mx)$,
        with domain $x<k/m$ and the asymptote $x=k/m$.
  \item \textbf{Logarithms with an inner factor.}
        $c\ln(ax+b)+d=c\ln\bigl(a(x+\tfrac{b}{a})\bigr)+d$: domain
        $ax+b>0$, vertical asymptote $x=-b/a$, range $\R$. For $a>0$ there
        is a second reading, $c\ln a+c\ln\bigl(x+\tfrac{b}{a}\bigr)+d$: the
        horizontal compression by $a$ is the same as a vertical shift by
        $c\ln a$.
  \item \textbf{Higher powers in a denominator.} $y=\dfrac{c}{(x-a)^{n}}+b$
        has the asymptotes $x=a$ and $y=b$. For \emph{even} $n$ the
        denominator is positive on both sides of the pole, so both branches
        lie on the \emph{same} side of $y=b$ (above it if $c>0$, below if
        $c<0$), and the range is $(b,\infty)$ or $(-\infty,b)$. For
        \emph{odd} $n$ the denominator changes sign at $x=a$, the branches
        lie on \emph{opposite} sides, as for $1/x$, and the range is
        $\R\setminus\{b\}$.
\end{itemize}
Factor powers too: $(2x-6)^{2}=4(x-3)^{2}$, so
$\dfrac{8}{(2x-6)^{2}}=\dfrac{2}{(x-3)^{2}}$, with the pole at $x=3$.
\end{rulebox}

\begin{gallery}
\begin{tikzpicture}
  \begin{axis}[fnaxis, xmin=-3.4, xmax=3.4, ymin=-1.2, ymax=4.6,
      xtick={-2,2}, ytick={2,4}, title={$y=2^{-x}=\bigl(\tfrac12\bigr)^{x}$}]
    \addplot[gray!60, dashed, thin, domain=-3.2:2.1]{2^x};
    \addplot[domain=-2.1:3.2]{2^(-x)};
  \end{axis}
\end{tikzpicture}\hfill
\begin{tikzpicture}
  \begin{axis}[fnaxis, xmin=-3.4, xmax=4.4, ymin=-1.2, ymax=4.6,
      xtick={-2,2,4}, ytick={2,4}, title={$y=\sqrt{4-2x}$}]
    \addplot[gray!60, dashed, thin, samples=101, domain=0:4.2]{sqrt(x)};
    \addplot[samples=201, domain=-3.2:2]{sqrt(4-2*x)};
  \end{axis}
\end{tikzpicture}\hfill
\begin{tikzpicture}
  \begin{axis}[fnaxis, xmin=-2.4, xmax=3.4, ymin=-3.2, ymax=4.6,
      xtick={-2,2}, ytick={-2,2,4},
      title={$\tfrac{1}{(x-1)^{2}}$ and $\tfrac{1}{(x-1)^{3}}$}]
    \addplot[samples=101, domain=-2.3:0.53]{1/(x-1)^2};
    \addplot[samples=101, domain=1.47:3.3]{1/(x-1)^2};
    \addplot[gray!60, dashed, thin, samples=101, domain=-2.3:0.31]{1/(x-1)^3};
    \addplot[gray!60, dashed, thin, samples=101, domain=1.62:3.3]{1/(x-1)^3};
  \end{axis}
\end{tikzpicture}
\gcap{A negative exponent reflects in the $y$-axis \quad\textbullet\quad
$\sqrt{-2(x-2)}$ starts at $x=2$ and runs left \quad\textbullet\quad even
power (solid): both branches above; odd power (dashed): opposite sides.}
\end{gallery}

\begin{rulebox}{Absolute values: two different operations}
\begin{itemize}[leftmargin=1.6em,itemsep=0.3ex]
  \item $y=|f(x)|$ acts on the \textbf{output}: everything below the
        $x$-axis is folded \emph{up}, everything above stays. Roots become
        points where the graph touches the axis and turns.
  \item $y=f(|x|)$ acts on the \textbf{input}: the part of the graph with
        $x\geq0$ is kept and \emph{mirrored} to the left, and whatever the
        function did for $x<0$ is discarded. The result is always even.
\end{itemize}
The two coincide only by accident. Note also $|ax+b|=|a|\,
\bigl|x+\tfrac{b}{a}\bigr|$, which turns any such expression into a shifted
and stretched V.
\end{rulebox}

\begin{gallery}
\begin{tikzpicture}
  \begin{axis}[fnaxis, xmin=-2.6, xmax=2.6, ymin=-2.6, ymax=2.6,
      xtick={-2,2}, ytick={-2,2}, title={$y=f(x)=x^{3}-3x$}]
    \addplot[domain=-2.1:2.1]{x^3-3*x};
  \end{axis}
\end{tikzpicture}\hfill
\begin{tikzpicture}
  \begin{axis}[fnaxis, xmin=-2.6, xmax=2.6, ymin=-2.6, ymax=2.6,
      xtick={-2,2}, ytick={2}, title={$y=|f(x)|$}]
    \addplot[domain=-2.1:2.1]{abs(x^3-3*x)};
  \end{axis}
\end{tikzpicture}\hfill
\begin{tikzpicture}
  \begin{axis}[fnaxis, xmin=-2.6, xmax=2.6, ymin=-2.6, ymax=2.6,
      xtick={-2,2}, ytick={-2,2}, title={$y=f(|x|)$}]
    \addplot[domain=-2.1:2.1]{abs(x)^3-3*abs(x)};
  \end{axis}
\end{tikzpicture}
\end{gallery}

\begin{rulebox}{Composing and decomposing}
$(f\circ g)(x)=f\bigl(g(x)\bigr)$: the \textbf{inner} function $g$ is
applied first. To \emph{decompose} a formula, ask what you would compute
first with a calculator; that is the inner function:
\[
  \begin{aligned}
    \sqrt{3x+1}\;&:\quad g(x)=3x+1, \quad f(u)=\sqrt{u},\\
    \mathrm{e}^{x^{2}}\;&:\quad g(x)=x^{2}, \quad f(u)=\mathrm{e}^{u},\\
    (\ln x)^{3}\;&:\quad g(x)=\ln x, \quad f(u)=u^{3}.
  \end{aligned}
\]
Decompositions are not unique, but the ``innermost bracket first'' reading
is the one the chain rule wants. Note that all four elementary moves are
themselves compositions with linear functions: $f(b(x-c))$ is $f$ composed
with a linear inner function.
\end{rulebox}

\begin{rulebox}{Inverse functions}
$f^{-1}$ undoes $f$: \;$f^{-1}\bigl(f(x)\bigr)=x$ and
$f\bigl(f^{-1}(y)\bigr)=y$. It exists exactly when $f$ is
\textbf{injective}; if it is not, restrict the domain first (as with
$x^{2}$ on $[0,\infty)$). Three consequences and one recipe:
\begin{itemize}[leftmargin=1.6em,itemsep=0.3ex]
  \item $D(f^{-1})=W(f)$ and $W(f^{-1})=D(f)$: domain and range swap.
  \item The graph of $f^{-1}$ is the graph of $f$ \textbf{mirrored in the
        line $y=x$}; the point $(p,q)$ becomes $(q,p)$.
  \item $f^{-1}$ is \emph{not} $1/f$.
  \item \textbf{Recipe:} write $y=f(x)$, solve for $x$, then rename the
        variables. \;$y=\mathrm{e}^{2x}\Rightarrow \ln y=2x\Rightarrow
        x=\tfrac12\ln y$, so $f^{-1}(x)=\tfrac12\ln x$ with
        $D(f^{-1})=(0,\infty)$.
\end{itemize}
This is the operation behind every \emph{inverse demand} function: the
demand $q(p)$ solved for $p$ gives $p(q)$, the same curve read along the
other axis.
\end{rulebox}

\begin{gallery}
\begin{tikzpicture}
  \begin{axis}[fnaxis, xmin=-3.2, xmax=4.4, ymin=-3.2, ymax=4.4,
      xtick={-2,2,4}, ytick={-2,2,4},
      title={$\mathrm{e}^{x}$ and $\ln x$}]
    \addplot[domain=-3:1.45]{exp(x)};
    \addplot[domain=0.05:4.2]{ln(x)};
    \addplot[gray!60, dashed, thin, samples=2, domain=-3:4.2]{x};
  \end{axis}
\end{tikzpicture}\hfill
\begin{tikzpicture}
  \begin{axis}[fnaxis, xmin=-0.8, xmax=4.4, ymin=-0.8, ymax=4.4,
      xtick={2,4}, ytick={2,4},
      title={$x^{2}$ on $[0,\infty)$ and $\sqrt{x}$}]
    \addplot[domain=0:2.05]{x^2};
    \addplot[domain=0:4.2]{sqrt(x)};
    \addplot[gray!60, dashed, thin, samples=2, domain=0:4.2]{x};
  \end{axis}
\end{tikzpicture}
\gcap{Each pair is symmetric about the dashed line $y=x$.}
\end{gallery}

\begin{rulebox}{Reverse engineering: choosing the parameters}
Given the shape you want, pick the family and fit the constants.
\begin{itemize}[leftmargin=1.6em,itemsep=0.3ex]
  \item \textbf{Prescribed vertex:} use $y=a(x-h)^{2}+k$ with the vertex
        $(h,k)$, and determine $a$ from one further point.
  \item \textbf{Prescribed asymptotes:} $y=\dfrac{k}{x-c}+d$ has the
        vertical asymptote $x=c$ and the horizontal one $y=d$; a rational
        $\dfrac{ax+b}{x-c}$ has $y=a$ as horizontal asymptote.
  \item \textbf{Prescribed saturation:} $y=S\bigl(1-\mathrm{e}^{-kx}\bigr)$
        starts at $0$ and rises to the ceiling $S$; $y=S\mathrm{e}^{-kx}$
        starts at $S$ and decays to $0$.
  \item \textbf{Two data points, exponential family:} $y=A\,q^{x}$ with
        $A=y(0)$, and $q$ from the ratio of the two values. Halving every
        $T$ units means $q=2^{-1/T}$.
  \item \textbf{Economic shifts:} a change in fixed cost shifts $C(q)$
        \emph{vertically}; a per-unit tax $t$ adds $tq$ and therefore
        changes the \emph{slope}, not the intercept; measuring the input in
        thousands instead of units rescales $x$ horizontally.
\end{itemize}
\end{rulebox}

\section{Worked Example}

\begin{worked}{Describe the moves taking $y=f(x)$ to $y=-3f(2x-4)+1$, apply them to $f(x)=\sqrt{x}$, and invert $g(x)=2\ln(x-1)+3$}
\wstep{Step 1. Factor the inside before saying anything about the shift.}
The inside is $2x-4$, and it must be written as
\[
  2x-4=2(x-2)
\]
before it can be read. In this form the horizontal instructions are
unambiguous: \textbf{compress by a factor $2$}, and \textbf{shift right by
$2$} --- not by $4$.
\wtrap{This is the chapter's central trap. $f(2x-4)$ shifts by $2$, because
what matters is the value of $x$ that makes the argument zero. Reading the
shift straight off the $-4$ is the most common error with
transformations.}

\wstep{Step 2. Separate inside from outside, and get the directions right.}
\begin{center}\footnotesize
\begin{tabular}{lll}
  \hline
  where & operation & effect on the graph\\
  \hline
  inside, $2(\cdot)$ & multiply $x$ by $2$ & horizontal compression by $2$ \\
  inside, $(\cdot-2)$ & subtract $2$ from $x$ & shift \emph{right} $2$ \\
  outside, $-3(\cdot)$ & multiply by $-3$ & vertical stretch $3$, flip in the $x$-axis\\
  outside, $+1$ & add $1$ & shift \emph{up} $1$ \\
  \hline
\end{tabular}
\end{center}
Inside changes act \textbf{opposite} to how they read (subtracting moves
right, multiplying compresses); outside changes act \textbf{as} they read.
The reason is that inside operations act on the input before $f$ sees it,
so to keep the output the same you must compensate.

\wstep{Step 3. Apply the moves in the right order.}
Inside: shift and scale, in the order that the factored form dictates ---
compress first, then shift, or equivalently shift the compressed picture by
$2$. Outside: stretch and reflect first, then translate. Applying the
vertical shift before the stretch would multiply the shift by $-3$ as well,
landing at $-3f(\cdot)-3$ instead.

\wstep{Step 4. Apply it to $f(x)=\sqrt{x}$.}
\[
  y=-3\sqrt{2x-4}+1=-3\sqrt{2(x-2)}+1 .
\]
The domain moves with the graph: $\sqrt{\cdot}$ needs $2x-4\ge0$, so
$x\ge2$. The starting point $(0,0)$ of $\sqrt{x}$ goes to $(2,1)$: $x=2$
because the shift is right by~$2$, and $y=1$ because $-3\cdot0+1=1$. From
there the curve falls (the reflection), so its range is $(-\infty,1]$.

\wstep{Step 5. Track one more point, as a check.}
The point $(4,2)$ on $y=\sqrt{x}$: solve $2x-4=4$ for the new abscissa,
giving $x=4$; the new ordinate is $-3\cdot2+1=-5$. So $(4,2)\mapsto(4,-5)$,
and substituting $x=4$ into the final formula gives
$-3\sqrt{4}+1=-5$. \checkmark

\wstep{Step 6. Now the inverse of $g(x)=2\ln(x-1)+3$: undo the operations in reverse order.}
Forward, $g$ does: subtract $1$, take $\ln$, double, add $3$. Backwards:
subtract $3$, halve, exponentiate, add $1$. In symbols, set $y=g(x)$ and
solve for $x$:
\[
  y-3=2\ln(x-1)
  \;\Longrightarrow\;
  \frac{y-3}{2}=\ln(x-1)
  \;\Longrightarrow\;
  \mathrm{e}^{(y-3)/2}=x-1
  \;\Longrightarrow\;
  x=1+\mathrm{e}^{(y-3)/2} .
\]
Renaming the variable,
\[
  \boxed{\;g^{-1}(x)=1+\mathrm{e}^{(x-3)/2}\;}
\]

\wstep{Step 7. State the domain and range, which swap.}
$g$ has domain $(1,\infty)$ and range $\R$. Therefore $g^{-1}$ has domain
$\R$ and range $(1,\infty)$ --- visible in the formula, since
$\mathrm{e}^{(x-3)/2}>0$ always, so $g^{-1}>1$ always.
\wtrap{The domain of the inverse is the \emph{range} of the original. This is
not bookkeeping: it is what makes $g^{-1}$ a function at all. And check
$g^{-1}(g(x))=x$ on one value --- $g(1+\mathrm{e}^{0})=g(2)=2\ln1+3=3$ and
$g^{-1}(3)=1+\mathrm{e}^{0}=2$. \checkmark}
\end{worked}

\section{Practice Problems}

In problems 1--10, describe in words how the graph of $y=f(x)$ has to be
moved, scaled or reflected.

\begin{problems}
  \item $y=f(x-3)$
  \item $y=f(x)+2$
  \item $y=-f(x)$
  \item $y=f(-x)$
  \item $y=2f(x)$
  \item $y=f(2x)$
  \item $y=f\!\left(\dfrac{x}{2}\right)+1$
  \item $y=-3f(x-1)+4$
  \item $y=f(x+2)-5$
  \item $y=\tfrac12 f(-x)$
\end{problems}

In problems 11--18, write down the formula.

\begin{problems}[resume]
  \item $y=x^{2}$ shifted $3$ to the right and $1$ up.
  \item $y=\sqrt{x}$ reflected in the $x$-axis and shifted $4$ up.
  \item $y=\mathrm{e}^{x}$ compressed horizontally by the factor $2$.
  \item $y=\ln x$ shifted $5$ to the left.
  \item $y=1/x$ shifted $2$ to the right and $3$ down.
  \item $y=|x|$ shifted $1$ to the left and stretched vertically by $3$.
  \item $y=x^{3}$ reflected in the $y$-axis and shifted $2$ up.
  \item A downward parabola with vertex $(-1,5)$, twice as steep as
        $y=x^{2}$.
\end{problems}

In problems 19--24, name the base function and the transformations, and
give the features asked for.

\begin{problems}[resume]
  \item $y=2(x+1)^{2}-8$: vertex, roots, range.
  \item $y=-\sqrt{x-2}+3$: domain, starting point, range.
  \item $y=3\mathrm{e}^{x+1}$: asymptote, value at $x=-1$.
  \item $y=\ln(2x)$: give \emph{two} readings of this graph, one
        horizontal and one vertical.
  \item $y=\dfrac{1}{x-4}+2$: both asymptotes.
  \item $y=-|x+3|+5$: vertex, roots, direction of opening.
\end{problems}

\begin{problems}[resume]
  \item For $f(x)=x^{2}-2x-3$, describe the graph of $y=|f(x)|$: which
        part moves, and what happens at the roots?
  \item Sketch $y=|2x-6|$ by first writing it as a stretched, shifted V.
  \item For $f(x)=\sqrt{x}$, describe $y=f(|x|)$: give its domain and its
        symmetry.
  \item With $f(u)=u^{2}$ and $g(x)=x-5$: give $(f\circ g)(x)$ and
        $(g\circ f)(x)$, and say which transformation of $u^{2}$ each one
        is.
  \item With $f(u)=\sqrt{u}$ and $g(x)=x^{2}+9$: give $(f\circ g)(x)$ and
        its domain.
  \item Decompose $h(x)=\sqrt{3x+1}$ into $f\circ g$.
  \item Decompose $h(x)=\mathrm{e}^{x^{2}}$ into $f\circ g$.
  \item Decompose $h(x)=(\ln x)^{3}$ into $f\circ g$.
  \item Decompose $h(x)=\dfrac{1}{x^{2}+1}$ into $f\circ g$.
\end{problems}

In problems 34--39, find the inverse together with its domain and range.

\begin{problems}[resume]
  \item $f(x)=3x-7$
  \item $f(x)=\dfrac{x+2}{x-1}$
  \item $f(x)=x^{2}$ on $D=[0,\infty)$
  \item $f(x)=\ln(x-1)$
  \item $f(x)=\sqrt{x+4}$
  \item $f(x)=\dfrac{2x+1}{x-3}$
  \item Verify for $f(x)=3x-7$ that $f\bigl(f^{-1}(x)\bigr)=x$ and
        $f^{-1}\bigl(f(x)\bigr)=x$.
  \item The demand function is $q(p)=50-2p$. Give the inverse demand
        function $p(q)$, and state on which interval of quantities it is
        economically meaningful.
\end{problems}

\begin{problems}[resume]
  \item Determine $a$ so that the parabola with vertex $(2,-3)$ passes
        through $(0,1)$, and write the function.
  \item An exponential $y=A\,q^{x}$ passes through $(0,5)$ and $(2,45)$.
        Determine $A$ and $q$.
  \item Write a rational function with vertical asymptote $x=2$,
        horizontal asymptote $y=3$, and a root at the origin.
  \item Total cost is $C(q)=5q+100$. Write the new cost function and name
        the transformation when (a) the fixed cost rises by $40$;
        (b) a tax of $2$ per unit is levied.
\end{problems}

\section{Solutions}

\begin{solutions}
  \item Shifted $3$ units to the \emph{right} (inside, against the sign).
  \item Shifted $2$ units up.
  \item Reflected in the $x$-axis: every $y$-value changes sign.
  \item Reflected in the $y$-axis: left and right are swapped.
  \item Stretched vertically by the factor $2$; roots and $x$-intercepts
        stay where they are.
  \item Compressed horizontally by the factor $2$: what happened at $x=4$
        now happens at $x=2$. The $y$-intercept is unchanged.
  \item Stretched horizontally by the factor $2$ (the inner factor
        $\tfrac12$ is a compression by $\tfrac12$, i.e. a stretch), then
        moved $1$ up.
  \item Right $1$, stretched vertically by $3$, reflected in the $x$-axis,
        then up $4$ --- inside first, outside afterwards.
  \item Left $2$ and down $5$.
  \item Reflected in the $y$-axis and compressed vertically to half its
        height.
  \item $y=(x-3)^{2}+1$.
  \item $y=-\sqrt{x}+4$.
  \item $y=\mathrm{e}^{2x}$ (a factor $2$ \emph{inside} compresses).
  \item $y=\ln(x+5)$, with the vertical asymptote moved to $x=-5$.
  \item $y=\dfrac{1}{x-2}-3$, asymptotes $x=2$ and $y=-3$.
  \item $y=3|x+1|$.
  \item $y=(-x)^{3}+2=-x^{3}+2$.
  \item $y=-2(x+1)^{2}+5$.
  \item Base $x^{2}$: left $1$, stretched by $2$, down $8$. Vertex
        $(-1,-8)$; roots from $2(x+1)^{2}=8\iff(x+1)^{2}=4\iff x=1$ or
        $x=-3$; range $[-8,\infty)$.
  \item Base $\sqrt{x}$: right $2$, reflected in the $x$-axis, up $3$.
        $D=[2,\infty)$, the curve starts at $(2,3)$ and falls,
        $W=(-\infty,3]$.
  \item Base $\mathrm{e}^{x}$: left $1$, stretched by $3$. Asymptote
        $y=0$; $y(-1)=3\mathrm{e}^{0}=3$. (Equivalently
        $3\mathrm{e}\cdot\mathrm{e}^{x}$: a horizontal shift of an
        exponential is also a vertical stretch.)
  \item Horizontally: $\ln x$ compressed by the factor $2$. Vertically:
        since $\ln(2x)=\ln 2+\ln x$, the same graph is $\ln x$ shifted up
        by $\ln2\approx0.69$. Both readings are correct --- a peculiarity
        of the logarithm.
  \item Base $1/x$: right $4$, up $2$. Vertical asymptote $x=4$,
        horizontal asymptote $y=2$.
  \item Base $|x|$: left $3$, reflected (so it opens downwards), up $5$.
        Vertex $(-3,5)$; roots where $|x+3|=5$, i.e. $x=2$ and $x=-8$.
  \item $f(x)=(x-3)(x+1)$ is negative exactly on $(-1,3)$, and that arc is
        folded up above the axis; outside it the graph is unchanged. At the
        roots $x=-1$ and $x=3$ the curve now touches the axis and turns
        instead of crossing it, and the old minimum $f(1)=-4$ becomes a
        local maximum of height $4$.
  \item $|2x-6|=2|x-3|$: a V with vertex $(3,0)$, twice as steep as
        $|x|$, i.e. with slopes $-2$ and $+2$; $y$-intercept $6$.
  \item $y=\sqrt{|x|}$: $D=\R$ now, since $|x|\geq0$ always. The right half
        is unchanged and mirrored to the left, so the function is even,
        with range $[0,\infty)$ and a sharp point at the origin.
  \item $(f\circ g)(x)=(x-5)^{2}$: the parabola shifted $5$ to the right.
        $(g\circ f)(x)=x^{2}-5$: the parabola shifted $5$ down. Inner
        shifts move horizontally, outer shifts vertically.
  \item $(f\circ g)(x)=\sqrt{x^{2}+9}$, and since $x^{2}+9\geq9>0$ for
        every $x$, $D=\R$. (Its range is $[3,\infty)$.)
  \item $g(x)=3x+1$, $f(u)=\sqrt{u}$.
  \item $g(x)=x^{2}$, $f(u)=\mathrm{e}^{u}$.
  \item $g(x)=\ln x$, $f(u)=u^{3}$.
  \item $g(x)=x^{2}+1$, $f(u)=\dfrac{1}{u}$.
  \item $y=3x-7\Rightarrow x=\dfrac{y+7}{3}$, so
        $f^{-1}(x)=\dfrac{x+7}{3}$, with $D=W=\R$.
  \item $y=\dfrac{x+2}{x-1}\Rightarrow y(x-1)=x+2\Rightarrow x(y-1)=y+2
        \Rightarrow x=\dfrac{y+2}{y-1}$. Hence
        $f^{-1}(x)=\dfrac{x+2}{x-1}=f(x)$: the function is its own
        inverse, and its graph is symmetric about $y=x$. Domain and range
        are both $\R\setminus\{1\}$.
  \item $f^{-1}(x)=\sqrt{x}$, with $D(f^{-1})=[0,\infty)=W(f)$ and
        $W(f^{-1})=[0,\infty)=D(f)$. The restriction to $[0,\infty)$ was
        needed to make $f$ injective in the first place.
  \item $y=\ln(x-1)\Rightarrow \mathrm{e}^{y}=x-1\Rightarrow
        x=\mathrm{e}^{y}+1$: $f^{-1}(x)=\mathrm{e}^{x}+1$, with $D=\R$ and
        $W=(1,\infty)$ --- exactly the domain $(1,\infty)$ of $f$.
  \item $y=\sqrt{x+4}\Rightarrow y^{2}=x+4\Rightarrow x=y^{2}-4$, and only
        $y\geq0$ occurs: $f^{-1}(x)=x^{2}-4$ with $D=[0,\infty)$ and
        $W=[-4,\infty)$. Without the restriction $D=[0,\infty)$ the formula
        $x^{2}-4$ would not be the inverse.
  \item $y=\dfrac{2x+1}{x-3}\Rightarrow y(x-3)=2x+1\Rightarrow
        x(y-2)=3y+1\Rightarrow x=\dfrac{3y+1}{y-2}$. So
        $f^{-1}(x)=\dfrac{3x+1}{x-2}$, with $D=\R\setminus\{2\}$ (the
        range of $f$) and $W=\R\setminus\{3\}$ (the domain of $f$). The two
        asymptotes have swapped roles.
  \item $f\bigl(f^{-1}(x)\bigr)=3\cdot\dfrac{x+7}{3}-7=x+7-7=x$, and
        $f^{-1}\bigl(f(x)\bigr)=\dfrac{(3x-7)+7}{3}=\dfrac{3x}{3}=x$.
  \item $q=50-2p\Rightarrow 2p=50-q\Rightarrow p(q)=25-\dfrac{q}{2}$. It
        is meaningful for $q\in[0,50]$: at $q=50$ the price has fallen to
        $0$, and larger quantities would require a negative price.
  \item $y=a(x-2)^{2}-3$ with $1=a\cdot4-3$, so $a=1$ and
        $y=(x-2)^{2}-3$.
  \item $A=y(0)=5$. Then $45=5q^{2}\Rightarrow q^{2}=9\Rightarrow q=3$
        (positive base), so $y=5\cdot3^{x}$.
  \item $y=\dfrac{3x}{x-2}$: the denominator gives the pole $x=2$, equal
        degrees give the horizontal asymptote $y=3$, and the numerator
        vanishes at $x=0$, so the graph passes through the origin.
  \item (a) $C(q)=5q+140$: a pure vertical shift by $40$; the curve keeps
        its slope. (b) $C(q)=7q+100$: the tax adds $2q$, so the
        \emph{slope} changes from $5$ to $7$ while the intercept stays.
        A vertical shift and a change of slope are different operations,
        and only the first one is a translation.
\end{solutions}

\section{Additional Problems}

\begin{problems}
  \item $y=(x+4)^{2}-2$: base function, transformations, vertex, range.
  \item $y=-2\sqrt{x}+1$: transformations, domain, range.
  \item $y=\mathrm{e}^{-(x-1)}$: transformations, asymptote, value at
        $x=1$.
  \item $y=\ln(x+3)$: transformation, vertical asymptote, domain.
  \item $y=\dfrac{3}{x+2}-1$: transformations and both asymptotes.
  \item $y=\tfrac12|x-6|$: transformations and vertex.
  \item Write: $y=1/x$ reflected in the $x$-axis and shifted $3$ up.
  \item Write: $y=\mathrm{e}^{x}$ compressed horizontally by the factor
        $3$ and shifted $1$ down.
  \item Write: $y=\sqrt{x}$ stretched vertically by $5$ and shifted $1$ to
        the right.
  \item Write: a downward parabola with vertex $(3,-4)$, twice as steep as
        $y=x^{2}$.
  \item Describe the graph of $y=|x^{2}-1|$: which arc is folded, and what
        does the graph do at $x=\pm1$?
  \item Describe the graph of $y=|\ln x|$: domain, roots, and behaviour as
        $x\to0^{+}$.
  \item For $f(x)=x^{3}-3x$, give the formula for $f(|x|)$ and state its
        symmetry.
  \item With $f(u)=u^{2}$ and $g(x)=x+3$: give both composites and say why
        they differ.
  \item Decompose $h(x)=\sqrt{1-x^{2}}$ into $f\circ g$ and give its
        domain.
  \item Decompose $h(x)=\mathrm{e}^{\sqrt{x}}$ into $f\circ g$.
  \item Decompose $h(x)=(2x+5)^{10}$ into $f\circ g$.
  \item Inverse of $f(x)=5-2x$.
  \item Inverse of $f(x)=x^{3}+1$.
  \item Inverse of $f(x)=\dfrac{1}{x-2}$, with domain and range.
  \item Inverse of $f(x)=10^{x}$, with domain and range.
  \item Inverse of $f(x)=\ln(3x)$, with domain and range.
  \item Why does $f(x)=x^{2}$ have no inverse on $\R$, and which two
        restrictions repair it?
  \item A quantity is $200$ at $x=0$ and halves every $5$ units. Write it
        both as $A\,q^{x}$ and in the form $A\,\mathrm{e}^{-kx}$.
  \item Cost is $C(q)=3q+40$. Fixed cost rises to $60$ and the unit cost
        falls to $2.5$. Write the new function and name both changes as
        transformations of the graph.
\end{problems}

In problems 26--36, rewrite the formula first (single base, factored
inside, factored power), then give the chain of moves from the base
function, the maximal domain, the range and every asymptote.

\begin{problems}[resume]
  \item $y=3-2\cdot5^{-x}$
  \item $y=1+3\cdot\bigl(\tfrac25\bigr)^{x}$
  \item $y=\dfrac{12^{x}}{3^{x}}-2$
  \item $y=4-\dfrac{2^{2x}}{8^{x}}$
  \item $y=2\sqrt{6-2x}+1$
  \item $y=2-\sqrt{5-x}$
  \item $y=3\ln(2x+6)-1$
  \item $y=1-2\ln(4-x)$
  \item $y=\dfrac{2}{(x-1)^{2}}-3$
  \item $y=1-\dfrac{1}{(x+2)^{3}}$
  \item $y=\dfrac{4}{(2x-6)^{2}}+1$
\end{problems}

\section{Additional Solutions}

\begin{solutions}
  \item Base $x^{2}$, shifted $4$ to the left and $2$ down; vertex
        $(-4,-2)$; range $[-2,\infty)$.
  \item Base $\sqrt{x}$: stretched by $2$, reflected in the $x$-axis, then
        up $1$. $D=[0,\infty)$, and since $-2\sqrt{x}\in(-\infty,0]$,
        $W=(-\infty,1]$.
  \item $\mathrm{e}^{-(x-1)}$: reflect $\mathrm{e}^{x}$ in the $y$-axis,
        then shift $1$ to the right (the inner expression is
        $-(x-1)$). Asymptote $y=0$; $y(1)=\mathrm{e}^{0}=1$.
  \item $\ln x$ shifted $3$ to the left; the asymptote moves with it to
        $x=-3$; $D=(-3,\infty)$.
  \item Base $1/x$: stretched by $3$, left $2$, down $1$. Vertical
        asymptote $x=-2$, horizontal asymptote $y=-1$.
  \item Base $|x|$: right $6$ and compressed vertically to half the height;
        vertex $(6,0)$, slopes $\pm\tfrac12$.
  \item $y=-\dfrac{1}{x}+3$.
  \item $y=\mathrm{e}^{3x}-1$.
  \item $y=5\sqrt{x-1}$.
  \item $y=-2(x-3)^{2}-4$.
  \item The arc between the roots, where $x^{2}-1<0$, i.e. over
        $(-1,1)$, is folded up; it becomes a hump of height $1$ at $x=0$.
        At $x=\pm1$ the graph touches the axis and turns, forming two
        corners.
  \item $D=(0,\infty)$; the single root is $x=1$; on $(0,1)$ the logarithm
        is negative and is folded up, so the graph falls from $+\infty$ to
        $0$ there and rises again afterwards, with a corner at $(1,0)$. As
        $x\to0^{+}$ it goes to $+\infty$.
  \item $f(|x|)=|x|^{3}-3|x|$. It is even: the right half of the original
        cubic is kept and mirrored into the left half-plane, so the old
        left branch is discarded.
  \item $(f\circ g)(x)=(x+3)^{2}$, the parabola shifted $3$ to the left;
        $(g\circ f)(x)=x^{2}+3$, the parabola shifted $3$ up. Composition
        is not commutative: an inner shift moves the graph horizontally, an
        outer one vertically.
  \item $g(x)=1-x^{2}$, $f(u)=\sqrt{u}$; $D=[-1,1]$ from $1-x^{2}\geq0$.
  \item $g(x)=\sqrt{x}$, $f(u)=\mathrm{e}^{u}$; $D=[0,\infty)$.
  \item $g(x)=2x+5$, $f(u)=u^{10}$.
  \item $y=5-2x\Rightarrow x=\dfrac{5-y}{2}$, so
        $f^{-1}(x)=\dfrac{5-x}{2}$.
  \item $y=x^{3}+1\Rightarrow x=\sqrt[3]{y-1}$, so
        $f^{-1}(x)=\sqrt[3]{x-1}$, with $D=W=\R$.
  \item $y=\dfrac{1}{x-2}\Rightarrow x-2=\dfrac1y\Rightarrow
        x=2+\dfrac1y$: $f^{-1}(x)=2+\dfrac{1}{x}$, with
        $D=\R\setminus\{0\}$ and $W=\R\setminus\{2\}$.
  \item $f^{-1}(x)=\log_{10}x$, with $D=(0,\infty)$ and $W=\R$.
  \item $y=\ln(3x)\Rightarrow \mathrm{e}^{y}=3x\Rightarrow
        x=\dfrac{\mathrm{e}^{y}}{3}$: $f^{-1}(x)=\dfrac{\mathrm{e}^{x}}{3}$,
        with $D=\R$ and $W=(0,\infty)$.
  \item It is not injective: $f(-2)=f(2)=4$, so a horizontal line meets the
        graph twice and the inverse assignment would not be unique.
        Restricting to $[0,\infty)$ gives the inverse $\sqrt{x}$;
        restricting to $(-\infty,0]$ gives $-\sqrt{x}$.
  \item Halving every $5$ units means $q=2^{-1/5}$, so
        $y=200\cdot2^{-x/5}$. Since $2^{-x/5}
        =\mathrm{e}^{-(\ln2/5)x}$, the same function is
        $y=200\,\mathrm{e}^{-kx}$ with
        $k=\dfrac{\ln2}{5}\approx0.1386$.
  \item $C(q)=2.5q+60$. Raising the fixed cost from $40$ to $60$ shifts the
        line $20$ units up; lowering the unit cost from $3$ to $2.5$
        flattens it, a change of slope rather than a translation. Only the
        first is a shift of the graph.
  \item Rewrite $5^{-x}=\bigl(\tfrac15\bigr)^{x}$. From $y=5^{x}$:
        reflect in the $y$-axis (the graph now falls), stretch vertically by
        $2$, reflect in the $x$-axis (it now rises again), shift up $3$.
        $D=\R$. Since $2\cdot5^{-x}$ takes every value in $(0,\infty)$,
        $-2\cdot5^{-x}$ takes every value in $(-\infty,0)$ and
        $W=(-\infty,3)$. Horizontal asymptote $y=3$, approached from below
        as $x\to\infty$; as $x\to-\infty$ the graph falls without bound.
        Check: $y(0)=3-2=1$. Root: $5^{-x}=\tfrac32$, so
        $x=-\dfrac{\ln1.5}{\ln5}\approx-0.25$.
  \item The base $\tfrac25$ lies between $0$ and $1$, so $(\tfrac25)^{x}$
        is a \emph{falling} exponential through $(0,1)$. Stretch vertically
        by $3$, shift up $1$. $D=\R$, $W=(1,\infty)$, horizontal asymptote
        $y=1$, approached from above as $x\to\infty$. $y(0)=1+3=4$.
        (Equivalently $3\cdot(\tfrac52)^{-x}+1$: the rising exponential
        $(\tfrac52)^{x}$ reflected in the $y$-axis.)
  \item Rewrite $\dfrac{12^{x}}{3^{x}}=\Bigl(\dfrac{12}{3}\Bigr)^{x}=4^{x}$,
        so $y=4^{x}-2$: the rising exponential $4^{x}$ shifted $2$ down.
        $D=\R$, $W=(-2,\infty)$, horizontal asymptote $y=-2$ as
        $x\to-\infty$. $y(0)=-1$; root $4^{x}=2$, i.e. $x=\tfrac12$.
  \item Rewrite $2^{2x}=4^{x}$, so $\dfrac{2^{2x}}{8^{x}}
        =\Bigl(\dfrac48\Bigr)^{x}=\bigl(\tfrac12\bigr)^{x}=2^{-x}$ and
        $y=4-2^{-x}$. From $2^{x}$: reflect in the $y$-axis, reflect in the
        $x$-axis, shift up $4$. $D=\R$, $W=(-\infty,4)$, horizontal
        asymptote $y=4$ as $x\to\infty$. $y(0)=3$; root $2^{-x}=4=2^{2}$,
        so $x=-2$.
  \item Factor out $-2$: $6-2x=-2(x-3)$, so $y=2\sqrt{-2(x-3)}+1$. From
        $\sqrt{x}$: reflect in the $y$-axis and compress horizontally by
        $2$ (giving $\sqrt{-2x}$), shift right $3$, stretch vertically by
        $2$, shift up $1$. Domain from $6-2x\ge0$: $D=(-\infty,3]$. The
        curve starts at $(3,1)$ and rises to the left without bound, so
        $W=[1,\infty)$; no asymptote. Check: $y(1)=2\sqrt4+1=5$.
  \item Factor out $-1$: $y=-\sqrt{-(x-5)}+2$. From $\sqrt{x}$: reflect in
        the $y$-axis, shift right $5$, reflect in the $x$-axis, shift up
        $2$. $D=(-\infty,5]$, starting point $(5,2)$, from which the curve
        falls to the left: $W=(-\infty,2]$. No asymptote. Root:
        $\sqrt{5-x}=2$, so $x=1$.
  \item Factor the inside: $2x+6=2(x+3)$, so $y=3\ln\bigl(2(x+3)\bigr)-1$.
        From $\ln x$: compress horizontally by $2$, shift left $3$, stretch
        vertically by $3$, shift down $1$. $D=(-3,\infty)$ from $2x+6>0$;
        vertical asymptote $x=-3$; $W=\R$. Second reading:
        $y=3\ln(x+3)+(3\ln2-1)$, i.e. $\ln(x+3)$ stretched by $3$ and
        shifted up by $3\ln2-1\approx1.08$. Root: $\ln(2x+6)=\tfrac13$,
        $x=\tfrac12\bigl(\mathrm{e}^{1/3}-6\bigr)\approx-2.30$.
  \item Factor out $-1$: $y=-2\ln\bigl(-(x-4)\bigr)+1$. From $\ln x$:
        reflect in the $y$-axis, shift right $4$, stretch vertically by
        $2$, reflect in the $x$-axis, shift up $1$. $D=(-\infty,4)$ from
        $4-x>0$; vertical asymptote $x=4$. As $x\to4^{-}$,
        $\ln(4-x)\to-\infty$, so $y\to+\infty$; the two reflections together
        leave an \emph{increasing} curve. $W=\R$. Root: $\ln(4-x)=\tfrac12$,
        $x=4-\sqrt{\mathrm{e}}\approx2.35$.
  \item Base $1/x^{2}$ (even power): stretch vertically by $2$, shift
        right $1$, down $3$. $D=\R\setminus\{1\}$, vertical asymptote $x=1$,
        horizontal asymptote $y=-3$. Since $\dfrac{2}{(x-1)^{2}}>0$ on both
        sides, both branches lie above $y=-3$: $W=(-3,\infty)$. Roots:
        $(x-1)^{2}=\tfrac23$, $x=1\pm\sqrt{2/3}\approx1\pm0.82$.
  \item Base $1/x^{3}$ (odd power): reflect in the $x$-axis, shift left
        $2$, up $1$. $D=\R\setminus\{-2\}$, asymptotes $x=-2$ and $y=1$.
        For $x>-2$ the cube $(x+2)^{3}$ is positive and $y<1$; for $x<-2$
        it is negative and $y>1$: the branches lie on opposite sides, and
        $W=\R\setminus\{1\}$. Root: $(x+2)^{3}=1$, $x=-1$; $y(0)=\tfrac78$.
  \item Factor the power: $(2x-6)^{2}=\bigl(2(x-3)\bigr)^{2}=4(x-3)^{2}$,
        so $y=\dfrac{1}{(x-3)^{2}}+1$. Base $1/x^{2}$: shift right $3$, up
        $1$. $D=\R\setminus\{3\}$, asymptotes $x=3$ and $y=1$, both branches
        above $y=1$: $W=(1,\infty)$. No roots.
\end{solutions}

\chapter{Inequalities in Two Variables: Regions in the Plane}

The algebra volume solved inequalities in one variable, where the answer is
a set of points on a line --- an interval, or a union of a few. With two variables
the same question has a two-dimensional answer: the solution set of
$2x+3y\le12$ is not a list of numbers but a \emph{region} of the plane, and
it is described by drawing it.

The mechanism is short. The corresponding \emph{equation} $2x+3y=12$ is a
curve, and a curve cuts the plane into pieces. On each piece the expression
$2x+3y-12$ cannot change sign without passing through zero, so its sign is
constant there. One test point per piece therefore settles the whole
region --- the same sign-table logic as in one variable, with curves in
place of points.

Two directions are worth drilling, and examinations favour the second. The
forward direction is: given inequalities, shade the region. The reverse is:
given a shaded region, write down the inequalities that describe it. The
reverse is harder and more useful, because it is what modelling actually
requires --- a feasible set is described in words or in a picture long before
anyone writes the constraints. It is also exactly the skill that the
graphical solution of a linear programme depends on.

\section{Core Rules}

\begin{rulebox}{The solution set is a region}
For an expression $F(x,y)$, the set $\{F(x,y)<0\}$ is bounded by the curve
$\{F(x,y)=0\}$. The recipe:
\begin{enumerate}[leftmargin=1.8em,itemsep=0.3ex]
  \item \textbf{Draw the boundary} $F(x,y)=0$. It is one of the curves of
        Chapter~2: a line, circle, ellipse, parabola or hyperbola.
  \item \textbf{Note which pieces it creates.} A line makes two half-planes;
        a circle makes an inside and an outside; a parabola makes an
        ``above'' and a ``below''.
  \item \textbf{Test one point per piece.} Substituting a single convenient
        point --- the origin whenever it is not on the boundary --- decides
        the sign for that whole piece.
  \item \textbf{Mark the boundary.} Solid if the inequality is
        $\le$ or $\ge$ (the boundary belongs to the set), dashed if it is
        $<$ or $>$ (it does not).
\end{enumerate}
\end{rulebox}

\begin{rulebox}{The standard regions}
\begin{itemize}[leftmargin=1.6em,itemsep=0.3ex]
  \item \textbf{Half-plane:} $ax+by\le c$. Draw the line, test the origin.
        With $c>0$ and the origin satisfying it, the region is the side
        containing the origin.
  \item \textbf{Disc and its complement:} $x^{2}+y^{2}\le r^{2}$ is the
        closed disc; $\ge$ gives everything outside. Shifted:
        $(x-x_{0})^{2}+(y-y_{0})^{2}\le r^{2}$.
  \item \textbf{Ellipse interior:} $\dfrac{x^{2}}{a^{2}}
        +\dfrac{y^{2}}{b^{2}}\le1$.
  \item \textbf{Above or below a parabola:} $y\ge x^{2}$ is the region
        \emph{inside} the parabola's opening; $y\le x^{2}$ is everything
        outside it. Rewriting an inequality as $y\ge\cdots$ or
        $y\le\cdots$ makes the side unambiguous: ``$y$ at least this'' means
        upwards.
  \item \textbf{Between hyperbola branches:} $xy\ge c$ with $c>0$ is the
        pair of regions beyond the two branches --- and on the positive
        quadrant alone it is the convex ``better set'' of Cobb--Douglas.
  \item \textbf{Strip:} $|x|\le a$ is the vertical strip $-a\le x\le a$;
        $|y|\le b$ a horizontal one.
\end{itemize}
\end{rulebox}

\begin{gallery}
\begin{tikzpicture}
  \begin{axis}[geoaxis, xmin=-1.4, xmax=7.4, ymin=-1.4, ymax=5.4,
      xtick={2,4,6}, ytick={2,4}, title={$2x+3y\le12$, $x,y\ge0$}]
    \addplot[draw=none, name path=lin, domain=0:6, samples=2]{(12-2*x)/3};
    \addplot[draw=none, name path=ax, domain=0:6, samples=2]{0};
    \region{lin}{ax}{}
    \addplot[domain=-1.2:7.2, samples=2]{(12-2*x)/3};
  \end{axis}
\end{tikzpicture}\hfill
\begin{tikzpicture}
  \begin{axis}[geoaxis, xmin=-3.4, xmax=3.4, ymin=-3.4, ymax=3.4,
      xtick={-2,2}, ytick={-2,2}, title={$x^{2}+y^{2}\le4$}]
    \addplot[draw=none, name path=up, domain=-2:2, samples=201]{sqrt(4-x^2)};
    \addplot[draw=none, name path=lo, domain=-2:2, samples=201]{-sqrt(4-x^2)};
    \region{up}{lo}{}
    \addplot[domain=0:360, samples=181]({2*cos(x)},{2*sin(x)});
  \end{axis}
\end{tikzpicture}\hfill
\begin{tikzpicture}
  \begin{axis}[geoaxis, xmin=-2.6, xmax=2.6, ymin=-0.8, ymax=4.6,
      xtick={-2,2}, ytick={2,4}, title={$x^{2}\le y\le 4$}]
    \addplot[draw=none, name path=par, domain=-2:2, samples=201]{x^2};
    \addplot[draw=none, name path=top, domain=-2:2, samples=2]{4};
    \region{par}{top}{}
    \addplot[domain=-2.15:2.15, samples=201]{x^2};
    \addplot[domain=-2.5:2.5, samples=2]{4};
  \end{axis}
\end{tikzpicture}
\gcap{Three regions: an intersection of three half-planes (a triangle), a
disc, and the region caught between a parabola and a line.}
\end{gallery}

\begin{rulebox}{Systems: intersect}
Several inequalities holding \emph{simultaneously} give the
\textbf{intersection} of their regions --- shade each one and keep only the
overlap. Two consequences:
\begin{itemize}[leftmargin=1.6em,itemsep=0.3ex]
  \item An intersection of half-planes is a \textbf{convex polygon}
        (possibly unbounded). This is precisely the feasible region of a
        linear programme, and the reason the optimum of an LP sits at a
        corner.
  \item An intersection of convex regions is always convex, so discs,
        ellipse interiors and half-planes may be combined freely and the
        result stays convex. A \emph{union}, by contrast, usually is not,
        and neither is the outside of a disc.
\end{itemize}
\end{rulebox}

\begin{rulebox}{The reverse direction: from picture to inequalities}
Given a shaded region, work boundary by boundary.
\begin{enumerate}[leftmargin=1.8em,itemsep=0.3ex]
  \item \textbf{Identify every boundary curve} and write its equation. A
        circle needs its centre and radius; a line two points or an
        intercept and a slope; a parabola its vertex and one further
        point.
  \item \textbf{Fix the direction of each} with a point that is clearly
        inside the shaded region: substitute it and keep the inequality
        sign that comes out true.
  \item \textbf{Join the conditions with ``and''}, one per boundary. A region
        bounded by three curves needs three inequalities.
  \item \textbf{Check a point just outside} each boundary: it must fail
        exactly one condition. If it fails none, a condition is missing.
\end{enumerate}
Read the boundary style too: a dashed boundary means a strict inequality.
\end{rulebox}

\begin{rulebox}{Which of these regions are convex}
One rule settles almost every case. For a function $f$ of one variable:
\begin{itemize}[leftmargin=1.6em,itemsep=0.3ex]
  \item $f$ \textbf{convex} $\Longrightarrow$ the region \emph{above} it,
        $\{y\ge f(x)\}$ (its \de{Epigraph}), is a convex set. The region
        \emph{below} it is then not convex.
  \item $f$ \textbf{concave} $\Longrightarrow$ the region \emph{below} it,
        $\{y\le f(x)\}$, is convex. The region above it is then not.
\end{itemize}
The proof is one line: if $y_{i}\ge f(x_{i})$ and $f$ is convex then
$ty_{1}+(1-t)y_{2}\ge tf(x_{1})+(1-t)f(x_{2})\ge f(tx_{1}+(1-t)x_{2})$.
So $y\ge x^{2}$ is convex and $y\le x^{2}$ is not, while $y\le4-x^{2}$
\emph{is} convex, because $4-x^{2}$ is concave. Add to this that
half-planes, discs and ellipse interiors are convex, that any
\textbf{intersection} of convex sets is convex, and that the
\emph{outside} of a disc and any \emph{union} of separated pieces is not,
and every region in this chapter is decided.
\end{rulebox}

\begin{rulebox}{Absolute values in two variables}
Split on the sign, exactly as in one variable, or recognise the standard
shapes:
\[
  |x|+|y|\le c \;\text{(a diamond, or square rotated }45^{\circ}),\qquad
  \max(|x|,|y|)\le c \;\text{(a square)} .
\]
For $|y|\le f(x)$ the region is squeezed between $y=f(x)$ and $y=-f(x)$,
and exists only where $f(x)\ge0$. For $y\ge|x|$ it is the wedge above both
lines $y=x$ and $y=-x$.
\end{rulebox}

\section{Worked Example}

\begin{worked}{Sketch $S=\{(x,y):x^{2}+y^{2}\le16,\ y\ge0,\ x\le2\}$, then decide whether $\{(x,y):y\ge\sqrt{x}\}$ is convex}
\wstep{Step 1. Take the constraints one at a time, and shade each.}
\begin{itemize}[leftmargin=1.6em,itemsep=0.2ex]
  \item $x^{2}+y^{2}\le16$: the closed disc of radius $4$ about the origin
        --- ``$\le$'' is the inside, because the origin satisfies it
        ($0\le16$).
  \item $y\ge0$: the upper half-plane, boundary included.
  \item $x\le2$: everything to the left of the vertical line $x=2$.
\end{itemize}
\wtrap{Decide inside-versus-outside by testing one point, not by staring at
the inequality sign. The origin is free and settles the disc immediately.
For a line not through the origin, any convenient point does.}

\wstep{Step 2. Intersect the three.}
The region is the part of the upper half-disc lying left of $x=2$: a
half-disc with a circular segment removed on the right. Its boundary
consists of three pieces --- the segment of the $x$-axis from $x=-4$ to
$x=2$, the vertical segment $x=2$ from $y=0$ up to $y=\sqrt{16-4}=2\sqrt3$,
and the arc of the circle from $(2,2\sqrt3)$ counter-clockwise round to
$(-4,0)$.

\wstep{Step 3. Find the corner points explicitly.}
Where $x=2$ meets the circle: $4+y^{2}=16$, so $y=2\sqrt3\approx3.46$ (the
negative root is cut off by $y\ge0$). Where $y=0$ meets the circle:
$x=\pm4$, of which $x=4$ is cut off by $x\le2$, leaving $(-4,0)$. And
$(2,0)$ where the two straight boundaries meet. Three corners:
$(-4,0)$, $(2,0)$, $(2,2\sqrt3)$.

\wstep{Step 4. Sketch, and check one interior point.}
\begin{center}
\begin{tikzpicture}
\begin{axis}[geoaxis, width=6.6cm, height=6.6cm,
             xmin=-5.2, xmax=5.2, ymin=-1.2, ymax=5.2,
             xtick={-4,-2,2,4}, ytick={2,3.464},
             yticklabels={$2$,$2\sqrt3$}]
  % the region: under the upper semicircle, above y=0, left of x=2
  \addplot[accent, very thick, domain=-4:2, samples=200,
           fill=regionfill, fill opacity=0.5] {sqrt(16-x^2)} \closedcycle;
  % the discarded part of the disc, for contrast
  \addplot[accent!35, thick, dashed, domain=2:4, samples=80]
           {sqrt(16-x^2)};
  \node[font=\scriptsize] at (axis cs:-1.2,1.4) {$S$};
  \node[font=\tiny, anchor=south west] at (axis cs:2.1,3.6) {$(2,2\sqrt3)$};
\end{axis}
\end{tikzpicture}
\end{center}
Test $(0,1)$: $0+1\le16$ \checkmark, $1\ge0$ \checkmark, $0\le2$
\checkmark. Inside, as drawn.

\wstep{Step 5. Note that $S$ is convex, and why you can say so instantly.}
Each of the three constraint sets is convex --- a disc, a half-plane, a
half-plane --- and \textbf{an intersection of convex sets is convex}. No
picture is needed for that conclusion, and no amount of picture-gazing
substitutes for it.

\wstep{Step 6. The second region: $\{y\ge\sqrt{x}\}$. Do not trust the eye.}
The boundary $y=\sqrt{x}$ is a smooth, gently bending curve, and the region
above it looks perfectly well behaved. Test the definition instead. Take the
two boundary points
\[
  P=(0,0),\qquad Q=(4,2),
\]
both in the set, since $0\ge\sqrt0$ and $2\ge\sqrt4$. Their midpoint is
$M=(2,1)$, and
\[
  \sqrt{2}\approx1.414>1 ,
\]
so $M$ violates $y\ge\sqrt{x}$ and lies \emph{outside}. A segment between
two points of the set leaves the set:
\[
  \boxed{\{(x,y):y\ge\sqrt{x}\}\ \text{is not convex.}}
\]

\wstep{Step 7. The principle behind it.}
$\sqrt{x}$ is \textbf{concave}. The set above the graph of a function --- its
epigraph --- is convex exactly when the function is \emph{convex}; the set
below --- the hypograph --- is convex exactly when the function is
\emph{concave}. So $\{y\ge\sqrt{x}\}$ fails while $\{y\le\sqrt{x}\}$
succeeds, and $\{y\ge x^{2}\}$ succeeds while $\{y\le x^{2}\}$ fails.
\wtrap{``Above'' is not a synonym for ``convex''. Match the direction of the
inequality to the curvature of the function: convex function with $\ge$, or
concave function with $\le$. Getting this backwards is the single most
common mistake in the chapter, and it recurs in Volume~7 where the feasible
set of a constrained problem must be convex for the first-order conditions
to be sufficient.}
\end{worked}

\section{Practice Problems}

Sketch the region described by each inequality in problems 1--12, and say
whether it is convex.

\begin{problems}
  \item $y\ge 2x-1$
  \item $x+y\le 3$
  \item $x^{2}+y^{2}\le 9$
  \item $x^{2}+y^{2}>4$
  \item $y\ge x^{2}$
  \item $y\le 4-x^{2}$
  \item $\dfrac{x^{2}}{9}+\dfrac{y^{2}}{4}\le1$
  \item $xy\ge 1$ on $x>0$
  \item $|x|\le 2$
  \item $|y|\le x$
  \item $(x-2)^{2}+(y+1)^{2}\le 9$
  \item $y>\dfrac1x$ on $x>0$
\end{problems}

In problems 13--20 sketch the region defined by the whole system.

\begin{problems}[resume]
  \item $x\ge0$, \;$y\ge0$, \;$2x+3y\le12$
  \item $x\ge0$, \;$y\ge0$, \;$x+y\le4$, \;$y\le2$
  \item $x^{2}+y^{2}\le16$, \;$y\ge0$
  \item $x^{2}+y^{2}\le16$, \;$y\ge0$, \;$x\le2$
  \item $x^{2}\le y$, \;$y\le x+2$
  \item $1\le x^{2}+y^{2}\le9$
  \item $x^{2}+y^{2}\le9$, \;$x\ge0$, \;$y\ge0$
  \item $y\ge x^{2}$, \;$y\le 4$, \;$x\ge0$
\end{problems}

Problems 21--28 refer to the eight shaded regions below. For each, give a
system of inequalities describing it.

\begin{gallery}
\begin{tikzpicture}
  \begin{axis}[geoaxis, xmin=-5.2, xmax=5.2, ymin=-5.2, ymax=5.2,
      xtick={-4,-2,2,4}, ytick={-4,-2,2,4}, title={(a)}]
    \addplot[draw=none, name path=up, domain=-4:2, samples=201]{sqrt(16-x^2)};
    \addplot[draw=none, name path=ax, domain=-4:2, samples=2]{0};
    \region{up}{ax}{}
    \addplot[domain=0:360, samples=181]({4*cos(x)},{4*sin(x)});
    \addplot[gray!60, dashed, thin, samples=2] coordinates {(2,-5)(2,5)};
  \end{axis}
\end{tikzpicture}\hfill
\begin{tikzpicture}
  \begin{axis}[geoaxis, xmin=-2.8, xmax=2.8, ymin=-0.9, ymax=5.1,
      xtick={-2,2}, ytick={2,4}, title={(b)}]
    \addplot[draw=none, name path=par, domain=-2:2, samples=201]{x^2};
    \addplot[draw=none, name path=top, domain=-2:2, samples=2]{4};
    \region{par}{top}{}
    \addplot[domain=-2.2:2.2, samples=201]{x^2};
    \addplot[domain=-2.6:2.6, samples=2]{4};
  \end{axis}
\end{tikzpicture}\hfill
\begin{tikzpicture}
  \begin{axis}[geoaxis, xmin=-3.9, xmax=3.9, ymin=-2.9, ymax=2.9,
      xtick={-2,2}, ytick={-2,2}, title={(c)}]
    \addplot[draw=none, name path=ell, domain=0:3, samples=201]{2*sqrt(1-x^2/9)};
    \addplot[draw=none, name path=ax, domain=0:3, samples=2]{0};
    \region{ell}{ax}{}
    \addplot[domain=0:360, samples=181]({3*cos(x)},{2*sin(x)});
  \end{axis}
\end{tikzpicture}
\end{gallery}

\begin{gallery}
\begin{tikzpicture}
  \begin{axis}[geoaxis, xmin=-2.4, xmax=3.4, ymin=-0.9, ymax=4.9,
      xtick={-2,2}, ytick={2,4}, title={(d)}]
    \addplot[draw=none, name path=par, domain=-1:2, samples=201]{x^2};
    \addplot[draw=none, name path=lin, domain=-1:2, samples=2]{x+2};
    \region{par}{lin}{}
    \addplot[domain=-2.1:2.2, samples=201]{x^2};
    \addplot[domain=-2.2:3.2, samples=2]{x+2};
  \end{axis}
\end{tikzpicture}\hfill
\begin{tikzpicture}
  \begin{axis}[geoaxis, xmin=-3.9, xmax=3.9, ymin=-3.9, ymax=3.9,
      xtick={-2,2}, ytick={-2,2}, title={(e)}]
    \addplot[draw=none, name path=o1, domain=-3:3, samples=201]{sqrt(9-x^2)};
    \addplot[draw=none, name path=o2, domain=-3:3, samples=201]{-sqrt(9-x^2)};
    \region{o1}{o2}{}
    \addplot[draw=none, name path=i1, domain=-1:1, samples=201]{sqrt(1-x^2)};
    \addplot[draw=none, name path=i2, domain=-1:1, samples=201]{-sqrt(1-x^2)};
    \addplot[white, forget plot, sharp plot] fill between[of=i1 and i2];
    \addplot[domain=0:360, samples=181]({3*cos(x)},{3*sin(x)});
    \addplot[domain=0:360, samples=181]({cos(x)},{sin(x)});
  \end{axis}
\end{tikzpicture}\hfill
\begin{tikzpicture}
  \begin{axis}[geoaxis, xmin=-1.4, xmax=7.4, ymin=-1.4, ymax=5.4,
      xtick={2,4,6}, ytick={2,4}, title={(f)}]
    \addplot[draw=none, name path=lin, domain=0:6, samples=2]{(12-2*x)/3};
    \addplot[draw=none, name path=ax, domain=0:6, samples=2]{0};
    \region{lin}{ax}{}
    \addplot[domain=-1.2:7.2, samples=2]{(12-2*x)/3};
  \end{axis}
\end{tikzpicture}
\end{gallery}

\begin{gallery}
\begin{tikzpicture}
  \begin{axis}[geoaxis, xmin=-1.2, xmax=5.4, ymin=-2.9, ymax=2.9,
      xtick={2,4}, ytick={-2,2}, title={(g)}]
    \addplot[draw=none, name path=up, domain=0:4, samples=201]{sqrt(x)};
    \addplot[draw=none, name path=lo, domain=0:4, samples=201]{-sqrt(x)};
    \region{up}{lo}{}
    \addplot[domain=0:4.6, samples=201]{sqrt(x)};
    \addplot[domain=0:4.6, samples=201]{-sqrt(x)};
    \addplot[gray!60, dashed, thin, samples=2] coordinates {(4,-2.8)(4,2.8)};
  \end{axis}
\end{tikzpicture}\hfill
\begin{tikzpicture}
  \begin{axis}[geoaxis, xmin=-3.4, xmax=3.4, ymin=-3.4, ymax=3.4,
      xtick={-2,2}, ytick={-2,2}, title={(h)}]
    \addplot[draw=none, name path=up, domain=-2:2, samples=101]{2-abs(x)};
    \addplot[draw=none, name path=lo, domain=-2:2, samples=101]{abs(x)-2};
    \region{up}{lo}{}
    \addplot[domain=-2:2, samples=101]{2-abs(x)};
    \addplot[domain=-2:2, samples=101]{abs(x)-2};
  \end{axis}
\end{tikzpicture}
\gcap{Regions (a)--(h) for problems 21--28. A dashed boundary means the
inequality there is strict.}
\end{gallery}

\begin{problems}[resume]
  \item Region (a).
  \item Region (b).
  \item Region (c).
  \item Region (d).
  \item Region (e).
  \item Region (f).
  \item Region (g).
  \item Region (h).
\end{problems}

\begin{problems}[resume]
  \item Which of the eight regions (a)--(h) are convex?
  \item Sketch $|x|+|y|\le 3$ and give its four vertices.
  \item Sketch $y\ge|x|$ and describe it in words.
  \item Sketch $|y|\le 4-x^{2}$ and give the $x$-values for which the
        region is non-empty.
  \item Write inequalities for the closed square with vertices
        $(\pm2,\pm2)$, using absolute values and then without them.
  \item A test point satisfies every inequality of a system but lies
        outside the intended region. What has gone wrong?
  \item Explain why the outside of a disc is never convex.
  \item Give a system of two inequalities whose solution set is empty, both
        of them satisfiable on their own.
  \item Show that the intersection of any two convex regions is convex.
  \item Is the union of the two half-planes $y\ge1$ and $y\le-1$ convex?
        Justify.
  \item Write inequalities for the set of points strictly inside the circle
        of radius $2$ about the origin and strictly above the line $y=x$.
  \item For which values of $c$ is the region $\{x^{2}+y^{2}\le c\}$ empty,
        a single point, and a disc?
\end{problems}

\begin{problems}[resume]
  \item \textbf{Budget set.} A consumer has income $60$ and faces prices
        $p_{x}=3$, $p_{y}=4$. Write the budget set as a system, sketch it,
        and give its three vertices.
  \item For the same consumer, write the set of bundles that are
        \emph{affordable and} give at least utility $12$ under
        $U(x,y)=xy$. Sketch it and say whether it is convex.
  \item \textbf{Production.} A firm can use at most $40$ hours of labour
        and $30$ units of capital; product A needs $2$ hours and $1$ unit,
        product B needs $1$ hour and $2$ units. Write the feasible set and
        sketch it.
  \item For the feasible set of the previous problem, give all its
        vertices, and say why a linear objective attains its maximum at one
        of them.
  \item \textbf{Tax wedge.} The set $\{(q,p): p\le 100-2q,\ p\ge 20+q,\
        q\ge0\}$ describes the quantities at which demand price exceeds
        supply price. Sketch it and find the largest $q$ in it.
\end{problems}

\section{Solutions}

\begin{solutions}
  \item The closed half-plane on or above the line $y=2x-1$. Testing
        $(0,0)$: $0\ge-1$ true, so the origin is in it. Convex.
  \item The closed half-plane on or below $x+y=3$ (through $(3,0)$ and
        $(0,3)$); the origin satisfies it. Convex.
  \item The closed disc of radius $3$ about the origin. Convex.
  \item Everything strictly outside the circle of radius $2$, boundary
        dashed and excluded. \textbf{Not} convex.
  \item The region on or above the parabola $y=x^{2}$ --- inside its
        opening. Convex.
  \item The region on or below the downward parabola $y=4-x^{2}$, which
        includes everything far below it. Since $4-x^{2}$ is
        \emph{concave}, the region below it is \textbf{convex}.
  \item The closed elliptical disc with semi-axes $3$ and $2$. Convex.
  \item The region on or beyond the hyperbola branch $y=1/x$ in the first
        quadrant. Convex (this is the quasiconcavity of $xy$).
  \item The closed vertical strip $-2\le x\le2$. Convex, being the
        intersection of two half-planes.
  \item Requires $x\ge0$ and $-x\le y\le x$: the wedge opening to the right
        between the lines $y=\pm x$. Convex.
  \item The closed disc of radius $3$ centred at $(2,-1)$. Convex.
  \item The region strictly above the hyperbola branch $y=1/x$ for $x>0$,
        boundary dashed. Convex.
  \item The triangle with vertices $(0,0)$, $(6,0)$, $(0,4)$ --- the
        intersection of three half-planes.
  \item The quadrilateral with vertices $(0,0)$, $(4,0)$, $(2,2)$,
        $(0,2)$: the triangle $x+y\le4$ in the first quadrant, with its top
        cut off by $y\le2$.
  \item The upper half-disc of radius $4$: the semicircular region above
        the $x$-axis.
  \item The same half-disc with everything right of $x=2$ removed --- the
        region shown in figure (a).
  \item The lens between the parabola and the line, from $x=-1$ to $x=2$
        (the boundaries meet where $x^{2}=x+2$, i.e. $x=-1,2$). Convex.
  \item The annulus between the circles of radius $1$ and $3$. Not convex.
  \item The quarter-disc of radius $3$ in the first quadrant. Convex.
  \item The region above the parabola, below $y=4$, right of the $y$-axis:
        the right half of figure (b).
  \item $x^{2}+y^{2}\le16$, \;$y\ge0$, \;$x<2$. The circle and the
        $x$-axis are solid, so those two are non-strict; the vertical
        boundary is dashed, so that one is strict.
  \item $y\ge x^{2}$ and $y\le4$. (Equivalently
        $x^{2}\le y\le 4$, which also forces $-2\le x\le2$ without needing
        a separate condition.)
  \item $\dfrac{x^{2}}{9}+\dfrac{y^{2}}{4}\le1$, \;$x\ge0$, \;$y\ge0$: the
        quarter of the elliptical disc in the first quadrant.
  \item $y\ge x^{2}$ and $y\le x+2$.
  \item $1\le x^{2}+y^{2}\le9$, that is $x^{2}+y^{2}\ge1$ \emph{and}
        $x^{2}+y^{2}\le9$.
  \item $2x+3y\le12$, \;$x\ge0$, \;$y\ge0$.
  \item $y^{2}\le x$ and $x<4$; equivalently
        $-\sqrt{x}\le y\le\sqrt{x}$ with $x<4$. The parabola opens to the
        right and is solid; the line $x=4$ is dashed.
  \item $|x|+|y|\le2$. Written without absolute values it is the four
        conditions $x+y\le2$, $x-y\le2$, $-x+y\le2$, $-x-y\le2$.
  \item All except (e). Each of (a), (b), (c), (d), (f), (g), (h) is an
        intersection of convex sets: discs and ellipse interiors,
        half-planes, the region above the convex $x^{2}$, the region
        \emph{right} of the parabola $x=y^{2}$ (the same rule with the axes
        exchanged, since $y\mapsto y^{2}$ is convex), and the diamond
        $|x|+|y|\le2$, which is an intersection of four half-planes.
        Only (e), the annulus, fails: it is the disc of radius $3$ with a
        hole, and the segment from $(-2,0)$ to $(2,0)$ runs straight
        through the hole.
  \item A diamond with vertices $(3,0)$, $(0,3)$, $(-3,0)$, $(0,-3)$: the
        square of side $3\sqrt2$ rotated by $45^{\circ}$.
  \item The wedge above both lines $y=x$ and $y=-x$ --- everything at or
        above the V of $y=|x|$, opening upwards with its vertex at the
        origin. Convex.
  \item Squeezed between $y=4-x^{2}$ and $y=-(4-x^{2})$, so it exists only
        where $4-x^{2}\ge0$, i.e. for $-2\le x\le2$. It is the
        lens-shaped region symmetric about both axes.
  \item With absolute values: $\max(|x|,|y|)\le2$, i.e. $|x|\le2$ and
        $|y|\le2$. Without: $-2\le x\le2$ and $-2\le y\le2$, four
        half-planes.
  \item The system describes a \emph{larger} set than intended --- a
        condition is missing. Every boundary of the intended region must
        contribute one inequality; a point outside that satisfies them all
        proves one boundary has been forgotten.
  \item Take two points diametrically opposite each other just outside the
        circle: the segment between them passes through the centre of the
        disc, which is not in the set. So the segment leaves the set and
        convexity fails.
  \item For instance $y\ge x+1$ and $y\le x-1$. Each is a satisfiable
        half-plane, but together they demand $x+1\le y\le x-1$, impossible
        since $x+1>x-1$. The two half-planes are on opposite sides of
        parallel lines.
  \item Let $C_{1},C_{2}$ be convex and take $P,Q\in C_{1}\cap C_{2}$.
        Since $P,Q\in C_{1}$ and $C_{1}$ is convex, the whole segment
        $[P,Q]$ lies in $C_{1}$; the same argument puts it in $C_{2}$.
        Hence $[P,Q]\subseteq C_{1}\cap C_{2}$, which is the definition of
        convexity. (The argument works for any number of sets.)
  \item No. The points $(0,1)$ and $(0,-1)$ both lie in the union, but the
        midpoint $(0,0)$ lies in neither half-plane. A union of disjoint
        convex sets is essentially never convex.
  \item $x^{2}+y^{2}<4$ \;and\; $y>x$. Both boundaries are dashed: the open
        half-disc lying above the diagonal. (It is convex, being an
        intersection of two convex sets.)
  \item Empty for $c<0$; the single point $(0,0)$ for $c=0$; the closed disc
        of radius $\sqrt{c}$ for $c>0$. The three cases are exactly the
        three things a level \emph{set} of $x^{2}+y^{2}$ can do, and they
        are why a sublevel set can be empty without the inequality being
        malformed.
  \item $3x+4y\le60$, \;$x\ge0$, \;$y\ge0$: the triangle with vertices
        $(0,0)$, $(20,0)$ and $(0,15)$.
  \item $3x+4y\le60$, \;$xy\ge12$, \;$x,y>0$: the part of the budget
        triangle lying beyond the indifference curve $y=12/x$. It is the
        intersection of a convex set (the triangle) with another convex set
        (the better set of a quasiconcave utility), hence \textbf{convex}.
  \item With $a$ units of A and $b$ of B: $2a+b\le40$ (labour),
        $a+2b\le30$ (capital), $a\ge0$, $b\ge0$. Four half-planes, so a
        convex polygon.
  \item Vertices: $(0,0)$, $(20,0)$, $(0,15)$, and the intersection of
        $2a+b=40$ with $a+2b=30$. Solving: $b=40-2a$ and
        $a+80-4a=30$, so $a=\tfrac{50}{3}$ and $b=\tfrac{20}{3}$. A linear
        objective has parallel straight level lines; pushing them across a
        convex polygon, the last point of contact is a vertex (or a whole
        edge, if the level lines happen to be parallel to it). Hence
        checking the vertices suffices --- the fundamental theorem of
        linear programming.
  \item The two lines meet where $100-2q=20+q$, i.e. $q=\tfrac{80}{3}$ with
        $p=\tfrac{140}{3}$. The region is the triangle with vertices
        $(0,20)$, $(0,100)$ and
        $\bigl(\tfrac{80}{3},\tfrac{140}{3}\bigr)$, so the largest feasible
        quantity is $q=\tfrac{80}{3}\approx26.7$ --- beyond it no price
        satisfies both buyer and seller, which is exactly the competitive
        equilibrium quantity.
\end{solutions}

\section{Additional Problems}

\begin{problems}
  \item Sketch $3x-2y>6$ and say which side of the line contains the
        origin.
  \item Sketch $y\le -x+5$ together with $x\ge1$.
  \item Sketch $(x+1)^{2}+(y-2)^{2}<4$ and give centre and radius.
  \item Sketch $x^{2}+y^{2}\ge 1$ and $x^{2}+y^{2}\le 4$ together.
  \item Sketch $y\ge x^{2}-4$.
  \item Sketch $y\le 2^{x}$.
  \item Sketch $y\ge\ln x$ and state the domain restriction it forces.
  \item Sketch $xy\le 4$ on the first quadrant.
  \item Sketch $|x-1|\le 2$.
  \item Sketch $|x|+|y|\ge 1$ and decide convexity.
  \item Sketch $y>x^{2}$ and $y<x+6$, and find the $x$-range of the
        region.
  \item Sketch $x^{2}+4y^{2}\le 16$ and $y\ge 0$.
  \item Write inequalities for the closed triangle with vertices $(0,0)$,
        $(4,0)$, $(0,3)$.
  \item Write inequalities for the closed disc of radius $2$ centred at
        $(1,1)$, intersected with the first quadrant.
  \item Write inequalities for the region above $y=x^{2}$ and below
        $y=2x+3$, and give the $x$-range.
  \item Write inequalities for the horizontal strip between $y=-1$ and
        $y=3$, once with an absolute value and once without.
  \item Write inequalities for everything outside the ellipse
        $x^{2}/4+y^{2}=1$ but inside the circle $x^{2}+y^{2}=9$.
  \item Is the region $\{y\ge x^{2}\}\cap\{y\le 4\}$ bounded? Give its
        extreme $x$- and $y$-values.
  \item Give a system of three inequalities whose solution set is a single
        point.
  \item Give a system of three inequalities whose solution set is an
        unbounded convex region with exactly two vertices.
  \item Show that $\{(x,y): y\ge \mathrm{e}^{x}\}$ is convex.
  \item Show that $\{(x,y): y\ge \sqrt{x},\ x\ge0\}$ is \emph{not}
        convex, and say which nearby set \emph{is}.
  \item \textbf{Two constraints, two goods.} Sketch
        $x+y\le 10$, \;$2x+y\le 16$, \;$x,y\ge0$, and give all vertices.
  \item \textbf{Diet problem.} A feed must supply at least $8$ units of
        protein and $12$ of energy. Grain gives $2$ protein and $2$ energy
        per kg; soy gives $1$ protein and $4$ energy per kg. Write the
        feasible set, sketch it, and say why it is unbounded.
  \item \textbf{Minimum-wage wedge.} Sketch
        $\{(L,w): w\ge 10,\ w\le 30-L,\ L\ge0\}$ and interpret the region
        and its largest $L$.
\end{problems}

\section{Additional Solutions}

\begin{solutions}
  \item The open half-plane strictly below the line $y=\tfrac32x-3$.
        Testing $(0,0)$: $0>6$ is false, so the origin is \emph{not} in the
        region; it lies on the other side.
  \item The intersection of the half-plane on or below $y=5-x$ with the
        half-plane $x\ge1$: an unbounded wedge opening down-right, with
        corner $(1,4)$.
  \item The \emph{open} disc of radius $2$ centred at $(-1,2)$; the
        boundary circle is dashed and excluded.
  \item The closed annulus $1\le x^{2}+y^{2}\le4$ between the two circles.
        Not convex.
  \item The region on or above the parabola $y=x^{2}-4$, whose vertex is
        $(0,-4)$ and which crosses the axis at $x=\pm2$. Convex.
  \item Everything on or below the increasing curve $y=2^{x}$, which
        passes through $(0,1)$ and hugs the $x$-axis to the left. Since
        $2^{x}$ is \emph{convex}, the region below it is \textbf{not}
        convex: $(-3,\tfrac18)$ and $(3,8)$ both satisfy the inequality,
        but their midpoint $(0,4.0625)$ does not, since $2^{0}=1<4.0625$.
  \item The region on or above $y=\ln x$; the logarithm's domain confines
        it to the open right half-plane $x>0$. Since $\ln x$ is
        \emph{concave}, the region above it is \textbf{not} convex: the
        boundary points $(0.1,\ln 0.1)$ and $(10,\ln 10)$ have midpoint
        $(5.05,0)$, and $0<\ln 5.05\approx1.62$, so the midpoint is
        outside.
  \item The part of the first quadrant on or below the hyperbola $y=4/x$
        --- the sublevel set $\{xy\le4\}$. \textbf{Not} convex: $(0.5,8)$
        and $(8,0.5)$ both lie on the boundary, but their midpoint
        $(4.25,4.25)$ has $xy\approx18.1>4$. (Equivalently: $4/x$ is convex
        on $x>0$, so the region below it is not.)
  \item $-2\le x-1\le2$, i.e. the closed vertical strip $-1\le x\le3$.
        Convex.
  \item Everything on or outside the diamond $|x|+|y|=1$. Not convex ---
        the segment from $(1,0)$ to $(-1,0)$ passes through the origin,
        where $|x|+|y|=0<1$.
  \item Strictly between the parabola and the line; they meet where
        $x^{2}=x+6$, i.e. $x=-2$ and $x=3$, so the region spans
        $-2<x<3$. Both boundaries dashed.
  \item $\dfrac{x^{2}}{16}+\dfrac{y^{2}}{4}\le1$ with $y\ge0$: the upper
        half of an elliptical disc with semi-axes $4$ and $2$.
  \item $x\ge0$, \;$y\ge0$, \;$3x+4y\le12$ (the line through $(4,0)$ and
        $(0,3)$).
  \item $(x-1)^{2}+(y-1)^{2}\le4$, \;$x\ge0$, \;$y\ge0$.
  \item $y\ge x^{2}$ and $y\le 2x+3$; the boundaries meet where
        $x^{2}=2x+3$, i.e. $x=-1$ and $x=3$, so $-1\le x\le3$.
  \item With an absolute value: $|y-1|\le2$. Without: $y\ge-1$ and
        $y\le3$.
  \item $\dfrac{x^{2}}{4}+y^{2}\ge1$ \;and\; $x^{2}+y^{2}\le9$: the disc of
        radius $3$ with the elliptical disc punched out. Not convex.
  \item Yes, bounded: $y$ ranges over $[0,4]$ and $x$ over $[-2,2]$, since
        $y\ge x^{2}$ together with $y\le4$ forces $x^{2}\le4$. The extreme
        points are $(\pm2,4)$ and $(0,0)$.
  \item For example $x\ge0$, \;$y\ge0$, \;$x+y\le0$: the first two force
        $x,y\ge0$ and the third then forces $x=y=0$, so the solution set is
        exactly $\{(0,0)\}$.
  \item $x\ge0$, \;$y\ge0$, \;$x+y\ge2$. The boundary consists of the
        ray $\{x=0,\ y\ge2\}$, the segment from $(0,2)$ to $(2,0)$, and the
        ray $\{y=0,\ x\ge2\}$, so there are exactly two corners, $(0,2)$
        and $(2,0)$. The region is unbounded (it contains every large
        point of the quadrant) and convex, being an intersection of three
        half-planes.
  \item $\mathrm{e}^{x}$ is a convex function, and the set lying \emph{on
        or above} the graph of a convex function --- its epigraph --- is a
        convex set: if $y_{1}\ge\mathrm{e}^{x_{1}}$ and
        $y_{2}\ge\mathrm{e}^{x_{2}}$ then for $t\in[0,1]$
        \[
          ty_{1}+(1-t)y_{2}\ \ge\ t\mathrm{e}^{x_{1}}+(1-t)\mathrm{e}^{x_{2}}
          \ \ge\ \mathrm{e}^{tx_{1}+(1-t)x_{2}},
        \]
        the last step being exactly convexity of $\mathrm{e}^{x}$.
  \item $\sqrt{x}$ is \emph{concave}, so the region \emph{above} it is
        not convex. Concretely, $(0,0)$ and $(4,2)$ both lie on the
        boundary and so belong to the set, but their midpoint $(2,1)$ does
        not, since $\sqrt2\approx1.414>1$. The set that \emph{is} convex is
        the one below the curve, $\{y\le\sqrt{x},\ x\ge0\}$ --- the
        hypograph of a concave function.
  \item The lines meet where $x+y=10$ and $2x+y=16$, i.e. $x=6$, $y=4$.
        Vertices: $(0,0)$, $(8,0)$, $(6,4)$, $(0,10)$.
  \item With $g$ kg of grain and $s$ kg of soy: $2g+s\ge8$,
        $2g+4s\ge12$, $g\ge0$, $s\ge0$. It is unbounded because the
        constraints are all ``at least'' --- buying arbitrarily much feed is
        always still feasible, so the region extends indefinitely up and to
        the right. A cost-minimising objective is nevertheless attained,
        at one of the two vertices on the lower-left boundary.
  \item The region is the triangle bounded below by $w=10$, above by
        $w=30-L$, and on the left by $L=0$; the two sloping boundaries
        meet at $L=20$, $w=10$. It shows the employment levels compatible
        with a wage floor of $10$ when labour demand is $w=30-L$: the
        largest feasible $L$ is $20$, the employment the floor permits.
        Any $L$ beyond that would require a wage below the minimum.
\end{solutions}

\chapter{Trigonometric Functions}

Trigonometry plays a small part in economics and a large part in the
prerequisites, for two reasons. The graphs of $\sin$, $\cos$ and $\tan$
belong to the catalogue of shapes one is expected to recognise, and they are
the standard test functions of calculus: many derivative and limit
exercises use them precisely because they oscillate and so expose careless
reasoning. Seasonality and business cycles supply the
occasional genuine application.

The chapter therefore covers what is actually needed and stops. The unit
circle and the handful of exact values; the three graphs with their periods,
amplitudes and asymptotes; the transformation form
$a\sin\bigl(b(x-c)\bigr)+d$, which is the machinery of \chref{functions}{3} applied to a
periodic function; the four identities that carry all the simplification
work; and how to write down \emph{every} solution of a trigonometric
equation rather than just the one a calculator returns.

\section{Core Rules}

\begin{rulebox}{Radians, and why}
An angle in \textbf{radians} is the arc length it cuts on the unit circle,
so a full turn is $2\pi$:
\[
  180^{\circ}=\pi,\qquad
  x_{\text{rad}}=x_{\deg}\cdot\frac{\pi}{180},\qquad
  x_{\deg}=x_{\text{rad}}\cdot\frac{180}{\pi} .
\]
Calculus uses radians exclusively, and for one reason: only in radians is
$\dfrac{\mathrm{d}}{\mathrm{d}x}\sin x=\cos x$. In degrees an unwanted
factor $\pi/180$ appears in every derivative.
\end{rulebox}

\begin{rulebox}{The unit circle definition}
For the point on the unit circle reached by turning through the angle $x$
from the positive $x$-axis, $\cos x$ is its first coordinate and $\sin x$
its second. Hence, immediately:
\[
  \sin^{2}x+\cos^{2}x=1,\qquad
  -1\le\sin x\le1,\qquad -1\le\cos x\le1,
\]
the first being Pythagoras. Also $\tan x=\dfrac{\sin x}{\cos x}$, the slope
of the ray, undefined where $\cos x=0$.

Signs by quadrant --- $\sin$ is the height, $\cos$ the width:
\[
  \begin{array}{c|cccc}
     & \text{I} & \text{II} & \text{III} & \text{IV}\\\hline
    \sin & + & + & - & -\\
    \cos & + & - & - & +\\
    \tan & + & - & + & -
  \end{array}
\]
\end{rulebox}

\begin{rulebox}{The values to know by heart}
\[
  \begin{array}{c|ccccc}
    x & 0 & \pi/6 & \pi/4 & \pi/3 & \pi/2\\\hline
    \sin x & 0 & \tfrac12 & \tfrac{\sqrt2}{2} & \tfrac{\sqrt3}{2} & 1\\[0.3ex]
    \cos x & 1 & \tfrac{\sqrt3}{2} & \tfrac{\sqrt2}{2} & \tfrac12 & 0\\[0.3ex]
    \tan x & 0 & \tfrac{1}{\sqrt3} & 1 & \sqrt3 & \text{---}
  \end{array}
\]
Everything else follows from three facts: $\sin$ is \textbf{odd}
($\sin(-x)=-\sin x$), $\cos$ is \textbf{even} ($\cos(-x)=\cos x$), and both
have period $2\pi$ while $\tan$ has period $\pi$. So
$\sin\tfrac{5\pi}{6}=\sin\tfrac{\pi}{6}=\tfrac12$ (second quadrant, $\sin$
positive) and $\cos\tfrac{4\pi}{3}=-\cos\tfrac{\pi}{3}=-\tfrac12$ (third
quadrant, $\cos$ negative).
\end{rulebox}

\begin{gallery}
\begin{tikzpicture}
  \begin{axis}[fnaxis, xmin=-0.6, xmax=7.0, ymin=-1.6, ymax=1.6,
      xtick={3.1416, 6.2832}, xticklabels={$\pi$, $2\pi$},
      ytick={-1,1}, title={$y=\sin x$}]
    \addplot[domain=-0.4:6.9, samples=201]{sin(deg(x))};
  \end{axis}
\end{tikzpicture}\hfill
\begin{tikzpicture}
  \begin{axis}[fnaxis, xmin=-0.6, xmax=7.0, ymin=-1.6, ymax=1.6,
      xtick={3.1416, 6.2832}, xticklabels={$\pi$, $2\pi$},
      ytick={-1,1}, title={$y=\cos x$}]
    \addplot[domain=-0.4:6.9, samples=201]{cos(deg(x))};
  \end{axis}
\end{tikzpicture}\hfill
\begin{tikzpicture}
  \begin{axis}[fnaxis, xmin=-3.4, xmax=3.4, ymin=-3.4, ymax=3.4,
      xtick={-1.5708, 1.5708}, xticklabels={$-\frac{\pi}{2}$,$\frac{\pi}{2}$},
      ytick={-2,2}, title={$y=\tan x$}]
    \addplot[domain=-1.32:1.32, samples=101]{tan(deg(x))};
    \addplot[domain=1.82:3.3, samples=101]{tan(deg(x))};
    \addplot[domain=-3.3:-1.82, samples=101]{tan(deg(x))};
    \addplot[gray!55, dashed, thin, samples=2, forget plot]
      coordinates {(1.5708,-3.3)(1.5708,3.3)};
    \addplot[gray!55, dashed, thin, samples=2, forget plot]
      coordinates {(-1.5708,-3.3)(-1.5708,3.3)};
  \end{axis}
\end{tikzpicture}
\gcap{$\sin$ and $\cos$ are the same wave shifted by $\pi/2$, with period
$2\pi$ and range $[-1,1]$. $\tan$ has period $\pi$, range $\R$, and a
vertical asymptote at every odd multiple of $\pi/2$.}
\end{gallery}

\begin{rulebox}{Reading $a\sin\bigl(b(x-c)\bigr)+d$}
The four constants are exactly the four moves of \chref{functions}{3}, and each has a
name here:
\begin{itemize}[leftmargin=1.6em,itemsep=0.3ex]
  \item $|a|$ is the \textbf{amplitude}: the wave runs from $d-|a|$ to
        $d+|a|$, so the range is $[d-|a|,\,d+|a|]$. Negative $a$ flips it.
  \item $b$ sets the \textbf{period}, $T=\dfrac{2\pi}{|b|}$ --- larger $b$
        means faster oscillation. (For $\tan$, $T=\pi/|b|$.)
  \item $c$ is the \textbf{phase shift}: the whole wave moves $c$ to the
        right. Factor before reading: $\sin(2x-\pi)=\sin\bigl(2(x-\pi/2)\bigr)$
        is a shift of $\pi/2$, not of $\pi$.
  \item $d$ is the \textbf{midline}, the horizontal line the wave
        oscillates about.
\end{itemize}
So $y=3\sin\bigl(2(x-\tfrac{\pi}{4})\bigr)+1$ has amplitude $3$, period
$\pi$, midline $y=1$, range $[-2,4]$, shifted $\pi/4$ right.
\end{rulebox}

\begin{rulebox}{The identities that do the work}
\[
  \sin^{2}x+\cos^{2}x=1,\qquad
  \tan x=\frac{\sin x}{\cos x},\qquad
  1+\tan^{2}x=\frac{1}{\cos^{2}x} .
\]
\[
  \sin(x\pm y)=\sin x\cos y\pm\cos x\sin y,\qquad
  \cos(x\pm y)=\cos x\cos y\mp\sin x\sin y .
\]
Setting $y=x$ in the addition formulas gives the \textbf{double-angle}
forms, which are the ones that appear in practice:
\[
  \sin 2x=2\sin x\cos x,\qquad
  \cos 2x=\cos^{2}x-\sin^{2}x=1-2\sin^{2}x=2\cos^{2}x-1 .
\]
Almost every simplification is one of: replace $1-\cos^{2}x$ by
$\sin^{2}x$, split a $\tan$ into $\sin/\cos$, or collapse a double angle.
\end{rulebox}

\begin{rulebox}{Solving trigonometric equations: give \emph{all} solutions}
A calculator returns one angle; periodicity supplies infinitely many.
\begin{itemize}[leftmargin=1.6em,itemsep=0.35ex]
  \item $\sin x=s$ with $|s|\le1$: two families,
        $x=x_{0}+2k\pi$ and $x=\pi-x_{0}+2k\pi$, where
        $x_{0}=\arcsin s$ and $k\in\mathbb{Z}$.
  \item $\cos x=s$: $x=\pm x_{0}+2k\pi$ with $x_{0}=\arccos s$.
  \item $\tan x=s$: one family, $x=\arctan s+k\pi$, since the period is
        $\pi$.
  \item If the argument is $bx$ rather than $x$, solve for $bx$ first and
        divide \emph{everything}, the $2k\pi$ included: $\sin 2x=1$ gives
        $2x=\tfrac{\pi}{2}+2k\pi$, hence
        $x=\tfrac{\pi}{4}+k\pi$.
\end{itemize}
\end{rulebox}

\begin{rulebox}{Inverse trigonometric functions}
None of the three is injective, so each is inverted only on a restricted
domain:
\[
  \arcsin:[-1,1]\to\Bigl[-\tfrac{\pi}{2},\tfrac{\pi}{2}\Bigr],\qquad
  \arccos:[-1,1]\to[0,\pi],\qquad
  \arctan:\R\to\Bigl(-\tfrac{\pi}{2},\tfrac{\pi}{2}\Bigr).
\]
So $\arcsin(\sin x)=x$ only for $x$ in that range: $\arcsin(\sin\pi)=0$,
not $\pi$. And $\arctan$ is bounded, with horizontal asymptotes
$y=\pm\pi/2$ --- the reason it appears as a saturating S-shaped curve.
\end{rulebox}

\section{Worked Example}

\begin{worked}{Solve $2\sin^{2}x-\sin x=0$ on $[0,2\pi]$ and in general, and read off the shape of $y=3\sin\!\bigl(2(x-\tfrac{\pi}{4})\bigr)+1$}
\wstep{Step 1. Factor. Do not divide.}
\[
  2\sin^{2}x-\sin x=0
  \quad\Longrightarrow\quad
  \sin x\,(2\sin x-1)=0 .
\]
\wtrap{Dividing both sides by $\sin x$ gives $2\sin x-1=0$ and loses every
solution with $\sin x=0$ --- three of the five on $[0,2\pi]$. You may divide
by a quantity only when you know it is non-zero, and $\sin x$ is exactly the
quantity in question. Factor and use the zero-product property instead.}

\wstep{Step 2. Solve each factor on the given interval.}
\emph{First factor:} $\sin x=0$ on $[0,2\pi]$ gives $x=0,\ \pi,\ 2\pi$.

\emph{Second factor:} $\sin x=\tfrac12$. The reference angle is
$\pi/6$, and the sine is positive in the first and second quadrants, so
\[
  x=\frac{\pi}{6}\qquad\text{and}\qquad x=\pi-\frac{\pi}{6}=\frac{5\pi}{6}.
\]
\wtrap{$\arcsin(\tfrac12)=\pi/6$ is one solution, not the solution. The
inverse sine returns a single value from $[-\pi/2,\pi/2]$; the second
solution in each period must be supplied by hand from the symmetry
$\sin(\pi-x)=\sin x$. For cosine the partner is $-x$ (or $2\pi-x$), for
tangent it is $x+\pi$.}

\wstep{Step 3. Collect the five solutions.}
\[
  \boxed{\;x\in\Bigl\{0,\ \tfrac{\pi}{6},\ \tfrac{5\pi}{6},\ \pi,\ 2\pi\Bigr\}\;}
\]
Check the smallest and one middle value: at $x=0$, $0-0=0$ \checkmark; at
$x=\pi/6$, $2\cdot\tfrac14-\tfrac12=0$ \checkmark.

\wstep{Step 4. Give the general solution, since ``all solutions'' usually means all of $\R$.}
\[
  x=k\pi,
  \qquad
  x=\frac{\pi}{6}+2k\pi,
  \qquad
  x=\frac{5\pi}{6}+2k\pi,
  \qquad k\in\mathbb{Z} .
\]
Note the two different periods: the zeros of $\sin$ repeat every $\pi$
(because $\sin(x+\pi)=-\sin x$, which is also zero), while the solutions of
$\sin x=\tfrac12$ repeat only every $2\pi$.
\wtrap{Adding $2k\pi$ to everything is safe but incomplete for the first
family; writing $k\pi$ for the second family is simply wrong. Work out the
period of each family separately from the equation it solves.}

\wstep{Step 5. Now the sinusoid: read the four parameters off $a\sin\bigl(b(x-c)\bigr)+d$.}
For $y=3\sin\bigl(2(x-\tfrac{\pi}{4})\bigr)+1$:
\begin{center}\footnotesize
\begin{tabular}{llll}
  \hline
  parameter & value & name & effect\\
  \hline
  $a$ & $3$ & amplitude & oscillates $3$ above and below the midline\\
  $b$ & $2$ & angular frequency & period $=2\pi/b=\pi$\\
  $c$ & $\pi/4$ & phase shift & shifted \emph{right} by $\pi/4$\\
  $d$ & $1$ & midline & centred on $y=1$, so range $[-2,4]$\\
  \hline
\end{tabular}
\end{center}

\wstep{Step 6. Sketch it from those four numbers.}
\begin{center}
\begin{tikzpicture}
\begin{axis}[fnaxis, width=8.2cm, height=4.4cm,
             domain=-0.4:3.6, samples=201,
             xmin=-0.5, xmax=3.7, ymin=-2.8, ymax=4.8,
             xtick={0.7854, 1.5708, 2.3562, 3.1416},
             xticklabels={$\frac{\pi}{4}$,$\frac{\pi}{2}$,$\frac{3\pi}{4}$,$\pi$},
             ytick={-2,1,4}]
  \addplot[accent!35, thick, dashed, samples=2, forget plot] {1};
  \addplot {3*sin(deg(2*(x-pi/4)))+1};
\end{axis}
\end{tikzpicture}
\end{center}
The curve crosses its midline $y=1$ going upward at $x=\pi/4$ --- that is
what the phase shift means --- peaks at $4$ a quarter-period later
($x=\pi/4+\pi/4=\pi/2$), and repeats after one full period $\pi$.
\wtrap{The period is $2\pi/b$, not $2\pi b$: a \emph{larger} $b$ gives a
\emph{shorter} period. And the phase shift is only readable once the inside
is factored as $b(x-c)$; from $\sin(2x-\tfrac{\pi}{2})$ the shift is
$\pi/4$, not $\pi/2$. It is the same trap as in \chref{functions}{3} on
transformations.}
\end{worked}

\section{Practice Problems}

Convert in problems 1--6.

\begin{problems}
  \item $60^{\circ}$ into radians.
  \item $270^{\circ}$ into radians.
  \item $\tfrac{3\pi}{4}$ into degrees.
  \item $\tfrac{5\pi}{3}$ into degrees.
  \item $1$ radian into degrees, to one decimal.
  \item Why does calculus insist on radians?
\end{problems}

Give the exact value in problems 7--16.

\begin{problems}[resume]
  \item $\sin\tfrac{\pi}{6}$
  \item $\cos\tfrac{\pi}{4}$
  \item $\tan\tfrac{\pi}{3}$
  \item $\sin\tfrac{5\pi}{6}$
  \item $\cos\tfrac{2\pi}{3}$
  \item $\tan\pi$
  \item $\sin\bigl(-\tfrac{\pi}{3}\bigr)$
  \item $\cos\tfrac{7\pi}{2}$
  \item $\sin\tfrac{9\pi}{4}$
  \item $\tan\tfrac{3\pi}{4}$
\end{problems}

In problems 17--24 give period, amplitude (where defined), range and any
asymptotes.

\begin{problems}[resume]
  \item $y=\sin x$
  \item $y=\cos x$
  \item $y=\tan x$
  \item $y=3\sin x$
  \item $y=\sin 2x$
  \item $y=\cos\tfrac{x}{2}$
  \item $y=2\sin x+1$
  \item $y=\tan 3x$
\end{problems}

\begin{problems}[resume]
  \item Describe $y=\sin\bigl(x-\tfrac{\pi}{2}\bigr)$ as a shift, and name
        the simpler function it equals.
  \item Give amplitude, period, midline and range of
        $y=3\sin\bigl(2(x-\tfrac{\pi}{4})\bigr)+1$.
  \item Give amplitude, period and phase shift of
        $y=\sin(2x-\pi)$ --- factor first.
  \item Give amplitude, period and range of $y=-2\cos 3x$.
  \item Write a sine function with amplitude $4$, period $\pi$ and midline
        $y=-1$.
  \item Write a cosine function with amplitude $2$, period $6\pi$ and no
        vertical shift.
  \item Monthly sales oscillate about $500$ units with a swing of $\pm80$
        and a twelve-month cycle, peaking in month $3$. Write $S(t)$.
  \item Why can no function of the form $a\sin(bx)+d$ ever be injective on
        $\R$?
\end{problems}

Simplify or verify in problems 33--38.

\begin{problems}[resume]
  \item $\dfrac{1-\cos^{2}x}{\sin x}$
  \item $\tan x\cdot\cos x$
  \item $\sin^{2}x-1$
  \item $\dfrac{\sin 2x}{2\cos x}$
  \item Verify $1+\tan^{2}x=\dfrac{1}{\cos^{2}x}$.
  \item Verify $\cos 2x=1-2\sin^{2}x$ from the addition formula.
\end{problems}

Give \emph{all} real solutions in problems 39--45.

\begin{problems}[resume]
  \item $\sin x=\tfrac12$
  \item $\cos x=0$
  \item $\tan x=1$
  \item $2\cos x+1=0$
  \item $\sin 2x=1$
  \item $\cos^{2}x=\tfrac14$
  \item $\sin x=\cos x$
\end{problems}

\section{Solutions}

\begin{solutions}
  \item $60\cdot\tfrac{\pi}{180}=\tfrac{\pi}{3}$.
  \item $270\cdot\tfrac{\pi}{180}=\tfrac{3\pi}{2}$.
  \item $\tfrac{3\pi}{4}\cdot\tfrac{180}{\pi}=135^{\circ}$.
  \item $\tfrac{5\pi}{3}\cdot\tfrac{180}{\pi}=300^{\circ}$.
  \item $\tfrac{180}{\pi}\approx57.3^{\circ}$.
  \item Because $\tfrac{\mathrm{d}}{\mathrm{d}x}\sin x=\cos x$ holds only
        in radians. In degrees the limit
        $\lim_{h\to0}\tfrac{\sin h}{h}$ equals $\tfrac{\pi}{180}$ instead
        of $1$, and that factor would appear in every derivative and integral
        of a trigonometric function.
  \item $\tfrac12$.
  \item $\tfrac{\sqrt2}{2}$.
  \item $\sqrt3$.
  \item Second quadrant, $\sin$ positive:
        $\sin\tfrac{5\pi}{6}=\sin\tfrac{\pi}{6}=\tfrac12$.
  \item Second quadrant, $\cos$ negative:
        $-\cos\tfrac{\pi}{3}=-\tfrac12$.
  \item $\tan\pi=\tfrac{\sin\pi}{\cos\pi}=\tfrac{0}{-1}=0$.
  \item $\sin$ is odd: $-\sin\tfrac{\pi}{3}=-\tfrac{\sqrt3}{2}$.
  \item $\tfrac{7\pi}{2}=\tfrac{3\pi}{2}+2\pi$, so
        $\cos\tfrac{7\pi}{2}=\cos\tfrac{3\pi}{2}=0$.
  \item $\tfrac{9\pi}{4}=\tfrac{\pi}{4}+2\pi$, so the value is
        $\sin\tfrac{\pi}{4}=\tfrac{\sqrt2}{2}$.
  \item Second quadrant, $\tan$ negative:
        $-\tan\tfrac{\pi}{4}=-1$.
  \item Period $2\pi$, amplitude $1$, range $[-1,1]$, no asymptotes.
  \item Period $2\pi$, amplitude $1$, range $[-1,1]$, no asymptotes.
  \item Period $\pi$, amplitude undefined (unbounded), range $\R$,
        vertical asymptotes at $x=\tfrac{\pi}{2}+k\pi$.
  \item Period $2\pi$, amplitude $3$, range $[-3,3]$.
  \item Period $\tfrac{2\pi}{2}=\pi$, amplitude $1$, range $[-1,1]$.
  \item Period $\tfrac{2\pi}{1/2}=4\pi$, amplitude $1$, range $[-1,1]$.
  \item Period $2\pi$, amplitude $2$, midline $y=1$, range $[-1,3]$.
  \item Period $\tfrac{\pi}{3}$, unbounded range $\R$, asymptotes where
        $3x=\tfrac{\pi}{2}+k\pi$, i.e. $x=\tfrac{\pi}{6}+\tfrac{k\pi}{3}$.
  \item A shift of $\tfrac{\pi}{2}$ to the right. It equals $-\cos x$,
        since $\sin\bigl(x-\tfrac{\pi}{2}\bigr)
        =\sin x\cos\tfrac{\pi}{2}-\cos x\sin\tfrac{\pi}{2}=-\cos x$.
  \item Amplitude $3$, period $\tfrac{2\pi}{2}=\pi$, midline $y=1$, range
        $[-2,4]$, shifted $\tfrac{\pi}{4}$ to the right.
  \item $\sin(2x-\pi)=\sin\bigl(2(x-\tfrac{\pi}{2})\bigr)$: amplitude $1$,
        period $\pi$, phase shift $\tfrac{\pi}{2}$ right. Reading the
        shift as $\pi$ from the unfactored form is the standard error.
  \item Amplitude $2$ (the minus sign reflects rather than changing the
        amplitude), period $\tfrac{2\pi}{3}$, range $[-2,2]$.
  \item Period $\pi$ needs $b=2$: $y=4\sin(2x)-1$.
  \item Period $6\pi$ needs $b=\tfrac13$:
        $y=2\cos\bigl(\tfrac{x}{3}\bigr)$.
  \item Midline $500$, amplitude $80$, period $12$ so
        $b=\tfrac{2\pi}{12}=\tfrac{\pi}{6}$, and a sine peaks a quarter
        period after its shift, so peaking at $t=3$ means $c=0$:
        \[
          S(t)=500+80\sin\!\left(\frac{\pi t}{6}\right).
        \]
        (Check: at $t=3$ the sine is $\sin\tfrac{\pi}{2}=1$, giving the
        peak $580$.)
  \item Because it is periodic: $f(x+T)=f(x)$ for $T=2\pi/|b|\neq0$, so two
        different inputs share a value. Injectivity fails everywhere, which
        is why the inverse functions need restricted domains.
  \item $\dfrac{\sin^{2}x}{\sin x}=\sin x$ (for $\sin x\neq0$).
  \item $\dfrac{\sin x}{\cos x}\cdot\cos x=\sin x$ (for
        $\cos x\neq0$).
  \item $-(1-\sin^{2}x)=-\cos^{2}x$.
  \item $\dfrac{2\sin x\cos x}{2\cos x}=\sin x$ (for $\cos x\neq0$).
  \item Divide $\sin^{2}x+\cos^{2}x=1$ by $\cos^{2}x$:
        $\tan^{2}x+1=\dfrac{1}{\cos^{2}x}$.
  \item $\cos2x=\cos^{2}x-\sin^{2}x$, and substituting
        $\cos^{2}x=1-\sin^{2}x$ gives $1-\sin^{2}x-\sin^{2}x
        =1-2\sin^{2}x$.
  \item $x=\tfrac{\pi}{6}+2k\pi$ or $x=\tfrac{5\pi}{6}+2k\pi$,
        $k\in\mathbb{Z}$.
  \item $x=\tfrac{\pi}{2}+k\pi$.
  \item $x=\tfrac{\pi}{4}+k\pi$.
  \item $\cos x=-\tfrac12$, so
        $x=\pm\tfrac{2\pi}{3}+2k\pi$.
  \item $2x=\tfrac{\pi}{2}+2k\pi$, hence
        $x=\tfrac{\pi}{4}+k\pi$.
  \item $\cos x=\pm\tfrac12$, giving
        $x=\pm\tfrac{\pi}{3}+2k\pi$ and $x=\pm\tfrac{2\pi}{3}+2k\pi$ ---
        more compactly $x=\tfrac{\pi}{3}+\tfrac{k\pi}{1}$ fails, so keep
        the four families, or write
        $x=\pm\tfrac{\pi}{3}+k\pi$.
  \item Dividing by $\cos x$ (legitimate, since $\cos x=0$ would force
        $\sin x=0$ too, impossible): $\tan x=1$, so
        $x=\tfrac{\pi}{4}+k\pi$.
\end{solutions}

\section{Additional Problems}

\begin{problems}
  \item $45^{\circ}$ and $150^{\circ}$ into radians.
  \item $\tfrac{7\pi}{6}$ and $\tfrac{11\pi}{6}$ into degrees.
  \item $\cos\tfrac{5\pi}{4}$
  \item $\sin\tfrac{4\pi}{3}$
  \item $\tan\tfrac{5\pi}{6}$
  \item $\sin\tfrac{11\pi}{6}$
  \item $\cos(-\tfrac{\pi}{6})$
  \item $\tan\tfrac{7\pi}{4}$
  \item Period, amplitude and range of $y=\tfrac12\cos 4x-2$.
  \item Period and asymptotes of $y=\tan\bigl(x-\tfrac{\pi}{4}\bigr)$.
  \item Amplitude, period, midline of
        $y=5\cos\bigl(\tfrac{\pi}{3}(x-2)\bigr)+3$.
  \item Write a cosine with maximum $7$, minimum $1$ and period $4$.
  \item Simplify $\dfrac{\cos^{2}x-1}{\cos x-1}$.
  \item Simplify $\sin x\cos x\cdot\dfrac{2}{\sin 2x}$.
  \item Simplify $\dfrac{1}{\cos^{2}x}-\tan^{2}x$.
  \item Verify $\sin(x+\pi)=-\sin x$.
  \item Verify $\cos\bigl(\tfrac{\pi}{2}-x\bigr)=\sin x$.
  \item Solve $\sin x=-\tfrac{\sqrt3}{2}$.
  \item Solve $\tan 2x=\sqrt3$.
  \item Solve $2\sin^{2}x-\sin x=0$.
  \item Solve $\cos 2x=\tfrac12$.
  \item Solve $\sin x=2$ and explain the answer.
  \item Give $\arcsin\tfrac12$, $\arccos(-1)$ and $\arctan(-1)$.
  \item Evaluate $\arcsin\bigl(\sin\tfrac{3\pi}{4}\bigr)$ and explain why
        the answer is not $\tfrac{3\pi}{4}$.
  \item Quarterly demand is $D(t)=200-40\cos\bigl(\tfrac{\pi t}{2}\bigr)$
        with $t$ in quarters. Give the period, the maximum and minimum
        demand, and the quarters in which each occurs.
\end{problems}

\section{Additional Solutions}

\begin{solutions}
  \item $\tfrac{\pi}{4}$ and $\tfrac{5\pi}{6}$.
  \item $210^{\circ}$ and $330^{\circ}$.
  \item Third quadrant, $\cos$ negative:
        $-\tfrac{\sqrt2}{2}$.
  \item Third quadrant, $\sin$ negative:
        $-\tfrac{\sqrt3}{2}$.
  \item Second quadrant, $\tan$ negative:
        $-\tfrac{1}{\sqrt3}$.
  \item Fourth quadrant, $\sin$ negative:
        $-\tfrac12$.
  \item $\cos$ is even: $\cos\tfrac{\pi}{6}=\tfrac{\sqrt3}{2}$.
  \item Fourth quadrant, $\tan$ negative: $-1$.
  \item Period $\tfrac{2\pi}{4}=\tfrac{\pi}{2}$, amplitude $\tfrac12$,
        midline $y=-2$, range $[-2.5,-1.5]$.
  \item Period $\pi$; asymptotes where
        $x-\tfrac{\pi}{4}=\tfrac{\pi}{2}+k\pi$, i.e.
        $x=\tfrac{3\pi}{4}+k\pi$.
  \item Amplitude $5$, period $\tfrac{2\pi}{\pi/3}=6$, midline $y=3$,
        shifted $2$ to the right; range $[-2,8]$.
  \item Midline $\tfrac{7+1}{2}=4$, amplitude $\tfrac{7-1}{2}=3$, period
        $4$ so $b=\tfrac{2\pi}{4}=\tfrac{\pi}{2}$:
        $y=3\cos\bigl(\tfrac{\pi x}{2}\bigr)+4$.
  \item $\dfrac{(\cos x-1)(\cos x+1)}{\cos x-1}=\cos x+1$ (for
        $\cos x\neq1$).
  \item $\dfrac{2\sin x\cos x}{\sin 2x}=\dfrac{\sin 2x}{\sin 2x}=1$.
  \item $\dfrac{1}{\cos^{2}x}-\tan^{2}x=1$ by the identity
        $1+\tan^{2}x=1/\cos^{2}x$.
  \item $\sin(x+\pi)=\sin x\cos\pi+\cos x\sin\pi=-\sin x$.
  \item $\cos\bigl(\tfrac{\pi}{2}-x\bigr)
        =\cos\tfrac{\pi}{2}\cos x+\sin\tfrac{\pi}{2}\sin x=\sin x$.
  \item $x=-\tfrac{\pi}{3}+2k\pi$ or
        $x=\pi+\tfrac{\pi}{3}+2k\pi=\tfrac{4\pi}{3}+2k\pi$.
  \item $2x=\tfrac{\pi}{3}+k\pi$, so
        $x=\tfrac{\pi}{6}+\tfrac{k\pi}{2}$.
  \item Factor: $\sin x\,(2\sin x-1)=0$, so $\sin x=0$ giving $x=k\pi$, or
        $\sin x=\tfrac12$ giving $x=\tfrac{\pi}{6}+2k\pi$ and
        $x=\tfrac{5\pi}{6}+2k\pi$.
  \item $2x=\pm\tfrac{\pi}{3}+2k\pi$, so
        $x=\pm\tfrac{\pi}{6}+k\pi$.
  \item No solution: $\sin$ never exceeds $1$, since it is a coordinate of
        a point on the unit circle. An equation can be perfectly
        well-formed and still have an empty solution set.
  \item $\tfrac{\pi}{6}$, \;$\pi$, \;$-\tfrac{\pi}{4}$.
  \item $\sin\tfrac{3\pi}{4}=\tfrac{\sqrt2}{2}$, and
        $\arcsin\tfrac{\sqrt2}{2}=\tfrac{\pi}{4}$. The answer is
        $\tfrac{\pi}{4}$, not $\tfrac{3\pi}{4}$, because $\arcsin$ returns
        values in $[-\tfrac{\pi}{2},\tfrac{\pi}{2}]$ only:
        $\arcsin(\sin x)=x$ holds just on that interval.
  \item Period $\tfrac{2\pi}{\pi/2}=4$ quarters, i.e. one year. Maximum
        $240$ where $\cos\bigl(\tfrac{\pi t}{2}\bigr)=-1$, i.e.
        $\tfrac{\pi t}{2}=\pi+2k\pi$ so $t=2,6,10,\dots$; minimum $160$
        where the cosine is $1$, i.e. $t=0,4,8,\dots$. Demand peaks in the
        third quarter of each year and troughs in the first.
\end{solutions}


\end{document}
